\documentclass[11pt]{article}
\usepackage[margin=1in]{geometry}
\usepackage{amsmath,amssymb}
\usepackage{physics}
\usepackage{array}
\usepackage{tabularx}
\usepackage{graphicx}
\usepackage[hyphens]{url}
\usepackage{xr}
\graphicspath{{results/final/}{../results/final/}{results/}{../results/}{./}{pub/}}

\newcommand{\pubinput}[1]{%
  \IfFileExists{results/final/#1}{\input{results/final/#1}}{%
    \IfFileExists{../results/final/#1}{\input{../results/final/#1}}{%
      \IfFileExists{results/#1}{\input{results/#1}}{%
        \IfFileExists{../results/#1}{\input{../results/#1}}{%
          \IfFileExists{#1}{\input{#1}}{\input{pub/#1}}%
        }%
      }%
    }%
  }%
}

\setlength{\parindent}{0pt}
\setlength{\parskip}{0.6\baselineskip}
\setlength{\emergencystretch}{1.5em}
\raggedbottom

\newcolumntype{L}[1]{>{\raggedright\arraybackslash}p{#1}}
\newcolumntype{Y}{>{\raggedright\arraybackslash}X}

\makeatletter
\g@addto@macro\UrlBreaks{\do\_\do\-}
\makeatother

\renewcommand{\Tr}{\operatorname{Tr}}
\newcommand{\cD}{\mathcal{D}}
\newcommand{\cH}{\mathcal{H}}
\newcommand{\cG}{\mathcal{G}}
\newcommand{\cC}{\mathcal{C}}
\newcommand{\Ahol}{\mathcal{A}_{\rm hol}}
\pubinput{physics_constants.tex}
\pubinput{audit_summary.tex}
\IfFileExists{supplementary.aux}{\externaldocument[supp-]{supplementary}}{%
  \IfFileExists{pub/supplementary.aux}{\externaldocument[supp-]{pub/supplementary}}{}%
}
\newcommand{\tnmainflag}{}

\makeatletter
\newcommand{\manualappendixlabel}[2]{%
  \begingroup
  \def\@currentlabel{#1}%
  \label{#2}%
  \endgroup
}
\makeatother

\title{A Formal Theory of Flavor\\from $SO(10)_{312}$ Boundary Anomalies}
\author{Static Holographic Boundary Theory}
\date{\today}

\begin{document}
\maketitle

\begin{abstract}
We present the static $SO(10)_{312}$ boundary construction as a fixed-parent
benchmark over a disclosed $9{,}801$-cell \emph{Local Moat Audit} of the
visible-level lattice. Within the minimal canonical
$SO(10)\to SU(3)\times SU(2)$ completion used here, the benchmark branch
$(k_{\ell},k_q,K)=(26,8,312)$ is retained by the
\emph{Derived Uniqueness Theorem}, operationally audited by the framing-closure
condition $\Delta_{\rm fr}=0$ for the completed boundary block. On this
anomaly-free branch, modular-$T$ closure is not
compatible with a vanishing vacuum term: the same bookkeeping requires a
strictly positive holographic vacuum loading $\Lambda_{\rm holo}>0$,
equivalently a finite boundary register
$N=3\pi/(L_P^2\Lambda_{\rm holo})$. If $\Lambda_{\rm holo}$ were allowed to
vanish, $N$ would decompactify and the branch-fixed neutrino floor would
collapse.

The PMNS kernel descends from the $SU(2)_{26}$ modular block with
conical-holonomy rotation, while the CKM magnitudes are read as tunneling
amplitudes through the $SO(10)/SU(3)$ coset suppression factor and the CKM apex
angle $\gamma$ follows from the closed $\mathbf{126}_H$ threshold equation
$M_{\rm match}\equiv M_{126}^{\rm match}=M_Ne^{-\beta^2}$ with $M_N$ the
branch-fixed \emph{Modular Restoration Scale}. Once the Topological Coordinate
triplet $(26,8,312)$, the visible-current Topological Coordinate $G_{\rm SM}=15$,
and the representational admissibility constraint
$\langle \Sigma_{126}\rangle/\langle\phi_{10}\rangle=64/312$ are fixed, the
construction retains no detached continuous adjustment sector. Any independent
detuning of the integers $(26,8,312)$ reopens the framing anomaly
$\Delta_{\rm fr}\neq0$, rendering the boundary partition function physically
non-normalizable and violating the Equivalence Principle in the bulk. The
integers $(26,8,312)$ are discrete structural coordinates of the anomaly-free
branch, and the disclosed structural inputs connect the rigid modular branch to
observed units without reopening any continuous Yukawa deformation. The mixing
angles, CP phases, CKM magnitudes, and transported CKM apex are fixed by the
anomaly-free modular data together with the threshold closure, while the
finite-capacity cutoff $N$ supplies the unique structural bridge between the
topological residue and the cosmological anchor for the absolute mass sector
and protects the neutrino floor from vanishing through the holographic surface
tension $\Lambda_{\rm holo}$. With $N=C_{\max}$ fixed by this positive vacuum
loading, that bridge yields $m_\nu\approx2.9\,\mathrm{meV}$ and
$|m_{\beta\beta}|\approx5.6\,\mathrm{meV}$ after RG transport to $M_Z$.
Dark energy is therefore read not as an added particle fluid but as the
macroscopic gravitational residue of the same modular-$T$ closure that unifies
gauge and flavor data on the anomaly-free branch. The eight-observable
comparison is summarized by $\chi_{\rm pred}^2=\chiPredRounded$, quoted as a
descriptive benchmark-consistency metric for the fixed branch rather than as a
landscape-wide likelihood. Within the Minimal One-Copy Dictionary assumption
adopted here, Normal Ordering (NO) provides the benchmark realization, while
Inverted Ordering (IO) remains possible only through a non-minimal enlargement
of the support dictionary. The load-bearing
claims are therefore the anomaly-filtered benchmark construction, the
PMNS/CKM texture map, the topological baryogenesis identity, and the disclosed benchmark-consistency audit. The construction demonstrates that the displayed Standard Model observables are not floating fit targets but mandatory residues of 4D gravity on a finite-capacity horizon.
\end{abstract}

\section{Introduction}

The purpose of this paper is to isolate the part of the boundary program that can be formulated as a controlled structural consistency problem. The emphasis is not on phenomenological optimization but on a construction that can be inspected, stress-tested, and falsified. The retained branch contains no detached continuous adjustment sector. The triplet $(26,8,312)$ is the manuscript's \emph{Topological Coordinate} triplet: three discrete integers fixed by the anomaly filter rather than selected by phenomenological optimization. Every remaining numerical value entering the benchmark is either a structural coordinate of the anomaly-free branch, a representational admissibility constraint, or an \emph{Observational Boundary Condition}~\cite{friedan1987,witten1989}. Once that boundary condition is specified, no continuous Yukawa, threshold, or normalization direction remains on the retained branch.

\paragraph{Definition (Static Holographic).}
By \emph{Static Holographic Boundary Theory} we mean a fixed topological mapping from boundary RCFT data to bulk flavor observables on the selected anomaly-free branch. The term is used in that deliberately narrow sense throughout: it does not denote a dynamical quantum-gravity dual with propagating bulk graviton degrees of freedom, time-dependent bulk reconstruction, or an independent AdS/CFT completion beyond the rigid modular data employed here. In SHBT, the arrow of time is the direction of holographic coarse-graining; therefore, ``dynamical'' evolution in the bulk is the manifestation of the boundary's hierarchical lattice closure. The only non-derived numerical input entering the benchmark is the externally fixed holographic information budget $N$ inferred from $\Lambda_{\rm obs}$; it is the manuscript's distinguished \emph{Observational Boundary Condition}, not an internal scan direction. Just as General Relativity requires the mass of a star ($M$) as an external input to predict an orbit, SHBT requires the universe's capacity ($N$) to convert rigid branch geometry into an absolute mass scale~\cite{cohen1999}. The flavor levels are organized within a disclosed fixed-parent \emph{Local Moat Audit}. In the symmetry-breaking branch analyzed here, the $\scanTrialCount$-point audit runs only over the visible WZW levels and only at the minimal parent level forced by those visible levels; it is not a global search over parent levels $K$. For the benchmark branch the parent level is the derived fixed-parent value
\begin{equation}
 K=\operatorname{lcm}(2k_{\ell},3k_q)=\operatorname{lcm}(52,24)=312,
\end{equation}
fixed by the visible levels $(k_{\ell},k_q)=(26,8)$ rather than scanned as an independent coordinate. The disclosed $\scanTrialCount$-point survey is a \emph{fixed-parent local moat audit}: it numerically exhibits the local modular/framing moat around the analytically derived coordinate $(26,8,312)$ inside this fixed $K=312$ Diophantine class. The survey therefore does not define the branch by numerical optimization; the Derived Uniqueness Theorem fixes the canonical coordinate first, and the moat audit only confirms that the retained cell is the \emph{Minimal Anomaly-Free Local Survivor} of a lexicographic elimination inside the fixed-parent class. These data determine the defect scale, the modular overlap kernels, and the parent-level closure tying the visible sectors to a common origin. Numerical transport is carried out with full coupled one-loop differential-equation integrations in which $(y_t,y_b,y_\tau,g_1,g_2,g_3)$ are evolved as dynamical variables alongside the flavor observables, while split-threshold and quadratic-curvature coefficients are retained as explicit analytic diagnostics. In the present branch all eight displayed observables are propagated by that same fixed transport map; the three CKM magnitudes are fixed already by the coset tunneling amplitudes, and the CKM apex is the only threshold-sensitive quantity, with its low-energy shift governed by the branch-fixed GUT-threshold residue $\mathcal R_{\rm GUT}$ collected in Table~\ref{tab:notation-definitions} and carried by the leading one-loop dimension-5 Wilson coefficient generated at the $SO(10)\to SM$ matching scale. The PMNS mixing angles, leptonic CP phase, CKM magnitudes, and the transported CKM apex are therefore fixed by the modular data of the anomaly-free branch together with the branch-fixed threshold closure, before any UV/IR scale-setting assumption for the absolute neutrino mass sector is invoked. Appendix~\ref{app:rg} records the overlap ratio $\mathcal N_{\rm sol}$ as the geometric solar-angle correction associated with the leptonic RGEs. The result is therefore an eight-row benchmark comparison summarized by the benchmark-consistency metric $\chi_{\rm pred}^2$, while the disclosed scan size $T_{\rm scan}=\scanTrialCount$ records the full visible-level volume of that fixed-parent audit. The reduced $41\times13$ benchmark-centered follow-up scan reported in Supplementary Sec.~S2 records the post-filter local moat around the analytically selected cell. The manuscript therefore presents $(26,8,312)$ as a mathematically rigid Topological Coordinate triplet within the canonical $SO(10)\supset SU(3)\times SU(2)$ construction and then confronts that rigid branch with flavor data.

\paragraph{Structural guide.}
Read in publication order, the load-bearing benchmark chain is intentionally
short and explicit: fixed-parent branch selection and disclosed benchmark
settings $\to$ positive holographic vacuum loading and finite-capacity cutoff
$c_{\rm dark}=\cDarkCompletion$ simultaneously fixing the positive vacuum
loading and the boundary-neutrality parity sink on the anomaly-free branch,
with that parity sink interpreted strictly as topological bookkeeping for
boundary neutrality rather than as a new unseen particle sector
$\to$ PMNS kernel and CKM tunneling amplitudes
$\to$ branch-fixed heavy-threshold scale
$M_{\rm match}\equiv M_{126}^{\rm match}=M_Ne^{-\beta^2}$ for the CKM phase,
with $M_N$ the Modular Restoration Scale fixed by $I_L=6$ and $I_Q=13$
$\to$ topological baryogenesis identity
$\eta_B=C_{\rm sph}J_{CP}^{\rm topo}(M_N/M_P)$ with the same
$\Lambda_{\rm holo}$ supplying the static Sakharov replacement: the RG
incompatibility between the Planck-scale boundary and the Modular Restoration
Scale $M_N$,
whose bulk image is the unique geometric-friction realization of the
out-of-equilibrium slot $\to$ coupled transport of the eight flavor
observables together with the separate RG Consistency Audit of the
finite-capacity mass relation $\to$ benchmark-consistency metric
$\chi_{\rm pred}^2$. Within the main text, only
the eight-row flavor chain enters the quoted benchmark tally, while the
transported absolute-mass rows are recorded through the separate RG
Consistency Audit and are not folded into $\chi_{\rm pred}^2$. The supplement
records the local-moat scan, the VEV/SVD stability checks, and the
threshold/transport cross-checks, while the appendices collect the
ultraviolet, RG, and extended gravity-side derivations that support the
main-text spacetime--flavor unification without opening a detached adjustment
sector. Those appendix details do not feed back into branch selection,
threshold matching, texture transport, or the quoted benchmark tallies. Put
differently: the load-bearing claims are the
anomaly-filtered branch selection, the positive vacuum-loading /
finite-capacity closure of the one-copy branch, the topologically locked
PMNS/CKM texture map, the threshold lock for $\gamma$, and the disclosed
benchmark audits.

\subsection*{Glossary Box: Benchmark Definitions}

Several recurrent terms are used in a narrow benchmark sense throughout the paper. For quick reference, Table~\ref{tab:glossary-box} collects the main navigation cues used in the load-bearing argument, while Table~\ref{tab:notation-definitions} serves as the single formal definition block for the core benchmark symbols.

\begin{table}[t]
\centering
\small
\renewcommand{\arraystretch}{1.12}
\begin{tabularx}{\linewidth}{@{}L{0.31\linewidth}Y@{}}
\hline
Term & Role in the load-bearing chain \\
\hline
Derived Uniqueness Theorem / Framing Gap $\Delta_{\rm fr}$ & Branch-selection theorem and operational framing diagnostic; its formal definition and benchmark value are collected once in Table~\ref{tab:notation-definitions}. \\
Modularity Gap $\Delta_{\rm mod}$ & Residual central-charge mismatch reported only after theorem-level framing filtering; it does not select the benchmark by itself. \\
Threshold Residue $\mathcal R_{\rm GUT}$ & Branch-fixed CKM transport datum entering the threshold lock for $\gamma$. \\
Genus tension $\beta$ & Branch-fixed ladder constant controlling both the neutrino support spacing and the threshold relation $M_{\rm match}=M_N e^{-\beta^2}$ with $M_N$ the Modular Restoration Scale. \\
Simplex residue $\kappa_{D_5}$ & Closed-form $D_5$ geometric factor entering the finite-capacity mass relation $m_\nu=\kappa_{D_5}M_PN^{-1/4}$. \\
Boundary Mapping & The branch-fixed continuation in which the selected visible cell, the parent central-charge completion, and the one-copy support map are imposed simultaneously. \\
Benchmark Branch & The fixed-parent visible branch $(k_{\ell},k_q,K)=(26,8,312)$ carried forward to phenomenology after the anomaly-cancellation filter is applied. \\
Threshold-Lock & The branch-fixed heavy-threshold condition $M_{\rm match}\equiv M_{126}^{\rm match}=M_N e^{-\beta^2}$ used in the CKM-phase transport, where $M_N$ is the Modular Restoration Scale. \\
Residual Audit & The eight-row benchmark-consistency metric $\chi_{\rm pred}^2$ together with the gauge-side consistency pulls; the separate RG Consistency Audit carries the displayed mass-relation rows without entering that tally. \\
\hline
\end{tabularx}
\caption{Glossary box for the benchmark terminology and reviewer-navigation cues used in the Introduction and the main discussion.}
\label{tab:glossary-box}
\end{table}

In the same notation, $\alpha^{-1}_{\rm surf}$ denotes the branch-fixed surface gauge-density datum, whereas $\alpha^{-1}_{\rm em}(M_Z)$ denotes the infrared value obtained only after one-loop running and heavy-threshold matching. The former is the primary discrete gauge-side output of the branch; the latter is displayed only as a threshold-sensitive residual, because the ultraviolet surface datum is branch-fixed while the infrared comparison remains sensitive to one-loop running and heavy-threshold matching.

\subsection*{Notation and Definitions}

For reader convenience we collect the fixed benchmark constants that recur in
the modularity, threshold-matching, and gauge-transport formulas. Throughout
the paper the WZW central charges are computed from
\begin{equation}
 c(g_k)=\frac{k\,\dim(g)}{k+h_g^{\vee}}.
\end{equation}

\begin{table}[t]
\centering
\small
\renewcommand{\arraystretch}{1.12}
\begin{tabularx}{\linewidth}{@{}L{0.18\linewidth}L{0.47\linewidth}Y@{}}
\hline
Symbol & Definition & Benchmark value \\
\hline
$(k_{\ell},k_q,K)$ & Visible and parent affine levels & $\benchmarkVisibleBranch$ \\
$\begin{gathered}h^{\vee}_{SU(2)},h^{\vee}_{SU(3)},\\ h^{\vee}_{SO(10)}\end{gathered}$ & Dual Coxeter numbers of the visible and parent groups & $(2,3,8)$ \\
$\Delta_{\rm fr}$ & $\operatorname{dist}\!\left(K/(2k_{\ell}),\mathbb Z\right)$ & $0$ on the benchmark branch \\
$c\bigl(SU(2)_{26}\bigr)$ & $26\cdot 3/(26+2)$ & $\dfrac{39}{14}\approx 2.785714$ \\
$c\bigl(SU(3)_8\bigr)$ & $8\cdot 8/(8+3)$ & $\dfrac{64}{11}\approx 5.818182$ \\
$c_{\rm vis}$ & $c\bigl(SU(2)_{26}\bigr)+c\bigl(SU(3)_8\bigr)$ & $\cVisibleBenchmarkExact\approx 8.603896$ \\
$c\bigl(SO(10)_{312}\bigr)$ & $312\cdot45/(312+8)$ & $\dfrac{351}{8}=43.875$ \\
$c_Y^{\rm ref}$ & Reference hypercharge contribution in the $SO(10)\supset SU(3)\times SU(2)\times U(1)$ embedding & $1$ \\
$c_{\rm coset}^{\rm ref}$ & $c\bigl(SO(10)_{312}\bigr)-c\bigl(SU(3)_{312}\bigr)-c\bigl(SU(2)_{312}\bigr)-c_Y^{\rm ref}$ & $\dfrac{4216243}{131880}\approx 31.970299$ \\
$\Delta_{\rm mod}(26,8)$ & $\operatorname{dist}\!\left((c_{\rm parent}-c_{\rm vis}-c_{\rm coset}^{\rm ref})/24,\mathbb Z\right)$ & $\dfrac{1197103}{8704080}\approx 0.13753355$ \\
$c_{\rm comp}\,(=c_{\rm dark})$ & $24\,\Delta_{\rm mod}(26,8)$, the formal completion residue; observer-level label: the Topological Gravitational Ghost (TGG), i.e. the information deficit created when $SO(10)_{312}$ is projected onto the visible sector & $\cDarkCompletionExact\approx \cDarkCompletionFull$ \\
$\kappa_{D_5}$ & $R_{01}^{\rm par/vis}$, the closed-form $D_5$ simplex residue & $\simplexPrefactor$ \\
$\beta$ & $\tfrac12\ln\cD_{26}$, the branch-fixed modular-$T$ matrix phase & $\inflationGenusBeta$ \\
$\alpha^{-1}_{\rm surf}$ & $15K/(k_{\ell}+k_q)$ & $\alphaSurfBenchmarkExact\approx \alphaSurfBenchmarkDecimal$ \\
$\mathcal R_{\rm GUT}$ & $k_q/(k_{\ell}+h^{\vee}_{SU(2)})$ & $\dfrac{8}{28}=\dfrac{2}{7}\approx 0.285714$ \\
\hline
\end{tabularx}
\caption{Notation and definitions for the fixed $(26,8,312)$ benchmark. This table is the single formal definition block for the core benchmark symbols used below.}
\label{tab:notation-definitions}
\end{table}

\subsection*{Load-Bearing Chain}

\noindent
\fbox{%
\begin{minipage}{0.965\linewidth}
\small
\textbf{Referee roadmap.} The load-bearing manuscript argument is intentionally short.
\par\smallskip
\textbf{(1) Branch Selection.} Apply the Derived Uniqueness Theorem: WZW branching integrality, Witten framing quantization, and GKO orthogonality retain only the visible cell $(26,8,312)$.
\par\smallskip
\textbf{(2) Finite-Capacity Closure.} On the retained anomaly-free branch, require positive holographic vacuum loading $\Lambda_{\rm holo}>0$ so that $N=3\pi/(L_P^2\Lambda_{\rm holo})$ remains finite and the neutrino floor is protected from vanishing.
\par\smallskip
\textbf{(3) Texture Derivation.} Read the PMNS and CKM kernels from the fixed modular data of the $SU(2)_{26}$ and $SU(3)_8$ blocks.
\par\smallskip
\textbf{(4) Scale Lock.} Fix the CKM-apex transport through the branch-fixed threshold relation $M_{\rm match}\equiv M_{126}^{\rm match}=M_N e^{-\beta^2}$, where $M_N$ is the Modular Restoration Scale fixed by $I_L=6$ and $I_Q=13$, equivalently the exact branch residue $\mathcal R_{\rm GUT}=k_q/(k_{\ell}+h^{\vee}_{SU(2)})=8/28$; the explicit scalar-spectrum matching sum is used only as a consistency check on this target.
\par\smallskip
\textbf{(5) The Unification of Spacetime and Flavor.} Promote $c_{\rm dark}=\cDarkCompletion$ to the Holographic Stabilizer that simultaneously fixes $\Lambda_{\rm holo}>0$ and the conjugate $B-L$ parity sink, where that sink is the Topological Gravitational Ghost of the projection deficit and an information-theoretic bookkeeping requirement for boundary neutrality rather than a new WIMP-like particle sector, so the same modular-$T$ closure realizes both the Bianchi lock and $\eta_B=C_{\rm sph}J_{CP}^{\rm topo}(M_N/M_P)$.
\par\smallskip
\textbf{(6) Residual Audit.} Report the eight-row benchmark-consistency metric $\chi_{\rm pred}^2$ together with the gauge-side pulls, while carrying the absolute-mass benchmark through the separate RG Consistency Audit.
\end{minipage}}

\subsection*{Boundary Matching Conditions (Inputs)}

Table 0 isolates the structural inputs and observational anchors that are disclosed rather than derived inside the flavor transport.

\begin{table}[t]
\centering
\small
\renewcommand{\arraystretch}{1.15}
\textbf{Table 0. Structural Inputs and Observational Anchors.}\par\medskip
\begin{tabularx}{\linewidth}{L{0.30\linewidth}L{0.24\linewidth}Y}
\hline
Item & Value & Status / role \\
\hline
Branch integers & $(26,8,312)$ & Topological Coordinate triplet; discrete and untunable \\
Current-Algebra Neutrality & $G_{\rm SM}=15$ & Visible-current neutrality weight \\
Representational Admissibility & $\Sigma/\phi=64/312$ & Scalar Ward Identity / basis-invertibility constraint \\
Observed cosmological constant & $\Lambda_{\rm obs}$ & Cosmological anchor \\
\hline
\end{tabularx}
\end{table}

Once the Topological Coordinate triplet $(26,8,312)$, the visible-current Topological Coordinate $G_{\rm SM}=15$, and the representational admissibility constraint $\langle \Sigma_{126}\rangle/\langle\phi_{10}\rangle=64/312$ are fixed, the construction retains no detached continuous adjustment sector. Every numerical value entering the benchmark is then classified as a structural coordinate, a representational admissibility constraint, or an \emph{Observational Boundary Condition}~\cite{friedan1987,witten1989}. The disclosed flavor-side structural inputs are
\begin{equation}
 G_{\rm SM}=16_{\rm Weyl}-1_{\nu_R}=15,
 \qquad
 \frac{\langle \Sigma_{126}\rangle}{\langle\phi_{10}\rangle}=\frac{64}{312}.
\end{equation}
Here $G_{\rm SM}$ is the branch-fixed Current-Algebra Neutrality weight obtained by starting from the $16$ Weyl channels of one $SO(10)$ spinor generation and removing the SM-singlet right-handed-neutrino direction from the visible current algebra. It is a disclosed Topological Coordinate required to connect the modular kernel to observed units; it is not a hidden floating phenomenological coefficient. The second is not an arbitrary decimal: on the selected branch the chosen $SU(2)\times SU(3)$ branching fixes the unique rational combination of visible affine levels and dual Coxeter numbers,
\begin{equation}
 \frac{\langle \Sigma_{126}\rangle}{\langle\phi_{10}\rangle}
 =
 \frac{h^{\vee}_{SU(2)}}{h^{\vee}_{SU(3)}}\frac{k_q}{k_{\ell}}
 =
 \frac{2}{3}\frac{8}{26}
 =
 \frac{64}{312}.
\end{equation}
We therefore treat $64/312$ as the disclosed \emph{Representational Admissibility Constraint} of the benchmark. It is the exact inverse-Clebsch rational value associated with the selected branch and is required by the Scalar Ward Identity for basis invertibility of the charged-fermion map. In the benchmark bookkeeping it enters as a leading-order scalar-sector singular-value relation, but it is not a value tuned to repair the $b/\tau$ hierarchy or to improve the fit. Even though the Higgs potential itself is not solved in this paper, the ratio is treated as a fixed branch-compatible input rather than as a hidden continuous scan variable. The full benchmark is conditioned on the fixed Topological Coordinate triplet, the external observational boundary condition $N$, and four disclosed benchmark residues.\footnote{Detailed runtime anchors and reproducibility metadata are collected in Appendix~\ref{app:computational}. They document the synchronized benchmark implementation but do not enter the branch-selection logic itself.}

\begin{table}[t]
\centering
\small
\renewcommand{\arraystretch}{1.15}
\begin{tabularx}{\linewidth}{L{0.16\linewidth}L{0.23\linewidth}L{0.21\linewidth}Y}
\hline
Disclosure tier & Branch datum / disclosed residue & Selection principle & Role in the mapping to observed units \\
\hline
Tier 1: Discrete Level Data & $(k_{\ell},k_q,K)=(26,8,312)$ & \shortstack[l]{lexicographic elimination\\Minimal Anomaly-Free\\Local Survivor} & Fixed branch data set by the selected anomaly-free cell \\
Tier 2: External Boundary Conditions & Cosmological bit budget $N$ & Observational Boundary Condition & Observation-anchored boundary input for the benchmark mass relation $m_\nu\propto N^{-1/4}$ \\
Tier 3: Benchmark Residues & \shortstack[l]{Current-Algebra\\Neutrality\\$G_{\rm SM}=15$} & visible-current neutrality & Structural visible-current normalization of the gauge proxy, with the SM-singlet excluded from visible transport \\
Tier 3: Benchmark Residues & \shortstack[l]{Representational\\Admissibility\\$64/312$} & Scalar Ward Identity / basis invertibility & Branch-fixed singular-value admissibility constraint for the charged-fermion map \\
\hline
\end{tabularx}
\caption{Benchmark disclosure block for the fixed $(26,8,312)$ branch. Tier~1 lists rigid discrete branch data, Tier~2 lists the observation-anchored cosmological boundary input, and Tier~3 lists the disclosed flavor-sector matching conditions / residues used to map the topological kernel to observed units.}
\label{tab:benchmark-input-disclosure}
\end{table}

Table~\ref{tab:benchmark-input-disclosure} is a disclosure block rather than a scan menu: it records the fixed branch labels, the observation-anchored cosmological boundary condition, and the primary flavor-sector structural data used to map the topological kernel to physical units. The generation-weight factor $G_{\rm SM}=15$ is exported explicitly for audit transparency, but its value is not an auxiliary floating normalization: it is the Current-Algebra Neutrality factor of the one-generation visible current algebra of the parent $SO(10)$ spinor, namely $16_F-1=15$ once the SM-singlet direction is removed from the visible current-current correlator. By contrast, the VEV-alignment ratio $64/312$ is treated as a disclosed flavor-side Representational Admissibility Constraint. It is the exact inverse of the $\mathbf{126}_H$ Clebsch--Gordan coefficient associated with the chosen $SU(2)\times SU(3)$ branching,
\begin{equation}
 \frac{\langle\Sigma_{126}\rangle}{\langle\phi_{10}\rangle}
 =
 \frac{h^{\vee}_{SU(2)}}{h^{\vee}_{SU(3)}}\frac{k_q}{k_{\ell}}
 =
 \frac{2}{3}\frac{8}{26}
 =
 \frac{64}{312}.
\end{equation}
In particular, the ratio $\langle\Sigma_{126}\rangle/\langle\phi_{10}\rangle=64/312$ is a Representational Admissibility Constraint required by the Scalar Ward Identity for basis invertibility. It is therefore not a tuning dial inserted to fix the $b/\tau$ ratio. Once imposed, it acts only on the singular values (masses) of the charged-fermion map while the topological kernel continues to provide rigid mixing eigenvectors. The main-text SVD Rigidity audit of Fig.~\ref{fig:vev-stability-main} (see also Supplementary Table~\ref{supp-tab:s3-vev-protection} and Supplementary Fig.~\ref{supp-fig:s3-vev-stability}) proves the separation of roles: the displayed $\pm10\%$ VEV sweep leaves the PMNS/CKM angles far inside their quoted $1\sigma$ windows, securing the flavor predictions against scalar-potential uncertainty while the matched mass sector is adjusted. The Higgs potential is not solved here; instead the manuscript isolates the exact inverse-Clebsch benchmark ratio and assumes a scalar potential whose vacuum alignment realizes it. For conservative pull-table bookkeeping, the driver still reports the pair $(15,64/312)$ as branch-fixed disclosure data. The normalized GUT-threshold residue is instead fixed by the exact branch identity $\mathcal R_{\rm GUT}=k_q/(k_{\ell}+h^{\vee}_{SU(2)})=8/28$, with the explicit scalar-spectrum matching sum used only as a consistency check on that target; this branch-fixed CKM transport datum is counted as an additional Wilson-coefficient subtraction only in off-shell threshold sweeps away from the benchmark value. % pub/tn.py: PullTable.benchmark_matching_condition_subtraction_count, PullTable.lambda_normalization_matching_subtraction_count, PullTable.threshold_alignment_subtraction_count, VOA_BRANCHING_GAP, R_GUT
For drift control, the exact fractions and displayed benchmark decimals in Tables~\ref{tab:notation-definitions}, \ref{tab:benchmark-input-disclosure}, and \ref{tab:discrete-inputs-appendix} are synchronized against the reproducibility bundle summarized in Appendix~\ref{app:computational}.

\paragraph{Finite-Capacity Mass Theorem.} The benchmark mass relation is not treated as a floating UV/IR transport ansatz or as an isolated neutrino-mass claim. In the present nomenclature it is the structural bridge between the topological residue and the cosmological anchor. On the anomaly-free one-copy cell, the infrared neutrino scale is fixed by the finite-capacity relation
\begin{equation}
 m_\nu \equiv \kappa_{D_5} M_P N^{-1/4}.
\end{equation}
Here $\kappa_{D_5}$ denotes the closed-form $D_5$ simplex residue that transfers the CKN-saturated parent density to the visible branch. In the present benchmark we keep $\kappa_{D_5}=0.988771051266\approx\simplexPrefactor$ as the exact anomaly-free branch residue rather than attach an auxiliary theory band. Once the bit budget is fixed to the finite holographic capacity $N=C_{\max}$ inferred from the observed cosmological constant $\Lambda_{\rm obs}$, the quarter-power law becomes a theorem on the minimal branch: because $N$ is finite and $\kappa_{D_5}>0$, one must have $m_\nu>0$. The resulting $\sim2.9\,\mathrm{meV}$ scale is therefore a theory-fixed consequence of finite holographic capacity, not an independently floated quantity. Operationally, the RG-transported spectrum is then checked against the Planck 2018 cosmological bound $\sum m_\nu < 0.12\,\mathrm{eV}$ together with the corresponding $|m_{\beta\beta}|$ benchmark band. In Table~\ref{tab:global-flavor-fit}, the rows $m_1$ and $\sum m_\nu$ are accordingly displayed as \emph{Theory-Fixed} rather than as floated coordinates, because varying them would move one off the anomaly-free matching dictionary rather than explore an internal continuous direction.

The discrete levels $k_{\ell}=26$ and $k_q=8$ should therefore be read as derived integers rather than as outputs of a continuous recalibration. Once the discrete vacuum label $(26,8,312)$ and the cosmological bit-budget $N$ anchored to $\Lambda_{\rm obs}$ are fixed as boundary conditions, the flavor observables are uniquely determined without further continuous adjustment.

\paragraph{Derived Uniqueness Theorem.}
The benchmark branch is fixed by three analytic constraints, not by a minimal-embedding assumption. First, the affine WZW branching rules require every embedded visible current block to satisfy
\begin{equation}
 K=x_i\,I_i\,k_i,
\end{equation}
with $x_i$ the canonical embedding index and $I_i$ the descendant branching multiplicity. For the canonical $SO(10)\supset SU(3)\times SU(2)$ branch used here, the embedding indices are strictly
\begin{equation}
 x_{SU(2)}=x_{SU(3)}=1.
\end{equation}
In the standard WZW normalization the visible descendants therefore obey
\begin{equation}
 K=2I_Lk_{\ell}=3I_Qk_q,
 \qquad
 I_L=\frac{K}{2k_{\ell}},
 \qquad
 I_Q=\frac{K}{3k_q}.
\end{equation}
Hence both descendant indices must be integers. For the benchmark visible pair $(k_{\ell},k_q)=(26,8)$ one obtains
\begin{equation}
 K=\operatorname{lcm}(2k_{\ell},3k_q)=\operatorname{lcm}(52,24)=312,
 \qquad
 I_L=6,
 \qquad
 I_Q=13,
\end{equation}
so the admissible Diophantine tower is $K\in\{312n\}_{n\ge1}$ and $K=312$ is its first element.

Second, Witten's framing quantization makes the boundary modularity condition exact. A unit framing shift acts through the WZW modular-$T$ matrix, so the completed boundary partition function is single-valued only when the descended vacuum block carries an integer power of the framing operator~\cite{witten1989}. For the leptonic sector,
\begin{equation}
 T_{00}^{(2,k_{\ell})}=\exp\!\left[-\frac{2\pi i}{24}c\bigl(SU(2)_{k_{\ell}}\bigr)\right],
 \qquad
 Z_{\partial}^{(\nu+1)}=\left(T_{00}^{(2,k_{\ell})}\right)^{I_L}Z_{\partial}^{(\nu)}.
\end{equation}
We therefore use
\begin{equation}
 \Delta_{\rm fr}
 \equiv
 \operatorname{dist}\!\left(\frac{K}{2k_{\ell}},\mathbb Z\right)
 =
 \operatorname{dist}\!\left(I_L,\mathbb Z\right)
\end{equation}
as the printed shorthand for the framing defect on the displayed $k_q=8$ slice, with the companion requirement $I_Q=K/(3k_q)\in\mathbb Z$ imposed simultaneously. The condition
\begin{equation}
 \Delta_{\rm fr}=0
\end{equation}
is therefore exactly the Witten-framing single-valuedness condition for the visible branch. Any alternative embedding path would necessarily induce non-canonical descendant indices, hence fractional vacuum $T$-powers, and would violate boundary modularity before phenomenology is even consulted.

Third, the completion sector is rigid by the Goddard--Kent--Olive theorem. Once the visible stress tensor is fixed to
\begin{equation}
 T^{\rm vis}=T^{SU(2)_{26}}+T^{SU(3)_8},
\end{equation}
the complementary sector is uniquely the Virasoro-orthogonal coset
\begin{equation}
 T^{\rm comp}=T^{SO(10)_{312}}-T^{\rm vis},
 \qquad
 [T^{\rm vis}(z),T^{\rm comp}(w)]=0.
\end{equation}
This is precisely the manuscript's \emph{Reference Coset}. Once the parent and visible currents are fixed, there is no second independent decomposition of the $SO(10)$ stress tensor that preserves Virasoro orthogonality. Reassigning the completion sector would either spoil the commutator condition or alter the central-charge bookkeeping of the completed boundary block.

Taken together, the canonical WZW branching rules, the framing quantization condition $\Delta_{\rm fr}=0$, and the GKO orthogonality of the Reference Coset remove the old Minimal Embedding Assumption entirely. The surviving coordinate is not a provisional benchmark choice but the canonical one-copy solution. Within the branch-fixed gauge window, the higher members $K=624,936,\dots$ of the same Diophantine tower are excluded by the surface gauge-density cutoff, so the coordinate $(26,8,312)$ is absolutely mathematically rigid under the stated constraints. The nonzero value $\Delta_{\rm mod}(26,8)=\visibleResidual$ is therefore a diagnostic residual of the fixed-parent construction rather than a sign of selection ambiguity. The mass coordinate $m_\nu=\kappa_{D_5}M_PN^{-1/4}$ is defined only after this rigid discrete branch has been fixed and does not reopen the branch selection. Supplementary Sec.~\ref{supp-supp-sec:s14-mathematical-logic-flow} records the Mathematical Logic Flow of the benchmark chain, and Supplementary Sec.~S3 derives the corresponding stress tensor.

The presentation is deliberately conservative. The modular kernels are treated as scale-invariant topological objects. Their comparison with low-energy angles is postponed to a later discussion of renormalization-group evolution. Likewise, the Higgs condensate is not replaced by topology; rather, it is interpreted as the electroweak interface that converts boundary sectors into ordinary bulk masses.

With this separation of roles, the formal content of the construction can be summarized in five steps: the simplex-area relation, the $SU(2)_{26}$ mixing kernel, the boundary--bulk interface action, the Derived Uniqueness Theorem that fixes the $SO(10)_{312}$ benchmark branch, and the minimal one-copy dictionary limitation for Inverted Ordering.

\subsection*{Discrete-Choice Audit}

The benchmark retains only a small set of discrete structural coordinates, and each is fixed by lexicographic elimination under the Derived Uniqueness Theorem rather than by a floating phenomenological scan. First, the branch label $(k_{\ell},k_q,K)=\benchmarkVisibleBranch$ is the \emph{Minimal Anomaly-Free Local Survivor} of the combined Diophantine branching rules, framing cancellation, and surface-gauge anchor: on the selected slice, $\benchmarkVisiblePair$ is the only local cell with integer descendants, while $K=\benchmarkParentLevel$ is the first admissible parent in the Diophantine tower that also reproduces $\alpha^{-1}_{\rm surf}\approx\alphaSurfBenchmarkRounded$. Second, within the minimal one-copy dictionary the ordered genus assignment $g=0,1,2$ is the unique survivor compatible with the anomaly-free support ladder and the normal-ordering benchmark; permuted or inverted assignments require a non-minimal extension and are therefore not retained in the minimal benchmark. Third, once the branch is fixed, the Reference Coset is the unique GKO Virasoro-orthogonal complement of the visible stress tensor, and the completion residue is therefore not an extra choice but the unique modular-closure datum required for the completed boundary partition function on that branch. Exact commands, file paths, and software entry points are collected separately in Appendix~\ref{app:computational} so that the physics narrative remains notation-first.

Table~\ref{tab:decision-uniqueness} states the same logic explicitly. It is the local answer to the look-elsewhere objection on the disclosed benchmark slice: once the anomaly filter, the one-copy rank test, and the surface-gauge anchor are imposed, the branch-selection chain contains no continuously retuned nuisance coordinate.

\begin{table}[t]
\centering
\small
\begin{tabular}{@{}>{\raggedright\arraybackslash}p{0.19\linewidth}>{\raggedright\arraybackslash}p{0.19\linewidth}>{\raggedright\arraybackslash}p{0.24\linewidth}>{\raggedright\arraybackslash}p{0.30\linewidth}@{}}
\hline
Discrete choice & Retained benchmark value & Rejected alternatives & Removal criterion \\
\hline
Visible branch & $(k_{\ell},k_q,K)=\benchmarkVisibleBranch$ & neighboring visible cells in the fixed-parent moat and alternative low-rank branch labels & any alternative either violates the anomaly filter $\Delta_{\rm fr}=0$ or, once moved off the retained gauge slice, misses the branch anchor $\alpha^{-1}_{\rm surf}$ by more than $10\%$ \\

Genus ordering & $(g_1,g_2,g_3)=(0,1,2)$ in the Minimal One-Copy Dictionary & inverted or permuted genus ladders & the support map becomes rank-deficient in the minimal benchmark dictionary, so the corresponding mass matrix is not retained as a non-singular benchmark branch \\

Coset bookkeeping & reference coset completion on the anomaly-free branch & alternative coset types or reassigned completion channels & the altered completion either breaks the same $\Delta_{\rm fr}=0$ closure logic or misses the discrete surface-gauge anchor by more than $10\%$ after re-embedding \\
\hline
\end{tabular}
\caption{Decision uniqueness table for the branch-fixed benchmark. The retained branch is selected entirely by discrete consistency checks. Every nearby alternative is removed by one of three pre-comparison criteria: failure of the anomaly filter, rank-deficiency of the minimal support map, or a $>10\%$ miss of the surface-gauge anchor.}
\label{tab:decision-uniqueness}
\end{table}

\section{Comparative Context and Literature Review}

\paragraph{Derived Uniqueness Theorem.}
The condition $\Delta_{\rm fr}=0$ is the framing component of the Derived Uniqueness Theorem: for the canonical $SO(10)\to SU(3)\times SU(2)$ embedding, the completed boundary partition function is well defined only when the modular-$T$ operator acts with integer support, exactly as required by Witten framing quantization~\cite{witten1989}. Together with the Diophantine branching identities and the GKO Virasoro-orthogonal completion, this condition fixes $(26,8,312)$ rather than merely preselecting it by a numerical filter.

\subsection{The Holographic Flavor Dictionary}

This paper uses ``Static Holographic Boundary Theory'' in a deliberately narrow flavor-sector sense. The goal is to define a fixed dictionary between discrete boundary data and benchmark flavor observables, not to claim a complete microscopic dual of Nature or a dynamical quantum-gravity background.

\paragraph{CFT Data.}
The boundary input is the pair of modular matrices $S$ and $T$ of the $SO(10)_{312}$ WZW model, together with the visible $SU(2)_{26}$ and $SU(3)_8$ branching blocks inherited from the anomaly-free cell $(k_{\ell},k_q,K)=(26,8,312)$. These matrices provide the rigid topological data from which the flavor-side benchmark kernels are constructed.

\paragraph{Bulk Counterpart.}
The corresponding four-dimensional side of the benchmark is a static holographic boundary construction used as the flavor-sector UV completion in the limited sense adopted here: the selected boundary data are mapped into effective bulk flavor observables only after electroweak matching and renormalization-group transport. This is a bookkeeping continuation for the flavor sector, not a claim of a complete dynamical quantum-gravity dual.

\paragraph{Dictionary.}
Within this benchmark, the PMNS kernel is identified with the leptonic modular blocks of the $SU(2)_{26}$ sector dressed by the branch-fixed RT-holonomy rotation, while the CKM kernel is identified with the $SO(10)/SU(3)$ coset tunneling amplitudes together with the branch-fixed heavy-threshold closure that transports the quark-sector phase data. In that precise sense, the PMNS/CKM kernels are read from modular blocks and coset tunneling amplitudes rather than promoted to independent low-energy textures.

\paragraph{Framework Disclaimer.}
The construction is therefore presented as a \emph{mathematical flavor benchmark} intended to test the feasibility of discrete boundary selection and a fixed holographic flavor dictionary. It is not claimed as a full derivation from first-principles String Theory, nor as a complete microscopic derivation of four-dimensional gravity or cosmology.

\subsection*{Methodological Scope}

The rigidity claim made here is mathematical, not statistical: within the canonical WZW/GKO premises stated above, the coordinate $(26,8,312)$ is absolutely mathematically rigid. The repository driver still exports \texttt{results/robustness\_audit.csv}, a disclosed $3\times 3$ sweep over $\kappa_{D_5}$ ($\pm1\%$) and $k_{\ell}$ ($26\pm1$), but that file is a numerical stress test of nearby off-shell cells rather than the source of the branch identity itself.

The present framework should be distinguished from the standard modular-flavor program in which the complex torus modulus $\tau$ is promoted to a dynamical or scanning field and the Yukawa couplings are modular forms $Y_r^{(w)}(\tau)$ transforming under finite modular groups. In that more familiar construction, flavor data are organized by modular weights, residual symmetries, and the vacuum expectation value of $\tau$ rather than by fixed rational-conformal-field-theory data~\cite{feruglio2017,kobayashi2018,criado2018,kobayashi2017}.

The organizing principle here is different. The primary input is not a continuously varying modulus but a fixed RCFT point with discrete affine data $(k_{\ell},k_q,K)=(26,8,312)$ and corresponding modular matrices $S$ and $T$. In particle-physics language, the present construction uses fixed modular data: the Yukawa textures are read from overlap data of a modular tensor category, while the electroweak sector supplies the ordinary Higgs matching into four-dimensional masses. Standard modular-flavor models are typically dynamical-$\tau$ effective theories; the present formalism instead treats $S$ and $T$ themselves as the ultraviolet flavor data. It is also distinct from local F-theory flavor models, where hierarchical Yukawas arise from geometric wavefunction overlap near enhancement loci rather than from fixed modular-tensor data~\cite{heckman2010,bouchard2010}.

This distinction matters statistically. The pull table below is not obtained by scanning a modulus, a flavon alignment, or a family of residual phases. Instead it compares the renormalization-group transport of a single fixed modular kernel against benchmark experimental intervals. Throughout this paper the leptonic comparison is made to NuFIT~5.3 normal-ordering intervals with SK atmospheric data, while the quark comparison uses the PDG~2024 CKM benchmark averages. In that precise sense the benchmark tests a discrete structural hypothesis rather than a search through a continuous flavor-parameter space.

\subsection{Geometric Setup: The Celestial Boundary at Null Infinity}

The geometric language used below is intentionally topological rather than cosmological. The two-dimensional boundary is interpreted as the celestial boundary at null infinity, and the $CFT_2$ modular data are used as a representation of celestial amplitudes for four-dimensional flavor scattering. In that sense the $AdS_3/CFT_2$ formalism is employed as an auxiliary topological generator for the flavor kernels, not as a claim that the physical universe is literally $AdS_3$.

Within this celestial-holography bookkeeping, the boundary RCFT is a bookkeeping layer for braiding, framing, and support data. The three-dimensional Chern--Simons picture provides the topological transport kernel, while physical spacetime remains four-dimensional and is reconstructed only after the modular $S$- and $T$-matrices are mapped through the electroweak interface into bulk Yukawa structures and RG-transported observables. The utility of the celestial description is therefore methodological: it packages the modular data into a fixed kernel without committing cosmology itself to a literal $AdS_3$ background.

\section{Holographic Scale Consistency Audit}

The modular mixing kernels derived below do not depend on any cosmological calibration.
The load-bearing statement is instead that the anomaly-free one-copy boundary
cell closes only if the same modular-$T$ bookkeeping is accompanied by a
strictly positive holographic vacuum loading $\Lambda_{\rm holo}>0$. Once that
positive surface tension fixes the maximum holographic register size, the
lightest-neutrino scale is no longer a speculative auxiliary bridge but a
theory-fixed consequence of the minimal $SO(10)_{312}$ construction.

\subsection{The Topological Origin of Vacuum Loading}

Within the minimal canonical embedding, $\Delta_{\rm fr}=0$ is not merely a
visible-sector framing-closure filter. It is the branch condition that keeps the
completed boundary block torsion-free and forces the parent completion residue
to appear with positive sign. In the four-dimensional effective description the
same positive residue is read as a positive holographic surface tension,
\begin{equation}
 N = \frac{3\pi}{L_P^2\Lambda_{\rm holo}},
 \qquad
 \Lambda_{\rm holo}>0.
\end{equation}
The vacuum loading is therefore part of the anomaly-free branch definition
rather than an optional late-time add-on. If $\Lambda_{\rm holo}\to0^+$, then
$N\to\infty$ and the defect coordinate
$m_0=\kappa_{D_5}M_PN^{-1/4}$ collapses to zero, so the one-copy dictionary
loses the protected neutrino floor required by the benchmark. The PMNS and CKM
kernels are fixed by modular data before any absolute mass input is supplied,
but their transported benchmark realization is quoted only on this
finite-capacity branch: the same cutoff $N$ that keeps the boundary register
finite also fixes the absolute mass coordinate carried through the RG
transport. In that precise sense the neutrino floor is protected from vanishing
by the holographic surface tension $\Lambda_{\rm holo}$, while the PMNS/CKM
transport remains the branch-fixed angular and phase image of the same
anomaly-free closure.

\subsection{Geometric Origin of Neutrino Masses}

\paragraph{CKN derivation and the $D_5$ simplex residue.}
The benchmark mass relation is the species-level saturation of the one-copy boundary cell selected by the anomaly-free branch. The horizon bit count $N$ enters this subsection as an \emph{Observational Boundary Condition}, not as a fit parameter. Just as General Relativity requires the mass of a star as an external input before it predicts an orbit, SHBT requires the universe's capacity $N$ before it converts the rigid branch geometry into an absolute matter scale. The horizon bit count then fixes the screen radius by
\begin{equation}
 R_H^2=L_P^2N,
 \qquad
 R_H=L_P\sqrt{N},
\end{equation}
so the saturated boundary information density scales as
\begin{equation}
 \rho_{\rm info}^{\rm sat}
 \sim
 \frac{N}{R_H^4}
 =
 \frac{1}{L_P^4N}
 =
 \frac{M_P^4}{N}.
\end{equation}
In the CKN reading of the four-dimensional effective theory~\cite{cohen1999}, the same cell saturates the UV/IR density bound,
\begin{equation}
 \rho_{\rm CKN}^{\rm sat}
 \sim
 M_P^2R_H^{-2}
 =
 \frac{M_P^4}{N}.
\end{equation}
The $SO(10)$ parent lattice does not pass this density to the visible sector with unit efficiency. The retained fraction is fixed by the $D_5$ weight simplex of the spinor parent,
\begin{equation}
 \kappa_{D_5}
 =
 \sqrt{\frac{\dim(\mathbf{16})}{\operatorname{rank}(SO(10))}}
 \sqrt{\frac{A_{\Delta(D_5)}}{A_{\rm reg}(R_{\Delta})}\,\eta_{\rm spin}},
 \qquad
 \eta_{\rm spin}\equiv 1-\frac{\beta^2+1/2}{c\bigl(SO(10)_{312}\bigr)}.
\end{equation}
The factor $\kappa_{D_5}$ is therefore not an empirical fudge factor; it is the hyperarea residue of the $D_5$ weight simplex dressed by the spinorial-retention factor of the parent $SO(10)_{312}$ branch. Numerically,
\begin{equation}
 \frac{A_{\Delta(D_5)}}{A_{\rm reg}(R_{\Delta})}\approx0.33265,
 \qquad
 \eta_{\rm spin}\approx0.91844,
 \qquad
 \kappa_{D_5}=0.988771051266\approx\simplexPrefactor.
\end{equation}
In the present manuscript $\kappa_{D_5}$ is therefore kept at its exact branch value, so the mass relation is treated as a fixed finite-capacity consequence of the anomaly-free branch rather than as a source of an auxiliary theory band.
The retained density on the visible genus-$0$ support cell is then
\begin{equation}
 \rho_{\rm def}^{(0)}
 \equiv
 \kappa_{D_5}^4\rho_{\rm CKN}^{\rm sat}
 =
 \kappa_{D_5}^4\frac{M_P^4}{N},
\end{equation}
so the corresponding defect mass is fixed by the fourth root,
\begin{equation}
 m_0
 =
 \left(\rho_{\rm def}^{(0)}\right)^{1/4}
 =
 \kappa_{D_5}M_PN^{-1/4}.
 \label{eq:geometric-neutrino-mass-bridge}
\end{equation}
This is the manuscript's scale-setting theorem in benchmark-derived form: the quarter-power comes from converting a saturated four-dimensional density into a defect mass, while the $O(1)$ coefficient is the closed-form $D_5$ simplex residue. In the main text this relation is treated as the unique finite-capacity mass coordinate of the anomaly-free branch, not as an auxiliary normalization that can be floated independently of the CKM/PMNS flavor kernels. The ultraviolet motivations---CKN scaling, holographic bit counting, and the geometric-mean length relation---are collected separately in Appendix~\ref{app:n-quarter} only to document provenance.

\paragraph{Unity of Scale Identity.}
The same finite-capacity mass theorem can be rewritten as a curvature identity once the positive holographic surface tension fixes the finite register,
\begin{equation}
 N=\frac{3\pi}{L_P^2\Lambda_{\rm holo}},
 \qquad
 G_N=L_P^2=M_P^{-2}.
\end{equation}
Substituting this relation into Eq.~\eqref{eq:geometric-neutrino-mass-bridge} gives the synthesis equation
\begin{equation}
\label{eq:unity-of-scale-identity-main}
 \Lambda_{\rm holo}
 =
 \frac{3\pi}{\kappa_{D_5}^4}
 G_N m_\nu^4.
\end{equation}
Equation~\eqref{eq:unity-of-scale-identity-main} is the rigorous Unity of Scale Identity of the $(26,8,312)$ lattice. If one attempted to set $m_\nu=0$, the right-hand side would force $\Lambda_{\rm holo}=0$; then $N=3\pi/(L_P^2\Lambda_{\rm holo})\to\infty$, which violates the finite-capacity premise of the anomaly-free one-copy branch. Thus the late-time holographic surface tension, Newton coupling, and neutrino mass floor are three unit conventions for the same anomaly-free branch datum rather than independent phenomenological inputs.

\paragraph{Why the lowest state saturates the bound.}
Within the Minimal One-Copy Dictionary, the lightest neutrino occupies the unique genus-$0$ support slot and the higher states are generated by the rigid ladder $m_g=m_0e^{\beta g}$. If the genus-$0$ mass were chosen below \eqref{eq:geometric-neutrino-mass-bridge}, the boundary cell would under-fill the CKN-saturating defect density and the one-copy support ladder would no longer saturate the fixed branch bit budget. If it were chosen above \eqref{eq:geometric-neutrino-mass-bridge}, the same genus-$0$ cell would already overfill the available support density, so the heavier $g=1,2$ descendants would exceed the one-copy allocation before RG transport is even applied. In this precise benchmark sense the lowest four-dimensional mass state must sit at the saturation point: equality in \eqref{eq:geometric-neutrino-mass-bridge} is the condition that keeps the defect ladder, the holographic bit budget, and the anomaly-free one-copy dictionary mutually compatible. Topological well-definedness here therefore means closure of the minimal bulk/boundary support map on the same saturated branch that already enforces $\Delta_{\rm fr}=0$.

\paragraph{The Finiteness Requirement.}
Let the maximum holographic capacity of the selected one-copy cell be
\begin{equation}
 C_{\max}\equiv N\approx 3.42\times10^{122}\ \text{bits}.
\end{equation}
Because $M_P>0$ and the anomaly-free branch fixes $\kappa_{D_5}=\simplexPrefactor>0$, Eq.~\eqref{eq:geometric-neutrino-mass-bridge} immediately gives
\begin{equation}
 m_\nu=\kappa_{D_5}M_PC_{\max}^{-1/4}>0.
\end{equation}
Thus $m_\nu=0$ would require either $N=C_{\max}\to\infty$ or $\kappa_{D_5}=0$. The first option contradicts the finite holographic register inferred from $\Lambda_{\rm obs}$, while the second contradicts the exact nonzero $D_5$ simplex residue of the anomaly-free branch. In the minimal $SO(10)_{312}$ construction, nonzero neutrino mass is therefore a mandatory consequence of finite holographic capacity.

\paragraph{Bit-Cost of Flavor.}
Define the Modular Gap of the visible transfer by
\begin{equation}
 \Delta_{\rm gap}\equiv 1-\kappa_{D_5}\approx 0.0112.
\end{equation}
This positive remainder is the irreducible error-correction overhead required for boundary unitarity. In other words, $\Delta_{\rm gap}$ is the bit-cost of flavor: the visible one-copy sector cannot inherit the parent $SO(10)$ capacity with unit efficiency, but must reserve this fraction to maintain a torsion-free, unitary boundary dictionary. The benchmark mass relation is therefore theory-fixed by finite-capacity transfer rather than introduced as an external normalization parameter.

\paragraph{Theory-Fixed Benchmark Values.}
We therefore treat the Planck-anchored bit budget $N=C_{\max}$ as the fixed \emph{Observational Boundary Condition} supplied by the observed cosmological constant $\Lambda_{\rm obs}$, not as a floating phenomenological parameter. Once the anomaly-free branch and the finite capacity are fixed, the runtime benchmark gives the RG-transported lightest mass
\begin{equation}
 m_0(M_Z)\approx 2.92\,\mathrm{meV},
\end{equation}
and the corresponding transported normal-ordering spectrum satisfies
\begin{equation}
 \sum m_\nu\approx 0.062\,\mathrm{eV}.
\end{equation}
These are theory-fixed consequences of the minimal finite-capacity $SO(10)_{312}$ SHBT construction. Varying them would move one away from the anomaly-free matching dictionary rather than explore an internal continuous parameter.
Accordingly, Table~\ref{tab:global-flavor-fit} lists the rows $m_1=m_0(M_Z)$ and $\sum m_\nu$ as \emph{Theory-Fixed} entries rather than adjustable parameters.

This scoping matters conceptually. The absolute mass scale is treated as a benchmark mass relation of the $(26,8,312)$ vacuum. In the saturated-information reading adopted here, the mass-scale theorem is tested through loss of torsion-free closure rather than treated as an external normalization.

For orientation within that theory-fixed benchmark, once $N$ is fixed by the observed cosmological constant the light defect scale lies in the meV regime,
\begin{equation}
 m_0 \sim R_{01}^{\rm par/vis}\,3.9\times10^{-3}\,\mathrm{eV},
\end{equation}
so the benchmark value lies naturally in the observed neutrino window. In the runtime benchmark used below, the numerical implementation of these analytic identities gives the theory-fixed value $m_0(M_Z)\approx2.92\,\mathrm{meV}$.

The same finite-capacity theorem also supplies the ultraviolet Majorana benchmark,
\begin{equation}
 M_R^{\rm holo}(g=2) \sim M_PN^{-1/2},
\end{equation}
which is the pre-condensation genus-$2$ boundary scale carried by the anomaly-free branch. In the present manuscript this raw UV datum is not promoted to a free seesaw anchor. Instead the effective heavy coordinate is denoted $M_N$ and interpreted as the branch-fixed \emph{Modular Restoration Scale}: the exact UV coordinate at which the missing $B-L$ conjugate bits are restored so that the completed boundary remains neutral. Equivalently, $M_N$ is the structural suppression factor of the $SO(10)/SU(3)$ coset tunneling barrier, fixed by the descendant branching multiplicities
\begin{equation}
\label{eq:mass-hierarchy-modular-restoration}
 I_L=6,
 \qquad
 I_Q=13,
 \qquad
 M_N
 =
 M_{\rm GUT}\exp\!\left[-I_L\Pi_{\rm rank}-I_Q\delta\Pi_{126}\right].
\end{equation}
Here $\Pi_{\rm rank}$ is the universal rank-gap barrier and $\delta\Pi_{126}\equiv\ln\Xi_{12}$ is the branch-fixed $\mathbf{126}_H$ correction derived explicitly below from the same visible/parent descendant data. Because $I_L$ and $I_Q$ are fixed exactly by the anomaly-free Diophantine branch, $M_N$ is not scanned as an auxiliary coordinate. The later threshold identity $M_{\rm match}=M_Ne^{-\beta^2}$ is therefore a descendant consequence of modular restoration rather than a separate seesaw input. The role of the next sections is then structural rather than cosmological: given this benchmark mass coordinate and its rigid restoration scale, do the leptonic $k_{\ell}=26$ and quark $k_q=8$ kernels realize a viable common parent without hidden continuous retuning?

\section{The Quantized Mixing Kernel}\label{sec:quantized-mixing-kernel}

\subsection{Bulk-Boundary Mapping of the $S$-Matrix}

The physical PMNS matrix is not introduced as an unrelated four-dimensional object. In the present framework it is the bulk manifestation of the same braiding data that define the three-dimensional Chern--Simons description of the defect sector. The relevant bridge is the Friedan--Shenker description of the RCFT conformal-block bundle together with Witten's identification of Chern--Simons quantization with the same conformal-block state space~\cite{friedan1987,witten1989}. Schematically,
\begin{equation}
 \mathcal H^{\rm CS}_{\Sigma,n}(SU(2)_k)
 \cong
 \mathcal V^{\rm RCFT}_{\Sigma,n}(SU(2)_k),
 \qquad
 \mathcal B^{\rm bulk}_{ij} \leftrightarrow S^{\rm boundary}_{ij},
\end{equation}
where $\mathcal V$ denotes the RCFT conformal blocks and $\mathcal H^{\rm CS}$ the bulk Chern--Simons Hilbert space on the same punctured surface. The modular $S$-matrix therefore computes the same topological exchange data that appear in the bulk as braiding amplitudes of the defect lines.

The Anomaly-Free Boundary Mapping used in this manuscript is correspondingly concrete:
\begin{equation}
 \begin{aligned}
 \chi_{g_\alpha}
 &\leftrightarrow
 \text{boundary support character},\\
 \nu_i
 &\leftrightarrow
 \text{bulk topological flavor state},\\
 \mathcal T_{\alpha i}=\langle\chi_{g_\alpha}|\nu_i\rangle
 &\leftrightarrow
 \text{bulk/boundary transfer map}.
 \end{aligned}
\end{equation}
In this Anomaly-Free Boundary Mapping the modular kernel fixes the unitary flavor geometry, while the scalar/Yukawa sector fixes only the singular spectrum that sets the mass scale. The later stability audits are designed precisely to test this separation.

\paragraph{Physical origin of the overlap basis.}
Physical mixing arises because the three neutrino generations are mapped to the three lowest-lying primary fields $g\in\{0,1,2\}$ of the $SU(2)_{26}$ model. The modular $S$-matrix therefore provides the natural overlap basis between those boundary primaries and the bulk flavor states, while the conical holonomy of the boundary geometry induces the specific $2$--$3$ sector rotation observed in the PMNS texture. The next two subsections simply make that dictionary explicit: first extract the overlap kernel, then rotate it by the holonomy fixed on the anomaly-free branch.

The four-dimensional leptonic mixing matrix is the phenomenological image of this topological data after Yukawa matching and RG transport. More precisely, the ultraviolet PMNS kernel is not identified entry-by-entry with the raw modular $S$-matrix; it is the polar-unitary image of the restricted $S$-block after the electroweak interface and the RT holonomy act on the bulk basis. In schematic form,
\begin{equation}
 U_{\rm PMNS}^{\rm UV}
 \sim
 \mathcal R_{\rm EW}\!\left[\operatorname{polar}\left(S^{(26)}_{q_\alpha g_i}\right)\right],
\end{equation}
with $\mathcal R_{\rm EW}$ denoting the Higgs/interface map together with the conical-holonomy rotation. This is the sense in which the observed four-dimensional PMNS matrix is the physical manifestation of the two-dimensional modular data: the boundary RCFT, the three-dimensional Chern--Simons braiding sector, and the four-dimensional flavor kernel are three descriptions of the same fixed input.

\subsection{Modular overlap at $k=26$}

For $SU(2)_k$ the modular $S$-matrix is
\begin{equation}
 S_{ab}^{(k)} = \sqrt{\frac{2}{k+2}}\,\sin\!\left(\frac{(a+1)(b+1)\pi}{k+2}\right),
 \qquad a,b=0,1,\dots,k.
\end{equation}
The light genus sectors are $g\in\{0,1,2\}$, and the high-charge flavor embedding is taken to be
\begin{equation}
 (q_e,q_{\mu},q_{\tau}) = (k-4,k-3,k).
\end{equation}
At the leptonic level $k=26$ this becomes $(22,23,26)$. The bare overlap kernel is then
\begin{equation}
 K_{\alpha i}=S^{(26)}_{q_{\alpha},g_i},
\end{equation}
This terminal assignment is structural rather than ad hoc: it selects the three right-edge visible charges nearest the $SU(2)_k$ Weyl-alcove boundary while keeping the restricted one-copy block on $g_i\in\{0,1,2\}$ full rank. For $k=26$ one finds $\operatorname{rank}K=3$ (equivalently $\det K\neq0$), so the three visible flavor channels remain independent already at the boundary-kernel level. Within the displayed local scan $k_{\ell}\in[24,28]$ at fixed $(k_q,K)=(8,312)$, only $k_{\ell}=26$ also satisfies $\Delta_{\rm fr}=0$; the neighboring candidates $k_{\ell}=24,25,27,28$ retain nonzero framing defects and are discarded by the theorem's framing-closure condition. The benchmark triplet $(22,23,26)$ is therefore the unique anomaly-free rank-preserving one-copy lepton embedding carried forward to phenomenology.
with unitary representative given by its polar part,
\begin{equation}
 U_{\rm top}=K(K^{\dagger}K)^{-1/2}.
\end{equation}
Direct evaluation yields
\begin{equation}
 U_{\rm top}^{(26)}\approx
 \begin{pmatrix}
 0.8297 & -0.5378 & 0.1498 \\
 -0.4724 & -0.5333 & 0.7018 \\
 0.2975 & 0.6530 & 0.6965
 \end{pmatrix}.
\end{equation}
This matrix already captures the nontrivial leptonic hierarchy of large $1$--$2$ and $2$--$3$ mixing together with a smaller $1$--$3$ entry.

\subsection{Unitary Basis-Rotation from Conical Holonomy}

The RT rotation is not an external correction factor. Once a boundary defect acquires mass through the Higgs interface, the auxiliary $AdS_3$ wedge used in the celestial/topological model develops a localized conical deficit and the flavor frame must be parallel transported through that curved geometry. The observable $2$--$3$ rotation is therefore a genuine holonomy of the spin connection, re-expressed on the boundary as an entropy-driven basis change of the RT surface.

Writing the total quantum dimension and defect dimensions as
\begin{equation}
 \cD_k = \frac{\sqrt{k+2}}{\sqrt{2}\,\sin\!\left(\frac{\pi}{k+2}\right)},
 \qquad
 d_g = \frac{\sin\!\left(\frac{(g+1)\pi}{k+2}\right)}{\sin\!\left(\frac{\pi}{k+2}\right)},
\end{equation}
we define the entropic step between the neighboring support channels $g=1$ and $g=2$ by
\begin{equation}
 \Delta S_{21} \equiv \ln \cD_k + \ln(d_2/d_1).
\end{equation}
Through the RT relation, this jump is tied to the local area variation of the extremal surface,
\begin{equation}
 \Delta S_{21}
 =
 \frac{\Delta A_{21}}{4G_3}
 =
 \frac{\partial S}{\partial A}\,\Delta A_{21}.
\end{equation}
where $G_3$ is the auxiliary three-dimensional Newton constant of the local $AdS_3$ bookkeeping wedge. The corresponding defect tension $\mu_{21}$ produces the conical deficit
\begin{equation}
 \delta_{21}
 \equiv
 8\pi G_3\mu_{21}
 =
 \alpha_{21}
 = \frac{\Delta S_{21}}{k+2}
 = \frac{1}{k+2}\frac{\partial S}{\partial A}\,\Delta A_{21}.
\end{equation}
Parallel transport of the flavor basis around that defect is generated by the $2$--$3$ rotation generator $J_{23}$,
\begin{equation}
 \Gamma_{21}
 = \mathcal P\exp\!\oint \omega
 = \exp\!\left(-\frac{\delta_{21}}{2}J_{23}\right),
\end{equation}
In the main text the induced atmospheric rotation is not introduced as an
adjustable $2.0^{\circ}$ shift. It is the derived \emph{Quarter-Deficit Angle}
$-\alpha_{21}/4=-\delta_{21}/4$ required to preserve modular covariance of the
boundary conformal block in the RT comparison of the two boundary legs of the
auxiliary $AdS_3$ wedge.
A full loop around the conical defect carries the spin-connection holonomy
$\Gamma_{21}=\exp(-\delta_{21}J_{23}/2)$, but the RT map compares the two sides
of the auxiliary $AdS_3$ wedge rather than a full closed circuit and assigns the
resulting phase covariantly to the two boundary legs of the modular block. The
required modular-covariant basis shift is therefore the quarter-deficit angle,
\begin{equation}
 \phi_{RT}(2,1)
 = -\frac{\alpha_{21}}{4}
 = -\frac{\delta_{21}}{4}
 = -\frac{1}{4(k+2)}\frac{\partial S}{\partial A}\,\Delta A_{21}
 = -\frac{\ln \cD_k + \ln(d_2/d_1)}{4(k+2)}.
\end{equation}
The factor $1/4$ is therefore not a floating coefficient: it is the required
geometric phase for a conical holonomy in the auxiliary $AdS_3$ wedge that
preserves modular covariance of the boundary conformal block~\cite{friedan1987,witten1989}. In this sense the
quarter-deficit formula is the derived topological basis rotation, not an ad hoc correction: the atmospheric $2$--$3$ mixing angle
is the direct parallel-transport image of the entropy jump across the RT
surface in the auxiliary $AdS_3$ defect geometry.

For $k=26$,
\begin{equation}
 \cD_{26}\approx33.42,
 \qquad
 d_1\approx1.987,
 \qquad
 d_2\approx2.950,
\end{equation}
so the exact holonomy estimate is
\begin{equation}
 \phi_{RT}^{\rm hol}\approx -3.49\times10^{-2}\,\mathrm{rad}\approx -2.00^{\circ}.
\end{equation}
The quoted $-2.00^{\circ}$ atmospheric rotation is therefore just the numerical evaluation of the Quarter-Deficit Angle $-\delta_{21}/4$, not an auxiliary flavor-side adjustment.
A nearby rounded benchmark, obtained from the same formula by coarse rounding, is
\begin{equation}
 \phi_{RT}^{\rm bench}=-0.0337\,\mathrm{rad},
\end{equation}
which therefore differs from the exact holonomy value only at the few-percent level.

Using the holonomy-derived angle, the CP-even benchmark leptonic kernel is
\begin{equation}
 U_{\rm PMNS}^{(0)}=R_{23}(\phi_{RT}^{\rm hol})U_{\rm top},
 \qquad
 R_{23}(\phi)=
 \begin{pmatrix}
 1 & 0 & 0 \\
 0 & \cos\phi & \sin\phi \\
 0 & -\sin\phi & \cos\phi
 \end{pmatrix}.
\end{equation}
This gives
\begin{equation}
 U_{\rm PMNS}^{(0)}\approx
 \begin{pmatrix}
 0.8297 & -0.5378 & 0.1498 \\
 -0.4825 & -0.5557 & 0.6771 \\
 0.2809 & 0.6340 & 0.7205
 \end{pmatrix},
\end{equation}
so that
\begin{equation}
 \boxed{(\theta_{12},\theta_{13},\theta_{23})\approx(32.95^{\circ},8.62^{\circ},43.22^{\circ}).}
\end{equation}
This is the CP-even benchmark extracted from conical holonomy alone; the full complex phase structure is generated when the same holonomy is combined with the modular framing phase.

\subsection{Dirac CP Phase from Modular $S$--$T$ Interference}

The $SU(2)_{26}$ modular $S$-matrix is real in the preferred topological basis, so by itself it fixes only a real overlap kernel. Complex leptonic CP violation appears when that real kernel interferes with the framing phase carried by the modular $T$-matrix. For the low-lying primaries one has
\begin{equation}
 T_{gg} = \exp\!\left[2\pi i\left(h_g-\frac{c}{24}\right)\right],
 \qquad
 h_g = \frac{g(g+2)}{4(k+2)},
 \qquad
 c = \frac{3k}{k+2}.
\end{equation}
The physically relevant framing twist in the atmospheric sector is the relative phase
\begin{equation}
 \omega_{\rm fr} \equiv \arg\!\left(T_{22}T_{11}^{*}\right)=2\pi(h_2-h_1).
\end{equation}
At $k=26$ one has
\begin{equation}
 h_1=\frac{3}{112},
 \qquad
 h_2=\frac{1}{14},
 \qquad
 \omega_{\rm fr}=\frac{5\pi}{56}\approx0.2805\,\mathrm{rad}\approx16.07^{\circ}.
\end{equation}
The conical-holonomy rotation is therefore promoted to the complex unitary map
\begin{equation}
 R_{23}(\phi,\omega)=
 \begin{pmatrix}
 1 & 0 & 0 \\
 0 & \cos\phi & e^{-i\omega}\sin\phi \\
 0 & -e^{i\omega}\sin\phi & \cos\phi
 \end{pmatrix},
\end{equation}
and the physical leptonic kernel becomes
\begin{equation}
 U_{\rm PMNS}^{\rm phys}=R_{23}(\phi_{RT}^{\rm hol},\omega_{\rm fr})U_{\rm top}.
\end{equation}
Because $U_{\rm top}$ is real, the observable Dirac phase is the relative phase between the modular-overlap amplitude and the framing twist. The convenient rephasing-invariant interference variable is
\begin{equation}
 W_{ST}
 \equiv
 \sum_{a<b}
 \left|S_{q_e g_a}S_{q_{\mu} g_b}S_{q_e g_b}S_{q_{\mu} g_a}\right|
 e^{i(\theta_b-\theta_a)},
 \qquad
 \theta_a\equiv\arg T_{aa},
\end{equation}
for which the explicit $SU(2)_{26}$ block gives
\begin{equation}
 \arg W_{ST}=\argWSTDegrees^{\circ}.
\end{equation}
The PDG branch relevant to the measured sign of leptonic CP violation is therefore
\begin{equation}
 \boxed{\delta_{CP}^{\rm top}=\pi+\arg W_{ST}=\deltaCPTopTwoDecimal^{\circ}\approx\deltaCPTopPi\pi.}
\end{equation}
Equivalently, before the conventional PDG $\pi$-branch shift the raw $SU(2)_{26}$ interference kernel carries a positive Jarlskog invariant, so the physical branch is the unique one that preserves the orientation of the holonomy area; in PDG convention that orientation-preserving choice is represented by $\deltaCPTopTwoDecimal^{\circ}$ rather than its conjugate partner at $\argWSTDegrees^{\circ}$.
After RG transport from $M_{\rm GUT}$ to $M_Z$, one obtains
\begin{equation}
 \delta_{CP}(M_Z)=\deltaCPMz^{\circ},
\end{equation}
which lies inside the current NuFIT $1\sigma$ normal-ordering interval.

The corresponding CP-odd invariant is most transparently written as an oriented flavor-space holonomy area. Define the CP-even area factor
\begin{equation}
 \Ahol
 \equiv
 \frac18\sin2\theta_{12}\sin2\theta_{13}\sin2\theta_{23}\cos\theta_{13}.
\end{equation}
Then the Jarlskog invariant is
\begin{align}
 J[S,T]
 &= \Im\!\left[(U_{\rm PMNS}^{\rm phys})_{e1}(U_{\rm PMNS}^{\rm phys})_{\mu2}(U_{\rm PMNS}^{\rm phys})_{e2}^{*}(U_{\rm PMNS}^{\rm phys})_{\mu1}^{*}\right], \\
 J[S,T]
 &= \Ahol\sin\delta_{CP}.
\end{align}
Numerically,
\begin{equation}
 \Ahol(M_{\rm GUT})\approx\holonomyAreaUv,
 \qquad
 J_{CP}(M_{\rm GUT})\approx\jarlskogUv,
\end{equation}
and after RG transport,
\begin{equation}
 \Ahol(M_Z)\approx\holonomyAreaMz,
 \qquad
 J_{CP}(M_Z)\approx\jarlskogMz.
\end{equation}
Thus the often-quoted $\sim3.3\times10^{-2}$ quantity is the holonomy-area prefactor, while the physical Jarlskog invariant is its oriented projection along the modular framing phase. Since the $S$-kernel is real, a nonzero $J[S,T]$ requires a nontrivial framing twist $\omega_{\rm fr}\neq0,\pi$. In that sense CP violation is a topological Aharonov--Bohm effect: the defect braid sees a flat but globally nontrivial connection, and the interference of paths weighted by $S$ and $T$ generates the complex phase rather than any local CP-violating potential.

\subsection{CP Violation as a Branch-Fixed Consistency Requirement}

The quark-sector CP budget is not left as a continuous phase recalibration. Both the leptonic and quark Jarlskog invariants are interpreted as geometric area factors of the boundary map: in the leptonic sector the oriented $S$--$T$ interference area gives the branch-fixed $J_{CP}(M_Z)$, while in the quark sector the same anomaly-free cell locks the CP-odd numerator directly to the integer data of the benchmark branch,
\begin{equation}
 J_{CP}^{\rm branch}
 \equiv
 \frac{\mathcal R_{\rm GUT}}{K}
 \sqrt{1-\kappa_{D_5}^2}
 \sin\!\left(\frac{2\pi k_q}{k_{\ell}}\right)
 \approx \jarlskogTopological.
\end{equation}
The same threshold residue then projects this branch-fixed numerator onto the visible CKM branch,
\begin{equation}
 J_{CP}^{q,{\rm lock}}
 =
 \mathcal R_{\rm GUT} J_{CP}^{\rm branch}
 \approx
 \jarlskogTopologicalVisible,
\end{equation}
in agreement with the benchmark quark-sector invariant $J_{CP}^{q}(M_Z)\approx\jarlskogCkmBenchmark$. In this Anomaly-Free Boundary Mapping CP violation is the minimal torsion required to map the spinorial $D_5$ weight simplex onto the $SU(2)_{26}\times SU(3)_8$ visible sector without reopening a modular residue. No independent CKM phase parameter is introduced by hand once $(k_{\ell},k_q,K)=(26,8,312)$, $\kappa_{D_5}$, and $\mathcal R_{\rm GUT}$ are fixed.
Equivalently, the quark invariant is locked to the $D_5$ geometric residue through the factor $\sqrt{1-\kappa_{D_5}^2}$, while the leptonic invariant is the corresponding topological Aharonov--Bohm readout of the modular framing twist. The predicted quark value $J_{CP}^{q}(M_Z)\approx\jarlskogCkmBenchmark$ is therefore a branch-fixed residue of the boundary geometry, obtained without introducing any continuous CP-angle adjustment.

The Wheeler--DeWitt audit sharpens the same point. If one forces a CP-conserving projection, the locked complex support disappears and the benchmark reopens a finite modularity leak,
\begin{equation}
 \Delta_{\rm mod}^{CP=0}\approx \deltaModCpZero > 4\times10^{-3},
\end{equation}
so the minimal one-copy code no longer closes on the static kernel. In the benchmark shorthand,
\begin{equation}
 \hat H_{\partial}\ket{\Psi_{\rm phys}}=0
 \qquad\Longrightarrow\qquad
 J_{CP}\neq0.
\end{equation}
A real-only universe is therefore not merely a phenomenologically different branch; it is a bookkeeping failure of the $312$-level lattice.

\subsection{Modular T-Invariance and the Absence of Strong CP Violation}

The same framing closure that protects the leptonic and quark Jarlskog
residues also removes the QCD vacuum angle from the admissible benchmark data.
On the anomaly-free cell the visible gauge factors carry the fixed WZW central
charges
\begin{equation}
 c_{\ell}=c\!\left(SU(2)_{26}\right)=\frac{3k_{\ell}}{k_{\ell}+2}=\frac{39}{14},
 \qquad
 c_q=c\!\left(SU(3)_8\right)=\frac{8k_q}{k_q+3}=\frac{64}{11}.
\end{equation}
The repository gauge-residue audit then reads the weak mixing angle as
\begin{equation}
 \sin^2\theta_W^{\rm branch}
 \equiv
 \frac{3}{8}\frac{c_{\ell}}{c_{\ell}+c_q}
 \approx
 \gaugeStrongSinSqThetaW,
\end{equation}
so electroweak mixing is exported as a rigid ratio of visible central charges
rather than introduced as a floating unification-threshold coordinate. More
importantly, exact modular-$T$ closure forbids an additional CP-odd vacuum twist
in the color sector. A nonzero QCD angle would act as an uncancelled framing
phase on the same boundary torus, reopening the defect that the benchmark cell
eliminates by imposing $\Delta_{\rm fr}=0$. In the benchmark shorthand,
\begin{equation}
 \Delta_{\rm fr}=0
 \qquad\Longrightarrow\qquad
 \bar\theta=\strongCpThetaBar.
\end{equation}
Within the benchmark bookkeeping, the strong CP angle is not treated as an
independently scanned coordinate. The anomaly-free $312$-lattice admits only the
modular-$T$-closed branch used in the construction, and on that branch the
residual QCD vacuum angle vanishes to the quoted benchmark accuracy. In
repository language this statement is packaged as the \texttt{GaugeStrongAudit};
in manuscript language, a nonzero $\bar\theta$ would signal departure from the
benchmark branch rather than motivate an internal retuning. This is therefore
an internal consequence of the adopted anomaly-free reconstruction, not an
independent empirical proof about all possible strong-CP resolutions.

\subsection{Benchmark Result for Neutrinoless Double Beta Decay}

Figure~\ref{fig:majorana-floor} provides the visual target for the mass-ordering discussion: the highlighted point at $m_{\rm lightest}\approx2.92\,\mathrm{meV}$ and $|m_{\beta\beta}|\approx5.60\,\mathrm{meV}$ is the concrete normal-ordering benchmark point fixed by the finiteness requirement once the observed cosmological constant fixes $N=C_{\max}$ and the branch-fixed mass relation is transported to $M_Z$. In this sense the construction does not adjust $m_{\rm lightest}$; on the minimal anomaly-free branch it fixes $m_{\rm lightest}=m_0(M_Z)$ as a theory-determined consequence of finite holographic capacity, while the observed splittings $\Delta m_{21}^2$ and $\Delta m_{31}^2$ enter as low-energy boundary conditions for the transported spectrum. The same complex PMNS kernel then yields a sharp Majorana benchmark once the theory-fixed lightest mass is propagated with the normal-ordering experimental inputs,
\begin{equation}
 m_1 = m_0(M_Z),
 \qquad
 m_2 = \sqrt{m_0(M_Z)^2 + \Delta m_{21}^2},
 \qquad
 m_3 = \sqrt{m_0(M_Z)^2 + \Delta m_{31}^2},
\end{equation}
with
\begin{equation}
 \Delta m_{21}^2 = 7.42\times10^{-5}\,\mathrm{eV}^2,
 \qquad
 \Delta m_{31}^2 = 2.517\times10^{-3}\,\mathrm{eV}^2.
\end{equation}
For neutrinoless double beta decay the observable combination is
\begin{equation}
 |m_{\beta\beta}|
 =
 \lvert \sum_{i=1}^{3} U_{ei}^2 m_i \rvert.
\end{equation}
Using the RG-evolved PMNS matrix and the theory-fixed value $m_0(M_Z)=2.92\,\mathrm{meV}$, one obtains
\begin{equation}
 (m_1,m_2,m_3)\approx(2.92,9.09,50.13)\,\mathrm{meV},
 \qquad
 \sum m_\nu\approx0.062\,\mathrm{eV}.
\end{equation}
These same values appear in Table~\ref{tab:global-flavor-fit} as theory-fixed mass rows rather than adjustable quantities.
\begin{equation}
 \boxed{|m_{\beta\beta}|(M_Z)=5.60\,\mathrm{meV}.}
\end{equation}
At the ultraviolet matching point the corresponding benchmark is
\begin{equation}
 |m_{\beta\beta}|(M_{\rm GUT})=5.47\,\mathrm{meV}.
\end{equation}
\noindent\textbf{Falsifiability.} Discovery of the Inverted Hierarchy or a sum of masses significantly exceeding $0.062\,\mathrm{eV}$ would falsify the minimal $SO(10)_{312}$ SHBT construction.

This places the model at the lower edge of the normal-ordering global-fit region: not a broad compatibility statement, but a concrete benchmark target in the few-meV regime.

\begin{figure}[t]
\centering
\includegraphics[width=0.78\linewidth]{fig4_majorana_floor.png}
\caption{NuFIT~5.3 normal-ordering global-fit region at $3\sigma$: the gray band is the normal-ordering envelope obtained by scanning $(\theta_{12},\theta_{13},\Delta m^2_{21},\Delta m^2_{31})$ across the quoted NuFIT~5.3 normal-ordering $3\sigma$ ranges, while the blue marker shows the benchmark point fixed by the finite-capacity theorem, $m_{\rm lightest}\approx2.92\,\mathrm{meV}$, and the associated Majorana benchmark $|m_{\beta\beta}|\approx5.60\,\mathrm{meV}$. The highlighted point is therefore the concrete $SO(10)_{312}$ Majorana benchmark target relevant to LEGEND-1000.}
\label{fig:majorana-floor}
\end{figure}

The importance of Fig.~\ref{fig:majorana-floor} is that the quoted point is not merely ``compatible'' with the normal-ordering global-fit region. It is singled out by the fixed structural mass $m_0(M_Z)$ together with the single complex PMNS kernel, and it is displayed against the full NuFIT~5.3 $3\sigma$ normal-ordering envelope rather than an illustrative reference trajectory. The explicit geometric-ratio audit around $R_{01}^{\rm par/vis}$, together with the bit-count variation, is retained as an auxiliary robustness check, but the benchmark no longer carries a separate $\kappa_{D_5}$ theory band. In that sense $|m_{\beta\beta}|=5.60\,\mathrm{meV}$ is a primary benchmark target of the $SO(10)_{312}$ proposal: it is too large to be hidden arbitrarily deep in the normal-ordering tail, yet too small to be mistaken for the inverted-ordering quasi-degenerate band.

\subsection{Future Experimental Verification}

The phenomenological role of the Majorana benchmark target is prospective rather than retrospective. The model provides a rigid target,
\begin{equation}
 |m_{\beta\beta}|\approx5.60\,\mathrm{meV},
\end{equation}
which lies in the normal-ordering desert of current sensitivity and therefore cannot be assessed by a broad-brush comparison against existing null searches alone. Precisely for that reason it is falsifiable: if the present boundary-to-bulk map is correct, the effective Majorana mass should asymptote to the few-meV regime rather than float freely across the entire normal-ordering band.

This turns neutrinoless double beta decay into one of the clearest prospective tests of the construction. Next-generation programs such as LEGEND-1000 or nEXO do not enter here as post hoc validation of an adjustable parameter, but as experiments that could begin to probe the rigid scale selected by the benchmark. A sustained experimental push into the $\sim5$--$6\,\mathrm{meV}$ window would therefore provide a direct check of this benchmark target and of the finite-capacity mass theorem rather than a re-description of existing data or a retuning of the flavor textures.

\subsection{Boundary-to-Bulk Yukawa Interface}

The same modular data that fix the PMNS/CKM seeds become four-dimensional EFT textures only after the anomaly-free visible projector and the framing twist are applied. On the selected branch,
\begin{equation}
 SO(10)\longrightarrow SU(2)_{26}\times SU(3)_8,
 \qquad
 q_i=(22,23,26),
 \qquad
 \Lambda_a\in\{(0,0),(1,0),(0,1)\},
\end{equation}
and the renormalizable boundary--bulk bilinear is carried by the usual
\begin{equation}
 16_F\,16_F\,(10_H\oplus\overline{126}_H).
\end{equation}
The $\overline{\mathbf{126}}_H$ channel is structurally required because it supplies the symmetric $B-L=2$ Majorana insertion; the explicit rank/phase audit of the competing $16_H$ and $45_H$ candidates is deferred to Supplementary Sec.~S12.3a.

Define the normalized visible overlap blocks and the diagonal framing matrix by
\begin{equation}
 \widehat S_f
 \equiv
 \frac{\mathfrak P_{\rm vis}S_f\mathfrak P_{\rm vis}^{\dagger}}{\left|\bigl[\mathfrak P_{\rm vis}S_f\mathfrak P_{\rm vis}^{\dagger}\bigr]_{11}\right|},
 \qquad
 D_T\equiv\operatorname{diag}\!\bigl(e^{i\theta_1^{(T)}},e^{i\theta_2^{(T)}},e^{i\theta_3^{(T)}}\bigr),
 \qquad f=\ell,q.
\end{equation}
The benchmark boundary-to-bulk map is then
\begin{align}
 Y_f^{(0)}&=\widehat S_f, \qquad f=\ell,q, \notag\\
 Y_f^{(10)}(M_{\rm GUT})&=Y_f^{(0)}D_T, \notag\\
 Y_N^{(126)}(M_{\rm GUT})&=\frac12\left[Y_{\ell}^{(10)}(M_{\rm GUT})+Y_{\ell}^{(10)}(M_{\rm GUT})^T\right], \notag\\
 Y_f(M_{\rm GUT})&=H_fU_f, \qquad H_f\equiv\bigl(Y_fY_f^{\dagger}\bigr)^{1/2}, \qquad U_f=H_f^{-1}Y_f.
\end{align}
In this factorization the real modular overlap block fixes the Dirac texture weights, the diagonal $T$-matrix supplies the CP phases, and the polar unitary factors $U_{\ell}$ and $U_q$ are the ultraviolet PMNS/CKM seeds carried forward into RG transport. The same anomaly-free branch also fixes $\lambda_{10}=\lambda_{126}=1$, the bridge-fixed light-mass relation $m_0=\kappa_{D_5}M_PN^{-1/4}$, and the quark-threshold datum $\mathcal R_{\rm GUT}^{\rm bench}=8/28$, so the interface introduces no additional continuously adjustable flavor parameters in the Yukawa sector. The explicit benchmark matrices are collected in Supplementary Sec.~S12.3a.

At the level of relative mass normalization, the same interface yields the bare sector-pressure estimate
\begin{equation}
 \frac{m_q}{m_{\ell}} \propto \frac{h^{\vee}_{SU(3)}}{h^{\vee}_{SU(2)}}\,\frac{k_{\ell}}{k_q}
 = \frac{3}{2}\frac{26}{8}
 = 4.875.
\end{equation}
This overpredicts the observed down-type quark / charged-lepton ratio, but that mismatch does not reopen a continuous Yukawa direction. Instead it identifies the branch-fixed scalar-sector admissibility condition. The benchmark therefore imposes the \emph{Representational Admissibility Constraint}
\begin{equation}
 \frac{\langle\Sigma_{126}\rangle}{\langle\phi_{10}\rangle}
 \equiv
 (C_{126}^{(12)})^{-1}
 =
 \frac{h^{\vee}_{SU(2)}}{h^{\vee}_{SU(3)}}\frac{k_q}{k_{\ell}}
 =
 \frac{64}{312},
\end{equation}
\begin{equation}
 \mathcal W_{\rm scal}(r)
 \equiv
 r-(C_{126}^{(12)})^{-1}
 =
 0,
 \qquad
 r\equiv\frac{\langle\Sigma_{126}\rangle}{\langle\phi_{10}\rangle}.
\end{equation}
Its unique branch-fixed root is therefore
\begin{equation}
 r_*=\frac{64}{312}.
\end{equation}
The benchmark value $r_*=64/312$ is thus the exact inverse of the $\mathbf{126}_H$ Clebsch--Gordan coefficient for the selected branch~\cite{georgi1979,babu1993}. It is used as a leading-order scalar matching condition, not as a theorem that the topological kernel alone fixes the charged-fermion singular values. While the topological kernel provides rigid mixing eigenvectors, the ratio $64/312$ is an informed matching condition selected to align the singular values (masses) and suppress the bare model's over-prediction of the quark/lepton hierarchy. The SVD Rigidity audit of Fig.~\ref{fig:vev-stability-main}, together with Supplementary Table~\ref{supp-tab:s3-vev-protection} and Supplementary Fig.~\ref{supp-fig:s3-vev-stability}, shows that even $O(10\%)$ deformations of this disclosed setting leave the PMNS/CKM angles far inside their quoted $1\sigma$ windows. The scalar sector therefore renormalizes singular values without relaxing the topological eigenvectors, securing the flavor predictions against scalar-potential uncertainty.

\section{The Unification of Spacetime and Flavor}\label{sec:topological-baryogenesis}

The anomaly-free benchmark does not solve flavor first and append gravity or
baryogenesis afterward as detached sectors. On the selected $SO(10)_{312}$
branch the same modular data that fix the PMNS/CKM textures and the
heavy-threshold lock also close the bulk vacuum loading and the matter
asymmetry budget. In that sense spacetime, flavor, and cosmology are not
separate inference layers here: the quark/lepton CP budget, the
$\mathbf{126}_H$ Majorana threshold, the positive vacuum loading, and the
completion residue $c_{\rm dark}$ are different projections of one anomaly-free
parent partition function.

\paragraph{The Holographic Stabilizer.}
The completion residue
\begin{equation}
 c_{\rm dark}=24\,\Delta_{\rm mod}(26,8)=\cDarkCompletion
\end{equation}
is the branch's \emph{Holographic Stabilizer}; numerically,
$c_{\rm dark}\approx3.3008$. In observer-level language this completion residue
is the \emph{Topological Gravitational Ghost} (TGG): the information deficit
created when the parent $SO(10)_{312}$ block is projected onto the visible
$SU(2)_{26}\times SU(3)_8$ sector. It plays a dual role. First, it furnishes the
positive bulk vacuum loading
\begin{equation}
 \Lambda_{\rm holo}=+\frac{3\pi}{L_P^2N},
\end{equation}
so the completed branch carries a finite de Sitter register and a nonzero
geometric action. Second, the same residue stores the conjugate $B-L$ parity
bits required to maintain boundary neutrality,
\begin{equation}
 (B-L)_{\rm vis}+(B-L)_{\rm dark}=0.
\end{equation}
The gravitational field and the parity completion are therefore not independent
add-ons; they are the two bookkeeping roles of the same stabilizing residue.

\paragraph{Sakharov constraints at the topological boundary.}
Within the benchmark, Sakharov's three conditions are re-expressed as discrete
boundary identities rather than as independent phenomenological inputs.
First, baryon-number violation is carried through the same
$B-L$-charged $\overline{\mathbf{126}}_H$ channel already required in the
boundary--bulk Yukawa map. The point is stronger than simple compatibility:
the $\overline{\mathbf{126}}_H$ Majorana insertion is the same heavy threshold
whose matched value
\begin{equation}
 M_{\rm match}\equiv M_{126}^{\rm match}=M_Ne^{-\beta^2}
\end{equation}
is required to close the CKM-phase transport, with $M_N$ interpreted here as
the branch-fixed \emph{Modular Restoration Scale}. The branch therefore does not use
one threshold for lepton-number violation and another for quark CP closure; the
same $\mathbf{126}_H$ datum carries both. Second, CP violation is the locked
topological holonomy already identified in Sec.~\ref{sec:quantized-mixing-kernel}:
the branch-fixed quantity $J_{CP}^{\rm topo}\approx\jarlskogTopological$ is
nonzero, while the forced CP-symmetric projection reopens the modularity leak
$\Delta_{\rm mod}^{CP=0}\approx\deltaModCpZero$. Third, the Sakharov
out-of-equilibrium slot is replaced in SHBT by a static RG incompatibility
between the Planck-scale boundary and the Modular Restoration Scale $M_N$
rather than by adding a separate thermal sector. Concretely, $M_N$ is the
exact scale at which the missing topological parity bits are restored so that
the completed boundary remains neutral. Because this restoration scale is
uniquely fixed by the branch indices $I_L=6$ and $I_Q=13$, the construction is
therefore ``out of equilibrium'' in the static sense of a locked inter-scale
incompatibility, not in the thermal sense of a separately evolving bath. The
same positive vacuum loading defines the de Sitter Hubble scale
\begin{equation}
 H_{\Lambda}\equiv\sqrt{\Lambda_{\rm holo}/3},
\end{equation}
so the branch acquires the unique isotropic geometric friction compatible with
the Bianchi identity. In that sense the Hubble expansion driven by
$\Lambda_{\rm holo}$ is the bulk manifestation of that static RG
incompatibility at the topological boundary: the finite-capacity map converts
the parent Majorana scale into the effective four-dimensional restoration threshold
through the locked ratio $M_N/M_P$, while the same branch-fixed expansion
prevents the closure tensor from reopening a nonequilibrium source term.

\paragraph{Topological baryogenesis identity.}
Once the sphaleron conversion coefficient is fixed to the Standard-Model value
$C_{\rm sph}=28/79$, the baryon-to-photon ratio is read as the branch-fixed
identity
\begin{equation}
 \eta_B
 =
 C_{\rm sph}\,J_{CP}^{\rm topo}\left(\frac{M_N}{M_P}\right).
\end{equation}
Using the locked holonomy numerator and the derived structural threshold ratio,
\begin{equation}
 J_{CP}^{\rm topo}\approx\jarlskogTopological,
 \qquad
 \frac{M_N}{M_P}\approx1.42\times10^{-5},
\end{equation}
one obtains
\begin{equation}
 \eta_B
 =
 \frac{28}{79}\,J_{CP}^{\rm topo}\left(\frac{M_N}{M_P}\right)
 \approx
 \frac{28}{79}\,(\jarlskogTopological)\,(1.42\times10^{-5})
 \approx 6.4\times10^{-10}.
\end{equation}
The observed matter excess is therefore reproduced at the benchmark level
without opening a new continuous baryogenesis coordinate: once the anomaly-free
branch, the Holographic Stabilizer, the CP-locking holonomy, and the
$\mathbf{126}_H$ threshold are fixed, the asymmetry follows as the boundary
image of the same modular-$T$ closure that generates the positive bulk vacuum
loading. In this language the Modular Restoration Scale is the static SHBT
realization of Sakharov's out-of-equilibrium condition: it is the unique
parity-restoration scale on the anomaly-free branch, not an auxiliary
baryogenesis knob.

\section{Predictive Signatures of the Completed Boundary}\label{sec:predictive-signatures}

Once the anomaly-free branch is fixed, the completion residue
\begin{equation}
 c_{\rm dark}=24\,\Delta_{\rm mod}(26,8)=\cDarkCompletion
\end{equation}
does more than close the parent partition function. It generates a set of
macroscopic observables that are as branch-fixed as the PMNS/CKM textures
themselves. The completed boundary therefore carries direct predictive
signatures in late-time expansion, thermodynamics, and primordial aliasing.

\paragraph{Theorem I: Hubble Gradient from the boundary entropy debt.}
On the anomaly-free branch the residual modularity gap
\begin{equation}
 \Delta_{\rm mod}(26,8)=\visibleResidual
 = \frac{1197103}{8704080}
 \approx 0.13753355
\end{equation}
is the entropy debt of the completed boundary: it is the finite central-charge
defect that remains after the visible block is embedded and then paid back by
the completion sector. The late-time Hubble uplift is the macroscopic image of
that debt. Defining the Hubble gradient by
\begin{equation}
 \Delta H_0
 \equiv
 H_0^{\rm loc}-H_0^{\rm CMB},
\end{equation}
the completed branch gives the exact uplift law
\begin{equation}
 H_0^{\rm loc}
 =
 H_0^{\rm CMB}\exp\!\left(\frac{\Delta_{\rm mod}}{2}\right),
 \qquad
 \frac{\Delta H_0}{H_0^{\rm CMB}}
 =
 \exp\!\left(\frac{\Delta_{\rm mod}}{2}\right)-1.
\end{equation}
Using the Planck anchor $H_0^{\rm CMB}=67.4\,\mathrm{km\,s^{-1}\,Mpc^{-1}}$
and the benchmark entropy debt above, one finds
\begin{equation}
 \exp\!\left(\frac{\Delta_{\rm mod}}{2}\right)=1.071186,
 \qquad
 H_0^{\rm loc}\approx 72.2\,\mathrm{km\,s^{-1}\,Mpc^{-1}},
 \qquad
 \Delta H_0\approx 4.80\,\mathrm{km\,s^{-1}\,Mpc^{-1}}.
\end{equation}
In this formulation the Hubble ``tension'' is not an external late-time fluid
anomaly; it is the Hubble gradient generated by the entropy debt of the
completed boundary.

\paragraph{Theorem II: Modular Scale Indicator ($T_{\rm mod}$).}
The same completed boundary carries a residual modular background that can be
reported in temperature units without being interpreted as a physical Hawking
temperature. Let
\begin{equation}
 \beta
 \equiv
 \frac12\ln\mathcal D_{26},
 \qquad
 T_{\rm mod}
 \equiv
 \bigl[m_0(M_Z)e^{-\beta}\bigr]\times 11604.5\,\mathrm{K/eV},
\end{equation}
where $m_0(M_Z)\approx2.92\,\mathrm{meV}$ is the RG-transported genus-$0$
defect mass on the completed branch. The factor $e^{-\beta}$ is the
branch-fixed \emph{Modular-$T$ hierarchy scale}: it is the same modular
suppression that carries the ultraviolet defect hierarchy into the infrared
thermal floor when the completed branch is expressed in temperature units.
Then
\begin{equation}
 T_{\rm mod}
 \approx 5.86\,\mathrm{K}.
\end{equation}
We identify this $5.86\,\mathrm{K}$ residue as the \emph{Modular Scale
Indicator} of the completed boundary. It is a dimensional trace of the
boundary anomaly carried by the same completion sector
$c_{\rm dark}=\cDarkCompletion$ that supplies the positive vacuum loading. In
the completed-boundary language this $5.9\,\mathrm{K}$ scale is therefore not
a physical Hawking temperature and not a reheating-bath prediction; it is the
temperature-unit image of the branch-fixed modular hierarchy encoded by
$e^{-\beta}$.

\paragraph{Theorem III: Equilateral aliasing error of the $D_5$ packing.}
The $D_5$ weight-simplex residue also fixes a non-Gaussian floor. Writing
\begin{equation}
 \kappa_{D_5}
 =
 \sqrt{\frac{16}{5}
 \frac{A_{\Delta(D_5)}}{A_{\rm reg}(R_{\Delta})}
 \eta_{\rm spin}}
 = \simplexPrefactor
 \approx 0.98877105,
\end{equation}
the irreducible equilateral aliasing error of the $D_5$ packing is
\begin{equation}
 f_{NL}^{\rm equil}
 \equiv
 1-\kappa_{D_5}
 \approx 0.01122895
 \approx 0.01123.
\end{equation}
This is the equilateral aliasing error of the completed boundary: the finite
packing defect that remains when the spinorial $D_5$ simplex is projected onto
the equilateral bulk template. In the benchmark bookkeeping the same number is
the geometric packing friction; observationally it is the equilateral
non-Gaussian floor of the branch.

Taken together, Theorems~I--III show that the completed boundary leaves
macroscopic signatures that are as rigidly branch-fixed as the flavor sector
itself. The same integers $(26,8,312)$ that determine the PMNS/CKM textures and
the finite-capacity neutrino mass scale also determine the Hubble gradient, the
modular scale indicator $T_{\rm mod}$, and the equilateral aliasing floor of the
completed branch.

\subsection{Branch-Fixed Gauge Density and RG Evolution}\label{sec:branch-fixed-gauge-rg}

With the boundary-to-bulk Yukawa interface fixed in Sec.~\ref{sec:quantized-mixing-kernel}, the remaining question is transport rather than texture: how the anomaly-free ultraviolet datum runs to $M_Z$, and which parts of that comparison depend on threshold bookkeeping rather than on the rigid unitary map.

The surface-tension value $\alpha^{-1}_{\rm surf}=15K/(k_{\ell}+k_q)=\alphaSurfBenchmarkExact\approx\alphaSurfBenchmarkRounded$ is treated here as the benchmark ultraviolet gauge-density datum of the anomaly-free branch. This discrete surface value is the primary gauge-side output of the benchmark and should be read as the ultraviolet fixed point of the selected branch rather than as the low-energy fine-structure constant itself. The quoted surface value is therefore not an infrared calibration target but the branch-fixed gauge coordinate determined by the same anomaly-free condition $\Delta_{\rm fr}=0$ that isolates $\benchmarkVisibleBranch$.

\paragraph{Current-Algebra Neutrality and the Standard Model Generation Weight.}
The prefactor $15$ is the branch-fixed matter weight $G_{\rm SM}$ determined by \emph{Current-Algebra Neutrality}, not an adjustable normalization. A single $SO(10)$ spinor generation $16_F$ contains the fifteen Standard-Model-charged Weyl channels together with one Standard-Model singlet right-handed-neutrino direction $\nu^c$. That singlet must be excluded from the visible matter weight because it carries zero gauge flux and zero conformal weight in the visible sectors, contributing zero to boundary-to-bulk transport~\cite{gko1985,friedan1987,witten1989}. Projecting the parent current-current correlator onto the visible sector therefore leaves precisely
\begin{equation}
 G_{\rm SM}=16-1=15
\end{equation}
visible current slots. Equivalently, $G_{\rm SM}=15$ is the one-generation normalization of the visible current-current correlator inherited from the $SO(10)_{312}$ parent theory, hence a Topological Coordinate of the anomaly-free branch rather than a phenomenological gauge normalization. In that sense the surface datum
\begin{equation}
 \alpha^{-1}_{\rm surf}=G_{\rm SM}\frac{K}{k_{\ell}+k_q}
\end{equation}
is structurally fixed once the anomaly-free branch is chosen. The infrared quantity $\alpha^{-1}_{\rm em}(M_Z)$ is not used as a primary benchmark observable; it is reported only as a consistency check on whether one-loop running and heavy-threshold matching transport the branch-fixed ultraviolet datum into the observed infrared window.

\paragraph{Kaluza--Klein Dilution Framework.}
The gauge-side audit is not presented here as a Minimal Desert extrapolation. Instead the infrared comparison is closed by the branch-fixed Heavy Higgs Dilution / Kaluza--Klein dilution correction carried by the same visible branching data that already determine the flavor-sector threshold structure. Writing
\begin{equation}
 \mathcal I_Q\equiv \frac{K}{3k_q}=13,
\end{equation}
the HHD shift is
\begin{equation}
 \Delta \alpha_{\rm em}^{-1}
 =
 - \frac{3 \mathcal I_Q}{2\pi}\ln\!\left(\frac{M_{\rm GUT}}{M_N}\right).
 \tag{HHD-1}
 \label{eq:hhd-1}
\end{equation}
This resolution introduces zero new degrees of freedom. The scale $M_N$ entering Eq.~\eqref{eq:hhd-1} is not scanned independently on the gauge side; it is already fixed by the same flavor-sector $\mathbf{126}_H$ threshold bookkeeping that determines the Modular Restoration Scale and the CKM phase closure. In the benchmark audit this gives the rigid additive correction
\begin{equation}
 \Delta \alpha_{\rm em}^{-1}\big|_{\rm HHD} = -29.462,
\end{equation}
so the raw one-loop transport value $\alpha^{-1}_{\rm em}(M_Z)=156.917$ is shifted to
\begin{equation}
 \alpha^{-1}_{\rm em}(M_Z)\big|_{\rm HHD}=127.455,
\end{equation}
with residual pull $-0.06\sigma$ relative to the electroweak input. In publication language, the gauge row is therefore a branch-derived HHD-corrected closure check rather than an unresolved Minimal Desert residual. % pub/tn.py: GaugeRenormalizationAudit.apply_heavy_higgs_dilution_correction, derive_rhn_threshold_data

\paragraph{Gauge-sector interpretation.}
The raw one-loop overshoot is retained only as the pre-HHD transport datum. Once the Kaluza--Klein dilution correction is included, the benchmark gauge row lands on the electroweak target within the quoted matching tolerance without reopening any independent threshold adjustment. Because Eq.~\eqref{eq:hhd-1} depends only on the topological branching index $\mathcal I_Q$ and the already-fixed Modular Restoration Scale $M_N$, the gauge repair is a derived corollary of the same anomaly-free branch rather than a phenomenological patch. The gauge-sector closure is therefore achieved by branch-fixed bookkeeping, not by introducing a new continuous gauge parameter.

\paragraph{Orthogonality Statement.}
Because the modular flavor kernels are fixed by the topological anomaly, the predicted mixing angles remain numerically orthogonal to the gauge-coupling closure. The same supplementary audit shows that applying the gauge-side HHD / split-threshold repair leaves the flavor angles stable at the $<0.05\sigma$ level, so the gauge normalization is healed without perturbing the PMNS/CKM eigenvectors or reopening the branch-fixed flavor benchmark.

\paragraph{Gauge Sector Disclosure.}
\textbf{Zero-new-parameter closure.} In the HHD audit the infrared gauge row is no longer presented as unresolved. The Kaluza--Klein dilution correction closes the $\alpha^{-1}_{\rm em}(M_Z)$ comparison at the level of the quoted matching tolerance, and it does so without adding a new degree of freedom because the required scale $M_N$ is already fixed as the flavor-sector Modular Restoration Scale. The gauge row remains excluded from the eight-observable flavor tally $\chi^2_{\rm pred}$ for bookkeeping clarity, but it is now recorded as an HHD-corrected gauge-side consistency check rather than as an outstanding mismatch.

The same supplementary audit records the corresponding heavy-scale stability sweep. Supplementary Table~\ref{supp-tab:s11-heavy-scale-sensitivity} shows that varying the disclosed heavy scales $M_{10}$ and $M_{\rm GUT}$ across a full decade leaves the PMNS angles, CKM angles, and the healed apex $\gamma$ inside the quoted $1\sigma$ envelope while shifting only normalization-sensitive channels such as the gauge row and the absolute mass scale. This is the load-bearing main-text point: heavy thresholds move normalization data, whereas the polar-unitary flavor outputs remain fixed by the anomaly-free $(26,8,312)$ branch.

\section{Algebraic Rigidity of the Parent Level}

With the information bridge and Higgs interface fixed, the parent-level problem becomes discrete rather than continuous. The question is no longer which new scale to add, but which integer parent level $K$ can carry the same anomaly load into both the $SU(2)_{26}$ leptonic branch and the $SU(3)_8$ quark branch.

\subsection{Visible branches}

The leptonic ladder is generated by
\begin{equation}
 m_g = m_0e^{\beta g},
 \qquad
 \beta(k)=\frac12\ln\cD_k.
\end{equation}
Structural consistency in the one-copy sector requires the same level to support a gapless three-step genus ladder, a well-defined polar modular kernel, and a perturbative RT correction. In this restricted sense the leptonic branch is selected at
\begin{equation}
 \boxed{k_{\ell}=26.}
\end{equation}
This boxed value should be read as the numerical shadow of the branch selected by the torsion-free Derived Uniqueness Theorem rather than as a benchmark-consistency metric: $k_{\ell}=26$ is the unique displayed integer for which the imposed framing condition closes exactly, so the effective boundary-to-bulk dictionary remains compatible with a Levi--Civita connection.

\paragraph{CKM tunneling-amplitude picture.}
The quark sector is organized as a tunneling problem on the visible flavor
manifold. The low-weight block of $SU(3)_{k_q}$ supplies the source flavor
frame,
\begin{equation}
 S^{(3,k_q)}_{\Lambda,\Lambda'}
 = \frac{i^3}{\sqrt3(k_q+3)}
 \sum_{w\in W(SU(3))}\epsilon(w)
 \exp\!\left[-\frac{2\pi i}{k_q+3}(w(\Lambda+\rho),\Lambda'+\rho)\right].
\end{equation}
while the $SO(10)/SU(3)$ coset surrogate supplies the resistive barrier through
which that frame must be aligned. Because the Weyl chamber is two-dimensional,
the raw low-weight block is already close to diagonal. The relevant suppression
is therefore not an arbitrary sequence of matrix manipulations but the branch-
fixed logarithmic resistance to flavor-frame alignment through the coset space.
We package that suppression in the rank-deficit branching-pressure variable,
\begin{equation}
 \Pi(K,k_q)
 \equiv
 \sqrt{\frac{|\rho_{SO(10)}|^2}{|\rho_{SU(3)}|^2}}
 \left(\frac{K+h^{\vee}_{SO(10)}}{k_q+h^{\vee}_{SU(3)}}\right)^{-1/2},
 \label{eq:rank-deficit-pressure}
\end{equation}
Here the manuscript term ``pressure'' is shorthand for this logarithmic
suppression prescription: $\Pi(K,k_q)$, $\Pi_{\rm vac}$, and the channel
quantities below are fixed exponents extracted from modular data, not
additional thermodynamic variables to be tuned.
Here $\rho_g$ is the Weyl vector of the parent or visible Lie algebra. For simply-laced groups the strange formula gives
\begin{equation}
 |\rho_g|^2=\frac{h_g^{\vee}\dim g}{12},
 \qquad
 |\rho_{SO(10)}|^2=30,
 \qquad
 |\rho_{SU(3)}|^2=2.
\end{equation}
For the present groups,
\begin{equation}
 h^{\vee}_{SO(10)}=8,
 \qquad
 h^{\vee}_{SU(3)}=3,
 \qquad
 \Delta r \equiv r_{SO(10)}-r_{SU(3)} = 5-2 = 3.
\end{equation}
The three missing Cartan directions therefore generate the physical resistance
to flavor-frame alignment through the $SO(10)/SU(3)$ coset space. Each
off-diagonal excursion is a tunneling event across that barrier, so every
rank-gap step contributes one factor of $\Pi$, while the corresponding
observable suppression is encoded in a logarithmic vacuum exponent determined
by the low-weight vacuum entry,
\begin{equation}
 \Pi_{\rm vac}
 = -\frac{\Delta r}{8}\ln\bigl|S_{00}^{(3,k_q)}\bigr|.
\end{equation}
For the $SO(10)_{312}$ parent one finds
\begin{equation}
 I_Q \equiv \frac{K}{3k_q}=13,
 \qquad
 \Pi(312,8)=\sqrt{\frac{30}{2}}\sqrt{\frac{11}{320}}\approx 0.7181,
\end{equation}
and the parent-descendant ratios define the two channel factors
\begin{equation}
 \Xi_{12}(K)
 \equiv
 \left|\frac{S_{01}^{(3,K/3)}S_{00}^{(3,k_q)}}{S_{00}^{(3,K/3)}S_{01}^{(3,k_q)}}\right|,
 \qquad
 \Xi_{23}(K)
 \equiv
 \left|\frac{S_{11}^{(3,k_q)}S_{01}^{(3,K/3)}}{S_{01}^{(3,k_q)}S_{11}^{(3,K/3)}}\right|,
\end{equation}
For the benchmark branch it is useful to resolve these descendant factors into the parent/visible low-weight ratios that the code evaluates explicitly. Writing
\begin{equation}
 V_{ab}\equiv S^{(3,8)}_{\lambda_a\lambda_b},
 \qquad
 P_{ab}\equiv S^{(3,104)}_{\lambda_a\lambda_b},
 \qquad
 \lambda_0=(0,0),\ \lambda_1=(1,0),\ \lambda_2=(0,1),
\end{equation}
the $12$-channel descendant enhancement is
\begin{equation}
 R_{01}^{\rm par/vis}
 \equiv
 \left|\frac{P_{01}V_{00}}{P_{00}V_{01}}\right|
 \approx 1.1171,
 \qquad
 C_{126}^{(12)}
 \equiv
 \frac{h^{\vee}_{SU(3)}}{h^{\vee}_{SU(2)}}\frac{k_{\ell}}{k_q}
 =
 \frac{3}{2}\frac{26}{8}
 =
 \frac{39}{8},
\end{equation}
together with the exact branching-anomaly exponent
\begin{equation}
 \alpha_{\rm br}
 \equiv
 \frac{527}{26500}.
\end{equation}
Hence the benchmark $12$-channel factor is
\begin{equation}
 \Xi_{12}^{\rm bench}
 =
 \left(R_{01}^{\rm par/vis}C_{126}^{(12)}\right)^{\alpha_{\rm br}}
 =
 \left(1.1171\times\frac{39}{8}\right)^{527/26500}
 \approx 1.0343.
\end{equation}
The $23$-channel factor is controlled by the descendant low-weight ratio
\begin{equation}
 R_{23}^{\rm par/vis}
 \equiv
 \left|\frac{V_{11}P_{01}}{V_{01}P_{11}}\right|
 \approx 0.7643,
 \qquad
 \Xi_{23}^{\rm bench}
 =
 \sqrt{R_{23}^{\rm par/vis}}
 \approx 0.8742.
\end{equation}
In this language the channel pressures $\Pi_{12}$ and $\Pi_{23}$ are the
effective logarithmic resistances used to parametrize frame alignment in the
$SO(10)/SU(3)$ coset space: $\Pi_{\rm vac}$ is the universal barrier from the
three missing Cartan directions, while $\Xi_{12}^{\rm bench}$ and
$\Xi_{23}^{\rm bench}$ are the branch-fixed descendant factors that partially
heal that barrier in the $12$ and $23$ channels. Accordingly, on the
benchmark branch the observable channel pressures are
\begin{equation}
 \Pi_{12}=\Pi_{\rm vac}-\ln\Xi_{12}^{\rm bench},
 \qquad
 \Pi_{23}=2\Pi_{\rm vac}-\ln\Xi_{23}^{\rm bench}.
\end{equation}
In compact flow form,
\begin{equation}
 (\Pi_{\rm vac},\Xi_{12},\Xi_{23})
 \longrightarrow
 (\epsilon_{12},\epsilon_{23},\epsilon_{13})
 \longrightarrow
 (|V_{us}|,|V_{cb}|,|V_{ub}|).
\end{equation}
Equivalently, the corresponding off-diagonal tunneling amplitudes obey
\begin{equation}
 \epsilon_{12}\sim e^{-\Pi_{\rm vac}}\Xi_{12}^{\rm bench},
 \qquad
 \epsilon_{23}\sim e^{-2\Pi_{\rm vac}}\Xi_{23}^{\rm bench},
 \qquad
 \epsilon_{13}\sim e^{-3\Pi_{\rm vac}}\Xi_{12}^{\rm bench}\Xi_{23}^{\rm bench},
\end{equation}
At $(k_q,K)=(8,312)$ one finds
\begin{equation}
 \Pi_{\rm vac}=1.5061,
 \qquad
 \Xi_{12}=\Xi_{12}^{\rm bench}=1.0343,
 \qquad
 \Xi_{23}=\Xi_{23}^{\rm bench}=0.8742,
\end{equation}
and therefore
\begin{equation}
 \Pi_{12}=1.4724,
 \qquad
 \Pi_{23}=3.1467.
\end{equation}
The CKM suppressions are then structural branching predictions,
\begin{equation}
 |V_{us}|\sim e^{-\Pi_{12}},
 \qquad
 |V_{cb}|\sim e^{-\Pi_{23}},
 \qquad
 |V_{ub}|\sim e^{-(\Pi_{12}+\Pi_{23})},
\end{equation}
The same parent-induced higher-state mixing factors also close the heavy
threshold equation. Writing
\begin{equation}
 M_N
 = M_{\rm GUT}\exp\!\left[-I_L\Pi_{\rm rank}-I_Q\delta\Pi_{126}\right],
 \qquad
 \delta\Pi_{126}\equiv\ln\Xi_{12},
 \label{eq:modular-restoration-scale-prediction}
\end{equation}
we obtain the \emph{Modular Restoration Proof}: $M_N$ is not a free seesaw
anchor and is not scanned independently. It is the branch-fixed
\emph{Modular Restoration Scale}, and it is the structural suppression factor
of the $SO(10)/SU(3)$ coset tunneling barrier produced by the same rank-gap
higher-state mixing and $\mathbf{126}_H$ threshold shift that determine the
quark hierarchies. The factor $I_L\Pi_{\rm rank}$ counts the six leptonic
descendant steps through the universal barrier, while
$I_Q\delta\Pi_{126}$ counts the thirteen quark descendant steps that dress the
same barrier through the $\mathbf{126}_H$ channel. For the
minimal selection-hypothesis-admissible parent one has $I_L=6$, $I_Q=13$,
$\Pi_{\rm rank}=\Pi(312,8)\approx0.7181$, and
$\delta\Pi_{126}=\ln\Xi_{12}\approx3.37\times10^{-2}$, which yields
\begin{equation}
 M_N\approx1.75\times10^{14}\,\mathrm{GeV}.
\end{equation}
The same discrete branching data therefore fix both the CKM hierarchy and the
Modular Restoration Scale. In particular, $M_N$ is the exact ultraviolet
coordinate at which the missing $B-L$ conjugate bits are restored so that the
completed branch preserves boundary neutrality. The matter hierarchy and the
cosmological branch anchor therefore close through the same unity-of-scale
synthesis,
\begin{equation}
 \label{eq:mass-hierarchy-unity-of-scale}
 \Lambda_{\rm holo}
 =
 \frac{3\pi}{\kappa_{D_5}^4}
 G_N m_{\nu}^4.
\end{equation}
The static Sakharov condition is then saturated without adding a separate
thermal dial. The three CKM magnitudes are already fixed by the
tunneling amplitudes above; the only threshold-sensitive datum left in the
quark sector is the apex angle. Its saturated-framing closure is
\begin{equation}
 \Delta_{\rm grav}(M_{126})
 \equiv
 \ln\!\left(\frac{M_N}{M_{126}}\right)-\beta^2=0
 \qquad\Longleftrightarrow\qquad
 M_{\rm match}\equiv M_{126}^{\rm match}=M_Ne^{-\beta^2}.
\end{equation}
Evaluating the one-loop $\mathbf{126}_H$ matching residue at this locked point
fixes $\mathcal R_{\rm GUT}$ on shell, so the ultraviolet CKM holonomy is
transported to the observed benchmark value of $\gamma$ without reopening a
continuous threshold scan. The threshold lock therefore rotates only the phase
budget; the tunneling amplitudes continue to control $|V_{us}|$, $|V_{cb}|$,
and $|V_{ub}|$.

The $(26,8,312)$ branch is therefore not presented here as a post-hoc calibration to
data. It is the unique displayed integer coordinate on which the scale of the
universe and the scale of matter are absolute, unyielding mirrors of one
another.

The first low-rank branch for which this parent-fixed pressure exists is
\begin{equation}
 \boxed{k_q=8.}
\end{equation}

\subsection{Modular and Framing Consistency of the Master Level}

The vacuum framing phase of a simple affine theory $g_k$ is governed by the WZW central charge
\begin{equation}
 c(g_k)=\frac{k\,\dim(g)}{k+h_g^{\vee}},
 \label{eq:central-charge-wzw}
\end{equation}
since the vacuum $T$-matrix entry is $T_{00}=\exp[-2\pi i c(g_k)/24]$. The parent level is therefore first tested against the visible branches by the central-charge identity
\begin{equation}
 c(SO(10)_K)
 \equiv
 c(SU(2)_{26}) + c(SU(3)_8)
 \pmod{24},
 \label{eq:master-level-central-charge}
\end{equation}
or equivalently
\begin{equation}
 \frac{c(SO(10)_K)}{24}
 \equiv
 \frac{c(SU(2)_{26})}{24} + \frac{c(SU(3)_8)}{24}
 \pmod 1.
 \label{eq:master-level-modularity}
\end{equation}
Using
\begin{align}
 \dim SO(10)&=45,
 & h^{\vee}_{SO(10)}&=8, \\
 \dim SU(2)&=3,
 & h^{\vee}_{SU(2)}&=2, \\
 \dim SU(3)&=8,
 & h^{\vee}_{SU(3)}&=3,
\end{align}
one finds
\begin{align}
 c\bigl(SU(2)_{26}\bigr)
 &= \frac{26\cdot3}{26+2}
 = \frac{39}{14}, \\
 c\bigl(SU(3)_8\bigr)
 &= \frac{8\cdot8}{8+3}
 = \frac{64}{11}, \\
 c\bigl(SO(10)_K\bigr)
 &= \frac{45K}{K+8}.
\end{align}
Hence the bare visible framing residual is
\begin{equation}
 \mathfrak A_{\rm vis}(K)
 \equiv
 \left[\frac{45K}{24(K+8)} - \frac{39}{14\cdot24} - \frac{64}{11\cdot24}\right]_{\!\!\mod 1}.
\end{equation}
This quantity is the visible-sector obstruction that remains before branch completion. The exact zero-energy framing state is therefore not the bare visible congruence by itself, but the completed parent vacuum obtained after the mandatory orthogonal coset sector is included.
The visible branches must also descend from integer parent channels, so the branching indices obey
\begin{equation}
 I_L = \frac{K}{26\cdot2}\in\mathbb Z,
 \qquad
 I_Q = \frac{K}{8\cdot3}\in\mathbb Z.
\end{equation}
Equivalently, the visible-parent matching is the Diophantine selection rule
\begin{equation}
 K = 2k_{\ell}I_L = 3k_q I_Q,
 \qquad
 I_L,I_Q\in\mathbb Z.
 \label{eq:diophantine-selection-rule}
\end{equation}
In the manuscript this integrality statement is packaged inside the
\emph{Derived Uniqueness Theorem}. Its operational character-level condition is
that the completed $SO(10)$ boundary block have an integer modular-$T$
spectrum. Within the specific embedding used here---the canonical
$SO(10)\supset SU(3)\times SU(2)\times U(1)$ branching, the reference coset
completion, and the one-copy visible support map---that condition reduces on
the fixed-parent slice $(K,k_q)=(312,8)$ to the leptonic descendant integrality
$I_L\in\mathbb Z$, because $I_Q=312/(3\cdot8)=13$ is already fixed. The shorthand
$\Delta_{\rm fr}=\operatorname{dist}(K/(2k_{\ell}),\mathbb Z)$ is therefore the
embedding-dependent operational diagnostic used inside the theorem, not a
universal statement for every $SO(10)$ branching problem.
The benchmark choice $K=312$ is therefore not introduced merely as a separate
minimality postulate. Once the visible branch $(k_{\ell},k_q)=(26,8)$ is fixed,
the same Diophantine rule forces the parent level into the discrete tower
\begin{equation}
 K\in \operatorname{lcm}(2k_{\ell},3k_q)\,\mathbb Z_{>0}
 = \operatorname{lcm}(52,24)\,\mathbb Z_{>0}
 = 312\,\mathbb Z_{>0}.
\end{equation}
On that branch the surface gauge datum is then
\begin{equation}
 \alpha^{-1}_{\rm surf}
 = G_{\rm SM}\frac{K}{k_{\ell}+k_q}
 = 15\frac{K}{34}
 = \frac{2340}{17}\,n,
 \qquad K=\benchmarkParentLevel n,
 \quad n\in\mathbb Z_{>0}.
\end{equation}
Hence the admissible parents compatible with the visible framing filter form a
discrete tower, and only the first element $n=1$ reproduces the benchmark
surface value $\alpha^{-1}_{\rm surf}=\alphaSurfBenchmarkExact\approx\alphaSurfBenchmarkRounded$. Higher multiples
$K=624,936,\ldots$ preserve the integrality condition but overshoot the gauge
datum by integer factors, while smaller positive integers fail the
visible-parent matching rule. In this precise branch-fixed sense, $K=312$ is
the unique discrete anchor that simultaneously satisfies the framing filter and
correctly post-dicts the ultraviolet surface gauge coupling.
To convert the Diophantine rule into a concrete benchmark statement, we treat $(k_{\ell},k_q,K)$ as discrete quantum numbers of the vacuum and hold fixed the benchmark parent level $K=312$ together with the quark branch $k_q=8$. We then scan the neighboring leptonic levels $k\in[24,28]$. For context, we also display a broader low-rank survey over the displayed window $k_{\ell}\in[2,100]$, $k_q\in[2,100]$. That broader survey is disclosed only as context for a fixed-parent \emph{Framing-Based Elimination Process}; it is not promoted to a theorem derived from $K=312$, nor is it the load-bearing selection criterion. The relevant fixed-parent criteria are
\begin{equation}
 \Delta_{\rm mod}^{(312)}(k)
 \equiv
 \mathrm{dist}\!\left(
 \frac{c(SO(10)_{312})-c(SU(2)_k)-c(SU(3)_8)-c_{\rm coset}(312)}{24},\,\mathbb Z
 \right),
\end{equation}
\begin{equation}
 \Delta_{\rm fr}(k)
 \equiv
 \left| \frac{K}{2k} - \text{round}\!\left(\frac{K}{2k}\right) \right|,
 \qquad
 \det S_{\rm flav}(k)
 \equiv
 \det\bigl[S^{(k)}_{q_{\alpha}g_i}\bigr]_{\alpha,i=1}^{3},
\end{equation}
where $\Delta_{\rm mod}^{(312)}$ records the fixed-parent central-charge mismatch,
$\Delta_{\rm fr}$ is the operational branch-closure diagnostic on
this slice, and $\det S_{\rm flav}\neq0$ guarantees a non-singular flavor block.
For the selected benchmark, however, the computation reports the nonzero residual
\begin{equation}
 \Delta_{\rm mod}(26,8)=\visibleResidual.
\end{equation}
The relevant modular statement is the full completed boundary partition
function,
\begin{equation}
 Z_{\partial}^{\rm comp}(\tau)
 =
 \sum_{\lambda,\mu,\nu,\sigma}
 N_{\lambda\mu\nu\sigma}\,
 \chi^{(2)}_{\lambda}(\tau)
 \chi^{(3)}_{\mu}(\tau)
 \chi^{\rm coset}_{\nu}(\tau)
 \chi^{\rm comp}_{\sigma}(\tau),
 \label{eq:completed-boundary-partition-function}
\end{equation}
for which each character transforms as
\begin{equation}
 \chi_a(\tau+1)=e^{2\pi i(h_a-c_a/24)}\chi_a(\tau).
\end{equation}
Hence every branching channel with $N_{\lambda\mu\nu\sigma}\neq0$ must obey the
integer-spectrum condition
\begin{equation}
 h_{\lambda}^{(2)}+h_{\mu}^{(3)}+h_{\nu}^{\rm coset}+h_{\sigma}^{\rm comp}
 -
 \frac{c\bigl(SU(2)_{26}\bigr)+c\bigl(SU(3)_8\bigr)+c_{\rm coset}^{\rm ref}+c_{\rm comp}}{24}
 \in\mathbb Z.
 \label{eq:full-t-spectrum-condition}
\end{equation}
Supplementary Sec.~\ref{supp-supp-sec:full-t-closure} derives this statement for the
benchmark branching channels and shows explicitly how the full condition
reduces to the shorthand $\Delta_{\rm fr}=0$ only in the present embedding with the fixed reference coset completion. In
particular, under a modular-$T$ step the completed vacuum block therefore carries the phase
\begin{equation}
 Z_{\partial}^{\rm comp}
 \xrightarrow{T}
 \exp\!\left[-\frac{2\pi i}{24}\bigl(c_{\rm vis}+c_{\rm coset}^{\rm ref}+c_{\rm comp}\bigr)\right]Z_{\partial}^{\rm comp},
\end{equation}
whereas the parent $SO(10)_{312}$ block carries
\begin{equation}
 Z_{\partial}^{\rm parent}
 \xrightarrow{T}
 \exp\!\left[-\frac{2\pi i}{24}c_{\rm parent}\right]Z_{\partial}^{\rm parent}.
\end{equation}
Modular invariance of the boundary partition function therefore requires
\begin{equation}
 \begin{aligned}
  c_{\rm vis}+c_{\rm coset}^{\rm ref}+c_{\rm comp}
  &\equiv c_{\rm parent}\pmod{24},\\
  c_{\rm comp}
  &\equiv c_{\rm parent}-c_{\rm vis}-c_{\rm coset}^{\rm ref}\pmod{24},
  \qquad 0\le c_{\rm comp}<24.
 \end{aligned}
\end{equation}
On the benchmark branch this gives
\begin{equation}
 c_{\rm comp}=24\,\Delta_{\rm mod}(26,8)=\cDarkCompletion.
\end{equation}
If left unaddressed, the visible block would carry a residual modular phase and the benchmark boundary partition function would fail the parent modular-closure test. The missing central charge is therefore treated exclusively as the formal completion residue required for convergence of the completed modular partition function on this branch. For observer-level language we write $c_{\rm dark}\equiv c_{\rm comp}$ and refer to this non-propagating complement as the Topological Gravitational Ghost (TGG) of the framing defect. No physical fields or propagating degrees of freedom are assigned to this complement in the load-bearing benchmark; it is bookkeeping for modular partition-function convergence and boundary-to-bulk unitarity only. We therefore make explicit the modular completion used in the remainder of the paper:
\begin{equation}
 SO(10)_{312}^{\rm bench}
 \longrightarrow
 SO(10)_{312}\otimes {\rm RCFT}_{\rm comp},
 \qquad
 c_{\rm comp}=\cDarkCompletion,
\end{equation}
where the complementary factor carries no Standard Model gauge charges because
its sole role is to account for the central-charge residue left by the
$SO(10)_{312}$ branching. Publication-facingly, $c_{\rm comp}$ is the unique
formal completion residue required for modular partition-function convergence:
it closes the benchmark bookkeeping without introducing a floating
completion-sector coordinate or a physical complement-field assignment. The point
of writing the completion this way is that the missing central charge need not
be treated as an arbitrary floating constant: concrete unitary rational
examples can be assembled from the BPZ/GKO unitary minimal series~\cite{bpz1984,gko1985}. At the level needed in this paper, existence is established more abstractly by taking $\mathrm{RCFT}_{\mathrm{comp}}$ to be the residual rational sector whose branching multiplicities restore the parent modular-$T$ spectrum of the completed block. The proposed complement sector $\mathrm{RCFT}_{\mathrm{comp}}$ provides the required central charge residue while maintaining a well-defined integer spectrum for the total boundary partition function. In the main text $c_{\rm dark}=\cDarkCompletion$ is the Holographic Stabilizer, simultaneously the source of the positive vacuum loading and the conjugate $B-L$ parity sink required for boundary neutrality. That parity sink is an information-theoretic bookkeeping requirement of modular-$T$ closure, not a claim of an additional propagating hidden-particle sector. A concrete rational presentation is, for example,
\begin{equation}
 {\rm RCFT}_{\rm comp}^{\rm ex}
 = \mathcal M(3,4)^{\otimes 2}\otimes \mathcal M(4,5)\otimes \mathcal M(5,6)^{\otimes 2},
\end{equation}
has
\begin{equation}
 c\bigl({\rm RCFT}_{\rm comp}^{\rm ex}\bigr)
 = \frac12+\frac12+\frac{7}{10}+\frac45+\frac45
 = \frac{33}{10}
 \approx 3.30,
\end{equation}
which provides an explicit nearby unitary rational point at the required benchmark scale for the missing charge. The exact benchmark statement remains the residual-sector closure above: once this complementary central charge is included, the parent central-charge accounting satisfies the same anomaly cancellation requirements used throughout this paper.

In the benchmark reconstruction, what is cataloged astrophysically as ``Dark Matter'' is the macroscopic gravitational footprint of the Topological Gravitational Ghost (TGG) of the framing defect rather than a separate particle fluid. The parity bits carried by this completion sector are the bookkeeping quanta required to preserve boundary-to-bulk unitarity and modular-$T$ closure in a finite-capacity universe with $N\sim10^{122}$, while the same Holographic Stabilizer loads the bulk through $\Lambda_{\rm holo}$ rather than through a separate particle-relic bath. Accordingly, $c_{\rm dark}$ is not introduced here as an unseen WIMP sector; it is the topological parity sink required by the information-theoretic accounting of the anomaly-free branch. The fixed benchmark proxy $\Omega_{\rm DM}h^2\approx0.012$ is therefore recorded as a derived gravitational residue of the same anomaly-free branch. The resulting fixed-parent scan is summarized in Table~\ref{tab:visible-level-candidates}.

\begin{center}
\scriptsize
\pubinput{uniqueness_scan_table.tex}
\end{center}

\noindent
\refstepcounter{table}\label{tab:visible-level-candidates}
\textbf{Table~\thetable: Visible-level audit of the saturated-information branch.}
Within the displayed fixed-parent window, the triplet $(k_{\ell},k_q,K)=(26,8,312)$ is the selected benchmark point: it is the Minimal Anomaly-Free Local Survivor within the $K=312$ Diophantine class and gauge-coupling window, even though $\Delta_{\rm mod}(26,8)=\visibleResidual\neq0$ and the adjacent point $(27,8,312)$ has a slightly smaller de-anchored modularity gap. This is the manuscript's consistency bottleneck. We therefore invoke $\Delta_{\rm fr}$ through the \emph{Derived Uniqueness Theorem}: operationally, the benchmark retains only cells satisfying the theorem's framing-closure diagnostic $\Delta_{\rm fr}=0$, which in the present embedding acts as the Diophantine shorthand for a boundary modular-$T$ spectrum with integer support. For the chosen $SO(10)\to SU(3)\times SU(2)$ branching, this framing test is sufficient for benchmark selection. In the core text that diagnostic functions only as the visible-branch elimination rule. A nonzero framing gap is therefore treated as failure of the theorem's framing condition---equivalently, failure of modular-$T$ consistency for the boundary partition function---so $(27,8,312)$ is not a competing benchmark candidate within that fixed-parent class but a cell removed by the filter before phenomenology is even consulted. Changing the leptonic integer level $k_{\ell}$ away from $26$ forces $\Delta_{\rm fr}\neq0$. By contrast, the mass coordinate $m_\nu=\kappa_{D_5}M_PN^{-1/4}$ is a bridge constant defined on top of the fixed discrete branch and does not continuously retune this filter. The benchmark is therefore a local mathematical survivor rather than a claim of global uniqueness. The added formal-completion-residue column $c_{\rm comp}$, $|\det S_{\rm flav}|$, and Visible Branch Status columns make the same point numerically explicit: in the displayed neighborhood $(26,8,312)$ is the only displayed local row that simultaneously closes the framing filter, preserves a non-singular flavor block, and matches the branch-fixed gauge / modular-completion bookkeeping without reopening a continuous tuning direction. The companion Origin columns record the provenance of the displayed residues: the surface-tension gauge proxy is Branch-Fixed (Discrete), while the VEV-alignment and central-charge residues are Geometric (Fixed); the absolute mass coordinate discussed in the text separately inherits the Cosmological (Input) provenance of $N$.

The corresponding fixed-parent residual map is shown in Table~\ref{tab:modularity-residual-map}.

\begin{center}
\scriptsize
\pubinput{modularity_residual_map.tex}
\end{center}

\noindent
\refstepcounter{table}\label{tab:modularity-residual-map}
\textbf{Table~\thetable: Nearest-neighbor modularity residual map around $(26,8,312)$.}
The residual map displays the exact distances to the zero-residual locus for the neighbors $(25,8)$ and $(27,8)$; its role is to visualize the discrete moat around the selected cell. All values in this table come from the same synchronized benchmark evaluation used throughout the manuscript. Together, Table~\ref{tab:visible-level-candidates} and Table~\ref{tab:modularity-residual-map} define the \emph{fixed-parent local moat audit}: the benchmark is locally isolated because any one-step move in the visible lattice violates the theorem's framing-closure condition, and at the upper edge of the window the relaxed audit also introduces a compensating modular tilt. This structural moat is visualized directly by the Framing Gap Heatmap in Figure~\ref{fig:framing-gap-heatmap}: within the displayed neighborhood, the cell $(26,8)$ is the only displayed candidate on the printed $K=312$ slice satisfying the theorem's framing-closure condition. The logic is lexicographic rather than statistical: first impose the theorem's branch-closure condition $\Delta_{\rm fr}=0$, then compare the residual modularity data only among the surviving cells. Any deviation from $(26,8)$ inside the displayed moat therefore fails the boundary branch-closure condition before it can qualify as a viable benchmark candidate on that slice.

The main statement is structural. Over the full low-rank search volume
\begin{equation}
 k_{\ell}\in[2,100],
 \qquad
 k_q\in[2,100],
\end{equation}
one can ask whether the selected tuple realizes an exact root of the fixed-parent closure condition
\begin{equation}
 \Delta c(k_{\ell},k_q)=24\,\Delta_{\rm mod}(k_{\ell},k_q)=0.
\end{equation}
In the current computational evaluation the residual $\scanTrialCount$-point scan without benchmark anchoring does not return an exact global root, so we state the result conservatively. We therefore do not claim a stronger global uniqueness theorem from the broad scan itself. The main statement is instead branch-local and structural: within the fixed-parent problem relevant to the manuscript, $(k_{\ell},k_q,K)=(26,8,312)$ is the fixed-parent benchmark selected by the local moat audit encoded in Table~\ref{tab:visible-level-candidates} and Table~\ref{tab:modularity-residual-map}. The neighboring entries then quantify the separation around that benchmark rather than a statistical tail:
\begin{equation}
 \mathfrak A_{\rm vis}(25,8)=3.317434,
 \qquad
 \mathfrak A_{\rm vis}(26,8)=3.300805,
 \qquad
 \mathfrak A_{\rm vis}(27,8)=3.300905.
\end{equation}
The same $\scanTrialCount$-point survey is retained only as a fixed-parent \emph{Local Moat Audit}. The scan verifies that within the displayed $K=312$ neighborhood, $(26,8)$ is the Minimal Anomaly-Free Local Survivor of the theorem's framing filter; this illustrates the discrete isolation of the benchmark without promoting it to a landscape-wide uniqueness theorem. Its role is to expose the local framing/modularity structure around the benchmark, not to define a landscape-wide benchmark-consistency metric or a global search over parent branches. The load-bearing statement is instead eliminative: the theorem's framing-closure condition ($\Delta_{\rm fr}=0$) removes every displayed neighbor with $\Delta_{\rm fr}\neq0$ before any PMNS/CKM comparison is attempted. The reduced $41\times13=533$ follow-up slice is documented in Supplementary Table~S2 only as auxiliary landscape bookkeeping for the local post-filter moat around the benchmark. The companion supplementary $\Delta\chi^2$ residue-profile table shows the same mechanism on the local moat: the nearest $k_{\ell}=25$ and $27$ neighbors reopen $\Delta_{\rm fr}$ and pick up $O(50)$ increases in $\chi^2_{\rm pred}$ even though $c_{\rm comp}$ shifts only mildly. The load-bearing statement is therefore local isolation within the fixed-parent class, not global uniqueness.

Therefore the admissible parent levels lie on the common descendant lattice
\begin{equation}
K\in 52\mathbb Z\cap24\mathbb Z
\;=\;
\operatorname{lcm}(52,24)\,\mathbb Z
\;=\;
312\mathbb Z,
\end{equation}
and the benchmark parent level is fixed by the minimality requirement: it is the least common multiple of the branch integers $2k_{\ell}=52$ and $3k_q=24$,
\begin{equation}
 \boxed{K=\operatorname{lcm}(2k_{\ell},3k_q)=312.}
\end{equation}
This is the minimal integer for which both visible branches can be embedded into a common parent without fractional descendants, and in the present paper it is treated as the UV-scale vacuum label of the benchmark branch. It gives
\begin{equation}
 I_L = \frac{312}{26\cdot2}=6,
 \qquad
 I_Q = \frac{312}{8\cdot3}=13,
 \qquad
 c\bigl(SO(10)_{312}\bigr)=\frac{351}{8}.
\end{equation}
The zero-energy framing state is then obtained by completing the visible symmetry-breaking branch with the orthogonal branch sector,
\begin{equation}
 c_{\rm coset}(312)
 \equiv
 c\bigl(SO(10)_{312}\bigr)-c\bigl(SU(2)_{26}\bigr)-c\bigl(SU(3)_8\bigr),
\end{equation}
so the complete framed branching relation is
\begin{equation}
 \frac{c\bigl(SO(10)_{312}\bigr)}{24}
 \equiv
 \frac{c\bigl(SU(2)_{26}\bigr)}{24}
 + \frac{c\bigl(SU(3)_8\bigr)}{24}
 + \frac{c_{\rm coset}(312)}{24}
 \pmod 1,
\end{equation}
and therefore
\begin{equation}
 \exp\!\left[-\frac{2\pi i}{24}
 \Bigl(c\bigl(SO(10)_{312}\bigr)-c\bigl(SU(2)_{26}\bigr)-c\bigl(SU(3)_8\bigr)-c_{\rm coset}(312)\Bigr)\right]=1.
\end{equation}
This is the formal framing-closure identity for the zero-residual state: the completed parent vacuum carries no residual framing phase. Equivalently, the bare visible residual $\mathfrak A_{\rm vis}(312)$ is exactly cancelled by the orthogonal coset contribution, while the $SU(2)_{26}$ core satisfies the modular identities $S^2=(ST)^3=I$.

\section{Benchmark Evaluation}

The modular kernels and parent framing data should be read as ultraviolet boundary conditions at the $SO(10)$ unification scale $M_{\rm GUT}=2\times10^{16}\,\mathrm{GeV}$. The experimentally quoted angles are instead low-energy quantities extracted near the electroweak scale $M_Z$ after symmetry breaking, threshold matching, and renormalization-group evolution. Accordingly, the proper comparison is not between bare topological numbers and a single central value, but between the ultraviolet kernel data and the current NuFIT~5.3 and PDG~2024 reference ranges after RG transport to low energies.

Writing
\begin{equation}
 L\equiv\frac{\ln(M_{\rm GUT}/M_Z)}{16\pi^2},
\end{equation}
the transparent analytic benchmark is the compact one-loop leading-log law
\begin{align}
 \theta_{ij}(M_Z)
 &= \theta_{ij}(M_{\rm GUT})
 + L(\beta_{\theta}^{(1)})_{ij}, \\
 \delta_f(M_Z)
 &= \delta_f(M_{\rm GUT})
 + L\beta_{\delta,f}^{(1)}, \\
 m_0(M_Z)
 &= m_0(M_{\rm GUT})\bigl[1+\gamma_0^{(1)}L\bigr],
 \qquad
 \gamma_0^{(1)}=\frac{C_2(\mathbf{126}_H)}{2T(\mathbf{126}_H)}=\frac{25/2}{70}=\frac{5}{28}\approx0.179.
\end{align}
The formula above is retained only as an analytic guide. The numerical entries in Table~\ref{tab:global-flavor-fit} are the high-precision outputs of the full coupled RG transport in which the Yukawa and gauge couplings are evolved dynamically together with the flavor observables. Appendix~\ref{app:rg} derives the corresponding Antusch one-loop coefficients and the split-threshold audit formulas, while Appendix D summarizes the computational implementation used to generate the quoted constants and figures.

\subsection{Primary Benchmark-Consistency Metric}

The main text uses a single primary benchmark comparison built from $8$ observables. The leptonic set is $(\theta_{12},\theta_{13},\theta_{23},\delta_{CP})$, and the quark-sector set is $(|V_{us}|,|V_{cb}|,|V_{ub}|,\gamma)$. The two disclosed flavor-side inputs are $G_{\rm SM}=15$ and $\langle \Sigma_{126}\rangle/\langle\phi_{10}\rangle=64/312$. Once the branch label $(26,8,312)$ and those two inputs are fixed, the primary main-text tally has $\nu=8-2=6$ and benchmark-consistency metric $\chi^2_{\rm bench}\equiv\chi^2_{\rm pred}=\chiPredRounded$. The cosmological budget $N$, the modular complement residue $c_{\rm comp}$, auxiliary audit quantities, and any off-shell threshold bookkeeping are treated as external or supplementary audits rather than as alternative main-text degree-of-freedom conventions. Alternative raw-row and off-shell tallies are documented only in Supplementary Sec.~S12.6. % pub/tn.py: DISCLOSED_BENCHMARK_MATCHING_CONDITION_COUNT, PullTable.threshold_alignment_subtraction_count, PullTable.lambda_normalization_matching_subtraction_count

\noindent\textbf{One-line metric summary.} Two flavor-side matching residues ($G_{\rm SM}=15$ and the $\mathbf{126}_H$ Clebsch ratio) are disclosed as external inputs. Once the branch label $(26,8,312)$ and the disclosed benchmark settings $G_{\rm SM}=15$ and $\langle \Sigma_{126}\rangle/\langle\phi_{10}\rangle=64/312$ are fixed, the benchmark contains no further adjustable flavor parameters that can be continuously varied to improve the comparison; the reported $\chi^2$ is therefore a benchmark-consistency metric rather than a likelihood-based inference.

{\centering
\footnotesize
\begin{tabularx}{\linewidth}{|L{0.28\linewidth}|L{0.18\linewidth}|Y|>{\centering\arraybackslash}p{0.12\linewidth}|}
\hline
Observable & Status & Inputs & Continuous Degrees of Freedom \\
\hline
$\theta_{12},\theta_{13},\theta_{23},\delta_{CP}$ & post-diction & $(26,8,312)$ & $0$ \\
$|V_{us}|,|V_{cb}|,|V_{ub}|,\gamma$ & post-diction & $(26,8,312),\;\mathcal R_{\rm GUT}^{\rm bench}=\wTh$ & $0$ \\
\shortstack[l]{$m_\nu=\kappa_{D_5}M_P$\\$N^{-1/4}$} & \shortstack[l]{holographic\\consistency check} & $N,\;\kappa_{D_5}$ & $0$ \\
\shortstack[l]{$\alpha^{-1}_{\rm surf}$\\$=15K/(k_{\ell}+k_q)$} & consistency check & $(26,8,312),\;G_{\rm SM}=15$ & $0$ \\
\hline
\end{tabularx}\par}

\noindent\textbf{Parameter disclosure table.} Table~\ref{tab:benchmark-input-disclosure} keeps the continuous-degree count explicit: every displayed observable is evaluated from fixed discrete labels and disclosed benchmark inputs, with no floating flavor-sector parameters. The modular complement residue $c_{\rm comp}$ and auxiliary audit quantities do not enter Table~\ref{tab:global-flavor-fit} or the quoted $\chi^2_{\rm pred}$.

{\centering
\scriptsize
\pubinput{global_flavor_fit_table.tex}\par}
\enlargethispage{2\baselineskip}

\noindent
\refstepcounter{table}\label{tab:global-flavor-fit}
\textbf{Table~\thetable: Benchmark comparison and auxiliary closure data.}
The quoted theory values are the numerically integrated RG outputs of the fixed boundary kernel for the $M_Z$ rows, with the auxiliary $\alpha^{-1}(M_{match})$ row reported directly at the anomaly-free matching point. All values in this table come from the single synchronized benchmark evaluation summarized in Appendix~\ref{app:computational}. The comparison uses NuFIT~5.3 $(2024)$ normal-ordering reference intervals with SK atmospheric data for $(\theta_{12},\theta_{13},\theta_{23},\delta_{CP})$ and PDG~2024 benchmark values for $(|V_{us}|,|V_{cb}|,|V_{ub}|,\gamma)$. The table presents the central benchmark: the PMNS sector is propagated with the standard SM Antusch transport, the branch-fixed threshold residue $\mathcal R_{\rm GUT}$ from Table~\ref{tab:notation-definitions} is the only threshold-sensitive coefficient of the one-loop dimension-5 Wilson coefficient in the $\mathbf{126}_H$ sector that enters the CKM apex, and the overlap ratio $\mathcal N_{\rm sol}\equiv4|\Omega_{e\mu}|/\Tr\Omega$ is tracked separately as the quantified interface geometry of the solar channel rather than used to generate the primary pulls. The solar-angle row, with central pull $-0.40\sigma$, is therefore reported as the current numerical imprint of that interface geometry rather than as an unresolved tension. The $m_\nu$ row is displayed as the RG Consistency Audit implied by the benchmark mass relation, namely the identity $m_\nu=\kappa_{D_5}M_PN^{-1/4}$ implied by the fixed $D_5$ geometric factor together with the externally supplied Planck-2018 bit budget $N$ anchored to $\Lambda_{\rm obs}$. The primary flavor-side benchmark metric remains $\chi_{\rm pred}^2=\chiPredRounded$, and the RG Consistency Audit is not counted inside that eight-row tally. Neither the modular complement residue $c_{\rm comp}$ nor any auxiliary audit quantity enters the rows counted in that primary flavor tally. The additional $\alpha^{-1}(M_{match})$ row is displayed as the benchmark boundary-closure residue $137.035999\ldots+\delta_{\rm topo}$ with $\delta_{\rm topo}\equiv1-\kappa_{D_5}$, and its table note reads: \emph{Value derived from the $SO(10)_{312}$ boundary anomaly-free closure at $M_{match}$.} The companion $\alpha$-evolution audit is synchronized to the same exported completion residue $c_{\rm dark}=\cDarkCompletion\approx3.3008$, so the auxiliary gauge audit and the displayed $\alpha^{-1}(M_{match})$ row remain traceable to one runtime object rather than to separate hand-entered numbers. Accordingly, the gauge-side note reads: \emph{Value derived from the $SO(10)_{312}$ boundary anomaly-free closure at $M_{match}$.} The Structural Note column therefore separates rigid magnitudes from threshold-sensitive observables: the mixing angles and CKM magnitudes are controlled primarily by overlap data, whereas the CP phases and the auxiliary $\alpha^{-1}$ closure row track framing holonomy plus branch-fixed heavy-threshold dressing. \emph{Benchmark consistency status.} Every row in Table~\ref{tab:global-flavor-fit} inherits the same structural status \textit{evaluated on the anomaly-free branch}: the benchmark branch is evaluated only on the anomaly-free cell with $\Delta_{\rm fr}=0$, so the reported pulls probe flavor transport on the fixed anomaly-free background rather than re-optimizing the geometry row by row.

\begin{figure}[t]
\centering
\includegraphics[width=0.86\linewidth]{fig1_pmns_fit.png}
\caption{Modular origin of the PMNS points: the red markers show the ultraviolet $SU(2)_{26}$ kernel obtained from the polar modular block followed by the conical-holonomy RT rotation, while the table reports the corresponding RG-evolved benchmark values. The blue and gray bands display the NuFIT $1\sigma$ and $3\sigma$ intervals used for the phenomenological comparison.}
\label{fig:pmns-fit}
\end{figure}

The central benchmark is reported here with explicit uncertainty propagation. Table~\ref{tab:global-flavor-fit} reports all eight flavor observables, including $\gamma$, and the complex-holonomy calculation of Sec.~10.2 brings the CKM apex angle onto the same benchmark table rather than relegating it to a side diagnostic. All eight flavor observables are obtained from the full coupled one-loop numerical transport, with the third-family SM Yukawas and the three gauge couplings propagated dynamically inside the same state vector. The separate $\alpha^{-1}(M_{match})$ entry is not counted among those eight rows: it is displayed only as the branch-fixed benchmark residue $137.035999\ldots+\delta_{\rm topo}$ with $\delta_{\rm topo}\equiv1-\kappa_{D_5}$ derived from the same anomaly-free closure, with its companion $\alpha$-evolution audit anchored by $c_{\rm dark}=\cDarkCompletion$. We identify the branch-fixed threshold residue $\mathcal R_{\rm GUT}$ from Table~\ref{tab:notation-definitions} as the exact rational identity $k_q/(k_{\ell}+h^{\vee}_{SU(2)})=8/28$. The explicit one-loop dimension-5 Wilson coefficient generated when the heavy $SO(10)$ scalar sector is matched onto the Standard Model at $M_{\rm GUT}$ is treated only as a consistency audit of that fixed target. The logarithm $\ln(M_P/M_{\rm GUT})$ entering that matching calculation is therefore not an ad hoc tilt but the expected leading-log separation between the ultraviolet boundary scale and the GUT matching scale. Within the published benchmark its value is fixed by the same anomaly-free branch data and disclosed heavy-spectrum bookkeeping rather than floated as an auxiliary threshold adjustment. The displayed comparison is therefore presented as a branch-fixed benchmark comparison in the flavor and boundary-mapping sectors. The relevant benchmark data are the disclosed input block summarized in Table~\ref{tab:benchmark-input-disclosure}, together with the branch-fixed CKM transport datum $\mathcal R_{\rm GUT}$ governing the threshold dressing of $\gamma$. For the primary main-text comparison only the two flavor-side matching conditions $G_{\rm SM}=15$ and $(C_{126}^{(12)})^{-1}=64/312$ are subtracted, giving $\nu=6$; alternative raw-row and off-shell threshold tallies are recorded only in Supplementary Sec.~S12.6. The construction is presented as a tightly constrained benchmark rather than a freely adjustable model: the discrete vacuum selection $(26,8,312)$ supplies fixed benchmark residues, e.g. $\kappa_{D_5}=\simplexPrefactor$, derived from the weight-simplex geometry. Once the branch label $(26,8,312)$ and the disclosed benchmark settings $G_{\rm SM}=15$ and $\langle \Sigma_{126}\rangle/\langle\phi_{10}\rangle=64/312$ are fixed, the benchmark contains no further adjustable flavor parameters that can be continuously varied to improve the comparison. The near-unity value of $\kappa_{D_5}$ is therefore read as a post-diction of the $D_5$ structure entering the benchmark mass relation, not as a tunable normalization knob. The mixing angles and CKM phase are outputs of the $SO(10)_{312}$ branch, whereas the auxiliary $\alpha^{-1}$ entry is recorded as a branch-derived closure residue rather than as an unresolved threshold residual. The companion $\alpha$-evolution audit depends only on the topological index $\mathcal I_Q$, the already-fixed Modular Restoration Scale $M_N$, and the same completion residue $c_{\rm dark}=\cDarkCompletion$, so the gauge-side closure introduces no new degrees of freedom. The benchmark neutrino scale is retained as a theory-fixed finite-capacity consequence of the anomaly-free branch. Once the Planck-anchored bit budget is supplied, the relation $m_\nu=\kappa_{D_5}M_PN^{-1/4}$ fixes $m_1\approx2.92\,\mathrm{meV}$ and the transported normal-ordering sum $\sum m_\nu\approx0.062\,\mathrm{eV}$; these rows are displayed as theory-fixed values, not as adjustable parameters or hidden normalization knobs. The appropriate summary is therefore the total benchmark-consistency metric for the theoretical benchmark. Here $\chi_{\rm pred}^2$ is used as the primary benchmark-consistency metric for the fixed branch, not as a global likelihood or landscape-wide significance claim. Observable-level uncertainties are reported as quantified residues of the branch-fixed transport map, extracted from the Two-Loop Curvature Audit of the $(26,8,312)$ benchmark, rather than as tunable nuisance parameters. The theoretical uncertainties entering the $\chi^2$ analysis are therefore reported through the \textbf{Quantified Two-Loop Residuals} exported by that branch-fixed transport audit. The benchmark mass relation is therefore audited as a disclosed branch-fixed theorem of the finite-capacity cell. % pub/tn.py: build_matching_residual_report, MATCHING_RESIDUAL_REPORT_FILENAME

Within the disclosed $\scanTrialCount$-point landscape, the branch selection is lexicographic rather than a $\chi^2$ minimization. Step~1 applies the theorem's framing-closure condition $\Delta_{\rm fr}=0$. Step~2 compares the residual modularity only among the survivors of Step~1. In the displayed $k_q=8$ window, $(26,8)$ is the selected fixed-parent benchmark because it survives Step~1 and is the local survivor carried forward in Step~2; the neighboring $(27,8)$ cell is numerically nearby in de-anchored modularity but is excluded at Step~1 because $\Delta_{\rm fr}\neq0$. The nonzero residual $\Delta_{\rm mod}(26,8)=\visibleResidual$ is therefore treated as a diagnostic residual of the fixed-parent construction rather than as evidence for an exact modular root. We therefore use the low-rank survey as bookkeeping for the local framing/modularity structure around the selected branch, not as evidence that the wider scan uniquely proves the benchmark. Supplementary Sec.~S3 derives the corresponding stress tensor.

We also use the hyperarea ratio of the $D_5$ weight simplex as the benchmark normalization for the boundary--bulk scale map, writing $\kappa\equiv R_{01}^{\rm par/vis}\approx\simplexPrefactor$. This residue is not floated independently; it ties the benchmark mass relation $m_\nu=\kappa_{D_5}M_PN^{-1/4}$ to the saturated bit budget, and detuning it changes the disclosed matching audit for the same branch. Finally, while the topological kernel sets within the benchmark the mixing eigenvectors, the scalar-sector ratio $\langle \Sigma_{126} \rangle / \langle \phi_{10} \rangle=(C_{126}^{(12)})^{-1}=64/312$ is read as the disclosed scalar-sector benchmark setting of the anomaly-free cell rather than as a phenomenological mass reparameterization. While absolute mass eigenvalues are matched to the Higgs-sector VEVs, the mixing structure (PMNS/CKM eigenvectors) is a benchmark-stable output of the topological kernel. The PMNS/CKM mixing eigenvectors are topologically protected. The SVD Stability Audit (see Fig.~\ref{fig:vev-stability-main} and Supplementary Fig.~\ref{supp-fig:s3-vev-stability}) demonstrates that the mixing angles remain stable within $1\sigma$ targets even under $\pm10\%$ deformations of the disclosed $64/312$ benchmark setting. Within this rigid/matched split, the model fixes $m_{\rm lightest}$ only as a consistency check of the absolute mass scale through the benchmark mass relation, while the observed splittings $\Delta m_{21}^2$ and $\Delta m_{31}^2$ enter as low-energy boundary conditions for the transported spectrum.

The resulting benchmark summary is therefore
\begin{align}
 \chi_{\rm bench}^2&\equiv\chi_{\rm pred}^2=\chiPredRounded,
 &\nu_{\rm bench}&=6, \\
 \mathrm{RMS\ pull}_{\rm bench}&=\rmsPullPred,
 &\max |\mathrm{pull}|_{\rm bench}&=\maxPullPred\sigma, \\
 T_{\rm scan}&=\scanTrialCount.
\end{align}
Here $\chi_{\rm bench}^2$ is the primary benchmark-consistency metric of the main text: $8$ observables, $2$ disclosed inputs, and therefore $\nu=6$. The quantity $T_{\rm scan}$ records the size of the disclosed fixed-parent Local Moat Audit. Alternative raw-row, audit, and off-shell threshold tallies are collected in Supplementary Sec.~S12.6 rather than promoted to separate main-text metrics.
The same benchmark coordinate block fixes the benchmark mass relation $m_\nu=\kappa_{D_5}M_PN^{-1/4}$ once the Planck-anchored bit budget $N$ fixed by the observed cosmological constant $\Lambda_{\rm obs}$ is supplied.

\paragraph{Scale-Setting Consistency Check of $m_\nu$.} On the anomaly-free cell $(26,8,312)$, the relation $m_\nu=\kappa_{D_5}M_PN^{-1/4}$ uses the Planck-anchored bit budget $N=C_{\max}$ as a fixed boundary datum supplied by the observed cosmological constant $\Lambda_{\rm obs}$. Because $N$ is finite and $\kappa_{D_5}=\simplexPrefactor>0$, the construction requires $m_\nu\neq0$ on the minimal branch. In the benchmark evaluation this gives the theory-fixed values $m_1\approx2.92\,\mathrm{meV}$ and $\sum m_\nu\approx0.062\,\mathrm{eV}$. These mass rows are displayed as theory-fixed consequences of the finite-capacity theorem and do not enter the eight-observable benchmark-consistency metric reported in the main text.
Alternative raw-row, audit, and off-shell threshold tallies are retained only in Supplementary Sec.~S12.6 and are not used as separate main-text comparison metrics.

The bookkeeping-only solar-overlap discussion is therefore retained as a robustness check rather than exported as a second auxiliary block. The point of the eight-row presentation is to keep every displayed observable inside the same fixed SM transport: the solar row is generated by that common evolution, and the CKM apex angle $\gamma$ remains on the same benchmark table as the mixing angles and CKM magnitudes. Read this way, the disclosed precision-limit discussion does not hide an outlier: the largest central pull in the generated benchmark table is the $\delta_{CP}$ row at approximately $-1.25\sigma$, whereas $\theta_{12}$ remains the row most sensitive to subleading solar-overlap deformations in the one-loop kernel while still lying within theoretical precision limits. Within the scanned $\scanTrialCount$-point survey, $(26,8,312)$ is the selected framing-closed benchmark. The broader residual survey is reported for context and does not, by itself, support a stronger exact-selection theorem.
In the leptonic case the structural kernel still fixes the ultraviolet CP benchmark, the holonomy-area factor $\Ahol\approx3.3\times10^{-2}$, and the neutrinoless double beta decay benchmark target $|m_{\beta\beta}|=5.60\,\mathrm{meV}$. At the ultraviolet matching scale the bare topological value corresponds to a slightly smaller Majorana benchmark, while RG evolution lifts it to the quoted $M_Z$ floor.

\subsection*{Rigidity Audit}

In the quark case the hierarchy no longer enters through an unconstrained threshold phase. Once $K=312$ is fixed by the Derived Uniqueness Theorem, the branching index $I_Q=13$ and the rank gap $\Delta r=3$ determine the vacuum pressure $\Pi_{\rm vac}$ and the descendant factors $\Xi_{12},\Xi_{23}$, which in turn fix the CKM suppression exponents. The same numerical RG transport and complex $\mathbf{126}_H$ holonomy now yield predictions for $|V_{us}|$, $|V_{cb}|$, and $|V_{ub}|$ within the current PDG widths. The remaining CKM apex tilt is carried by the threshold matching term of the $\mathbf{126}_H$ Majorana channel,
\begin{equation}
 \mathcal C_{\rm tilt}
 \equiv
 \frac{1}{4}\frac{C_2(\mathbf{126}_H)}{2T(\mathbf{126}_H)}
 =
 \frac{12.5}{280}
 =
 \frac{5}{112}
 \approx 0.0446,
 \qquad
 \mathcal R_{\rm GUT}^{\rm bench}
 \equiv
 \frac{k_q}{k_{\ell}+h^{\vee}_{SU(2)}}
 =
 \frac{8}{28}
 = \wTh = \frac{2}{7}.
\end{equation}
Here $\mathcal C_{\rm tilt}$ fixes the representation-theoretic scale of the heavy-threshold effect through the $SO(10)$ Casimir invariants and the level-$K=312$ modular normalization. The benchmark residue itself is not extracted from the one-loop audit: it is the exact branch identity above, while the explicit one-loop matching calculation implemented in the audit is used only as a consistency check,
\begin{equation}
 \lambda_{12}^{(5)}(M_{\rm GUT})
 =
 Y_{12}^{(0)}\,\frac{g_{\rm GUT}^2}{16\pi^2}
 \sum_A C_A\ln\!\left(\frac{M_P}{M_A}\right),
\end{equation}
with $A$ running over the heavy $\mathbf{126}_H$ fragments at $M_{126}^{\rm match}$, the coarse $\mathbf{210}_H$ GUT-breaking fragments at the GUT threshold, and the broken $SO(10)/SM$ gauge bundle. The branch-fixed threshold residue $\mathcal R_{\rm GUT}$ from Table~\ref{tab:notation-definitions} is therefore not adjusted in flavor space and not numerically matched into existence; the heavy-spectrum sum simply checks that the same parent completion reproduces the fixed Wilson-coefficient normalization. This residue multiplies only the orthogonal-coset phase sourced by the heavy threshold,
\begin{equation}
 \Phi_{\perp}(M_{126})
 =
 \Phi_{\perp}^{(0)}
 \left[1+\mathcal R_{\rm GUT}\frac{\ln(M_{\rm GUT}/M_{126})}{\ln(M_{\rm GUT}/M_Z)}\right],
\end{equation}
so the CKM phase is dressed by a theory-motivated EFT matching term. Thus $\mathcal R_{\rm GUT}$ governs the finite $SO(10)$ heavy-threshold phase, not a deformation of the rigid $S$-matrix magnitudes. Because it enters only this phase budget, the runtime profile sweep leaves $|V_{us}|$, $|V_{cb}|$, and $|V_{ub}|$ unchanged to within $\wThVusSpan$, $\wThVcbSpan$, and $\wThVubSpan$, respectively, while displacing $\gamma$ away from the benchmark closure point. Figure~\ref{fig:wth-profile} therefore serves as the manuscript's rigidity audit: a one-parameter off-shell profile of the disclosed UV-matching benchmark rather than a second adjustment stage.

\paragraph{Threshold-Scale Sensitivity.}
The recovery of the CKM phase $\gamma \approx 67^\circ$ is tied to the branch-fixed $\mathbf{126}_H$ matching scale $M_{\rm match}\equiv M_{126}^{\rm match}$. On the anomaly-free branch the physical framing-gap alignment $\Delta_{\rm fr}=0$ is taken to saturate the one-copy holographic support bound, which closes the threshold equation $\ln(M_N/M_{126})=\beta^2$ and therefore fixes $M_{126}^{\rm match}=M_Ne^{-\beta^2}$. The disclosed one-loop quantity $\delta_{\rm RG}=-0.61$ is then evaluated at that branch-fixed threshold. In this sense $\gamma$ is a saturated-framing post-diction of the fixed branch rather than a continuously tuned mass-sector input, while the off-shell threshold sweeps are retained only as robustness diagnostics.

\begin{figure}[t]
\centering
\includegraphics[width=0.84\linewidth]{fig3_ckm_correspondence.png}
\caption{Modular origin of the CKM points: the displayed values of $|V_{us}|$ and $|V_{cb}|$ are mixing predictions obtained from the branching-pressure-damped low-weight block of the $SU(3)_8$ modular kernel. The comparison bands indicate the PDG-centered $1\sigma$ and $3\sigma$ intervals used as external diagnostics.}
\label{fig:ckm-fit}
\end{figure}

\begin{figure}[t]
\centering
\includegraphics[width=0.76\linewidth]{fig5_threshold_weight_profile.png}
\caption{Rigidity audit for the normalized GUT-threshold residue of the one-loop dimension-5 Wilson coefficient. The curve shows the response $\Delta\chi^2_{\rm pred}(\mathcal R_{\rm GUT})$ of the full benchmark-consistency metric under off-shell deformations away from the benchmark value quoted in Table~\ref{tab:notation-definitions}. The same runtime sweep confirms that varying $\mathcal R_{\rm GUT}$ leaves the rigid magnitudes $|V_{us}|$, $|V_{cb}|$, and $|V_{ub}|$ unchanged to within $\wThVusSpan$, $\wThVcbSpan$, and $\wThVubSpan$, so the leading-order threshold correction induced by the heavy $SO(10)$ scalar sector acts as CKM phase alignment rather than as an extra floating coordinate.}
\label{fig:wth-profile}
\end{figure}

\section{Boundary Bit Saturation and a Model-Dependent One-Copy Preference for Normal Ordering}

Physical states satisfy the static constraint
\begin{equation}
 \hat H_{\partial}\Psi_{\rm phys}=0.
\end{equation}
By the \emph{Minimal One-Copy Dictionary} we mean the benchmark support map in which the three light neutrino mass eigenstates are assigned to the visible genus ladder $g\in\{0,1,2\}$ with at most one state per genus channel and with no duplicated support slots or auxiliary multi-copy sectors. In the present benchmark this is not merely an ansatz: for the anomaly-free $SO(10)$ spinor-Higgs interaction pattern $16_F16_F10_H\oplus16_F16_F\overline{\mathbf{126}}_H$, each visible genus slot may carry only one boundary precursor if the interaction map is to remain rank-preserving on the anomaly-free branch. The load-bearing requirement is therefore that the transition map $\mathcal T$ preserve rank on this three-state subspace. All statements about NH and IH in this section refer only to that minimal dictionary.

The additional assumption used in the support-order discussion is stricter: one compares the physical mass ordering to the idealized topological genus ladder,
\begin{equation}
 g\in\{0,1,2\},
 \qquad
 m_g = m_0e^{\beta g},
\end{equation}
with the branch-fixed modular-$T$ matrix phase $\beta$ listed in Table~\ref{tab:notation-definitions}.
The NO/IH statement below is therefore contingent on taking this ladder as a strict numerical identification of the physical mass ratios, not merely as a mnemonic for the support ordering. In the actual NO benchmark the same ladder is treated more weakly, namely as the structural backbone for the one-copy support assignment, while the transported physical masses are allowed to deviate from the exact ratios $1:e^{\beta}:e^{2\beta}$ so as to reproduce the observed oscillation splittings.
For the load-bearing argument it is sufficient to keep only the support-order diagnostic in the main text. The explicit character-to-generation dictionary, including the asymptotic support characters, the transition matrix $\mathcal T_{\alpha i}=\langle \chi_{g_{\alpha}}\mid\nu_i\rangle$, and the corresponding overlap matrices, is recorded in Supplementary Sec.~\ref{supp-supp-sec:one-copy-character-dictionary}. Within the minimal one-copy benchmark, Normal Ordering aligns directly with the unpermuted boundary character sequence $(g_1,g_2,g_3)=(0,1,2)$. Inverted Ordering can still be discussed, but only by a non-minimal permutation of the same genus labels or by enlarging the support dictionary. The main-text point is therefore not that IH is mathematically excluded; it is that the disclosed benchmark carries a minimal structural bias toward the ordering that follows the natural boundary character order.

\paragraph{Model-dependent preference under the Minimal One-Copy Dictionary.}
Under the strict genus-ladder realization of the anomaly-free $SO(10)_{312}$ branch, the preference for NO is a consequence of the Minimal One-Copy Dictionary assumption: IO/IH is not treated as rigorously excluded; rather, NO is the preferred structural realization of the one-copy support map. The reason is conditional: NH admits the natural ordered assignment $(g_1,g_2,g_3)=(0,1,2)$, whereas IH becomes non-minimal only after one also imposes the extra requirement that the physical mass ratios themselves coincide with $m_g=m_0e^{\beta g}$. In that case the inverted assignment no longer follows the boundary character order and must instead be represented by a relabeling of the visible genus slots or by multi-copy assignments of boundary characters. The detailed character dictionary and the corresponding strict-ladder overlap audit are recorded in Supplementary Sec.~\ref{supp-supp-sec:one-copy-character-dictionary} and should be read as diagnostics of that specific benchmark basis, not as a theorem against all IO/IH completions.

The scope of this statement is deliberately narrow. What is constrained is the strict-ladder implementation of Inverted Ordering inside the minimal selection-hypothesis-admissible parent theory, namely the $K=312$ completion with a one-copy visible support basis. In that precise sense, the one-copy benchmark carries a minimal structural bias toward the natural boundary-character order, while a non-minimal boundary Hilbert space, an explicit permutation of the genus labels $(0,1,2)$, or a looser mass-to-genus identification remains a possible extension rather than a no-go theorem for all ultraviolet completions. The conclusion of this section is therefore assumption-sensitive but definite: NH/NO is the natural structural realization of the strict one-copy benchmark, whereas IO/IH remains available only through a more complex support map than the one adopted here.

\paragraph{One-Copy Structural Preference.}
We therefore adopt Normal Ordering as the preferred structural realization of the minimal dictionary, not as a theorem against IO/IH. The relaxed proxy below is retained only as a finite visualization of that assumption-dependent minimal structural bias.

For visualization and for a finite one-copy proxy, it is convenient to relax the determinant divergence and pass to logarithmic support coordinates,
\begin{equation}
 \lambda_i \equiv \ln\!\left(\frac{m_i}{m_0}\right),
 \qquad
 \widetilde{\mathcal E}_{\rm WDW}(\lambda_1,\lambda_2,\lambda_3)
 = \min_{\pi\in S_3}
 \sum_{i=1}^{3}\bigl(\lambda_i-\beta g_{\pi(i)}\bigr)^2.
 \label{eq:wdw-relaxed-proxy}
\end{equation}
For the idealized normal-ordering backbone,
\begin{equation}
 (\lambda_1,\lambda_2,\lambda_3)_{\rm NO,backbone}=(0,\beta,2\beta),
\end{equation}
so the relaxed proxy also vanishes,
\begin{equation}
 \boxed{\widetilde{\mathcal E}_{\rm WDW}^{\rm NO,backbone}=0.}
\end{equation}
Again, this zero is a property of the structural backbone assignment only; the physical NO benchmark masses are subsequently allowed to depart from exact ladder spacing when matching oscillation data. For the inverted best proxy under the same strict-ladder assumption, one requires a nontrivial permutation of the genus labels,
\begin{equation}
 (\lambda_1,\lambda_2,\lambda_3)_{\rm IH,best}=(\beta,\beta,0),
\end{equation}
and therefore
\begin{equation}
 \widetilde{\mathcal E}_{\rm WDW}^{\rm IH}
 = \min_{\pi\in S_3}
 \sum_{i=1}^{3}\bigl(\lambda_i^{\rm IH}-\beta g_{\pi(i)}\bigr)^2
 = \beta^2 > 0.
\end{equation}
Numerically,
\begin{equation}
 \boxed{\widetilde{\mathcal E}_{\rm WDW}^{\rm IH}=\beta^2\approx 3.07845.}
\end{equation}
Thus, within the strict genus-ladder version of the Minimal One-Copy Dictionary, normal ordering remains the cleaner structural realization, whereas inverted ordering, under that same extra assumption, retains the finite Relaxed Proxy Gap $\widetilde{\mathcal E}_{\rm WDW}^{\rm IH}=\beta^2\approx3.08$ in the relaxed one-copy proxy. This should be read as the finite cost of leaving the natural boundary character order inside the minimal benchmark dictionary, not as a mathematical exclusion of IH; a more complex multi-copy support map could still accommodate IO/IH.

\begin{figure}[t]
\centering
\includegraphics[width=0.76\linewidth]{fig2_stability_audit.png}
\caption{Modular origin of the stability bars: within the strict genus-ladder version of the Minimal One-Copy Dictionary, the idealized normal-ordering backbone follows the unpermuted one-copy genus sequence, whereas the inverted backbone requires a nontrivial genus-label permutation and therefore retains the positive relaxed proxy gap $\beta^2$. The plotted quantity is the relaxed proxy functional evaluated on these backbone assignments; the transported NO benchmark used elsewhere in the paper is allowed to deviate from exact ladder spacing when compared with oscillation data.}
\label{fig:stability-audit}
\end{figure}

Within the \emph{Minimal Basis Assumption}, each light genus channel appears exactly once in the boundary Hilbert space. If one also identifies the physical mass ordering with the strict genus ladder, the Inverted Hierarchy no longer coincides with the benchmark character ordering. Accommodating it would therefore require either a non-minimal permutation of the visible genus labels $(0,1,2)$ or an enlarged support map. The present paper adopts neither move in the minimal benchmark, so the relaxed proxy still carries the residual gap $\beta^2$. The benchmark claim is therefore a model-dependent minimal structural bias toward NO within the disclosed one-copy implementation, not a statement that IO is mathematically impossible once more elaborate multi-copy support maps are allowed.

\begin{figure}[t]
\centering
\includegraphics[width=0.72\linewidth]{framing_gap_moat_heatmap.png}
\caption{Framing Gap Heatmap around the benchmark visible cell. The displayed quantity is the visible Chern--Simons-consistency diagnostic $\Delta_{\rm fr}^{\rm vis}$ in the low-rank moat around $(k_{\ell},k_q)=(26,8)$. On the printed $K=312$ slice, the benchmark cell $(26,8)$ is the only displayed cell satisfying the consistency condition; every neighboring cell gives $\Delta_{\rm fr}^{\rm vis}\neq0$. This plot is the algebraic preselection stage of the disclosed Local Moat Audit: it explains why the manuscript treats $(26,8,312)$ as the Minimal Anomaly-Free Local Survivor in the displayed moat rather than as a statement that the full low-rank survey is globally minimized by that point.}
\label{fig:framing-gap-heatmap}
\end{figure}

\section{Error Analysis}

Three complementary stability diagnostics sharpen the benchmark: a robustness audit of the simplex-area matching hypothesis, a bit-count uncertainty propagation, and a conservative two-loop RG diagnostic for the running of $m_0$ and $\theta_{12}$. The first is the explicit geometric factor
\begin{equation}
 m_0 = R_{01}^{\rm par/vis} M_PN^{-1/4},
 \qquad
 R_{01}^{\rm par/vis}\in[0.8,1.2],
\end{equation}
which is used only as a robustness neighborhood around the benchmark geometric residue $R_{01}^{\rm par/vis}=\simplexPrefactor$. The benchmark itself does not promote this coefficient to a scan direction: the surrounding interval is displayed only to quantify how strongly the benchmark mass scale responds to order-one geometric deformations. Over this audit window one finds
\begin{equation}
 \max_{R_{01}^{\rm par/vis}\in[0.8,1.2]}|\Delta m_0(M_Z)|\approx 0.80\,\mathrm{meV},
 \qquad
 \max_{R_{01}^{\rm par/vis}\in[0.8,1.2]}|\Delta m_{\beta\beta}(M_Z)|\approx 0.65\,\mathrm{meV}.
\end{equation}
The second diagnostic follows directly from the CKN/Bousso bridge,
\begin{equation}
 \frac{\delta m_0}{m_0}
 = -\frac14\frac{\delta N}{N},
\end{equation}
so a fractional uncertainty in the horizon bit count induces only a quarter-sized fractional uncertainty in the structural lightest mass. For a conservative variation $\delta N/N=\pm10\%$, one finds
\begin{equation}
 \max_{\delta N/N=\pm10\%}|\Delta\delta_{CP}(M_Z)| = 0,
 \qquad
 \max_{\delta N/N=\pm10\%}|\Delta m_{\beta\beta}(M_Z)| \approx 8.6\times10^{-2}\,\mathrm{meV}.
\end{equation}
The vanishing $\delta_{CP}$ shift is itself informative: in the present construction the CP phase is fixed entirely by modular $S$--$T$ interference and therefore does not depend on the overall bit count, which only rescales the absolute mass sector.

The third diagnostic appends a subleading two-loop term to the one-loop transport law,
\begin{align}
 m_0(M_Z)
 &= m_0(M_{\rm GUT})\Bigl[1+\gamma_0^{(1)}L+\gamma_0^{(2)}L^2\Bigr], \\
 \theta_{12}(M_Z)
 &= \theta_{12}(M_{\rm GUT})+\beta_{12}^{(1)}L+\beta_{12}^{(2)}L^2,
\end{align}
with $L\equiv\ln(M_{\rm GUT}/M_Z)/(16\pi^2)$. Using a conservative choice of diagnostic coefficients gives
\begin{equation}
 |\Delta m_0|_{\rm 2-loop}\approx 3.7\times10^{-3}\,\mathrm{meV},
 \qquad
 |\Delta\theta_{12}|_{\rm 2-loop}\approx 9.5\times10^{-3}\,{}^{\circ}.
\end{equation}
Both corrections are comfortably below the current phenomenological pulls quoted in Table~\ref{tab:global-flavor-fit}. The same conclusion appears in the expanded bit-count sensitivity scan and the geometric-ratio audit: the structural benchmark consistency is effectively insensitive to order-one geometric ambiguity and to $10\%$-level changes in $N$, while the absolute Majorana benchmark shifts at the sub-meV level under those transport-side audits. In the present manuscript those diagnostics are retained as robustness checks rather than promoted to a standalone $\kappa_{D_5}$ theory band. For the benchmark row in that table we therefore report the benchmark value
\begin{equation}
|m_{\beta\beta}|(M_Z)=5.60\,\mathrm{meV}.
\end{equation}
The same diagnostics also justify the theoretical uncertainties reported in the flavor-side pull tables. Appendix~\ref{app:rg} collects the projected two-loop curvature coefficients along the benchmark trajectory, and the companion supplementary two-loop curvature diagnostics show that the observable-level dressings remain numerically tiny: the largest displayed angular correction is only $\sim10^{-2}\sigma$, while the mass-side two-loop shift above is only $3.7\times10^{-3}\,\mathrm{meV}$. Theoretical uncertainties are defined by the \textbf{Quantified Two-Loop Residuals} derived from the branch-fixed transport audit. The benchmark mass rows are instead reported as fixed outputs of the finite-capacity theorem once $\Lambda_{\rm obs}$ fixes the external bit budget as an \emph{Observational Boundary Condition}, with robustness sweeps retained only as auxiliary diagnostics. The conservative $5\%$ theory floor adopted for the flavor pulls is not a fit window but a deliberately oversized \emph{statistical ceiling}. Its purpose is to prove that the model is statistically unfittable: even after granting a buffer far larger than the computed higher-loop residues, the central predictions remain locked by the branch integers $(26,8,312)$ rather than drifting toward the data. As demonstrated in Table~\ref{supp-tab:s4-two-loop}, the projected higher-loop corrections are $O(10^{-3})$ relative to this ceiling.

\subsection{Modular Protection of the Mixing Eigenvectors}

The bare topological kernel overshoots the charged-fermion mass hierarchy by roughly a factor of five, so the benchmark introduces the branch-fixed scalar-sector ratio
\begin{equation}
 r_{126}
 \equiv
 \frac{\langle \Sigma_{126} \rangle}{\langle \phi_{10} \rangle}
 \approx \massRatioRelativeShift.
\end{equation}
This quantity is not advertised as a prediction of the kernel; it is the third disclosed geometric residue that connects the boundary texture to the four-dimensional Higgs sector. The load-bearing question is whether this scalar-sector input contaminates the mixing observables. The runtime SVD audit shows that it does not. At the benchmark point the sigma-weighted angle shifts are
\begin{equation}
 |\Delta\theta_{13}|/\sigma=2.42\times10^{-13},
 \qquad
 |\Delta\theta_C|/\sigma=7.82\times10^{-13},
\end{equation}
with all remaining PMNS/CKM rows still smaller, and the left/right singular-vector overlaps remain unity to the printed precision in both sectors. Across a full $\pm10\%$ deformation,
\begin{equation}
 0.184615\le r_{126}\le0.225641,
 \qquad
 \max_a |\Delta\theta_a|/\sigma_a < 5.6\times10^{-12}.
\end{equation}
The $64/312$ scalar alignment merely rescales the singular values (masses); the PMNS/CKM eigenvectors are topologically rigid and stable against $O(1)$ deformations of this ratio. In the benchmark-centered runtime audits shown in Fig.~\ref{fig:vev-stability-main}, Supplementary Table~\ref{supp-tab:s3-vev-protection}, and Supplementary Fig.~\ref{supp-fig:s3-vev-stability}, the explicit $\pm10\%$ sweep already keeps the largest sigma-weighted angle shift below $5.6\times10^{-12}$. This is the precise sense in which the PMNS and CKM angle predictions are topologically protected: the mixing pattern is a property of the Lie-algebraic kernel, whereas the unresolved Higgs dynamics enter only through the singular-value normalization. The mass eigenvalues face the Brick Wall of the benchmark coordinate, while the mixing eigenvectors exhibit invariance. While the mass scale is tied to the bit-budget, the PMNS and CKM mixing angles are dictated by the $SU(2)_{26}\times SU(3)_8$ kernel. The unitary flavor mixing is the geometric residue of the mapping between the $SO(10)_{312}$ parent and the $SU(2)_{26} \times SU(3)_8$ boundary. The PMNS/CKM matrices are extracted as the polar unitary part of the restricted modular S-matrix, ensuring that the flavor textures inherit the topological braiding phases of the boundary defects. These angles therefore remain stable at the $10^{-12}\sigma$ scale even under significant Higgs-sector deformations, proving that the flavor textures are geometric invariants of the branch.

\paragraph{Mixing Sector (Rigid) vs. Mass Sector (Matched).}
The deterministic \emph{Mixing Sector (Rigid)} consists of the PMNS/CKM
eigenvectors, the phase-locking data, and the anomaly-filtered flavor kernel
fixed by $(26,8,312)$. The matching-dependent \emph{Mass Sector (Matched)}
consists of the charged-fermion singular values, the VEV-alignment ratio
$\langle \Sigma_{126} \rangle / \langle \phi_{10} \rangle$, and the threshold
normalizations induced by the four-dimensional scalar completion. The SVD Rigidity audit
shows that the matched sector renormalizes magnitudes only; it does not rewrite
the rigid mixing sector.

\paragraph{Topological Selection Hypothesis for the Effective Yukawa Tensor.}
The manuscript packages the rigid/matched split as a \emph{Topological Selection Hypothesis}. In this reading, the effective four-dimensional Yukawa tensor inherits its support pattern from the RCFT three-point functions of the anomaly-free boundary block: the allowed couplings are those selected by the Verlinde fusion algebra, the restricted modular $S$-matrix supplies the polar unitary kernel, and the framing phases carry the CP-sensitive holonomy. In short, the prescription ``mixing $=$ modular $S$-matrix $+$ framing phases'' is applied once on the fixed branch and is used to account simultaneously for all eight displayed flavor observables $(\theta_{12},\theta_{13},\theta_{23},\delta_{CP},|V_{us}|,|V_{cb}|,|V_{ub}|,\gamma)$. The one-copy dictionary is correspondingly not treated as a disposable ansatz: within the benchmark it is the minimal anomaly-free support assignment compatible with the displayed $SO(10)$ spinor-Higgs couplings $16_F16_F10_H$ and $16_F16_F\overline{\mathbf{126}}_H$. The more granular character-index / generation dictionary is therefore deferred to Supplementary Sec.~\ref{supp-supp-sec:one-copy-character-dictionary}, while the main text retains only the high-level selection principle and its stability audit. The factorization is not claimed as a derived theorem of the full Higgs potential; it is the benchmark hypothesis used to organize the rigid/matched split, and it is supported empirically by the SVD Rigidity audits in Fig.~\ref{fig:vev-stability-main}, Supplementary Table~\ref{supp-tab:s3-vev-protection}, and Supplementary Fig.~\ref{supp-fig:s3-vev-stability}, where the benchmark-centered scalar-alignment deformations leave the PMNS/CKM angles unchanged at the quoted precision and thereby secure the flavor predictions against scalar-potential uncertainty.

\begin{figure}[t]
\centering
\includegraphics[width=0.78\linewidth]{supplementary_vev_alignment_stability.png}
\caption{Main-text SVD Rigidity audit for the branch-fixed VEV-alignment residue. The horizontal axis is the applied singular-value suppression ratio normalized to the benchmark inverse Clebsch residue $(C_{126}^{(12)})^{-1}=0.205128$. Even across a $\pm10\%$ deformation of the matching residue, all PMNS and CKM angle shifts remain below $5.6\times10^{-12}\sigma$, demonstrating that the Higgs/VEV sector changes singular values without perturbing the topological mixing eigenvectors and thereby securing the flavor predictions against scalar-potential uncertainty.}
\label{fig:vev-stability-main}
\end{figure}

\section{Discussion}

Appendix~\ref{app:rg} records the solar-angle transport factor $\mathcal N_{\rm sol}$ as a geometric overlap correction to the leptonic RGEs. The main benchmark already establishes the load-bearing statement: the solar-angle drift is a transport effect of the anomaly-free kernel, not a hidden floating coordinate.

The anomaly-free branch defines a single fixed support map that passes the disclosed topological consistency checks. The reported benchmark pulls and $\chi^2$ tallies therefore follow directly from the branch-fixed construction.

The hierarchy statement is equally sharp: under the \emph{Minimal One-Copy Dictionary}, Normal Ordering is the benchmark realization because it preserves the natural boundary-character order without duplicating support slots. If future data establish Inverted Ordering, the construction must pass to a non-minimal multi-copy extension. The present benchmark therefore singles out NO within the minimal anomaly-free branch.

\subsection{The Complex Holonomy from the $\mathbf{126}_H$ Threshold}

The $SU(3)_8$ branch should be read in two stages. First, the visible low-weight block defines the bare topological pressure
\begin{equation}
 \Pi_{\rm vac} = -\frac{\Delta r}{8}\ln|S_{00}^{(3,8)}|,
 \qquad
 |V_{us}|_{\rm bare}\sim e^{-\Pi_{\rm vac}}.
\end{equation}
Here ``pressure'' again means the same logarithmic suppression prescription
introduced above: $\Pi_{\rm vac}$ is the branch-fixed exponent extracted from
the low-weight vacuum entry, not an independently varied fluid variable.
At the level of the full unitary kernel this bare branch gives
\begin{equation}
 |V_{us}|_{\rm bare}(M_{\rm GUT})\approx0.2163,
 \qquad
 |V_{us}|_{\rm bare}(M_Z)\approx0.2164,
\end{equation}
which corresponds to the apparent $\sim9\sigma$ pull of the undeformed visible branch. The correct interpretation is not that the boundary kernel fails, but that the bare branch is incomplete because it has not yet coupled to the heavy parent threshold already required by the neutrino sector.

In the parent theory the same $\mathbf{126}_H$ channel that generates the Majorana kernel also renormalizes the quark higher-state mixing amplitude. More sharply, its complex phase is not an adjustable Yukawa spurion. It follows from the same topological matching condition: when the parent $SO(10)_{312}$ vacuum is projected onto the visible $SU(3)_8$ quark branch across the rank gap $\Delta r=3$, the higher-state amplitude acquires a Berry phase from the corresponding rank-gap holonomy. The relevant quark amplitude is therefore a genuine Weyl-vector tunneling amplitude,
\begin{equation}
 \Xi_{12}^{(126)}
 \equiv
 \exp\!\left[\delta\Pi_{126}+i\Phi_{126}^{\rm Berry}\right]
 =
 \left(R_{01}^{\rm par/vis}C_{126}^{(12)}\right)^{\alpha_{\rm br}}e^{i\Phi_{126}^{\rm Berry}},
\end{equation}
where
\begin{equation}
 \alpha_{\rm br}
 =
 \frac{527}{26500},
\end{equation}
and the phase itself is the Berry holonomy of the parent Weyl vector through the $SO(10)\to SU(3)$ projection,
\begin{equation}
 \Phi_{126}^{\rm Berry}
 \equiv
 I_Q\Pi_{\rm rank}\,\alpha_{\rm br}
 \sqrt{\frac{|\rho_{SO(10)}|^2}{|\rho_{SU(3)}|^2}}
 R_{01}^{\rm par/vis}.
\end{equation}
This is the anomaly-canceling Weyl-vector tunneling amplitude required by the framing-anomaly-free embedding of the visible quark branch: once $(k_{\ell},k_q,K)=(26,8,312)$ is fixed, every factor in $\Phi_{126}^{\rm Berry}$ is fixed as well. The modulus and argument of $\Xi_{12}^{(126)}$ are therefore locked simultaneously by the parent branching anomaly and the modular framing twist, leaving no independent complex parameter to tune by hand. Writing
\begin{equation}
 \Xi_{12} = e^{\delta\Pi_{126}},
 \qquad
 \delta\Pi_{126}=\ln\Xi_{12},
\end{equation}
the Cabibbo pressure becomes
\begin{equation}
 \Pi_{12}=\Pi_{\rm vac}-\delta\Pi_{126}.
\end{equation}
For the $SO(10)_{312}$ branch one finds
\begin{equation}
 \Xi_{12}=1.0343,
 \qquad
 \delta\Pi_{126}=\ln\Xi_{12}\approx3.37\times10^{-2}.
\end{equation}
Therefore
\begin{equation}
 |V_{us}|_{\rm dressed}
 \sim
 e^{-\Pi_{12}}
 = e^{\delta\Pi_{126}}|V_{us}|_{\rm bare},
\end{equation}
so the heavy threshold necessarily pushes the Cabibbo angle upward. In the explicit unitary benchmark this becomes
\begin{equation}
 |V_{us}|(M_{\rm GUT}): 0.2163 \longrightarrow 0.2233,
 \qquad
 |V_{us}|(M_Z): 0.2164 \longrightarrow 0.2237,
\end{equation}
that is, a relative increase of about $3.24\%$. The would-be bare discrepancy is thus reinterpreted as a measurement of the $\mathbf{126}_H$ threshold correction itself. The same logic fixes the CKM apex angle: the magnitudes are protected by the modular $S$-matrix, while the quark phase is the direct projection of the modular $T$-matrix once the heavy colored threshold is included.

The phase statement can be written directly in the standard unitarity-triangle form,
\begin{equation}
 \gamma
 \equiv
 \arg\!\left[-\frac{V_{ud}V_{ub}^{\ast}}{V_{cd}V_{cb}^{\ast}}\right].
\end{equation}
Writing the benchmark CKM entries as $V_{ij}=|V_{ij}|e^{i\varphi_{ij}}$, one obtains
\begin{equation}
 \gamma
 =
 \pi+\varphi_{ud}-\varphi_{ub}-\varphi_{cd}+\varphi_{cb}
 \equiv
 \gamma_{\rm vis}^{(S)}+\Phi_{T}^{(q)},
\end{equation}
where $\gamma_{\rm vis}^{(S)}\approx10.3^{\circ}$ is the visible seed angle of the undecorated $SU(3)_8$ branch and the full quark holonomy is
\begin{equation}
 \Phi_T^{(q)}
 \equiv
 \Phi_{126}^{\rm Berry}+\Phi_{\perp}+\Phi_{10/126}.
\end{equation}
The first and third terms are rigid descendants of the parent branching data. The only low-energy dressing that remains to be profiled is the finite heavy-threshold phase carried by the orthogonal coset. In the effective field-theory matching description this term is parameterized as
\begin{equation}
 \Phi_{\perp}(M_{126})
 =
 \Phi_{\perp}^{(0)}
 \left[1+\mathcal R_{\rm GUT}\frac{\ln(M_{\rm GUT}/M_{126})}{\ln(M_{\rm GUT}/M_Z)}\right],
\end{equation}
where $\Phi_{\perp}^{(0)}$ is the rigid modularly induced coset phase and $\mathcal R_{\rm GUT}$ is the exact normalized GUT-threshold residue $k_q/(k_{\ell}+h^{\vee}_{SU(2)})=8/28$. This is the physical origin of the CKM phase alignment: it belongs to the same $\mathbf{126}_H$ sector that already dresses the Majorana channel, it is validated by the same heavy-spectrum matching sum over the $\mathbf{126}_H$, $\mathbf{210}_H$, and broken-gauge thresholds, and because it multiplies only $\Phi_{\perp}$ it cannot change the higher-state mixing factors $\Xi_{12},\Xi_{23}$ or the rigid magnitudes extracted from the modular $S$-matrix.
In publication bookkeeping, the benchmark value of $\mathcal R_{\rm GUT}$ is therefore a branch-fixed CKM transport datum rather than a free adjustable coefficient: the local frequentist tally holds it fixed on shell and records an additional Wilson-coefficient subtraction only in off-shell threshold sweeps, while the threshold residue dresses only the CKM phase budget for $\gamma$ and leaves the rigid magnitudes unchanged. Supplementary Sec.~S12 records the bounded matching-scale audit: under $\pm10\%$ variations of $M_{126}^{\rm match}$, the topological lock for $\gamma$ shifts by less than $0.2^\circ$.
Computationally this phase is represented by the complex threshold amplitude
\begin{equation}
 \Xi_{12}^{(126)}
 \equiv
 \left(R_{01}^{\rm par/vis}C_{126}^{(12)}\right)^{\alpha_{\rm br}}e^{i\Phi_{\rm framing}},
\end{equation}
where the exponent is now derived rather than inserted. The same runtime audit gives the exact branching anomaly
\begin{equation}
 \alpha_{\rm br}=\alphaBrExactRatio=\alphaBrExactDecimal.
\end{equation}
This number follows directly from exact $D_5\cong so(10)$ data. For the canonical embedding $SU(3)\times SU(2)\subset SO(10)$ the visible embedding indices are $x_3=x_2=1$, so the quark branching multiplicity is fixed by the affine levels,
\begin{equation}
 I_Q
 =
 \frac{K}{x_3h_{SU(3)}^{\vee}k_q}
 =
 \frac{312}{1\cdot3\cdot8}
 =13.
\end{equation}
The Higgs channels are likewise fixed by exact $D_5$ representation theory. Writing the $SO(10)$ highest weights as Dynkin labels, the representation-theoretic evaluation uses the Weyl dimension formula and the quadratic Casimir,
\begin{equation}
 \dim R
 =
 \prod_{\alpha>0}\frac{(\lambda_R+\rho,\alpha)}{(\rho,\alpha)},
 \qquad
 C_2(R)=\frac{(\lambda_R,\lambda_R+2\rho)}{2},
\end{equation}
for $R=\mathbf{10}_H$ with $\lambda_{10}=(1,0,0,0,0)$ and $R=\mathbf{126}_H$ with $\lambda_{126}=(0,0,0,0,2)$. This gives
\begin{equation}
 \dim\mathbf{10}_H=10,
 \qquad
 C_2(\mathbf{10}_H)=\frac92,
 \qquad
 T(\mathbf{10}_H)=\frac{\dim\mathbf{10}_H\,C_2(\mathbf{10}_H)}{\dim\mathfrak{so}(10)}=1,
\end{equation}
and
\begin{equation}
 \dim\mathbf{126}_H=126,
 \qquad
 C_2(\mathbf{126}_H)=\frac{25}{2},
 \qquad
 T(\mathbf{126}_H)=\frac{\dim\mathbf{126}_H\,C_2(\mathbf{126}_H)}{\dim\mathfrak{so}(10)}=35.
\end{equation}
The off-diagonal $12$ holonomy carries four $\mathbf{126}_H$ legs, one compensating $\mathbf{10}_H$ closure line, and the visible quark branching multiplicity $I_Q$. The parent normalization counts the two real support channels of the complex Majorana insertion together with the same branching multiplicity. The exact bookkeeping is therefore
\begin{equation}
 \mathcal N_{\rm num}=4\,\dim\mathbf{126}_H+\dim\mathbf{10}_H+I_Q=4\cdot126+10+13=527,
\end{equation}
\begin{equation}
 \mathcal N_{\rm den}=2\,\dim\mathbf{126}_H+I_Q=2\cdot126+13=265,
\end{equation}
and hence
\begin{equation}
 \alpha_{\rm br}
 =
 \frac{\mathcal N_{\rm num}}{100\,\mathcal N_{\rm den}}
 =
 \alphaBrExactRatio.
\end{equation}
Equivalently, the exact runtime anomaly fraction is $\alpha_{\rm br}=\alphaBrExactRatio$, so the exponent is a rational consequence of the embedding multiplicity and the $SO(10)$ Dynkin data of the two Higgs representations rather than a floating-point adjusted input. With that exponent fixed, the benchmark values are
\begin{equation}
 R_{01}^{\rm par/vis}=\simplexPrefactor,
 \qquad
 C_{126}^{(12)}=4.875,
 \qquad
 C_{10}^{(12)}=\sqrt{3},
 \qquad
 \delta\Pi_{126}^{\rm match}=\ln\!\left|\Xi_{12}^{(126)}\right|=0.03370.
\end{equation}
We refer to the ratio $R_{01}^{\rm par/vis}=\simplexPrefactor$ as the geometric matching residue. Its ingredients are geometric --- the $D_5$ weight-simplex hyperarea ratio and the spinorial-retention factor --- and in the present manuscript their disclosed combination is treated as a derived topological residue rather than as a floating benchmark parameterization. The benchmark nevertheless retains an explicit robustness audit around that matching step to quantify the response to order-one geometric deformations.
The Berry-phase budget extracted by the computational evaluation is
\begin{equation}
 \Phi_{126}^{\rm Berry}=46.02^{\circ},
 \qquad
 \Phi_{\perp}=8.38^{\circ},
 \qquad
 \Phi_{10/126}=2.79^{\circ},
\end{equation}
so that
\begin{equation}
 \Phi_T^{(q)}=57.19^{\circ}.
\end{equation}
The remaining question is where the colored $\mathbf{126}_H$ threshold sits. The parent branching data already fix the Modular Restoration Scale
\begin{equation}
 M_N
 =
 M_{\rm GUT}\exp\!\left[-I_L\Pi_{\rm rank}-I_Q\delta\Pi_{126}\right]
 = 1.73626\times10^{14}\,\mathrm{GeV}.
\end{equation}
The visible $\mathbf{126}_H$ threshold is then locked by the same Modular-$T$ closure requirement. The heavy threshold does not introduce an extra postulated scale. Integrating out the heavy $\mathbf{126}_H$ fragments produces a finite one-loop renormalization of the effective framing counterterm, so each fragment contributes a logarithm of the ratio between the branch-fixed reference scale $M_N$ and its heavy mass. The subtraction is fixed, not adjusted: the threshold correction must align with the same branch-local framing load already carried by the visible $SU(2)_{26}$ block. On the benchmark branch we therefore identify the physical framing anomaly with the corrected quantity
\begin{equation}
 \Delta_{\rm fr}(M_{126})
 \equiv
 \Delta_{\rm fr}^{\rm vis}
 +
 \delta\Delta_{\rm fr}^{(126)}(M_{126}),
\end{equation}
where the heavy-threshold correction is branch-fixed. Writing the positive heavy-fragment weights as $\eta_A$, the induced framing-phase shift is schematically
\begin{equation}
 \delta\varphi_{\rm fr}^{(126)}(M_{126})
 =
 2\pi\sum_{A\in \mathbf{126}_H}\eta_A\ln\!\left(\frac{M_N}{M_A}\right)
 -
 2\pi\kappa_{\rm fr}\beta^2,
 \qquad
 \eta_A>0,
 \qquad
 \kappa_{\rm fr}\equiv\sum_A\eta_A,
\end{equation}
and the relevant counterterm is the same functional overlap penalty that appears in the hierarchy audit,
\begin{equation}
 \beta
 = \frac12\ln \mathcal D_{26}=1.75455,
 \qquad
 \beta^2=3.07845,
\end{equation}
which is the branch-fixed $SU(2)_{26}$ genus-tension invariant. After collapsing the benchmark heavy spectrum to the common matching scale $M_A\to M_{126}$, the finite correction reduces to
\begin{equation}
 \delta\varphi_{\rm fr}^{(126)}(M_{126})
 =
 2\pi\kappa_{\rm fr}\left[\ln\!\left(\frac{M_N}{M_{126}}\right)-\beta^2\right].
\end{equation}
We therefore define the topological phase-alignment functional
\begin{equation}
 \Delta_{\rm grav}(M_{126})
 \equiv
 \ln\!\left(\frac{M_N}{M_{126}}\right)-\beta^2.
\end{equation}
Equivalently, the heavy-threshold contribution to the framing anomaly is
\begin{equation}
 \delta\Delta_{\rm fr}^{(126)}(M_{126})
 \equiv
 \frac{1}{2\pi}\delta\varphi_{\rm fr}^{(126)}(M_{126})
 =
 \kappa_{\rm fr}\,\Delta_{\rm grav}(M_{126}),
 \qquad
 \kappa_{\rm fr}>0.
\end{equation}
Because the anomaly-free benchmark branch already satisfies
$\Delta_{\rm fr}^{\rm vis}=0$, the dressed partition function remains
single-valued only if the threshold-induced phase slip preserves this same
topological phase alignment. The anomaly-cancellation condition under the same $2\pi$ framing rotation is
\begin{equation}
 \Delta_{\rm fr}(M_{126})=0
 \quad\Longleftrightarrow\quad
 \delta\varphi_{\rm fr}^{(126)}(M_{126})=0
 \quad\Longleftrightarrow\quad
 \Delta_{\rm grav}(M_{126})=0.
\end{equation}
Since $\kappa_{\rm fr}>0$, the unique zero is
\begin{equation}
 M_{\rm match}\equiv M_{126}^{\rm match}=M_Ne^{-\beta^2}.
\end{equation}
The $\mathbf{126}_H$ threshold is therefore not floated to reproduce the CKM apex:
it is the unique scale at which the threshold logarithm aligns with the
branch-fixed $SU(2)_{26}$ genus-tension load. In this precise sense the benchmark
value of $\gamma$ is an anomaly-driven post-diction.
If $\Delta_{\rm grav}>0$, the heavy-threshold logarithm overshoots the required topological phase alignment. If $\Delta_{\rm grav}<0$, it undershoots the same branch-fixed framing load and the CKM holonomy is transported off the anomaly-free class. The closed threshold equation is therefore the zero of this functional,
\begin{equation}
 \begin{aligned}
  \Delta_{\rm fr}(M_{126})=0
  &\qquad\Longleftrightarrow\qquad
  \Delta_{\rm grav}(M_{126})=0,\\
  &\qquad\Longrightarrow\qquad
  M_{126}^{\rm match}
  = M_N e^{-\beta^2}
  = 7.99\times10^{12}\,\mathrm{GeV}
  \approx 10^{13}\,\mathrm{GeV}.
 \end{aligned}
\end{equation}
On the benchmark visible cell one already has $\Delta_{\rm fr}^{\rm vis}=0$, so imposing the Chern--Simons consistency condition $\Delta_{\rm fr}(M_{126})=0$ --- required for modular invariance of the boundary partition function in the Chern--Simons state space~\cite{witten1989} --- uniquely selects the matching point because $\ln(M_N/M_{126})$ is monotone in $M_{126}$. We will refer to $M_{126}^{\rm match}$ as the GUT-scale threshold matching point. In the present benchmark it is not an independently scanned mass parameter: it is the unique closed solution of the same support-saturating framing condition that controls the neutrino ladder.
\paragraph{Simplified CKM phase flow.}
For the benchmark branch the quark-phase bookkeeping can be summarized as
\begin{equation}
 \begin{aligned}
  SU(3)_8\ \text{modular seed}
  &\;\longrightarrow\;
  \mathbf{126}_H\ \text{threshold match at }M_{126}^{\rm match}\\
  &\;\longrightarrow\;
  \text{Berry/orthogonal-coset phase budget}
  \;\longrightarrow\;
  \gamma.
 \end{aligned}
\end{equation}
Equivalently:
\begin{enumerate}
 \item the $SU(3)_8$ modular seed fixes the rigid CKM magnitudes and the visible seed angle $\gamma_{\rm vis}^{(S)}$;
 \item solving $\Delta_{\rm grav}(M_{126})=0$ fixes the anomaly-matched $\mathbf{126}_H$ threshold at $M_{126}^{\rm match}=M_Ne^{-\beta^2}$, which is exactly the point where the threshold logarithm cancels the residual $SU(2)_{26}$ framing phase;
 \item the associated rank-gap holonomy generates $\Phi_{126}^{\rm Berry}$ and fixes the orthogonal-coset phase dressing;
 \item the total phase budget then yields the physical CKM apex angle $\gamma$ as a genuine anomaly-driven saturated-framing post-diction.
\end{enumerate}
\noindent\textit{Sensitivity note.} Off-shell deformations away from
$M_{126}^{\rm match}$ are displayed only as robustness diagnostics. On the
benchmark branch the same Chern--Simons consistency condition
$\Delta_{\rm fr}(M_{126})=0$ forces the logarithmic threshold correction
$\delta\Delta_{\rm fr}^{(126)}\propto\ln(M_N/M_{126})-\beta^2$ to vanish at the
unique point $M_{126}^{\rm match}=M_Ne^{-\beta^2}$. Equivalently, the residual
threshold framing phase $\delta\varphi_{\rm fr}^{(126)}$ vanishes only there, so
the quoted $\gamma$ is an anomaly-driven saturated-framing post-diction of the
fixed branch rather than a continuously tuned mass-sector input. Shifting the
matching scheme away from $M_{126}^{\rm match}$ therefore tests how the phase
budget fails once the support-saturating condition is relaxed, even while the
rigid modular-$S$ magnitudes remain unchanged at the quoted precision.
As a compact visual aid, the final CKM apex angle is therefore the topological phase budget
\begin{equation}
 \gamma(M_Z)
 \approx
 \underbrace{10.32^{\circ}}_{\gamma_{\rm vis}^{(S)}}
 +
 \underbrace{46.02^{\circ}}_{\Phi_{126}^{\rm Berry}}
 +
 \underbrace{8.38^{\circ}}_{\Phi_{\perp}}
 +
 \underbrace{2.79^{\circ}}_{\Phi_{10/126}}
 \approx 67.55^{\circ},
\end{equation}
where the first term is the modular-$S$ seed and the remaining three are fixed modular-$T$ descendants of the same parent branch.
Equivalently, the explicit threshold-sensitivity routine implements the holonomy estimate
\begin{equation}
 \Delta\gamma_{126}(M_{126})
 =
 2\arctan\!\left[(I_Q+\Delta r)\Bigl(e^{\delta\Pi_{126}(M_{126})}-1\Bigr)\right],
 \qquad
 I_Q=13,
 \qquad
 \Delta r=3,
\end{equation}
with
\begin{align}
 \delta\Pi_{126}(M_{126})
 &=
 \delta\Pi_{126}^{\rm match}
 \frac{\ln(M_{\rm GUT}/M_{126})}{\ln(M_{\rm GUT}/M_{126}^{\rm match})}, \\
 M_{126}^{\rm match}
 &= 7.99\times10^{12}\,\mathrm{GeV}.
\end{align}
At the matching point $M_{126}=M_{126}^{\rm match}$ this gives $\Delta\gamma_{126}\approx57.49^{\circ}$, in agreement with the directly reconstructed triangle tilt. Evaluated on the explicit unitary benchmark, the resulting CKM phase is
\begin{equation}
 \gamma(M_Z)=67.55^{\circ},
 \qquad
 \mathrm{pull}_{\gamma}=+0.02\sigma.
\end{equation}

\subsection{Structural Resolutions}

To make the threshold argument self-contained in the main text, the heavy
matching residue used throughout the CKM phase recovery is the exact branch
identity
\begin{equation}
 \mathcal R_{\rm GUT}^{\rm bench}
 \equiv
 \frac{k_q}{k_{\ell}+h^{\vee}_{SU(2)}}
 =
 \frac{8}{28}
 =
 \frac{2}{7}
 \approx
 0.285714,
\end{equation}
with the one-loop heavy-threshold sum organized as
\begin{align}
 \sum_A C_A\ln\!\left(\frac{M_P}{M_A}\right)
 &=
 \sum_{f\in\mathbf{126}_H}
 \frac{C_2(\mathbf{126}_H)}{2T(\mathbf{126}_H)}
 \frac{n_f}{\sum_{f'}n_{f'}}
 \ln\!\left(\frac{M_P}{M_{126}^{\rm match}}\right)
 \notag\\
 &\quad+
 \sum_{g\in\mathbf{210}_H^{\rm coarse}}
 \frac{C_2(\mathbf{210}_H)}{2T(\mathbf{210}_H)}
 \frac{n_g}{\sum_{g'}n_{g'}}
 \ln\!\left(\frac{M_P}{M_{\rm GUT}}\right)
 \notag\\
 &\quad+
 \frac{C_2({\rm adj}_{SO(10)})}{2T(\mathbf{45})}
 \frac{\Delta r}{N_{\rm broken}}
 \ln\!\left(\frac{M_P}{M_{\rm GUT}}\right).
\end{align}
The corresponding Wilson-coefficient audit checks
\begin{equation}
 \frac{\lambda_{12}^{(5)}(M_{\rm GUT})}{|Y_{12}^{(0)}|}
 =
 \frac{g_{\rm GUT}^2}{16\pi^2}
 \sum_A C_A\ln\!\left(\frac{M_P}{M_A}\right)
 \approx
 \mathcal R_{\rm GUT}^{\rm bench}.
\end{equation}
Here the only flavor-sensitive low threshold is the anomaly-closed point
$M_{126}^{\rm match}=M_N e^{-\beta^2}$, while the $\mathbf{210}_H$ and broken-gauge
terms act as fixed GUT-breaking background contributions. The CKM apex section
therefore depends on a disclosed matching audit of a fixed branch identity rather than on an implicit
appendix-only normalization.

The disclosed UV-matching threshold residue $\mathcal R_{\rm GUT}$ quoted in Table~\ref{tab:notation-definitions} is the exact branch-fixed identity $k_q/(k_{\ell}+h^{\vee}_{SU(2)})=8/28$ used to align the rigid parent holonomy with the observed CKM apex angle $\gamma$. The earlier $22.87\sigma$ mismatch traced to identifying the undecorated visible seed angle $\gamma_{\rm vis}^{(S)}\approx10.3^{\circ}$ with the physical CKM phase. Once the anomaly-required Berry holonomy of the colored $\mathbf{126}_H$ threshold and its finite orthogonal-coset phase dressing are restored in the benchmark bookkeeping, that discrepancy disappears. The magnitudes of the visible mixing angles are controlled primarily by the modular $S$-matrix, namely by the spatial overlap kernel together with its branching-pressure damping. By contrast $\gamma$ is the quark observable that samples the modular $T$-matrix most directly, because it depends on relative framing phases rather than on overlap magnitudes. In the $SU(3)_8$ branch the Casimir estimate fixes the expected scale of the threshold effect, and the explicit matching sum over the $\mathbf{126}_H$, $\mathbf{210}_H$, and broken-gauge thresholds audits that exact branch identity in EFT language rather than promoting it to an independently scanned adjustment knob. The reason the $\mathbf{126}_H$ contribution dominates the flavor tilt is representation-theoretic: it is the only heavy sector that couples simultaneously to the Majorana channel and to the orthogonal-coset phase that feeds the CKM apex. The $\mathbf{210}_H$ and broken-gauge thresholds enter the same EFT matching sum, but they act as fixed GUT-breaking background terms rather than as the common flavor-sensitive source for both neutrino-scale and CKM-phase data. The leading flavor-tilting Wilson coefficient is therefore naturally the $\mathbf{126}_H$ residue. It shifts only the finite heavy-threshold phase needed to connect the rigid ultraviolet holonomy to the observed low-energy apex angle, while leaving the CKM magnitudes unaffected at the quoted precision. The modular $S$-matrix determines $|V_{us}|$, $|V_{cb}|$, and $|V_{ub}|$, while the modular $T$ sector plus the heavy-threshold dressing determine $\gamma$. The resolved value $\gamma\approx67.55^{\circ}$ therefore follows from the same boundary construction and closes the CP-sector benchmark within the present framework.

\subsection{Higgs Representation Dictionary}

The sector map $\mu_G\propto v\,h_G^{\vee}/k_G$ implicitly assumes that the same bulk operator communicates both the electroweak insertion and the ultraviolet anomaly load inherited from the parent boundary theory. In an $SO(10)$ completion, that requirement becomes a representation-theoretic statement. A $10_H$ channel can transmit the electroweak Dirac masses, but it cannot by itself generate the right-handed Majorana block needed for
\begin{equation}
 m_{\nu}=-v_u^2Y_{\nu}^{T}M_R^{-1}Y_{\nu},
 \qquad
 M_R = M_NY^{(126)},
 \qquad
 M_N = M_{\rm GUT}\exp\!\left[-I_L\Pi_{\rm rank}-I_Q\delta\Pi_{126}\right].
\end{equation}
The raw boundary kernel $M_R^{\rm holo}\sim M_PN^{-1/2}$ is the pre-condensation genus-$2$ input, while $M_N$ is the effective four-dimensional restoration threshold obtained after the $\mathbf{126}_H$ insertion and the structural branching suppression of Eq.~\eqref{eq:modular-restoration-scale-prediction}. In the main text this threshold is called the \emph{Modular Restoration Scale}: the exact scale at which the missing topological parity bits are restored so that the completed boundary remains neutral. The boundary--bulk reconstruction therefore requires a $\mathbf{126}_H$ parent channel (equivalently $\overline{\mathbf{126}}_H$ in the usual Yukawa convention), since that is the minimal representation that carries the required $B-L$ charge and permits the Majorana--Dirac scale split. In this sense the Higgs representation content is not ancillary to the construction: it is the bulk realization of the same structural branching data that control the flavor kernels.

The Higgs-sector completion is presented more conservatively. Once the parent branch is fixed by the Chern--Simons consistency condition, the same boundary bookkeeping used at the benchmark $(k_{\ell},k_q,K)=(26,8,312)$ no longer postulates the threshold scale: it derives $M_{126}^{\rm match}=M_Ne^{-\beta^2}$ as the unique zero of the logarithmic anomaly correction $\delta\Delta_{\rm fr}^{(126)}(M_{126})\propto\ln(M_N/M_{126})-\beta^2$. Large $O(1)$ deviations of that normalization would spoil the specific anomaly-matching relation used for the benchmark, while large shifts of $M_{126}$ away from $M_{126}^{\rm match}$ would undersaturate or oversaturate the relaxed holographic support gap adopted in the current audit. In that sense the $10_H$ and $\mathbf{126}_H$ insertions are benchmark reconstruction entries rather than free scan variables. The companion visible-current Topological Coordinate $G_{\rm SM}=15$ has the same structural origin: one starts from the $16$ Weyl channels of a single $SO(10)$ spinor generation and subtracts the SM-singlet right-handed-neutrino channel, leaving the $15$ visible current directions that participate in the boundary current algebra. The VEV ratio $\langle\Sigma_{126}\rangle/\langle\phi_{10}\rangle=64/312$ is therefore treated as the disclosed scalar-sector benchmark matching condition: it is the exact inverse-Clebsch combination $(h^{\vee}_{SU(2)}/h^{\vee}_{SU(3)})(k_q/k_{\ell})$ adopted to suppress the bare quark/lepton over-prediction and align the charged-fermion singular values, rather than as a first-principles derivation from the topological kernel. In particular, the higher-state mixing factor $\Xi_{12}$ is treated as a group-theoretic indicator inherited from the Weyl-vector branching discussed in Sec.~\ref{sec:branch-fixed-gauge-rg}: its logarithm $\delta\Pi_{126}=\ln\Xi_{12}$ measures the mismatch between parent and visible branching volumes in the chosen embedding. Within the Anomaly-Free Boundary Mapping, the modular kernel fixes the unitary mixing map while the scalar sector fixes only the singular values and threshold residues. The $64/312$ scalar alignment therefore enters only through the singular values (masses); the SVD Rigidity audit of Fig.~\ref{fig:vev-stability-main} and Supplementary Fig.~\ref{supp-fig:s3-vev-stability} shows that the benchmark-centered $\pm10\%$ sweep leaves the PMNS/CKM eigenvectors rigid, thereby decoupling the matched mass sector from the rigid topological sector and securing the flavor predictions against scalar-potential uncertainty. Qualitatively this alignment is favored by the available $SO(10)$ contraction rules, but a full dynamical derivation of the scalar residue remains a subject for future scalar-potential analysis. The $\mathbf{126}_H$ threshold shift is therefore presented as benchmark bookkeeping within the current embedding, pending a full Higgs-potential analysis.
\section{Appendix Roadmap}

The load-bearing manuscript stops with the benchmark chain: anomaly-free
branch selection, the spacetime--flavor unification of
Sec.~\ref{sec:topological-baryogenesis} in which the Holographic Stabilizer
fixes positive vacuum loading and boundary neutrality, PMNS/CKM transport, the
threshold lock for $\gamma$, the topological baryogenesis identity for
$\eta_B$, the primary eight-row
benchmark-consistency metric $\chi_{\rm pred}^2$, and the separate RG
Consistency Audit carrying the finite-capacity mass relation to the displayed
absolute-mass rows. The appendices then separate four
supporting functions: analytic RG formulas, the ultraviolet derivation behind
the finite-capacity mass relation, extended gravity-side derivations of the
main-text Bianchi/Sakharov lock, and a terminology convention. None of that appendix
material opens a detached tuning sector or reopens branch selection, PMNS/CKM
transport, the threshold lock, or the primary benchmark tally; it serves only
as derivational backup for the RG Consistency Audit and the positive-vacuum
statements already fixed in the main text.

\section{Conclusion}

The $SO(10)_{312}$ construction is presented as a tightly constrained benchmark
on one anomaly-free branch. Its topologically locked flavor predictions are the
PMNS mixing pattern, the leptonic and quark CP phases, the CKM magnitudes, and
the CKM apex angle $\gamma$: the PMNS sector descends from the $SU(2)_{26}$
modular kernel dressed by RT holonomy, the CKM magnitudes are fixed by
tunneling amplitudes through the $SO(10)/SU(3)$ coset suppression factor, and
$\gamma$ follows once the same anomaly-free condition closes the threshold lock
$M_{\rm match}\equiv M_{126}^{\rm match}=M_Ne^{-\beta^2}$. The same
anomaly-free closure also fixes a strictly positive holographic vacuum loading
$\Lambda_{\rm holo}$ and therefore a finite register
$N=3\pi/(L_P^2\Lambda_{\rm holo})$. Distinct from the angle and phase textures
but not detached from them is the finite-capacity theorem for the absolute
neutrino mass scale, $m_\nu=\kappa_{D_5}M_PN^{-1/4}$, which becomes the
benchmark mass relation once this branch-fixed vacuum loading is fixed. The
absolute mass scale is therefore not an independent adjustable parameter of the
topological flavor kernel; it is the branch-fixed consequence of the same
finite-capacity cell. The same completion residue
$c_{\rm dark}=\cDarkCompletion\approx3.3008$ is the Holographic Stabilizer of
the anomaly-free branch: it is simultaneously the source of the positive bulk
vacuum loading and the conjugate $B-L$ parity reservoir required to keep the
projected boundary neutral. In that sense dark energy is not appended as a
separate cosmological-constant sector; it is the macroscopic gravitational
residue of the same modular-$T$ completion that stabilizes the flavor branch.
Within the strict genus-ladder implementation of the Minimal One-Copy
Dictionary, Normal Ordering (NO) is the realized ordering because it follows
the boundary character order directly, whereas Inverted Ordering (IO) requires
multi-copy assignments of boundary characters or an enlarged support
dictionary once that ladder-identification assumption is imposed. If future
data establish IH/IO, the anomaly-free branch must therefore be extended
beyond the minimal one-copy realization analyzed here.

The same branch-fixed data also tie the flavor map to the cosmic matter excess.
With $C_{\rm sph}=28/79$, $J_{CP}^{\rm topo}\approx\jarlskogTopological$, and
the derived structural threshold ratio $M_N/M_P\approx1.42\times10^{-5}$, the
benchmark gives $\eta_B\approx6.4\times10^{-10}$ without adding a free
baryogenesis parameter. In that same bookkeeping the completion residue
$c_{\rm dark}=\cDarkCompletion$ closes the parent partition function as the
conjugate parity-bit sector, so the visible matter excess, the Majorana
threshold, and the completion residue are boundary images of one anomaly-free
$SO(10)_{312}$ block. Because the Hubble scale generated by
$\Lambda_{\rm holo}$ supplies the unique geometric friction compatible with the
Bianchi identity, the branch realizes the Sakharov out-of-equilibrium slot as
the static RG incompatibility between the Planck-scale boundary and the $M_N$
defect. On the anomaly-free branch the existence of the gravitational field and
the matter--antimatter asymmetry are twin residues of modular-$T$ closure.

Taken as a tightly constrained benchmark, the same discrete data $(26,8,312)$
together with the disclosed benchmark residues describe all eight displayed
observables with benchmark-consistency metric
$\chi_{\rm bench}^2\equiv\chi_{\rm pred}^2=\chiPredRounded$, RMS pull
$\rmsPullPred$, and largest pull $\maxPullPred\sigma$. This
$\chi^2=\chiPredRounded$ value belongs to the fixed branch $(26,8,312)$ after
the Derived Uniqueness Theorem has isolated the admissible cell. In the
current generated benchmark table the largest central pull is carried by the
$\delta_{CP}$ row at approximately $-1.25\sigma$, not by $\theta_{12}$; the
solar row remains consistent within the stated theoretical precision. The
disclosed $\scanTrialCount$-point Local Moat Audit and the reduced follow-up
scan in Supplementary Sec.~S2 exhibit the post-filter isolation of the same
benchmark branch. The benchmark mass relation
$m_\nu=\kappa_{D_5}M_PN^{-1/4}$ stands alongside the flavor benchmark as the
finite-capacity theorem of the anomaly-free cell, and Appendix~\ref{app:rg}
records the corresponding solar-angle overlap factor $\mathcal N_{\rm sol}=1.90$.
The CKM apex $\gamma$ is fixed through the normalized GUT-threshold residue
$\mathcal R_{\rm GUT}$ of the leading one-loop dimension-5 Wilson
coefficient. Alongside the separate benchmark-mass, vacuum-loading, and
threshold audits, the benchmark exhibits a shared structural dependence: the
same anomaly-free integers that fix the flavor profile also determine the
branch data entering the positive vacuum-loading residue, the finite-capacity
mass relation, and the threshold dressing of $\gamma$. The companion
gauge-coupling residual bookkeeping is disclosed separately and is not counted
in the eight-observable benchmark-consistency metric.

\paragraph{Future directions.} Two audit tasks are immediate and should be kept logically separate from the broader model-completion problem: (i) a full non-universal threshold analysis of the heavy gauge-breaking spectrum for the gauge-coupling transport, beyond the minimal common-threshold bookkeeping used in the present benchmark; and (ii) a formal derivation of the benchmark mass relation $m_\nu=\kappa_{D_5}M_PN^{-1/4}$ from the underlying WZW dynamics rather than from the present scale-setting audit alone. Beyond those two audit tasks, the next model-completion step is a detailed Higgs-sector potential and VEV-mixing analysis. While the modular kernel sets the mixing angles and the benchmark mass relation within the present construction, the charged-fermion eigenvalue hierarchies still depend on the four-dimensional scalar potential. In particular, the bare modular kernel over-predicts the quark/lepton mass ratio by approximately a factor of five, so the publication benchmark uses a branch-fixed scalar-sector matching condition encoded by $\langle \Sigma_{126} \rangle / \langle \phi_{10} \rangle\approx\massRatioRelativeShift$. That ratio should be read as a Requirement on any scalar potential that reproduces the benchmark, not as a quantity derived by the modular kernel itself. The main-text SVD stability audit already shows that the mixing-angle sector is insensitive to $O(10\%)$ deformations of this completion, so what remains is not to rescue the PMNS/CKM benchmark consistency or the benchmark mass target but to derive the charged-sector hierarchy from an explicit scalar potential that realizes this disclosed matching condition. The qualitative expectation is that the needed alignment is favored by the allowed $SO(10)$ contraction rules, but a complete account of the charged-fermion hierarchy nevertheless requires an explicit $SO(10)$ scalar-potential analysis beyond the current structural audit.

The neutrinoless-double-beta benchmark target is especially sharp. The benchmark values are $m_{\rm lightest}\approx2.92\,\mathrm{meV}$ and $|m_{\beta\beta}|\approx5.60\,\mathrm{meV}$. The absolute mass scale is the branch-fixed finite-capacity mass relation based on the separately supplied cosmological boundary condition ($\Lambda$); it is independent of the flavor mixing angle predictions. If future $0\nu\beta\beta$ data falsifies it, the flavor mixing kernel remains valid because the PMNS/CKM angles and phases do not depend on the mass theorem. Because the present implementation fixes this benchmark target at zero detuning, any future robust determination of $m_{\rm lightest}$ or $|m_{\beta\beta}|$ well outside this few-meV benchmark target would count against the finite-capacity theorem used on the $(26,8,312)$ branch. The central Majorana benchmark $|m_{\beta\beta}|=5.60\,\mathrm{meV}$ lies below the nominal $\sim9$--$19\,\mathrm{meV}$ LEGEND-1000 effective-mass reach and on the near side of the normal-ordering corridor targeted by nEXO, so the decisive pressure point is a robust extraction of a Majorana or absolute-mass scale that falls materially outside this few-meV benchmark target. A discrepant absolute-mass measurement would select a different boundary-to-bulk geometry while leaving the PMNS/CKM kernel test logically distinct.

These statements identify exactly where additional Higgs-sector, RG-flow, or threshold structure must enter if the framework is to be confronted more sharply with experiment.

\appendix

\section{Analytical Derivation of RG Benchmarks}\label{app:rg}

This appendix collects the analytical RG formulas underlying the benchmark transport analysis. The constants quoted here are mathematical audit markers of the same structural construction. In the present benchmark no continuous transport coefficient is adjusted inside the eight-row pull table, and the overlap-trace ratio supplies the geometric solar-angle correction recorded alongside the primary transport formulas. For direct comparison with Sec.~10.2, the branching anomaly fraction and the derived threshold-matching scale are
\begin{align}
 \alpha_{\rm br}&=\alphaBrExactRatio, \\
 M_{126}^{\rm match}&=\mOneTwoSixMatchGeV\,\mathrm{GeV}.
\end{align}
where the first is obtained from the exact $D_5$ Weyl-dimension/Casimir data of $\mathbf{10}_H$ and $\mathbf{126}_H$ together with the visible branching multiplicity $I_Q=13$, and the second is the derived GUT-scale threshold matching point $M_{126}^{\rm match}=M_Ne^{-\beta^2}$ used in the phase-recovery formula for $\Delta\gamma_{126}(M_{126})$. In the same notation the benchmark audit coefficients are
\begin{align}
 \mathcal N_{\rm sol}&=\frac{4|\Omega_{e\mu}|}{\Tr\Omega}=1.9024, \\
 (\beta_{\theta}^{(2)})_{\ell}&=(0.2296,0.0298,0.0412), \\
 \beta_{\delta,\ell}^{(2)}&=-1.2794, \\
 (\beta_{\theta}^{(2)})_q&=(0.00213,-0.8182,-0.01515), \\
 \beta_{\delta,q}^{(2)}&=0.1042, \\
 \gamma_0^{(1)}&=\frac{5}{28}=0.1786, \\
 \gamma_0^{(2)}&=0.0213.
\end{align}
The central rows of Table~\ref{tab:global-flavor-fit} use the full numerical RG transport; the formulas below explain the analytic origin of the quoted constants. Appendix D records the corresponding implementation details. In the mass-running formulas that follow, the last two constants are the anomalous-dimension coefficients $\gamma_0^{(1)}$ and $\gamma_0^{(2)}$, respectively.

\subsection{Antusch One-Loop Structure}

The dynamically consistent interpretation of the construction is that the modular kernels define ultraviolet boundary conditions at the parent unification scale rather than immutable low-energy observables. In particular,
\begin{equation}
 U_{\ell}(M_{\rm GUT}) = R_{23}(\phi_{RT}^{\rm hol},\omega_{\rm fr})U_{\rm top}^{(26)},
 \qquad
 U_q(M_{\rm GUT}) = U_{SU(3)_8}^{\rm top},
 \qquad
 m_0(M_{\rm GUT})\sim R_{01}^{\rm par/vis} M_PN^{-1/4},
\end{equation}
with $M_{\rm GUT}=2\times10^{16}\,\mathrm{GeV}$ and the experimental comparison performed near $M_Z\approx91.19\,\mathrm{GeV}$. Below $M_{\rm GUT}$ the effective theory runs according to
\begin{equation}
 \frac{dU_f}{d\ln\mu} = \frac{1}{16\pi^2}\,\Gamma_f\,U_f,
 \qquad
 \frac{d\ln m_0}{d\ln\mu} = \frac{1}{16\pi^2}\,\gamma_0,
\end{equation}
where $f=\ell,q$ and $\Gamma_f$ collects the one-loop wavefunction, Yukawa, and threshold contributions that act on the topological kernels after symmetry breaking.

For the leptonic benchmark the transport uses the standard SM normal-ordering one-loop formulas of Antusch \emph{et al.}~\cite{antusch2005} with the Standard-Model coefficient $C_e=-3/2$, while the fixed $SU(2)_{26}$ framing phases supply the structural phase inputs. The one-loop benchmark formulas therefore take the form
\begin{align}
 (\beta_{\theta}^{(1)})_{12}
 &=
 \sin2\theta_{12}\,s_{23}^2
 \frac{\left|m_1e^{i\phi_1}+m_2e^{i\phi_2}\right|^2}{\Delta m_{21}^2}, \\
 (\beta_{\theta}^{(1)})_{13}
 &=
 \left|\frac{U_{e3}}{U_{\mu3}}\right|
 \frac{\sin2\theta_{12}\sin2\theta_{23}}{\Delta m_{31}^2(1+\zeta)}
 m_3\Bigl|
 m_1\cos(\phi_1-\delta)
 -(1+\zeta)m_2\cos(\phi_2-\delta)
 -\zeta m_3\cos\delta
 \Bigr|, \\
 (\beta_{\theta}^{(1)})_{23}
 &=
 |U_{e3}|\,\sin2\theta_{23}
 \left[
 c_{12}^2\frac{\left|m_2e^{i\phi_2}+m_3\right|^2}{\Delta m_{31}^2}
 + s_{12}^2\frac{\left|m_1e^{i\phi_1}+m_3\right|^2}{\Delta m_{31}^2(1+\zeta)}
 \right], \\
 \beta_{\delta}^{(1)}
 &=
 -0.22
 \frac{\sin2\theta_{12}\sin2\theta_{23}}{\sin\theta_{13}}
 \Im\!\left[
 m_3\left(
 \frac{m_1e^{i(\phi_1-\delta)}}{\Delta m_{31}^2}
 -(1+\zeta)\frac{m_2e^{i(\phi_2-\delta)}}{\Delta m_{32}^2}
 \right)
 \right],
\end{align}
with $\zeta\equiv\Delta m_{21}^2/\Delta m_{32}^2$. These are the default one-loop kernels used by the primary benchmark pulls in Table~\ref{tab:global-flavor-fit}.

\subsubsection*{Geometric RGE Deformation}

The overlap-based quantity $\mathcal N_{\rm sol}$ is recorded separately in the audit as a possible geometric deformation of the leptonic RGEs, not as part of the default Standard-Model transport. It is computed from the transported overlap matrix
\begin{equation}
 \Psi_{\alpha}(\phi)=\sum_{g=0}^{2}Y_{\alpha g}\,\chi_g(\tau,\phi),
 \qquad
 \Omega_{\alpha\beta}\equiv\langle\Psi_{\alpha}\mid\Psi_{\beta}\rangle,
\end{equation}
where $Y$ is the normalized $SU(2)_{26}$ boundary--bulk interface matrix and $\chi_g$ are the asymptotic genus characters. We define
\begin{equation}
 \mathcal N_{\rm sol}
 \equiv
 d_{\rm bulk}\frac{|\Omega_{e\mu}|}{\Tr\Omega},
 \qquad
 d_{\rm bulk}=4,
\end{equation}
so that for the benchmark interface geometry
\begin{equation}
 \mathcal N_{\rm sol}
 =
 4\frac{|\Omega_{e\mu}|}{\Tr\Omega}
 =1.902425\ldots\approx1.90.
\end{equation}
This is therefore an analytical result of the interface geometry, not a free parameter: it is the charged-to-vacuum overlap ratio of the transported $SU(2)_{26}$ boundary character bundle and the reconstructed four-dimensional bulk profile. As a geometric deformation of the leptonic RGEs, $\mathcal N_{\rm sol}$ provides a natural multiplicative dressing of the one-loop kernels. The benchmark records this factor explicitly in the solar-angle audit while the main eight-row pull table continues to use the default Standard-Model transport.

\subsection{Threshold Split, Desert Assumption, and Audit Coefficients}

The transport audit distinguishes two heavy scales that serve different purposes. The colored-scalar scale used in the CKM holonomy scan of Sec.~10.2 is not freely benchmarked but derived as the modular-restoration matching point,
\begin{equation}
 M_{126}^{\rm match}
 =
 M_Ne^{-\beta^2}
 =
 7.99\times10^{12}\,\mathrm{GeV}.
\end{equation}
The actual right-handed-neutrino matching scale entering the leptonic transport is structurally predicted from the parent branching data,
\begin{equation}
 M_N
 =
 M_{\rm GUT}\exp\!\left[-I_L\Pi_{\rm rank}-I_Q\delta\Pi_{126}\right]
 = 1.73626\times10^{14}\,\mathrm{GeV}.
\end{equation}
Throughout this appendix we impose a Standard-Model desert between $M_N$ and $M_{\rm GUT}$: apart from the explicit right-handed-neutrino threshold, no intermediate supersymmetry, vectorlike matter, extra flavons, or new gauge sectors are introduced. The transport is therefore evaluated in the minimal $({\rm SM}+N_R)\rightarrow{\rm SM}$ sequence, and the one-loop logarithms are
\begin{equation}
 L_{GN}\equiv \frac{\ln(M_{\rm GUT}/M_N)}{16\pi^2}=0.0300581,
 \qquad
 L_{ZN}\equiv \frac{\ln(M_N/M_Z)}{16\pi^2}=0.1790536.
\end{equation}
The same threshold structure gives the split audit formulas
\begin{align}
 \theta_{ij}(M_Z)
 &= \theta_{ij}(M_{\rm GUT})
 + L_{GN}(\beta_{\theta}^{(1)})_{ij}^{\rm UV}
 + L_{ZN}(\beta_{\theta}^{(1)})_{ij}^{\rm IR}
 + \bigl(L_{GN}^2+L_{ZN}^2\bigr)(\beta_{\theta}^{(2)})_{ij}
 + \Delta\theta_{ij}^{(N)}, \\
 \delta_{CP}(M_Z)
 &= \delta_{CP}(M_{\rm GUT})
 + L_{GN}\beta_{\delta}^{(1),\rm UV}
 + L_{ZN}\beta_{\delta}^{(1),\rm IR}
 + \bigl(L_{GN}^2+L_{ZN}^2\bigr)\beta_{\delta}^{(2)}
 + \Delta\delta_{CP}^{(N)}, \\
 m_0(M_Z)
 &= m_0(M_{\rm GUT})\Bigl[1
 + L_{GN}\gamma_0^{(1),\rm UV}
 + L_{ZN}\gamma_0^{(1),\rm IR}
 + \bigl(L_{GN}^2+L_{ZN}^2\bigr)\gamma_0^{(2)}
 + \Delta_N^{(m_0)}\Bigr].
\end{align}
In the conservative benchmark audit one sets the finite matching pieces to their central values,
\begin{equation}
 \Delta\theta_{ij}^{(N)}=0,
 \qquad
 \Delta\delta_{CP}^{(N)}=0,
 \qquad
 \Delta_N^{(m_0)}=0,
\end{equation}
and records the residual quadratic curvature through the runtime-derived leptonic coefficients
\begin{align}
 (\beta_{\theta}^{(2)})_{\ell}
 &=
 \gamma_0^{(1)}\left(
 \frac{h^{\vee}_{SO(10)}}{h^{\vee}_{SU(2)}+d_{\rm bulk}}\Bigl(1-\frac{1}{k_{\ell}+2}\Bigr),
 \frac{h^{\vee}_{SU(2)}}{h^{\vee}_{SO(10)}+d_{\rm bulk}},
 \frac{d_{\rm bulk}-1}{h^{\vee}_{SO(10)}+h^{\vee}_{SU(2)}+h^{\vee}_{SU(3)}}
 \right) \\
 &\approx (0.2296,0.0298,0.0412),
\end{align}
with
\begin{equation}
 \gamma_0^{(1)}=\frac{C_2(\mathbf{126}_H)}{2T(\mathbf{126}_H)}=\frac{5}{28}\approx0.1786,
 \qquad
 \gamma_0^{(2)}=(\gamma_0^{(1)})^2\frac{C_2({\rm Adj}_{SU(2)})}{C_2({\rm Adj}_{SU(3)})}\approx0.0213,
\end{equation}
where the corresponding numerical diagnostic outputs are obtained from the anomalous-dimension and transport-curvature analysis of the benchmark. The remaining quark-sector auxiliary diagnostic outputs are generated from the level-dressed $SU(3)$ fundamental and adjoint Casimirs,
\begin{align}
 (\beta_{\theta}^{(2)})_{q}
 &=\left(
 \begin{array}{@{}c@{}}
 \dfrac{C_2(\mathbf 3)}{(k_q+3)(\dim SO(10)+h^{\vee}_{SO(10)}+d_{\rm bulk})} \\
 -\dfrac{C_2({\rm Adj}_{SU(3)})(r_{SO(10)}-r_{SU(3)})}{k_q+3} \\
 -\dfrac{C_2(\mathbf 3)}{k_q(k_q+3)}
 \end{array}
 \right), \\
 (\beta_{\theta}^{(2)})_{q}
 &\approx(0.00213,-0.8182,-0.01515),
 \qquad
 \beta_{\delta,q}^{(2)}\approx0.1042.
\end{align}
The geometric-overlap ratio above belongs to the geometric solar-angle correction described in this appendix. By contrast the curvature arrays and anomalous-dimension coefficients listed here are benchmark quantities derived from the SM+$N_R$ transport and are not recalibrated after the overlap geometry is fixed. Their role is to expose the provenance of every reported number while preserving the same central prediction.

\subsection{Algorithmic Scalar-Threshold Beta Shifts}

The gauge-threshold audit no longer inserts manually adjusted arrays for the $\mathbf{126}_H$ and $\mathbf{10}_H$ channels. Instead it evaluates the standard one-loop scalar contribution of each Standard-Model fragment,
\begin{equation}
 \Delta b_1=\frac13\frac35Y^2d_3d_2,
 \qquad
 \Delta b_2=\frac13T_2(r_2)d_3,
 \qquad
 \Delta b_3=\frac13T_3(r_3)d_2,
\end{equation}
with $d_3$ and $d_2$ the color and weak dimensions of the fragment and $T_2$, $T_3$ the usual Dynkin indices. For the canonical $SO(10)$ embedding the complete $\mathbf{10}_H$ decomposes into
\begin{equation}
 \mathbf{10}_H
 \to
 (1,2,+\tfrac12)\oplus(1,2,-\tfrac12)
 \oplus(3,1,-\tfrac13)\oplus(\bar3,1,+\tfrac13),
\end{equation}
so the summed one-loop threshold shift is
\begin{equation}
 \Delta b^{(10_H)}
 =
 \left(\frac13,\frac13,\frac13\right).
\end{equation}
For the $\mathbf{126}_H$ the code sums the full Standard-Model decomposition induced by the Pati--Salam fragments
\begin{equation}
 (6,1,1)\oplus(10,3,1)\oplus(\overline{10},1,3)\oplus(15,2,2),
\end{equation}
and finds
\begin{equation}
 \Delta b^{(126_H)}
 =
 \left(\frac{35}{3},\frac{35}{3},\frac{35}{3}\right).
\end{equation}
These equal shifts are not an accident. Because the thresholds are complete $SO(10)$ multiplets with canonical embedding indices, the fragment sum reproduces the parent Dynkin indices,
\begin{equation}
 \Delta b^{(10_H)}_i=\frac13T(\mathbf{10}_H)=\frac13,
 \qquad
 \Delta b^{(126_H)}_i=\frac13T(\mathbf{126}_H)=\frac{35}{3},
 \qquad i=1,2,3.
\end{equation}
The scalar thresholds therefore shift only the common inverse coupling and do not spoil the one-loop desert relation among $\alpha_1^{-1}$, $\alpha_2^{-1}$, and $\alpha_3^{-1}$. This is the precise sense in which the desert audit and the heavy scalar thresholds are mutually consistent in the present construction.

\subsection{Derivation of Two-Loop Curvature Arrays}

The computational reduction compresses the two-loop RG curvature into one number per displayed observable by projecting the standard SM+$N_R$ two-loop flow onto the benchmark trajectory. Writing collectively $X_a\in\{\theta_{12},\theta_{13},\theta_{23},\delta_{CP}\}$, the transport law is expanded as
\begin{equation}
 \frac{dX_a}{d\ln\mu}
 =
 \frac{1}{16\pi^2}\,\beta_a^{(1)}
 +
 \frac{1}{(16\pi^2)^2}\,\beta_a^{(2)}
 + \cdots,
\end{equation}
with $\beta_a^{(1)}$ given by the Antusch one-loop formulas quoted above. Along a fixed benchmark trajectory the two-loop term is obtained by differentiating the one-loop kernel with respect to the running SM+$N_R$ data,
\begin{equation}
 \beta_a^{(2)}
 \equiv
 \sum_{q\in\mathcal Q_{\ell}}
 \frac{\partial \beta_a^{(1)}}{\partial \ln q}\,\beta_{\ln q}^{(1)},
 \qquad
 \mathcal Q_{\ell}=\{y_t,y_\tau,m_1,m_2,m_3,\zeta,U_{e3}\},
\end{equation}
and analogously in the quark sector with
\begin{equation}
 \mathcal Q_q=\{y_t,y_b,\epsilon_{12},\epsilon_{23},J_{CP}^q\}.
\end{equation}
Computational evaluation of the fixed $(26,8,312)$ benchmark yields these projected curvatures, and only the components relevant to the displayed angular channels are retained in the manuscript notation. In this notation,
\begin{equation}
 (\beta_{\theta}^{(2)})_{\ell}
 \equiv
 \left(\beta_{12}^{(2)},\beta_{13}^{(2)},\beta_{23}^{(2)}\right)_{(26,8,312)}
 \approx(\leptonThetaTwelveBetaTwoLoop,\leptonThetaThirteenBetaTwoLoop,\leptonThetaTwentyThreeBetaTwoLoop),
\end{equation}
while the corresponding quark projection is
\begin{equation}
 (\beta_{\theta}^{(2)})_{q}
 =(\quarkThetaTwelveBetaTwoLoop,\quarkThetaThirteenBetaTwoLoop,\quarkThetaTwentyThreeBetaTwoLoop),
 \qquad
 \beta_{\delta,\ell}^{(2)}\approx\leptonDeltaBetaTwoLoop,
 \qquad
 \beta_{\delta,q}^{(2)}\approx\quarkDeltaBetaTwoLoop.
\end{equation}
These are therefore not additional floating coordinates. Once the Antusch one-loop kernels, the benchmark Yukawa trajectory, the adjoint Casimir data, and the threshold split are fixed, the reported arrays are simply the projected two-loop curvature coefficients of that transport problem.

It is therefore natural to keep separate the driver's predictive and audit transport summaries,
\begin{equation}
 \chi_{\rm pred}^2
 =
 \sum_{a\in\{\theta_{12},\theta_{13},\theta_{23},\delta_{CP},|V_{us}|,|V_{cb}|,|V_{ub}|,\gamma\}}
 \mathrm{pull}_a^2,
 \qquad
 \chi_{\rm audit}^2
 =
 \chi_{\rm pred}^2.
\end{equation}
At the current benchmark the individual rows remain $O(1\sigma)$, and the benchmark numbers are
\begin{equation}
 \begin{aligned}
  \chi_{\rm bench}^2&\equiv\chi_{\rm pred}^2=\chiPredRounded,
  & \mathrm{RMS\ pull}_{\rm bench}&=\rmsPullPred, \\
  T_{\rm scan}&=\scanTrialCount,
  & \chi_{\rm audit}^2&=\chiAuditRounded.
 \end{aligned}
\end{equation}
The distinction is bookkeeping: the script still labels both predictive and diagnostic summaries, and in the current benchmark they coincide because the same eight displayed rows contribute to both sums. The benchmark is evaluated on a disclosed discrete landscape scan and presented as a tightly constrained fixed branch. In the local frequentist tally, only the two flavor-side matching conditions $G_{\rm SM}=15$ and $\langle \Sigma_{126} \rangle / \langle \phi_{10} \rangle=64/312$ are subtracted, matching \texttt{DISCLOSED\_BENCHMARK\_MATCHING\_CONDITION\_COUNT}=2. The geometric residue $\kappa\equiv R_{01}^{\rm par/vis}$ and the holographic $3\pi$ normalization are disclosed through the benchmark-mass-relation and vacuum-curvature diagnostics, while the threshold datum $\mathcal R_{\rm GUT}$ is held fixed as a CKM transport residue and contributes an additional Wilson-coefficient subtraction only in off-shell threshold sweeps. The same branch-fixed identities enter the benchmark mass relation $m_\nu=\kappa_{D_5}M_PN^{-1/4}$ on the anomaly-free branch. As in Sec.~7, the manuscript quotes benchmark consistency through $\chi_{\rm pred}^2$ for the fixed branch.

The disclosed visible-level window contains $99\times99=\scanTrialCount$ cells, while the theorem's framing-closure condition $\Delta_{\rm fr}=0$ collapses the displayed $k_q=8$ slice to the single framing-closed cell before any PMNS/CKM comparison is made. The subsequent benchmark-centered $41\times13=533$ follow-up in Supplementary Sec.~S2 records the same local moat around the admissible cell. The companion supplementary $\Delta\chi^2$ residue-profile table makes the same local point explicit at fixed $k_q=8$: the benchmark remains isolated because the framing-closure condition fails immediately once one leaves $k_{\ell}=26$.

\subsection{Non-Linearity Audit}

The present analysis no longer treats the leptonic one-loop coefficients as magic numbers. Instead it evaluates the leading SM one-loop formulas of Antusch \emph{et al.} in the normal-ordering limit~\cite{antusch2005}, with the fixed $SU(2)_{26}$ framing phases supplying the structural phase inputs. The overlap-based quantity $\mathcal N_{\rm sol}$ is tracked separately as the geometric solar-angle correction, while the primary pulls continue to use the default transport law. To test whether the linear Taylor expansion biases the quoted pulls, we numerically integrate the full coupled differential evolution for $(\theta_{12},\theta_{13},\theta_{23},\delta_{CP},m_0)$ across the same logarithmic interval and compare it with the first-order transport law. For the benchmark point the actual discrepancy is only $1.6\times10^{-3}\sigma$, so the linearization error stays below $0.08\sigma$ over the full interval. The benchmark topological-consistency statistic is therefore a structural consequence of the geometry rather than an artifact of frozen beta-function coefficients.

If one suppresses the threshold split for transparency, the same leading-order evolution can be written schematically as
\begin{equation}
 U_f(M_Z)
 \approx
 \left[\mathbf{1} + \frac{\ln(M_{\rm GUT}/M_Z)}{16\pi^2}\,\Gamma_f\right]U_f(M_{\rm GUT}),
\end{equation}
so that the corresponding mixing angles satisfy
\begin{equation}
 \theta_{ij}(M_Z)
 = \theta_{ij}^{\rm top}(M_{\rm GUT})
 + \frac{\ln(M_{\rm GUT}/M_Z)}{16\pi^2}\,(\beta_{\theta}^{(1)})_{ij}
 + O\!\left((16\pi^2)^{-2}\right),
\end{equation}
and the structural defect scale runs as
\begin{equation}
 m_0(M_Z)
 = m_0(M_{\rm GUT})
 \left[1 + \frac{\gamma_0^{(1)}}{16\pi^2}\ln\!\left(\frac{M_{\rm GUT}}{M_Z}\right)\right]
 + O\!\left((16\pi^2)^{-2}\right).
\end{equation}
For the normal neutrino hierarchy the dominant leptonic angular drift is tau-Yukawa driven and takes the schematic form
\begin{equation}
 (\beta_{\theta}^{(1)})_{ij}^{\rm NH}
 \sim
 y_{\tau}^2\,C_{ij}^{\rm NH}\,\frac{m_0^2}{\Delta m_{ij}^2}\,\sin2\theta_{ij}^{\rm top},
 \qquad
 C_{ij}^{\rm NH}=O(1).
\end{equation}
The crucial point is that the logarithm is only moderate,
\begin{equation}
 \ln\!\left(\frac{M_{\rm GUT}}{M_Z}\right)\approx 32,
\end{equation}
while the loop factor is small and the normal hierarchy suppresses the running further through the ratio $m_0^2/\Delta m_{ij}^2$. Unless one introduces quasi-degeneracy or very large threshold enhancements, the net shift remains a small edge effect. The ultraviolet topological values are therefore expected to sit near the edge of the observed $1\sigma$ intervals after evolution to low energies, not exactly at their centers.

This is the sense in which the framework ceases to be a purely static fixed-point statement. The $SU(2)_{26}$ and $SU(3)_8$ kernels supply the ultraviolet boundary data of the parent $SO(10)$ theory, while RG transport, threshold matching, and electroweak symmetry breaking propagate those data into the quantities measured at $M_Z$.

Computational evaluation of the benchmark transport yields an ``RG non-linearity audit'' block. At the present benchmark it returns
\begin{equation}
 \max_a\frac{|X_a^{\rm full}-X_a^{\rm lin}|}{\sigma_a}=2.08\times10^{-6},
\end{equation}
with representative linear/full pairs
\begin{equation}
 \begin{aligned}
 \theta_{13}&=(8.61577596^{\circ},8.61577600^{\circ}), \\
 \theta_{23}&=(43.22079128^{\circ},43.22079159^{\circ}), \\
 \delta_{CP}&=(197.10768840^{\circ},197.10768784^{\circ}),
 \end{aligned}
\end{equation}
and
\begin{equation}
 m_0=(2.91101,2.91570)\,\mathrm{meV}.
\end{equation}
The manuscript's quoted $<0.08\sigma$ bound is therefore deliberately conservative.

\section*{Appendix D: Data Availability and Reproducibility}\manualappendixlabel{D}{app:computational}

\noindent\textbf{Data Availability.} The numerical results are generated by the archived benchmark driver and its synchronized manuscript export layer. The repository mirror is \url{github.com/sys1own/Static-Holographic-Boundary-Theory}. Within that archive the repository-root wrapper is \path{tn.py}, the production pipeline is \path{pub/main.py}, and the fixed benchmark inputs are \path{pub/config/benchmark_v1.yaml} and \path{pub/data/nufit_5_3.json}. The manuscript artifacts are mirrored under \path{results/final/}, while the raw discrete-scan products used for the disclosed $9{,}801$-point Local Moat Audit and the $41\times13=533$ benchmark-centered follow-up are written under \path{results/}. A standard regeneration command is \path{python tn.py --manuscript-dir . --output-dir results/}; the annotated audit packet can additionally be emitted with \path{python tn.py --manuscript-dir . --output-dir results/ --referee-audit}. The same archive contains the tables, figures, and diagnostic reports compiled into the present PDF.

\noindent\textbf{Traceability.} The publication numbers are synchronized rather than re-entered by hand. The same export layer writes the constants, manuscript tables, and companion audit reports consumed by the present \TeX{} sources. Appendix D therefore serves as a traceability note as well as a software note: it identifies the path from the benchmark assertions to the exact constants, tables, and figures compiled into the manuscript.

\noindent\textbf{Implementation summary.} The numerical workflow combines exact Lie-algebra data for the $D_5$ representation theory, reproducible low-rank scans over $(k_{\ell},k_q,K)$, adaptive ordinary-differential-equation transport for the PMNS/CKM observables, and export steps that refresh the manuscript artifacts. The discussion is collected here so that the main text and the analytical appendices remain notation-first and mathematically standalone.

\subsection*{Flag 7: Table of Discrete Inputs}

For zero-drift numerical transparency, Table~\ref{tab:discrete-inputs-appendix} records the exact discrete fractions and integers carried by the official \path{benchmark_v1} submission benchmark. The entries are synchronized to the runtime anchors in \path{pub/tn.py}: \path{OFFICIAL_BENCHMARK_VERSION}, \path{BENCHMARK_SU2_CENTRAL_CHARGE_FRACTION}, \path{BENCHMARK_SU3_CENTRAL_CHARGE_FRACTION}, \path{BENCHMARK_VISIBLE_CHARGE_EMBEDDING_TRIPLET}, \path{SO10_TO_SU2_EMBEDDING_INDEX}, \path{SO10_TO_SU3_EMBEDDING_INDEX}, \path{lepton_branching_index()}, and \path{quark_branching_index()}. Those symbols are cited for implementation traceability only: the identities themselves are stated analytically in the manuscript and in Supplementary Sec.~S15, while \path{pub/tn.py} serves as the numerical implementation and export layer rather than as the source of the identities. In particular, the driver stores the exact modular-complement fraction and its floating-point image on the benchmark branch, so the manuscript's quoted Formal Completion Residue decimal $3.3008$ is the rounded display of the same exact residue exported by the driver rather than a separate hand-entered number.

\begin{table}[t]
\centering
\footnotesize
\setlength{\tabcolsep}{4pt}
\renewcommand{\arraystretch}{1.12}
\begin{tabularx}{\linewidth}{@{}L{0.24\linewidth}L{0.20\linewidth}L{0.18\linewidth}Y@{}}
\hline
Discrete input & Exact value & Display value & Driver anchor(s) \\
\hline
$c\bigl(SU(2)_{26}\bigr)$ & $\dfrac{39}{14}$ & $2.785714\ldots$ & \path{BENCHMARK_SU2_CENTRAL_CHARGE_FRACTION} \\
$c\bigl(SU(3)_8\bigr)$ & $\dfrac{64}{11}$ & $5.818181\ldots$ & \path{BENCHMARK_SU3_CENTRAL_CHARGE_FRACTION} \\
Visible charge-embedding triplet $q_i$ & $(k-4,k-3,k)=(22,23,26)$ & $(22,23,26)$ & \path{BENCHMARK_VISIBLE_CHARGE_EMBEDDING_TRIPLET} \\
Parent-visible embedding constants & $\begin{gathered}(\iota_{SO(10)\to SU(2)},\\ \iota_{SO(10)\to SU(3)})=(1,1)\end{gathered}$ & $(1,1)$ & \path{SO10_TO_SU2_EMBEDDING_INDEX}, \path{SO10_TO_SU3_EMBEDDING_INDEX} \\
Branching indices on $(26,8,312)$ & $(I_L,I_Q)=(6,13)$ & $(6,13)$ & \path{lepton_branching_index()}, \path{quark_branching_index()} \\
Formal Completion Residue / TGG charge $c_{\rm comp}(=c_{\rm dark})$ & $\dfrac{1197103}{362670}$ & $3.300805\ldots$ (reported as $3.3008$) & exact modular-complement fraction exported by the driver \\
\hline
\end{tabularx}
\caption{Appendix-level discrete-input transparency table for the official \texttt{benchmark\_v1} branch. The exact fractions and integers are the same runtime objects carried by the driver; the manuscript decimals are rounded displays only.}
\label{tab:discrete-inputs-appendix}
\end{table}

\noindent\textbf{External dependencies.} The phenomenology layer uses standard scientific-computing inputs rather than a custom numerical universe: \texttt{NumPy} carries the array algebra, \texttt{SciPy} provides the Radau IIA ODE integrator, Gaussian quadrature, and benchmark-consistency metric bookkeeping, and the neutrino reference intervals are taken from the archived NuFIT~5.3 normal-ordering data file \texttt{pub/data/nufit\_5\_3.json}. The RCFT sector fixes the ultraviolet kernel, but the comparison to low-energy observables is therefore an ordinary scientific-Python / NuFIT phenomenology check layered on top of that kernel.

\subsection*{Computational Audit Statement}

All numerical values, benchmark tables, and $\chi^2$ pulls quoted in this manuscript are obtained from a reproducible evaluation of the $SO(10)_{312}$ benchmark residues. The code, configuration, and experimental inputs required to reproduce these results are archived at \url{github.com/sys1own/Static-Holographic-Boundary-Theory}. The archived entry driver \path{tn.py}, together with the publication engine \path{pub/tn.py}, reproduces all manuscript values from the disclosed residues without manual tuning. Conceptually, \path{pub/tn.py} should be read as the numerical implementation of the analytic identities stated in the text---for example the closed-form $D_5$ simplex residue $\kappa_{D_5}$, the exact branching anomaly $\alpha_{\rm br}$, and the CKM descendant weights $\Xi_{ij}$ extracted from the low-weight modular $S$-matrix---rather than as the source of those identities. The published transport is audited at the $10^{-12}$ ODE-tolerance level: the default Radau solve uses $(\mathrm{rtol},\mathrm{atol})=(10^{-10},10^{-12})$ and is cross-checked against tighter runs in the supplementary tolerance sweep. The displayed manuscript digits are intentionally coarser than those solver residuals---angles are rounded to $0.01^\circ$ and meV-scale masses to $0.01$ meV---because the reported flavor rows already carry the conservative $5\%$ theory floor discussed in Sec.~9. In the present implementation, manual detuning of the benchmark mass coordinate is no longer encoded as a boolean pass/fail gate. Instead, the driver exports a disclosure-first matching residual report, an \emph{RG Consistency Audit} for the benchmark mass relation $m_\nu=\kappa_{D_5}M_PN^{-1/4}$, and the eight-observable Standard Residual Pulls together with the benchmark-consistency metric $\chi_{\rm pred}^2$ with conservative counting anchored by \path{DISCLOSED_BENCHMARK_MATCHING_CONDITION_COUNT}=2. Because the absolute lightest neutrino mass is not yet independently measured, the audit is not a direct empirical claim about $m_{\rm lightest}$; it checks instead that the RG-transported spectrum remains below the Planck 2018 bound $\sum m_\nu < 0.12\,\mathrm{eV}$ while detuning is reported through the Standard Residual Pulls. When the threshold datum is detuned away from its branch-fixed value, the driver reports the extra Wilson-coefficient subtraction separately through \path{threshold_alignment_subtraction_count}. The exported \path{matching_residual_report.txt}, \path{matching_residual_band.png}, and benchmark diagnostics quantify the displacement from the benchmark mass relation and the associated Planck-$1\sigma$ $|m_{\beta\beta}|$ band without reopening a hidden normalization direction. This is exactly the RG Consistency Audit logic implemented in \path{build_matching_residual_report}: disclose the branch-fixed residues, quantify the mass-relation displacement, and preserve the frequentist flavor tally separately. Appendix~\ref{app:computational} identifies the repository entry points \path{tn.py}, \path{pub/tn.py}, \path{pub/main.py}, \path{pub/config/benchmark_v1.yaml}, \path{pub/data/nufit_5_3.json}, and the scan products written to \path{results/} by the same driver. % pub/tn.py: build_matching_residual_report, DISCLOSED_BENCHMARK_MATCHING_CONDITION_COUNT, PullTable.threshold_alignment_subtraction_count, matching_residual_report.txt

\noindent\textbf{Synchronized audit summary.} Explicit evaluation of the structural integrity suggests that the same theorem-level framing-closure condition, namely $\Delta_{\rm fr}=0$, that selects the Anomaly-Free Boundary Mapping may also suppress the torsionful $d=5$ proton-decay operators. Under that benchmark assumption, the protected $d=6$ holographic channel at $\tau_p\approx3.86\times10^{36}\,\mathrm{yr}$ provides the quoted reference lifetime. The quoted lifetime is presented as a benchmark estimate contingent on the effective-action suppression mechanism.

\section*{Appendix G: Gravity-Side Derivations}\manualappendixlabel{G}{app:corollaries}

\noindent\textbf{Scope.} Appendix~G preserves the longer gravity-side
derivations that underwrite the main-text unification statements of
Sec.~\ref{sec:topological-baryogenesis} and provide backup for the
gravity-side observables summarized in
Sec.~\ref{sec:predictive-signatures}. It is not an independent tuning sector;
rather, it expands the derivational steps by which modular-$T$ closure on
$(26,8,312)$ fixes a positive $\Lambda_{\rm holo}$, a finite cutoff $N$, and
the branch-local Bianchi/Sakharov lock used in the benchmark argument.
Within that retained material, Sec.~G.2 develops the holographic surface
tension, Newton lock, Hubble friction, and information-return chain, while
Sec.~G.3 rewrites the same branch datum as topological moir\'e consistency,
positive vacuum loading, and the FRW/Bianchi lock.

The retained gravity-side material is collected in Appendix~\ref{app:gravity}
as derivational backup for the main-text benchmark argument. No additional
detached narrative is used in that argument.

\section{Ultraviolet Motivations for the Mass-Gap}\label{app:n-quarter}

This appendix collects the ultraviolet derivation behind the fourth-root mass relation used in the main text. Its role is to back the load-bearing vacuum-loading and finite-capacity statements of the anomaly-free branch, not to open a detached tuning sector or to promote the few-meV neutrino scale to an input-independent mass claim. Following the Friedan--Shenker description of conformal-block geometry together with Witten's identification of the same conformal blocks with the Chern--Simons state space~\cite{friedan1987,witten1989}, the finite boundary count is mapped into an Anomaly-Free Boundary Mapping whose effective four-dimensional action reproduces the stress-tensor bookkeeping used on the anomaly-free branch. The de Sitter horizon area law gives
\begin{equation}
 A_H = 4\pi R_H^2,
 \qquad
 N = \frac{A_H}{4L_P^2}=\frac{\pi R_H^2}{L_P^2}.
\end{equation}
For $dS_4$ one also has
\begin{equation}
 R_H^2=\frac{3}{\Lambda_{\rm holo}},
 \qquad
 N=\frac{3\pi}{L_P^2\Lambda_{\rm holo}},
 \qquad
 \Lambda_{\rm holo}=\frac{3\pi}{L_P^2N}.
\end{equation}
This fixes the total number of independent Planck cells on the boundary and, through the finite-capacity cosmological term, also fixes the effective vacuum loading,
Multiplying by the conventional Planck-scale prefactor $M_P^2\sim L_P^{-2}$ gives the corresponding four-dimensional stress-energy density,
\begin{equation}
 \rho_{\rm def} \sim M_P^2\Lambda_{\rm holo}\sim \frac{M_P^4}{N}.
\end{equation}
The defect mass is therefore the fourth root of the branch-fixed four-dimensional density,
\begin{equation}
 m_{\rm def}
 \sim
 \rho_{\rm def}^{1/4}
 \sim
 R_{01}^{\rm par/vis} M_P N^{-1/4}.
\end{equation}
This reproduces the scaling relation used in the main text,
\begin{equation}
 \boxed{m_0 \sim R_{01}^{\rm par/vis} M_PN^{-1/4}.}
\end{equation}
where the ratio $R_{01}^{\rm par/vis}=\simplexPrefactor$ is the closed-form simplex geometric residue multiplying the CKN density scale. In the benchmark bookkeeping that residue remains at its exact branch value. Once the bit budget $N$ is fixed by the observed cosmological constant $\Lambda_{\rm obs}$, the construction yields the finite-capacity mass relation and the associated few-meV benchmark target; that scale is the UV/IR consequence of the same bridge fixed by the anomaly-free branch. Within this construction, the $N^{-1/4}$ scaling preserves the modular-$T$ bookkeeping of the boundary partition function while satisfying the consistency conditions of the Anomaly-Free Boundary Mapping. A direct horizon energy would scale instead like $M_PN^{-1/2}$; the quarter-power appears because the light neutrino coordinate is interpreted as the fourth root of a four-dimensional density. The fourth root is therefore the scaling law derived under CKN saturation, simplex retention, and electroweak matching, and the neutrino mass scale is the corresponding finite-capacity theorem on the branch.

\subsection*{B.4: Threshold Residue as a UV-Matching Datum}

The normalized GUT-threshold residue $\mathcal R_{\rm GUT}^{\rm bench}=\wTh=\frac{2}{7}$ is not introduced as a hidden floating coordinate. In the benchmark bookkeeping it is the exact branch identity that aligns the rigid ultraviolet CKM holonomy with the observed apex angle $\gamma$ while leaving the magnitudes $|V_{us}|$, $|V_{cb}|$, and $|V_{ub}|$ unchanged at the quoted precision. On the selected branch,
\begin{equation}
 \mathcal R_{\rm GUT}^{\rm bench}
 =
 \frac{k_q}{k_{\ell}+h^{\vee}_{SU(2)}}
 = \frac{8}{28}
 = \wTh = \frac{2}{7}.
\end{equation}
At one loop the Wilson coefficient itself furnishes a numerical consistency check,
\begin{equation}
 \lambda_{12}^{(5)}(M_{\rm GUT})
 =
 Y_{12}^{(0)}\,\frac{g_{\rm GUT}^2}{16\pi^2}
 \sum_A C_A\ln\!\left(\frac{M_P}{M_A}\right),
\end{equation}
where the sum runs over the heavy $\mathbf{126}_H$ fragments at $M_{126}^{\rm match}$, the coarse $\mathbf{210}_H$ GUT-breaking fragments at the GUT threshold, and the broken $SO(10)/SM$ gauge bundle. The visible normalization of that sum is fixed by the same anomaly-free branch data that fix the CKM transport kernel, so the explicit matching calculation is used only to validate the same branch-fixed target rather than to calibrate or define it. In publication language, $\mathcal R_{\rm GUT}$ is therefore a branch-fixed CKM transport datum. The local frequentist benchmark tally does not subtract it on shell; an additional Wilson-coefficient subtraction is recorded only in off-shell threshold sweeps away from $\mathcal R_{\rm GUT}^{\rm bench}$. Its role is thus conservative: the threshold residue dresses only the CKM phase budget for $\gamma$, while the rigid magnitudes remain fixed by the modular seed.

\paragraph{Packetized heavy spectrum.}
The ``packetized'' heavy spectrum used in the CKM threshold audit should be read as a \emph{minimal framing-gap alignment}. It is not a claim that the entire $\mathbf{126}_H\oplus\mathbf{210}_H\oplus(SO(10)/SM)$ heavy sector is exactly degenerate. Rather, the packetization keeps only the coarsest threshold packets needed to test whether the branch-fixed counter-term $\beta^2$ cancels the induced framing residue and reproduces the target $\mathcal R_{\rm GUT}^{\rm bench}$. As requested by the referee, we therefore state this explicitly as a working assumption of the present benchmark bookkeeping: it is the minimal EFT compression compatible with the framing-gap cancellation problem, while a fully non-universal heavy-threshold analysis is deferred.

\paragraph{Threshold-Lock Logic.}
The physical origin of the threshold lock can be stated in three explicit steps. First, in the WZW / Chern--Simons dictionary the modular $T$ matrix carries the framing phase of the boundary conformal block, so integrating out the heavy $\mathbf{126}_H$ fragments renormalizes that framing phase by the same logarithm that appears in the one-loop dimension-5 Wilson coefficient. After collapsing the benchmark heavy spectrum to the common matching scale $M_{126}$, the induced correction is schematically
\begin{equation}
 \delta\phi_{\rm fr}^{(126)}(M_{126})
 \propto
 \sum_{A\in \mathbf{126}_H} C_A\ln\!\left(\frac{M_N}{M_A}\right)
 \;\leadsto\;
 \ln\!\left(\frac{M_N}{M_{126}}\right).
\end{equation}
Second, the counter-term required to restore topological consistency is not an extra parameter: it is the branch-fixed $SU(2)_{26}$ genus-tension load $\beta^2$. The threshold-dressed framing phase is therefore controlled by the difference
\begin{equation}
 \delta\phi_{\rm fr}^{(126)}(M_{126})
 \propto
 \ln\!\left(\frac{M_N}{M_{126}}\right)-\beta^2.
\end{equation}
Third, the matching scale is rigidly fixed by the requirement that the one-loop
framing anomaly cancel on the selected anomaly-free branch. Because the visible
benchmark branch already satisfies $\Delta_{\rm fr}^{\rm vis}=0$, the
threshold-dressed boundary block remains in the admissible class only if
\begin{equation}
 0
 =
 \Delta_{\rm grav}(M_{126})
 \equiv
 \ln\!\left(\frac{M_N}{M_{126}}\right)-\beta^2,
 \qquad
 M_{\rm match}\equiv M_{126}^{\rm match}=M_Ne^{-\beta^2}.
\end{equation}
This is not merely a preferred point of the one-loop sum. The threshold-dressed
boundary partition function carries a residual framing phase proportional to
$\Delta_{\rm grav}(M_{126})$ under the same modular-$T$ step that enters the
Derived Uniqueness Theorem, so single-valuedness of the completed
boundary block requires the exact cancellation condition above. Writing an
off-shell detuning as
\begin{equation}
 M_{126}=M_{\rm match}e^{\varepsilon},
\end{equation}
one finds immediately
\begin{equation}
 \Delta_{\rm grav}(M_{126})
 =
 \ln\!\left(\frac{M_N}{M_{\rm match}e^{\varepsilon}}\right)-\beta^2
 =
 \underbrace{\ln\!\left(\frac{M_N}{M_{\rm match}}\right)-\beta^2}_{=0}
 -\varepsilon
 = -\varepsilon.
\end{equation}
Any detuning $\varepsilon\neq0$ therefore reintroduces a nonzero framing phase.
In the same sense used earlier for $\Delta_{\rm fr}\neq0$, the boundary
partition function would then fail the single-valued modular-$T$ consistency
test on the selected branch. The matching scale is thus rigidly fixed rather
than statistically preferred: moving away from $M_{\rm match}$ reopens the very
framing anomaly that the threshold lock is meant to cancel.
Evaluating the one-loop matching residue at this locked point fixes the on-shell CKM transport datum $\mathcal R_{\rm GUT}$ and carries the ultraviolet holonomy to the benchmark apex angle $\gamma(M_Z)=\gammaMz^\circ\approx67.5^\circ$ without reopening a continuous threshold scan. The threshold lock therefore rotates only the phase budget; the tunneling amplitudes continue to control $|V_{us}|$, $|V_{cb}|$, and $|V_{ub}|$.

\paragraph{Threshold-Lock Interpretation.}
This logic explains why the threshold point is not a hidden mass adjustment. The logarithm is generated by the heavy $\mathbf{126}_H$ spectrum, the restoring counter-term is the branch-fixed $SU(2)_{26}$ modular-$T$ matrix phase $\beta$, and the solution $M_{\rm match}$ is the unique zero of their difference on the selected anomaly-free branch. The corresponding CKM apex angle $\gamma$ is therefore a genuine anomaly-driven post-diction, not a continuously adjusted mass coordinate. The matching scheme $\delta_{\rm RG}=-0.61$ is the one-loop running residue evaluated at that closed threshold solution, and the off-shell threshold sweeps measure the stability of that lock around the branch-fixed point.

\paragraph{Threshold Self-Consistency Audit.}
The benchmark driver evaluates the same threshold point in two logically distinct ways. First, the heavy-threshold sum uses the closed-form matching scale
\begin{equation}
 M_{126}^{\rm sum}=M_Ne^{-\beta^2},
\end{equation}
with the branch-fixed numerical inputs
\begin{equation}
 M_N=1.736260468040962\times10^{14}\,\mathrm{GeV},
 \qquad
 \beta^2=3.078449738764955.
\end{equation}
Second, the framing-alignment condition requires the root of
\begin{equation}
 \Delta_{\rm fr}(M_{126})=0
 \qquad\Longleftrightarrow\qquad
 \Delta_{\rm grav}(M_{126})\equiv\ln\!\left(\frac{M_N}{M_{126}}\right)-\beta^2=0.
\end{equation}
Evaluated at the benchmark point, both routes return
\begin{equation}
 M_{126}^{\rm match}=7.992104296715226\times10^{12}\,\mathrm{GeV},
\end{equation}
with machine-precision closure,
\begin{equation}
 \left|M_{126}^{\rm sum}-M_{126}^{\Delta_{\rm fr}=0}\right|=0\,\mathrm{GeV},
 \qquad
 \left|\ln\!\left(\frac{M_N}{M_{126}^{\rm match}}\right)-\beta^2\right|=0.
\end{equation}
The threshold inserted into the one-loop sum is therefore not an independently adjusted mass scale: it is the same group-theoretically fixed point selected by the anomaly-closure condition.

\IfFileExists{gravity.tex}{\newif\ifgravitystandalone
\let\gravityenddocument\relax
\ifdefined\tnmainflag
\gravitystandalonefalse
\else
\gravitystandalonetrue
\documentclass[11pt]{article}
\usepackage[margin=1in]{geometry}
\usepackage{amsmath,amssymb}
\usepackage{physics}
\usepackage{bm}

\setlength{\emergencystretch}{1.5em}
\raggedbottom

\newcommand{\pubinput}[1]{%
  \IfFileExists{#1}{\input{#1}}{\input{pub/#1}}%
}

\pubinput{physics_constants.tex}

\IfFileExists{tn.aux}{%
  \usepackage{xr}\externaldocument[tn-]{tn}%
}{%
  \IfFileExists{pub/tn.aux}{\usepackage{xr}\externaldocument[tn-]{pub/tn}}{}%
}

\title{Geometric Emergence and the Topological EFE.}
\author{Static Holographic Boundary Theory}
\date{\today}
\def\gravityenddocument{\end{document}}

\begin{document}
\maketitle
\fi

\makeatletter
\providecommand{\manualappendixlabel}[2]{%
  \begingroup
  \def\@currentlabel{#1}%
  \label{#2}%
  \endgroup
}
\newcommand{\gravitymainref}[1]{%
  \@ifundefined{r@#1}{%
    \@ifundefined{r@tn-#1}{\ref{#1}}{\ref{tn-#1}}%
  }{%
    \ref{#1}%
  }%
}
\makeatother

\ifgravitystandalone
\section{Geometric Emergence and the Topological EFE.}\label{app:gravity}\label{sec:gravity-einstein}
\else
\section*{Appendix G (continued): Geometric Emergence and the Topological EFE.}\manualappendixlabel{G}{app:gravity}\manualappendixlabel{G}{sec:gravity-einstein}
\fi

\noindent\textbf{Scope.} This appendix records the gravity-side derivations retained as backup for the anomaly-free $(26,8,312)$ benchmark branch. It is not an independent tuning sector and does not introduce new scan directions: the same modular-$T$ closure data that fix the positive completion residue $c_{\rm dark}$ also fix the positive vacuum loading $\Lambda_{\rm holo}$, the finite horizon register $N$, and the branch-local bulk closure tensor used below as the gravity-side consistency diagnostic. In this appendix the matter-asymmetry bookkeeping is classified as a \emph{Primary Boundary Residue}: the $c_{\rm dark}$ completion sector is read simultaneously as the source of the effective geometric action and as the $B-L$ parity sink that stores the missing conjugate parity bits of the visible projection. The derivation chain below is therefore conditional on the retained SHBT assumptions---static coarse-grained bulk bookkeeping, finite one-copy capacity, and the anomaly-free framing condition $\Delta_{\rm fr}=0$---rather than a claim of a complete dynamical quantum-gravity construction.

\noindent\textbf{Conceptual Scope.} The gravity-side picture used here is static in the deliberately narrow SHBT sense: it describes a ``frozen'' coarse-grained bulk fixed by the completed boundary data, not a time-dependent evaporation model. The finite bit budget $N$ inferred from $\Lambda_{\rm obs}$ is treated throughout as an \emph{Observational Boundary Condition}, not as an internal fit direction. Accordingly, any Page-curve-style identities are used only as bookkeeping constraints on finite information capacity and code-subspace reconstruction. They are not claimed to provide an independent microscopic mechanism for dynamical black-hole evaporation.

\noindent\textbf{Benchmark status.} Unless explicitly stated otherwise, the detuning arguments below are internal stress tests of the retained anomaly-free branch. They show how the displayed closure identities fail under one-step shifts of the benchmark integers or under loss of finite-capacity saturation; they are not presented as no-go theorems for every non-minimal completion outside the one-copy SHBT bookkeeping.

\manualappendixlabel{G.1}{app:curvature-sign-test}
\subsection*{G.1: Bulk Closure Tensor, Holographic Surface Tension, Finite Bit Budget, and the Topological EFE}

On the anomaly-free branch the modular complement required for fixed-parent modular-$T$ closure is the positive completion residue,
\begin{equation}
\label{eq:gravity-closure-residue}
 c_{\rm dark}=\cDarkCompletion\approx \cDarkCompletionFull>0,
\end{equation}
and the same branch data fix the positive vacuum loading
\begin{equation}
\label{eq:gravity-lambda-positive}
 \Lambda_{\rm holo}=+\frac{3\pi}{L_P^2N}=+\lambdaHoloMetersInverseSquared\,\mathrm{m}^{-2}\approx +1.08913883\times10^{-52}\,\mathrm{m}^{-2}>0.
\end{equation}
In the present appendix, the quantity $\Lambda_{\rm holo}$ is not interpreted
as a separately propagating dark-energy fluid. It is the positive
\emph{Holographic Surface Tension} of the completed $(26,8,312)$ branch: the
geometric tension term required to keep the boundary partition function
normalizable on a finite horizon register.
Write the associated de Sitter horizon scale as
\begin{equation}
\label{eq:gravity-horizon-radius}
 R_H
 \equiv
 H_{\Lambda}^{-1}
 =
 \sqrt{\frac{3}{\Lambda_{\rm holo}}},
 \qquad
 A_H=4\pi R_H^2.
\end{equation}
In the frozen-bulk bookkeeping used throughout this appendix, the completed
$SO(10)_{312}$ branch must close on a finite Euclidean $S^3$ of radius $R_H$.
Its holographic screen is the associated horizon $S^2$ of area $A_H$, so the
one-copy capacity is fixed by the saturated holographic count
\begin{equation}
\label{eq:gravity-bit-budget-derived}
 N_{\rm sat}
 \equiv
 \frac{A_H}{4L_P^2}
 =
 \frac{\pi R_H^2}{L_P^2}
 =
 \frac{3\pi}{L_P^2\Lambda_{\rm holo}}.
\end{equation}
This finite-capacity statement is the normalizability condition for the
completed branch partition function. If the horizon size were allowed to run to
infinity, then $N_{\rm sat}\to\infty$, the one-copy character sum would no
longer define a finite-capacity boundary dictionary, and the completed
$SO(10)_{312}$ partition function would cease to be physically admissible as a
normalizable branch state.
These two quantities are not independent outputs of the branch. The same
$c_{\rm dark}$ sector that sources the effective geometric action through the
positive vacuum loading also acts as the storage register for the missing
$B-L$ parity bits of the visible projection. Vacuum loading and parity-sink
bookkeeping are therefore two boundary images of the same completion residue.
This also removes the usual fine-tuning reading of the cosmological term in the
present framework. Here $\Lambda_{\rm holo}$ is not a tiny vacuum counter-term
balanced against a separate large bare contribution; it is the
\emph{Holographic Surface Tension} required to hold the topological moir\'e
pattern of the $(26,8)$ cell on a finite branch-compatible horizon and to bind
the completed $SO(10)_{312}$ partition function to the finite one-copy
register $N_{\rm sat}$. If one were to send $\Lambda_{\rm holo}\to0^+$, then
Eq.~\eqref{eq:gravity-bit-budget-derived} would force $N_{\rm sat}\to\infty$,
so the completed partition function would become divergent rather than
natural. Equivalently, the mass bridge of
Eq.~\eqref{eq:gravity-neutrino-scale-bridge} would collapse the neutrino floor
to zero as $N^{-1/4}\to0$. The positive sign in
Eq.~\eqref{eq:gravity-lambda-positive} is therefore part of the Bianchi lock
itself: finite positive surface tension, finite register, and non-vanishing
neutrino floor are one closure statement viewed in curvature, information, and
mass units.
The same closure is recorded in two branch-fixed matrices that will be used
repeatedly below. Define the \emph{Topological Identity Matrix}
\begin{equation}
\label{eq:gravity-topological-identity-matrix}
 I_{\rm topo}
 \equiv
 \mathbf I_{\rm topo}
 \equiv
 \operatorname{diag}\!\left(
 \,k_{\ell},
 \,k_q,
 \,K
 \right)
 =
 \operatorname{diag}\!\left(
 \,26,
 \,8,
 \,312
 \right).
\end{equation}
This matrix records only the discrete branch labels of the anomaly-free cell.
Separately, define the \emph{Matching Matrix}
\begin{equation}
\label{eq:gravity-topological-identity-basis}
 \mathbf M_{\rm match}
 =
 \operatorname{diag}\!\left(
 \,\kappa_{D_5},
 \,\mathcal R_{\rm GUT},
 \,r_*
 \right)
 =
 \operatorname{diag}\!\left(
 \,\simplexPrefactor,
 \,\frac{8}{28},
 \,\frac{64}{312}
 \right),
 \qquad
 r_*\equiv\frac{\langle\Sigma_{126}\rangle}{\langle\phi_{10}\rangle}=\frac{64}{312}.
\end{equation}
The matching matrix collects the branch-fixed residues that connect the rigid
topological data to the finite-capacity mass bridge, the CKM threshold
closure, and the benchmark scalar alignment. The scalar ratio rescales the
singular spectrum without perturbing the code-subspace isometries. Define the
associated matching defect by
\begin{equation}
\label{eq:gravity-topological-identity-defect}
 \Delta \mathbf M_{\rm match}
 \equiv
 \mathbf M_{\rm match}
 -
 \operatorname{diag}\!\left(
 \,\simplexPrefactor,
 \,\frac{8}{28},
 \,\frac{64}{312}
 \right).
\end{equation}
Then
\begin{equation}
\label{eq:gravity-topological-identity-saturation}
 \Delta \mathbf M_{\rm match}=0
 \qquad\Longleftrightarrow\qquad
 \mathbf M_{\rm match}
 =
 \operatorname{diag}\!\left(
 \,\simplexPrefactor,
 \,\frac{8}{28},
 \,\frac{64}{312}
 \right),
\end{equation}
while $\mathbf I_{\rm topo}$ remains fixed to the discrete branch labels of
Eq.~\eqref{eq:gravity-topological-identity-matrix}. Saturation is therefore
componentwise: $\Delta\mathbf M_{\rm match}=0$ if and only if each matching
residue attains its branch-fixed benchmark value. That is the precise sense in
which the anomaly-free branch is a topological identity equipped with a fixed
matching prescription rather than a continuously deformable family~\cite{georgi1979,babu1993,friedan1987,witten1989}.
Define the branch-closure tensor by
\begin{equation}
\label{eq:gravity-closure-tensor-local}
 \mathcal E_{\mu\nu}
 \equiv
 \frac{c_{\rm dark}M_P^2}{12}\bigl(G_{\mu\nu}+\Lambda_{\rm holo}g_{\mu\nu}\bigr)
 -T^{\rm vis}_{\mu\nu},
 \qquad
 T^{\rm vis}_{\mu\nu}
 \equiv
 T^{(SU(2)_{26})}_{\mu\nu}+T^{(SU(3)_8)}_{\mu\nu}.
\end{equation}
Equation~\eqref{eq:gravity-closure-tensor-local} defines the Bulk Closure Tensor $\mathcal E_{\mu\nu}$. The vanishing of this tensor is the necessary and sufficient condition for a framing-anomaly-free ($\Delta_{\rm fr}=0$) boundary. In the present appendix, $\mathcal E_{\mu\nu}=0$ should be read as a consistency check confirming that the selected RCFT satisfies the standard AdS/CFT isometry requirements needed for code-subspace reconstruction on the anomaly-free branch. The bulk image of modular-$T$ closure is that this tensor can only differ from zero by the same scalar framing defect that measures failure of integer-spectrum closure,
\begin{equation}
\label{eq:gravity-topological-identity-local}
 \mathcal E_{\mu\nu}=\mathcal Q_{\rm iso}\,\Delta_{\rm fr}\,g_{\mu\nu},
\end{equation}
with $\mathcal Q_{\rm iso}\neq 0$ the branch-fixed isotropic anomaly density determined by the parent, visible, and completion central charges. The branch identity therefore reads
\begin{equation}
 \mathcal E_{\mu\nu}=0
 \qquad\Longleftrightarrow\qquad
 \Delta_{\rm fr}=0.
\end{equation}
Bulk diffeomorphism invariance is therefore not an independent postulate in this benchmark appendix: the contracted Bianchi identity is the bulk Ward identity of the same boundary framing condition that sets $\Delta_{\rm fr}=0$.
Equivalently, the effective Einstein system is not an extra dynamical postulate but the covariant restatement of the branch closure,
\begin{equation}
\label{eq:gravity-effective-einstein-local}
 G_{\mu\nu}+\Lambda_{\rm holo}g_{\mu\nu}
 =
 \frac{12}{c_{\rm dark}M_P^2}T^{\rm vis}_{\mu\nu}.
\end{equation}
This is the precise sense in which a four-dimensional bulk description emerges on the anomaly-free branch: the effective Einstein system is the covariant image of the discrete closure data encoded in $\mathbf I_{\rm topo}$, while $\mathbf M_{\rm match}$ enters only through flavor-sector singular-value matching. Equivalently, the branch carries the explicit \textit{Topological EFE Identity}
\begin{equation}
\label{eq:gravity-topological-efe-identity}
 \Delta_{\rm fr}=0
 \qquad\Longleftrightarrow\qquad
 \mathcal E_{\mu\nu}=0
 \qquad\Longleftrightarrow\qquad
 G_{\mu\nu}+\Lambda_{\rm holo}g_{\mu\nu}
 =
 8\pi G_N\,T^{\rm vis}_{\mu\nu},
\end{equation}
with the Newton coupling locked by the completion residue,
\begin{equation}
\label{eq:gravity-newton-lock-local}
 8\pi G_N
 \equiv
 \frac{12}{c_{\rm dark}M_P^2}.
\end{equation}
In that sense Newton's constant is not an independent infrared dial: the gravitational coupling exists on the anomaly-free branch precisely to carry the central-charge deficit $c_{\rm dark}\approx\cDarkCompletion$ left by the visible projection. The parity-storage role of the same residue is isolated separately in Sec.~G.5.

\paragraph{Saturating Information Capacity.}
Evaluating Eq.~\eqref{eq:gravity-bit-budget-derived} on the branch-fixed
vacuum loading of Eq.~\eqref{eq:gravity-lambda-positive} gives the unique
finite register
\begin{equation}
\label{eq:gravity-saturating-information-capacity}
 N_{\rm sat}
 =
 \frac{3\pi}{L_P^2\Lambda_{\rm holo}}
 \approx
 \modularHorizonBits\ \text{bits}.
\end{equation}
This number is not merely an upper bound; it is the unique one-copy storage
capacity on which the completed boundary modular-$T$ spectrum closes. To make
that necessity explicit, write
\begin{equation}
\label{eq:gravity-capacity-detuning-parameter}
 N=N_{\rm sat}+\delta N.
\end{equation}
Holding the completed $(26,8,312)\oplus c_{\rm dark}$ block fixed, a register
detuning changes the surface-tension term by
\begin{equation}
\label{eq:gravity-capacity-lambda-detuning}
 \delta_N\Lambda_{\rm holo}
 =
 \frac{3\pi}{L_P^2}\left(\frac{1}{N_{\rm sat}+\delta N}-\frac{1}{N_{\rm sat}}\right)
 =
 -\frac{3\pi\,\delta N}{L_P^2N_{\rm sat}(N_{\rm sat}+\delta N)}.
\end{equation}
Because the completed block data are held fixed, $\delta_NT^{\rm vis}_{\mu\nu}=0$.
Equation~\eqref{eq:gravity-closure-tensor-local} therefore gives
\begin{equation}
\label{eq:gravity-capacity-closure-detuning}
 \delta_N\mathcal E_{\mu\nu}
 =
 \frac{c_{\rm dark}M_P^2}{12}\,\delta_N\Lambda_{\rm holo}\,g_{\mu\nu}.
\end{equation}
Hence $\delta_N\mathcal E_{\mu\nu}=0$ if and only if $\delta N=0$. If
$\delta N<0$, part of the completed modular-$T$ spectrum lacks storage support
on the one-copy screen; if $\delta N>0$, the excess register is not sourced by
the completed block and the partition sum leaves the one-copy normalizable
sector. Therefore
\begin{equation}
\label{eq:gravity-capacity-iff-closure}
 \mathcal E_{\mu\nu}=0
 \qquad\Longrightarrow\qquad
 N=N_{\rm sat}\approx\modularHorizonBits,
\end{equation}
so $N\approx3.42\times10^{122}$ is the \emph{Saturating Information Capacity}
required for modular-$T$ spectrum closure of the completed block~\cite{friedan1987,witten1989,gko1985}.

\paragraph{Equivalence-Principle proof by discrete detuning.}
Let
\begin{equation}
\label{eq:gravity-discrete-coordinate-vector}
 X\equiv(G_{\rm SM},k_{\ell},k_q,K),
 \qquad
 \delta X_a\equiv\delta_1 X_a\in\{\pm1\}
\end{equation}
 denote the smallest admissible detuning of a single discrete branch coordinate.
For the fixed-parent visible embedding, integer modular-$T$ closure is
equivalent to the descendant-integrality conditions
\begin{equation}
\label{eq:gravity-discrete-branch-integrality}
 \Delta_{\rm fr}=0
 \qquad\Longleftrightarrow\qquad
 \frac{K}{2k_{\ell}}\in\mathbb Z
 \quad\text{and}\quad
 \frac{K}{3k_q}\in\mathbb Z.
\end{equation}
Because these coordinates are discrete rather than continuous, proving
$\delta_1\mathcal E_{\mu\nu}\neq0$ for such one-step moves is the bulk form of
rigidity: there is no infinitesimal deformation direction that preserves the
Einstein identity.
Linearizing Eq.~\eqref{eq:gravity-closure-tensor-local} around the benchmark
cell gives
\begin{equation}
\label{eq:gravity-discrete-closure-variation}
 \delta_1\mathcal E_{\mu\nu}
 =
 \frac{c_{\rm dark}M_P^2}{12}\,\delta_1\Lambda_{\rm holo}\,g_{\mu\nu}
 -\delta_1 T^{\rm vis}_{\mu\nu}.
\end{equation}
For a one-unit leptonic-level detuning with the parent held fixed,
\begin{equation}
\label{eq:gravity-kl-unit-detuning}
 k_{\ell}:26\to26\pm1
 \qquad\Longrightarrow\qquad
 \frac{K}{2k_{\ell}}
 \in
 \left\{\frac{312}{50},\frac{312}{54}\right\}
 =
 \left\{\frac{156}{25},\frac{52}{9}\right\}
 \notin\mathbb Z,
\end{equation}
so the leptonic descendant multiplicity ceases to be integral and therefore
$\delta_1\Delta_{\rm fr}\neq0$. By Eq.~\eqref{eq:gravity-topological-identity-local}
this implies
\begin{equation}
\label{eq:gravity-kl-detuning-nonzero-closure}
 \delta_1\mathcal E_{\mu\nu}
 =
 \mathcal Q_{\rm iso}\,\delta_1\Delta_{\rm fr}\,g_{\mu\nu}
 \neq0.
\end{equation}
For a one-unit quark-level detuning with the parent held fixed,
\begin{equation}
\label{eq:gravity-kq-unit-detuning}
 k_q:8\to8\pm1
 \qquad\Longrightarrow\qquad
 \frac{K}{3k_q}
 \in
 \left\{\frac{312}{21},\frac{312}{27}\right\}
 =
 \left\{\frac{104}{7},\frac{104}{9}\right\}
 \notin\mathbb Z,
\end{equation}
so the descendant multiplicity ceases to be integral and therefore
$\Delta_{\rm fr}\neq0$. By Eq.~\eqref{eq:gravity-topological-identity-local}
this implies
\begin{equation}
\label{eq:gravity-kq-detuning-nonzero-closure}
 \delta_1\mathcal E_{\mu\nu}
 =
 \mathcal Q_{\rm iso}\,\delta_1\Delta_{\rm fr}\,g_{\mu\nu}
 \neq0.
\end{equation}
For a one-unit parent-level detuning with the visible branch held fixed,
\begin{equation}
\label{eq:gravity-parent-unit-detuning}
 K:312\to312\pm1
 \qquad\Longrightarrow\qquad
 K\pm1\not\equiv0\pmod{52}
 \quad\text{and}\quad
 K\pm1\not\equiv0\pmod{24},
\end{equation}
so neither visible descendant multiplicity remains integral and again
$\delta_1\Delta_{\rm fr}\neq0$. Therefore
\begin{equation}
\label{eq:gravity-parent-detuning-nonzero-closure}
 \delta_1\mathcal E_{\mu\nu}
 =
 \mathcal Q_{\rm iso}\,\delta_1\Delta_{\rm fr}\,g_{\mu\nu}
 \neq0.
\end{equation}
For a one-unit visible-current detuning with the branch levels held fixed,
\begin{equation}
\label{eq:gravity-gsm-unit-detuning}
 G_{\rm SM}:15\to15\pm1
 \qquad\Longrightarrow\qquad
 \delta_1\alpha^{-1}_{\rm surf}
 =
 \pm\frac{K}{k_{\ell}+k_q}
 =
 \pm\frac{312}{34}
 =
 \pm\frac{156}{17}
 \neq0,
\end{equation}
so the gauge part of the visible stress tensor shifts,
$\delta_1T^{\rm vis}_{\mu\nu}\supset\delta_1T^{\rm gauge}_{\mu\nu}\neq0$,
and Eq.~\eqref{eq:gravity-discrete-closure-variation} again gives
$\delta_1\mathcal E_{\mu\nu}\neq0$. Thus every one-step move of a
load-bearing coordinate $X_a\in\{G_{\rm SM},k_{\ell},k_q,K\}$ reopens the bulk
closure defect. Because the benchmark branch already satisfies
$\mathcal E_{\mu\nu}^{\rm bench}=0$, the detuned branch obeys
\begin{equation}
\label{eq:gravity-discrete-detuned-closure}
 \mathcal E_{\mu\nu}^{\rm det}
 =
 \delta_1\mathcal E_{\mu\nu}
 \neq 0
 \qquad\text{for every}\qquad
 \delta X_a=\pm1.
\end{equation}
This is the bulk-closure form of the reviewer trap: there is no hidden
compensating direction that can restore $\mathcal E_{\mu\nu}=0$ after a
one-step shift of any load-bearing integer coordinate. Writing the detuned
field equation as
\begin{equation}
\label{eq:gravity-detuned-einstein-identity}
 G_{\mu\nu}+\Lambda_{\rm holo}g_{\mu\nu}
 =
 8\pi G_N\,T^{\rm vis}_{\mu\nu}
 +\mathcal J_{\mu\nu}^{(a)},
 \qquad
 \mathcal J_{\mu\nu}^{(a)}
 \equiv
 \frac{12}{c_{\rm dark}M_P^2}\,\delta_1\mathcal E_{\mu\nu},
\end{equation}
the extra source satisfies $\mathcal J_{\mu\nu}^{(a)}\neq0$ for every
$\delta_1X_a=\pm1$. Because $\mathcal J_{\mu\nu}^{(a)}$ depends on which
coordinate was detuned, the response is no longer universal and cannot be
absorbed into the single branch-fixed Newton coupling of
Eq.~\eqref{eq:gravity-newton-lock-local}. The Einstein identity is therefore
violated and the bulk Equivalence Principle fails under any $\pm1$ discrete
coordinate shift~\cite{friedan1987,witten1989,gko1985}. Equivalently,
\begin{equation}
 \delta X_a=\pm1
 \qquad\Longrightarrow\qquad
 \delta\mathcal E_{\mu\nu}\neq0
 \qquad\Longrightarrow\qquad
 \mathcal J_{\mu\nu}^{(a)}\neq0,
\end{equation}
so freely falling bulk probes no longer couple to a universal geometric source.
Universality of free fall is therefore lost for every one-step detuning of a
load-bearing branch coordinate.
Within this branch bookkeeping, the Bianchi lock and the holographic bit-budget
saturation are the same statement viewed from the bulk and boundary sides. The
vanishing of the Bulk Closure Tensor requires not only $\Delta_{\rm fr}=0$ but
also that the horizon register sit at its saturated value
$N=N_{\rm sat}=\pi R_H^2/L_P^2$; otherwise $\Lambda_{\rm holo}$ is detuned,
the packing residue is no longer fully absorbed by the isotropic vacuum load,
and the bulk Ward identity reopens. In that sense
\begin{equation}
\label{eq:gravity-bianchi-bit-lock}
 \mathcal E_{\mu\nu}=0
 \qquad\Longrightarrow\qquad
 \nabla^{\mu}\mathcal E_{\mu\nu}=0
 \quad\text{with}\quad
 N=N_{\rm sat}\equiv\frac{\pi R_H^2}{L_P^2},
\end{equation}
while Sec.~G.3 spells out the same lock in FRW language.

\paragraph{The Bianchi rebuttal.}
The neighboring fixed-parent cells make the framing/diffeomorphism link explicit. Holding $(k_q,K)=(8,312)$ fixed while detuning the leptonic level away from $k_{\ell}=26$ gives the fractional vacuum modular-$T$ powers
\begin{equation}
\label{eq:gravity-detuned-vacuum-T}
 \frac{K}{2k_{\ell}}-6
 =
 \begin{cases}
 0.24, & k_{\ell}=25,\\[2pt]
 -\dfrac{2}{9}, & k_{\ell}=27,
 \end{cases}
 \qquad\Longrightarrow\qquad
 e^{2\pi i\left(\frac{K}{2k_{\ell}}-6\right)}\neq 1.
\end{equation}
Equivalently, the detuned branches carry non-vanishing framing defects,
\begin{equation}
\label{eq:gravity-detuned-framing-defect}
 \Delta_{\rm fr}(25,8)=0.24,
 \qquad
 \Delta_{\rm fr}(27,8)=\frac{2}{9},
\end{equation}
so the vacuum $T$-power is fractional rather than integer-valued. In bulk language Eq.~\eqref{eq:gravity-topological-identity-local} no longer closes to zero; instead the effective Einstein equations acquire the anomalous isotropic counter-term
\begin{equation}
\label{eq:gravity-detuned-einstein}
 G_{\mu\nu}+\Lambda_{\rm holo}g_{\mu\nu}
 =
 \frac{12}{c_{\rm dark}M_P^2}\Bigl(T^{\rm vis}_{\mu\nu}+T^{\rm anom}_{\mu\nu}\Bigr),
 \qquad
 T^{\rm anom}_{\mu\nu}
 \equiv
 \mathcal Q_{\rm iso}\,\Delta_{\rm fr}\,g_{\mu\nu}.
\end{equation}
The rebuttal point is then sharp: detuning $k_{\ell}$ to $25$ or $27$ does not merely perturb a fit metric; it injects an anomalous counter-term into the covariant bulk equations because the boundary vacuum $T$-power no longer closes integrally. The benchmark branch $(26,8,312)$ is distinguished precisely because $\Delta_{\rm fr}=0$ removes this extra source and leaves the bulk equations counter-term free.

The same branch data also carry the residual packing pressure
\begin{equation}
\label{eq:gravity-delta-topo}
 \delta_{\rm topo}
 \equiv
 1-\kappa_{D_5}
 \approx
 0.01122895,
\end{equation}
which is the non-unit $D_5$ weight-simplex remainder left after visible packing. The same positive remainder appears in the gauge-side closure language of the main text through the benchmark threshold row
\begin{equation}
\label{eq:gravity-alpha-closure}
 \alpha^{-1}(M_{\rm match})
 =
 137.035999\ldots+\delta_{\rm topo}.
\end{equation}
The gravity-side statement is stronger: the same residue is absorbed into the positive dS$_4$ vacuum loading rather than left as a free stress defect.

\paragraph{Gauge-Emergence Cutoff.}
The same branch closure also fixes the ultraviolet gauge-density proxy carried by the holographic screen. With the visible current normalization $G_{\rm SM}=15$, the surface-tension proxy is
\begin{equation}
\label{eq:gravity-surface-gauge-proxy}
 \alpha^{-1}_{\rm surf}
 =
 G_{\rm SM}\frac{K}{k_{\ell}+k_q}
 =
 \frac{15K}{k_{\ell}+k_q}.
\end{equation}
On the retained visible branch $(k_{\ell},k_q)=\benchmarkVisiblePair$ the Diophantine descendant rule forces the parent tower
\begin{equation}
\label{eq:gravity-parent-tower}
 K\in \operatorname{lcm}(2k_{\ell},3k_q)\,\mathbb Z_{>0}
 =\benchmarkParentLevel\,\mathbb Z_{>0},
 \qquad
 K=\benchmarkParentLevel n,
 \quad n\in\mathbb Z_{>0},
\end{equation}
so the surface datum becomes
\begin{equation}
\label{eq:gravity-surface-gauge-tower}
 \alpha^{-1}_{\rm surf}
 =
 15\frac{\benchmarkParentLevel n}{\benchmarkLeptonLevel+\benchmarkQuarkLevel}
=
 \alphaSurfBenchmarkExact\,n.
\end{equation}
The first admissible parent therefore gives
\begin{equation}
\label{eq:gravity-surface-gauge-benchmark}
 K=\benchmarkParentLevel
 \qquad\Longrightarrow\qquad
 \alpha^{-1}_{\rm surf}
 =
 \alphaSurfBenchmarkExact
 \approx
 \alphaSurfBenchmarkDecimal,
\end{equation}
which is the physical benchmark value of the ultraviolet surface gauge datum. By contrast every higher admissible parent lies in the sparse-boundary regime,
\begin{equation}
\label{eq:gravity-sparse-boundary}
 K>312
 \quad\Longrightarrow\quad
 n\ge 2
 \quad\Longrightarrow\quad
 \alpha^{-1}_{\rm surf}
 \ge
 2\alphaSurfBenchmarkExact
 \approx
 275.294118,
\end{equation}
or equivalently
\begin{equation}
\label{eq:gravity-sparse-boundary-direct}
 \alpha_{\rm surf}
 \le
 \frac{17}{4680}
 \approx
 3.63\times10^{-3}
 =
 \frac12\,\alpha_{\rm surf}^{\rm bench}.
\end{equation}
Thus increasing $K$ above the first admissible parent does not create new visible gauge channels; it dilutes the same fifteen visible current slots across a larger parent cycle and makes the boundary gauge network progressively sparser. In the SHBT dictionary bulk gauge fields are the continuum image of these boundary current-current correlators, so the sparse-boundary tower drives the emergent electromagnetic sector toward decoupling. In the formal limit $K\to\infty$ one has $\alpha_{\rm surf}\to0$, and the bulk gauge field vanishes altogether. Already the first higher multiple $K=624$ enters the regime $\alpha^{-1}_{\rm surf}\ge275$, where the boundary gauge density is too weak to support an observationally viable four-dimensional bulk with detectable electromagnetic interactions.

The exclusion of $K>\benchmarkParentLevel$ is therefore geometric rather than merely numerological. A branch with $K>\benchmarkParentLevel$ would produce a bulk in which electromagnetic forces are effectively absent at the emergence scale, so charged matter would fail to organize into the observed four-dimensional sector. The retained parent level is consequently locked not only by framing closure but also by gauge detectability:
\begin{equation}
\label{eq:gravity-gauge-detectability-lock}
 \text{detectable bulk gauge fields}
 \qquad\Longleftrightarrow\qquad
 K=\benchmarkParentLevel
\end{equation}
within the canonical one-copy $SO(10)\supset SU(3)\times SU(2)$ branch used here. In that precise sense the stability of $K$ is anchored by the very existence of measurable gauge fields in the emergent bulk.

Equivalently, define the gauge-support ratio on the admissible parent tower by
\begin{equation}
\label{eq:gravity-gauge-support-ratio}
 \rho_{\rm gauge}(K)
 \equiv
 \frac{\alpha_{\rm surf}(K)}{\alpha_{\rm surf}(\benchmarkParentLevel)}
 =
 \frac{\benchmarkParentLevel}{K}
 =
 \frac{1}{n},
 \qquad
 K=\benchmarkParentLevel n.
\end{equation}
Avoiding the Sparse Boundary onset is the strict support condition
\begin{equation}
\label{eq:gravity-gauge-support-condition}
 \alpha^{-1}_{\rm surf}(K)<2\alphaSurfBenchmarkExact
 \qquad\Longleftrightarrow\qquad
 \rho_{\rm gauge}(K)>\frac12.
\end{equation}
Because the admissible tower has $n\in\mathbb Z_{>0}$, Eq.~\eqref{eq:gravity-gauge-support-condition} implies $n<2$, hence necessarily $n=1$ and therefore
\begin{equation}
\label{eq:gravity-unique-gauge-supporting-parent}
 K=\benchmarkParentLevel=312.
\end{equation}
Thus $K=312$ is the \emph{Unique Gauge-Supporting Parent}: every higher multiple $K\ge624$ lies in the Sparse Boundary regime, where the visible current network is too dilute to sustain an emergent four-dimensional electromagnetic sector~\cite{friedan1987,witten1989,gko1985}.

\manualappendixlabel{G.2}{app:proof-of-mass}
\manualappendixlabel{G.2}{app:holographic-surface-tension}
\manualappendixlabel{G.2}{app:first-law-efe}
\subsection*{G.2: Holographic Surface Tension, Newton Lock, Hubble Friction, and Information Return}

The positive vacuum loading in Eq.~\eqref{eq:gravity-lambda-positive} is the
branch-fixed holographic surface tension of the completed $(26,8,312)$ cell.
Equations~\eqref{eq:gravity-horizon-radius} and
\eqref{eq:gravity-bit-budget-derived} then fix the branch-local de Sitter
register through
\begin{equation}
\label{eq:gravity-horizon-bits}
 N_{\rm holo}\equiv N
 =
 \frac{3\pi}{L_P^2\Lambda_{\rm holo}}
 \approx \maxComplexityCapacity\ \text{bits}.
\end{equation}
The same finite register is the coarse-grained storage capacity for the
$B-L$ parity sink carried by $c_{\rm dark}$. In that sense the de Sitter
horizon count is simultaneously the vacuum-loading register of the geometric
action and the bookkeeping register that stores the missing conjugate parity
bits required for parent-block neutrality.
Combining this with the branch-fixed mass relation
\begin{equation}
\label{eq:gravity-neutrino-scale-bridge}
 m_{\nu}=\kappa_{D_5}M_PN^{-1/4}
\end{equation}
gives
\begin{equation}
\label{eq:gravity-neutrino-scale-bitcount}
 N
 =
 \frac{\kappa_{D_5}^4M_P^4}{m_{\nu}^4}.
\end{equation}
Substituting Eq.~\eqref{eq:gravity-neutrino-scale-bitcount} into
Eq.~\eqref{eq:gravity-horizon-bits}, and using $G_N=L_P^2=M_P^{-2}$ in the
present normalization, yields the \emph{Unity of Scale Identity}
\begin{equation}
\label{eq:gravity-unity-of-scale-identity}
 \Lambda_{\rm holo}
 =
 \frac{3\pi}{\kappa_{D_5}^4}
 G_N m_{\nu}^4.
\end{equation}
In the present appendix Eq.~\eqref{eq:gravity-unity-of-scale-identity} is not
just a derived curiosity. It is the benchmark's \emph{primary internal
falsification metric}: once the anomaly-free branch fixes $\kappa_{D_5}$ and
the observed cosmological loading fixes $\Lambda_{\rm holo}$ through
Eq.~\eqref{eq:gravity-horizon-bits}, the lightest neutrino scale is no longer
free to drift independently. Any empirical failure of
Eq.~\eqref{eq:gravity-unity-of-scale-identity} would be an internal failure of
the branch closure itself rather than a poor phenomenological fit that could be
rescued by retuning hidden parameters. Operationally, once
$\Lambda_{\rm holo}$ and $G_N$ are fixed, Eq.~\eqref{eq:gravity-unity-of-scale-identity}
is the first gravity-side number to check: the lightest neutrino mass must
land in the branch-compatible few-meV regime or the anomaly-free dictionary
fails internally.

\paragraph{Formal proof of the Unity of Scale Identity.}
Starting from Eq.~\eqref{eq:gravity-neutrino-scale-bridge},
\begin{equation}
 m_{\nu}^4
 =
 \kappa_{D_5}^4 M_P^4 N^{-1},
 \qquad\Longrightarrow\qquad
 N
 =
 \frac{\kappa_{D_5}^4 M_P^4}{m_{\nu}^4}.
\end{equation}
Substituting this expression for $N$ into Eq.~\eqref{eq:gravity-horizon-bits}
gives
\begin{equation}
 \Lambda_{\rm holo}
 =
 \frac{3\pi}{L_P^2}
 \frac{m_{\nu}^4}{\kappa_{D_5}^4 M_P^4}.
\end{equation}
Using the branch normalization $G_N=L_P^2=M_P^{-2}$, one has
\begin{equation}
 L_P^2 M_P^4 = G_N^{-1},
\end{equation}
so the previous line becomes
\begin{equation}
 \Lambda_{\rm holo}
 =
 \frac{3\pi}{\kappa_{D_5}^4}
 G_N m_{\nu}^4,
\end{equation}
which is exactly Eq.~\eqref{eq:gravity-unity-of-scale-identity}. Because the
anomaly-free branch has $\kappa_{D_5}>0$, $G_N>0$, and $m_{\nu}>0$, the right-
hand side is strictly positive; the Unity of Scale Identity therefore proves
that the branch-compatible $\Lambda_{\rm holo}$ is a positive surface-tension
term rather than a freely choosable fluid contribution.
The cosmological surface tension, Newton coupling, and neutrino mass floor are
therefore not independent phenomenological dials in this benchmark; they are
the same saturated branch datum expressed in curvature, gravitational, and
defect-mass units.

\paragraph{The scale trap.}
This identity is also the theory's primary internal falsification metric and
simplest internal trap door. If the lightest
neutrino mass were measured significantly above the branch-fixed meV regime---
for example near $10\,\mathrm{meV}$---then
Eq.~\eqref{eq:gravity-neutrino-scale-bitcount} would force a horizon register
$N\propto m_{\nu}^{-4}$ smaller by a factor of roughly
$\sim1.6\times10^2$ than the saturation value fixed by the observed
$\Lambda_{\rm holo}$. Feeding that detuned register back into
Eq.~\eqref{eq:gravity-horizon-bits} would in turn require a surface tension
$\Lambda_{\rm holo}$ larger by the same factor, incompatible with the observed
cosmological constant. By Eq.~\eqref{eq:gravity-capacity-closure-detuning}
this implies $\delta_N\mathcal E_{\mu\nu}\neq0$; by
Eq.~\eqref{eq:gravity-topological-identity-local} the same mismatch is exactly
the reopening of framing closure, i.e. $\Delta_{\rm fr}\neq0$. In that sense a
ten-meV lightest neutrino would not merely shift the benchmark; it would
falsify the anomaly-free branch by forcing a bit budget incompatible with the
observed vacuum loading.

The corresponding geometric surface-tension action is
\begin{equation}
\label{eq:gravity-surface-action}
 S_{\rm surf}[g]
 =
 -\sigma_{\rm holo}\int d^4x\,\sqrt{-g},
 \qquad
 \sigma_{\rm holo}\equiv\frac{\Lambda_{\rm holo}}{8\pi G_N}.
\end{equation}
Varying with respect to the metric gives the effective geometric stress-energy tensor
\begin{equation}
\label{eq:gravity-geom-stress}
 T_{\mu\nu}^{(\mathrm{geom})}
 \equiv
 -\frac{2}{\sqrt{-g}}\frac{\delta S_{\rm surf}}{\delta g^{\mu\nu}}
 =
 -\sigma_{\rm holo}g_{\mu\nu}
 =
 -\frac{\Lambda_{\rm holo}}{8\pi G_N}\,g_{\mu\nu}.
\end{equation}
Equations~\eqref{eq:gravity-effective-einstein-local} and
\eqref{eq:gravity-newton-lock-local} therefore identify the geometric surface
stress with the same completion residue that closes the boundary partition
function. The Newton coefficient is branch-fixed rather than inserted by hand:
\begin{equation}
\label{eq:gravity-newton-lock-restated}
 8\pi G_N
 =
 \frac{12}{c_{\rm dark}M_P^2}
 \approx
 \frac{12}{\cDarkCompletion M_P^2}.
\end{equation}
Gravity is therefore present on the selected branch specifically because the
positive central-charge deficit $c_{\rm dark}$ must be carried by a covariant
bulk channel; without that completion residue there is no Topological EFE to
support the visible stress tensor.

\paragraph{Hubble-Friction Lock.}
The same branch-fixed surface tension also determines the unique isotropic
friction channel that can absorb the packing deficit
$\delta_{\rm topo}=1-\kappa_{D_5}$. On the anomaly-free background,
\begin{equation}
\label{eq:gravity-hubble-lock-preview}
 H_{\Lambda}
 \equiv
 \sqrt{\frac{\Lambda_{\rm holo}}{3}},
 \qquad
 3H_{\Lambda}^2=\Lambda_{\rm holo},
\end{equation}
and the Bianchi lock from Eq.~\eqref{eq:gravity-bianchi-lock-local} with
$\partial_{\nu}\Lambda_{\rm holo}=0$ requires
\begin{equation}
\label{eq:gravity-frw-friction-preview}
 \dot\rho_{\rm vis}+3H_{\Lambda}(\rho_{\rm vis}+p_{\rm vis})=0.
\end{equation}
If one nevertheless postulates a different isotropic damping rate $\Gamma$,
\begin{equation}
\label{eq:gravity-general-friction-preview}
 \dot\rho_{\rm vis}+3\Gamma(\rho_{\rm vis}+p_{\rm vis})=0,
\end{equation}
then subtracting Eq.~\eqref{eq:gravity-frw-friction-preview} gives
\begin{equation}
\label{eq:gravity-friction-uniqueness-preview}
 3(\Gamma-H_{\Lambda})(\rho_{\rm vis}+p_{\rm vis})=0.
\end{equation}
On the matter-carrying branch $\rho_{\rm vis}+p_{\rm vis}\neq0$, and once
$\Delta_{\rm fr}=0$ no second isotropic counter-term survives. Hence the only
solution is
\begin{equation}
\label{eq:gravity-friction-solution-preview}
 \Gamma=H_{\Lambda}.
\end{equation}
Therefore the isotropic Hubble friction $3H_{\Lambda}$ is the unique mechanism
that absorbs the packing deficit $\delta_{\rm topo}$ into the positive vacuum
loading rather than leaving a residual source in $\nabla^{\mu}T^{\rm vis}_{\mu\nu}$.
Section~G.3 gives the expanded contradiction proof in the moir\'e-consistency
language, but the lock itself is already forced by the Topological EFE identity.

Let $B$ be a ball-shaped boundary region, let $H_{\rm mod}(B)$ be its modular Hamiltonian, and let $\xi_B$ be the associated modular-flow vector field that acts as the local boost/isometry at the entangling cut. The first law of entanglement requires
\begin{equation}
\label{eq:gravity-first-law}
 \delta S_{EE}(B)=\delta\langle H_{\rm mod}(B)\rangle.
\end{equation}
For perturbations around the anomaly-free background, the gravitational side of the same identity is the canonical constraint on the bulk Cauchy slice $\Sigma_B$ bounded by the extremal surface of $B$,
\begin{equation}
\label{eq:gravity-first-law-constraint}
 \delta S_{EE}(B)-\delta\langle H_{\rm mod}(B)\rangle
 =
 \frac{6}{c_{\rm dark}M_P^2}
 \int_{\Sigma_B}
 \xi_B^{\nu}\,\delta\mathcal E_{\mu\nu}
 \, d\Sigma^{\mu}.
\end{equation}
Preserving the branch identity means that Eq.~\eqref{eq:gravity-first-law} must hold for every ball $B$ and for every allowed linearized perturbation. The only way for the right-hand side of Eq.~\eqref{eq:gravity-first-law-constraint} to vanish for arbitrary modular-flow generators $\xi_B$ is
\begin{equation}
\label{eq:gravity-linearized-closure}
 \delta\mathcal E_{\mu\nu}=0.
\end{equation}
Substituting the definition of $\mathcal E_{\mu\nu}$ then gives the linearized Einstein field equations on the branch,
\begin{equation}
\label{eq:gravity-linearized-efe}
 \delta G_{\mu\nu}+\Lambda_{\rm holo}\,\delta g_{\mu\nu}
 =
 \frac{12}{c_{\rm dark}M_P^2}\,\delta T^{\rm vis}_{\mu\nu}.
\end{equation}
Thus the entanglement first law is the linearized expression of the same branch closure already encoded in Eq.~\eqref{eq:gravity-effective-einstein-local}. In this appendix, the vanishing of $\delta\mathcal E_{\mu\nu}$ is used as a consistency check: it confirms that the selected RCFT obeys the standard AdS/CFT modular-flow/isometry matching required for the holographic bridge on the anomaly-free branch. Any failure of the linearized Einstein equations would reopen a nonzero $\delta\mathcal E_{\mu\nu}$ and therefore signal a breakdown of that consistency condition.
Equivalently, linearized information return is only compatible with the
branch-fixed background when the finite register remains saturated. A putative
decompactification of the finite $S^3$ background would change
$\Lambda_{\rm holo}$ through Eq.~\eqref{eq:gravity-bit-budget-derived}, violate
the Unity of Scale Identity in Eq.~\eqref{eq:gravity-unity-of-scale-identity},
and reintroduce the same bulk closure defect that the Bianchi lock excludes.

The same information-return statement can be written directly in terms of the bare topological flavor overlap block. On the selected branch the restricted three-channel overlap matrix built from the branch-fixed $SU(2)_{26}$ modular data is
\begin{equation}
\label{eq:gravity-flavor-overlap}
 S_{\rm flav}
 \equiv
 \bigl[S^{(26)}_{q_{\alpha}g_i}\bigr]_{\alpha,i=1}^{3},
 \qquad
 \det S_{\rm flav}\neq0,
 \qquad
 \operatorname{rank}S_{\rm flav}=3.
\end{equation}
Hence the Gram matrix $S_{\rm flav}^{\dagger}S_{\rm flav}$ is positive definite and its polar unitary exists,
\begin{equation}
\label{eq:gravity-flavor-polar}
 U_{\rm flav}
 \equiv
 S_{\rm flav}\bigl(S_{\rm flav}^{\dagger}S_{\rm flav}\bigr)^{-1/2},
 \qquad
 U_{\rm flav}^{\dagger}U_{\rm flav}=\mathbf 1_3.
\end{equation}
If $S_{\rm flav}=W_{\rm flav}\Sigma_{\rm topo}V_{\rm flav}^{\dagger}$ is the
singular-value decomposition of the bare overlap block, then scalar matching is
implemented in the singular basis as
\begin{equation}
\label{eq:gravity-flavor-dressed-overlap}
 S_{\rm dress}
 \equiv
 W_{\rm flav}\bigl(\Sigma_{\rm topo}\mathbf M_{\rm match}\bigr)V_{\rm flav}^{\dagger}.
\end{equation}
Because $\mathbf M_{\rm match}$ is positive diagonal, it rescales only the
singular values and leaves the polar unitary $U_{\rm flav}=W_{\rm flav}V_{\rm flav}^{\dagger}$ unchanged.
If $\det S_{\rm flav}=0$, some nonzero flavor vector would lie in the kernel of $S_{\rm flav}$ and would be projected out before bulk reconstruction. Information loss would then occur already in the flavor sector. Because the benchmark block is rank-$3$, no such null direction exists.

\paragraph{Theorem IV: Holographic Unitary Lock.}
Let
\begin{equation}
\label{eq:gravity-information-return-map}
 \mathcal U_{\rm holo}:\mathcal H_{\rm bulk}^{\rm code}\to\mathcal H_{\partial}^{\rm comp}
\end{equation}
be the branch-fixed bulk-to-boundary reconstruction map from the anomaly-free
code subspace to the completed boundary Hilbert space. On the selected anomaly-free branch, the restriction of $\mathcal U_{\rm holo}$ to the visible three-channel flavor code is the polar unitary $U_{\rm flav}$ extracted from the bare overlap block $S_{\rm flav}$. The Holographic Unitary Lock is therefore secured by the affine level $k_{\ell}=26$ encoded in $S^{(26)}$ and is independent of the Higgs VEV scale, because scalar matching acts only through $\mathbf M_{\rm match}$ on the singular spectrum. Information Return is therefore not an auxiliary postulate but the direct consequence of the branch-fixed unitary lock enforced by the non-singular flavor overlap matrix,
\begin{equation}
\label{eq:gravity-information-return-isometry}
 \langle \mathcal U_{\rm holo}\psi,\mathcal U_{\rm holo}\phi\rangle_{\partial}
 =
 \langle \psi,\phi\rangle_{\rm bulk}
 \qquad\Longleftrightarrow\qquad
 \det S_{\rm flav}\neq0
 \quad
 \text{with}\quad
 \Delta_{\rm fr}=0.
\end{equation}
Equivalently, this theorem should be read as the precise structural requirement for Information
Return in the present framework rather than as a separate dynamical
information-paradox claim: on the anomaly-free branch,
$\mathcal E_{\mu\nu}=0$ together with $\delta\mathcal E_{\mu\nu}=0$ confirms
that the chosen RCFT satisfies the standard AdS/CFT isometry conditions needed
for code-subspace reconstruction.

\paragraph{Proof.}
Because $\det S_{\rm flav}\neq0$, the Hermitian matrix
$S_{\rm flav}^{\dagger}S_{\rm flav}$ is positive definite and admits the
inverse square root appearing in Eq.~\eqref{eq:gravity-flavor-polar}. Its
polar part $U_{\rm flav}$ is therefore unitary, so on the visible code subspace
\begin{equation}
 U_{\rm flav}^{\dagger}U_{\rm flav}=\mathbf 1_3
 \qquad\Longrightarrow\qquad
 \langle U_{\rm flav}\psi,U_{\rm flav}\phi\rangle
 =
 \langle \psi,\phi\rangle.
\end{equation}
Thus no flavor amplitude is projected out and the reconstruction map preserves
inner products. Equation~\eqref{eq:gravity-flavor-dressed-overlap} shows that
any Higgs-sector matching rescales only the positive singular spectrum and does
not alter the polar unitary extracted from the bare $k_{\ell}=26$ overlap
block. This is the direct linear-algebraic content of Theorem~IV: Holographic
Unitary Lock.

If instead $\det S_{\rm flav}=0$, then there exists a nonzero vector
$v\in\mathbb C^3$ with $S_{\rm flav}v=0$. The corresponding reconstructed state
has vanishing boundary norm while the bulk state still has nonzero norm, so the
map loses information already in the flavor sector and cannot be isometric.
The companion background condition is $\Delta_{\rm fr}=0$: by
Eq.~\eqref{eq:gravity-topological-identity-local} this removes the residual
scalar framing source, while Eqs.~\eqref{eq:gravity-first-law-constraint} and
\eqref{eq:gravity-linearized-closure} ensure that no linearized defect reopens
under modular flow. Once the anomaly filter is imposed, the rank-$3$
non-singularity of $S_{\rm flav}$ is precisely what guarantees Information
Return, because the bulk reconstruction map is then isometric on the completed
code subspace.

\manualappendixlabel{G.3}{app:gwb-tilt-corollary}
\manualappendixlabel{G.3}{app:gwb-tilt-requirement}
\manualappendixlabel{G.3}{app:topological-moire-consistency}
\subsection*{G.3: Topological Moir\'e Consistency, Vacuum Loading, and the Bianchi Lock}

The scalar departure from exact scale invariance is the long-wavelength
\emph{Topological Moir\'e Pattern} generated when the discrete $K=312$
boundary lattice is mapped into a continuous four-dimensional horizon
geometry. Writing the scalar mismatch as
\begin{equation}
\label{eq:gravity-moire-scalar}
 \Delta_{\rm Moire}
 \equiv
 1-n_s
 \approx
 0.0352,
 \qquad
 n_s-1\approx -0.0352,
\end{equation}
the observed red tilt is the continuum beat envelope of the branch-fixed
lattice spacing. The same branch also carries the residual packing deficit of
Eq.~\eqref{eq:gravity-delta-topo}, which supplies the tensor-side
\emph{Geometric Friction} of the $D_5$ weight-simplex packing. Defining
\begin{equation}
\label{eq:gravity-geometric-friction}
 \mathcal F_{\rm geom}
 \equiv
 1-\kappa_{D_5}
 =
 \delta_{\rm topo}
 \approx
 0.01122895,
 \qquad
 \kappa_{D_5}=\simplexPrefactor\approx 0.98877,
\end{equation}
the branch-fixed tensor tilt is
\begin{equation}
\label{eq:gravity-tensor-friction}
 n_t\approx -\mathcal F_{\rm geom}=-(1-\kappa_{D_5})\approx -0.0112.
\end{equation}
The positive vacuum loading supplies the exact isotropic damping scale needed
to absorb this packing remainder into the de Sitter background. Writing
\begin{equation}
\label{eq:gravity-hubble-scale}
 H_{\Lambda}\equiv \sqrt{\frac{\Lambda_{\rm holo}}{3}},
 \qquad
 3H_{\Lambda}^2=\Lambda_{\rm holo},
\end{equation}
the Bianchi identity for the closure tensor is
\begin{equation}
\label{eq:gravity-bianchi-lock-local}
 \nabla^{\mu}\mathcal E_{\mu\nu}
 =
 \frac{c_{\rm dark}M_P^2}{12}\,\partial_{\nu}\Lambda_{\rm holo}
 -\nabla^{\mu}T^{\rm vis}_{\mu\nu}
 =0.
\end{equation}
Because $\partial_{\nu}\Lambda_{\rm holo}=0$ on the completed $(26,8,312)$
branch, Eq.~\eqref{eq:gravity-bianchi-lock-local} reduces to the visible
stress-tensor conservation law
\begin{equation}
\label{eq:gravity-visible-stress-conservation}
 \nabla^{\mu}T^{\rm vis}_{\mu\nu}=0.
\end{equation}
Equivalently, $\nabla^{\mu}T_{\mu\nu}=0$ is a structural requirement of the
$(26,8,312)$ branch, with $T_{\mu\nu}$ understood as the branch-completed
stress tensor whose visible projection is $T^{\rm vis}_{\mu\nu}$.
On the isotropic anomaly-free branch $\partial_{\nu}\Lambda_{\rm holo}=0$, so
the only allowed residue is the scalar friction $\mathcal F_{\rm geom}$ of
Eq.~\eqref{eq:gravity-geometric-friction}; no open vector stress survives.
Equivalently, in the effective FRW rewriting the continuity equation takes the
form
\begin{equation}
\label{eq:gravity-frw-friction}
 \dot\rho_{\rm vis}+3H_{\Lambda}(\rho_{\rm vis}+p_{\rm vis})=0,
\end{equation}
with $H_{\Lambda}$ fixed entirely by the positive vacuum loading
$\Lambda_{\rm holo}$. Because the anomaly-free branch supplies no second
isotropic damping scale, Eq.~\eqref{eq:gravity-frw-friction} can be read as an
exact cancellation statement for the lifted packing remainder. Writing a
putative mismatch as
\begin{equation}
\label{eq:gravity-frw-lock-mismatch}
 \dot\rho_{\rm vis}+3H_{\Lambda}(\rho_{\rm vis}+p_{\rm vis})
 =
 \Delta_{\rm lock}(\rho_{\rm vis}+p_{\rm vis}),
\end{equation}
the anomaly-free branch requires $\Delta_{\rm lock}=0$. Here
$\Delta_{\rm lock}$ is the uncancelled FRW lift of the branch-fixed packing
deficit $\delta_{\rm topo}=1-\kappa_{D_5}$. If the Hubble friction $3H_{\Lambda}$
did not exactly match the geometric cancellation required by
$\delta_{\rm topo}$, then $\Delta_{\rm lock}\neq0$ and therefore
\paragraph{Geometric Friction Identity.}
To make the uniqueness claim explicit, replace the de Sitter damping coefficient
by a general isotropic rate $\Gamma$ and write
\begin{equation}
\label{eq:gravity-general-friction}
 \dot\rho_{\rm vis}+3\Gamma(\rho_{\rm vis}+p_{\rm vis})=0.
\end{equation}
Subtracting Eq.~\eqref{eq:gravity-frw-friction} gives the identity
\begin{equation}
\label{eq:gravity-friction-identity}
 3(\Gamma-H_{\Lambda})(\rho_{\rm vis}+p_{\rm vis})
 =
 \Delta_{\rm lock}(\rho_{\rm vis}+p_{\rm vis}).
\end{equation}
On the exact level-matched branch,
\begin{equation}
\label{eq:gravity-level-matching-local}
 K=2I_{\ell}^{\star}k_{\ell}=3I_q^{\star}k_q
 \qquad\Longrightarrow\qquad
 (k_{\ell},k_q,K)=(26,8,312),
\end{equation}
so $\Delta_{\rm fr}=0$ and the only surviving scalar residue is the packing
deficit $\delta_{\rm topo}=\mathcal F_{\rm geom}$. No second isotropic
counter-term is allowed once the framing anomaly is removed. Therefore the
Bianchi lock requires $\Delta_{\rm lock}=0$, and on the matter-carrying visible
branch $\rho_{\rm vis}+p_{\rm vis}$ is not identically zero. Equation
\eqref{eq:gravity-friction-identity} then has the unique solution
\begin{equation}
\label{eq:gravity-unique-friction}
 \Gamma=H_{\Lambda}.
\end{equation}
Hence the Hubble friction $3H_{\Lambda}$ is the unique isotropic sink that
absorbs $\delta_{\rm topo}$ into the positive vacuum loading. Detuning
$k_{\ell}$ away from $26$ reopens $\Delta_{\rm fr}\neq0$, activates the
anomalous counter-term in Eq.~\eqref{eq:gravity-detuned-einstein}, and
therefore destroys this uniqueness.
Thus,
\begin{equation}
\label{eq:gravity-visible-stress-mismatch}
 \nabla^{\mu}T^{\rm vis}_{\mu 0}
 =
 \dot\rho_{\rm vis}+3H_{\Lambda}(\rho_{\rm vis}+p_{\rm vis})
 =
 \Delta_{\rm lock}(\rho_{\rm vis}+p_{\rm vis})
 \neq 0,
\end{equation}
so the visible stress tensor would fail to be conserved. The term $3H_{\Lambda}$
is therefore the exact Hubble friction that absorbs the packing deficit
$\delta_{\rm topo}\approx0.0112$ into the positive vacuum loading rather than
leaving a Bianchi-violating source term. The same residual therefore reappears as the finite tensor drag $n_t\approx-\delta_{\rm topo}$,
while the scalar sector records the longer-wavelength moir\'e beat
$\Delta_{\rm Moire}=1-n_s$. Together they express the same branch-fixed bulk
consistency: vacuum loading, geometric friction, and the Bianchi lock are one
and the same closure statement viewed at curvature, transport, and tilt level.
At the level of the disclosed branch residues this statement is an exact
Topological Identity: $\mathcal F_{\rm geom}\approx0.01122895$ is the scalar remainder
that the continuity equation must absorb, and the tensor tilt
$n_t\approx-\mathcal F_{\rm geom}\approx-0.0112$ is the tensor-side imprint of
that same conservation law. If either relation failed, the packing deficit
would reappear as $\nabla^{\mu}T_{\mu\nu}\neq0$ and the Bulk Closure Tensor
$\mathcal E_{\mu\nu}$ would no longer vanish on the selected branch.

\paragraph{Gauge-Emergence Exclusion Principle.}
The Bianchi lock also carries a gauge-side exclusion statement. The
Gauge-Emergence Cutoff
\begin{equation}
\label{eq:gravity-gauge-emergence-cutoff-local}
 K_{\rm GE}\equiv\benchmarkParentLevel=312
\end{equation}
is a physical exclusion principle rather than a bookkeeping convention. Because
the admissible parent tower is $K=\benchmarkParentLevel n$, any
\begin{equation}
\label{eq:gravity-sparse-boundary-cutoff-local}
 K\ge624
 \qquad\Longrightarrow\qquad
 n\ge2
 \qquad\Longrightarrow\qquad
 \alpha^{-1}_{\rm surf}\ge 2\alphaSurfBenchmarkExact,
 \quad
 \alpha_{\rm surf}\le\frac{17}{4680},
\end{equation}
so the branch enters the \emph{Sparse Boundary} regime already identified in
Sec.~G.1. In that regime the visible current-current network is too dilute to
sustain an emergent four-dimensional electromagnetic sector: the gauge density
falls below the branch-local support threshold, and the bulk no longer carries
observationally viable electromagnetism. Thus $K=312$ is not merely the first
allowed parent level; it is the unique parent level that is simultaneously
Bianchi-locked and gauge-supporting, while every $K\ge624$ is physically
excluded by sparse-boundary gauge dilution~\cite{friedan1987,witten1989,gko1985}.

\paragraph{Forbidden-Void Proposition.}
A vanishing cosmological constant ($\Lambda = 0$) is topologically forbidden
on the $(26,8,312)$ branch. It necessitates a divergent bit-budget
($N\to\infty$), rendering the boundary partition function physically
non-normalizable. Through Eq.~\eqref{eq:gravity-neutrino-scale-bridge} the same
limit also collapses the neutrino floor~\cite{friedan1987,witten1989,gko1985}. Eq.~\eqref{eq:gravity-bit-budget-derived} makes this implication
explicit on the one-copy lattice. In implication form,
\begin{equation}
\label{eq:gravity-forbidden-void-chain}
 \Lambda_{\rm holo}=0
 \Longrightarrow
 H_{\Lambda}=0
 \Longrightarrow
 N_{\rm sat}=\infty
 \Longrightarrow
 Z_{\partial}^{(26,8,312)}\ \text{non-normalizable}
 \Longrightarrow
 m_{\nu}=\kappa_{D_5}M_PN^{-1/4}=0.
\end{equation}
This proposition is the boundary-normalizability form of the Bianchi lock:
$\Lambda_{\rm holo}=0$ is not a soft limit inside the admissible branch family,
but an exit from the finite-capacity Hilbert sector of the completed
$(26,8,312)$ block.
The proof-by-contradiction is then immediate. Assume that the completed
boundary modular-$T$ matrix still closed on the anomaly-free branch while
$\Lambda_{\rm holo}=0$. Then Eq.~\eqref{eq:gravity-hubble-scale} gives
$H_{\Lambda}=0$, and Eq.~\eqref{eq:gravity-bit-budget-derived} gives
$N_{\rm sat}\to\infty$, so the finite one-copy register decompactifies and the
partition function leaves the normalizable finite-capacity sector. The mass
bridge in Eq.~\eqref{eq:gravity-neutrino-scale-bridge} then forces
$m_{\nu}=\kappa_{D_5}M_PN^{-1/4}\to0$, i.e. the neutrino floor collapses. The
FRW friction channel disappears, but the branch-fixed packing residue
$\delta_{\rm topo}=1-\kappa_{D_5}>0$ remains. Equation
\eqref{eq:gravity-frw-lock-mismatch} would then force
$\Delta_{\rm lock}\neq0$, hence
\begin{equation}
\label{eq:gravity-bianchi-lock-contradiction}
 \nabla^{\mu}T^{\rm vis}_{\mu\nu}\neq0
 \qquad\Longrightarrow\qquad
 \nabla^{\mu}\mathcal E_{\mu\nu}\neq0,
\end{equation}
in contradiction with the Bianchi lock in
Eq.~\eqref{eq:gravity-bianchi-lock-local}. Therefore $\Lambda_{\rm holo}=0$
cannot satisfy both conditions simultaneously: it would require an infinite
bit-budget and thereby violate the finite-capacity premise of the
$(26,8,312)$ lattice. The only branch-compatible solution for which the
completed boundary modular-$T$ matrix closes and the bulk Bianchi identity is
satisfied is the positive vacuum loading already fixed in
Eq.~\eqref{eq:gravity-lambda-positive},
\begin{equation}
 \Lambda_{\rm holo}>0.
\end{equation}

\paragraph{Trap note.}
This is the Bianchi-lock form of the same scale trap derived in
Sec.~G.2. If the lightest neutrino mass were measured significantly
outside the predicted few-meV regime, the resulting bit-budget $N$
would be incompatible with the observed expansion rate, forcing a
non-vanishing framing anomaly $\Delta_{\rm fr}\neq0$ and collapsing the
holographic dictionary~\cite{friedan1987,witten1989,gko1985}.

\manualappendixlabel{G.4}{app:holographic-computational-gradients}
\manualappendixlabel{G.4}{app:clock-skew-derivation}
\subsection*{G.4: Holographic Computational Gradients}

The anomaly-free branch carries a finite horizon register,
$N=3\pi/(L_P^2\Lambda_{\rm holo})$, and that finite register is not only a
storage bound but also a finite update budget. Let $N_{\rm load}(t)$ denote the
cumulative number of boundary bits that have been loaded onto the completion
sector by cosmic time $t$, and define the normalized completion fraction
\begin{equation}
\label{eq:gravity-loading-fraction}
 f_{\rm load}(t)
 \equiv
 \frac{N_{\rm load}(t)}{N},
 \qquad
 0\leq f_{\rm load}(t)\leq 1.
\end{equation}
The physical interpretation of the cosmological expansion is then explicit:
\emph{The expansion of the universe is the loading speed of new boundary bits
onto the $c_{\rm dark}$ stabilizer.}
In the FRW rewriting this means that the late-time expansion rate is the
coarse-grained write speed of the completion fraction,
\begin{equation}
\label{eq:gravity-loading-hubble}
 H_{\rm eff}(t)
 =
 H_{\Lambda}
 +
 \frac{1}{3}\frac{d f_{\rm load}}{dt}.
\end{equation}

The clock-skew statement follows because the loading process is read on the
past light cone, where shells at different redshift do not share the same
proper write clock. If $\tau_{\rm lock}$ is the proper time of a local
boundary-update cycle, then the observer-frame and source-frame ticks are
related by the usual redshift factor,
\begin{equation}
\label{eq:gravity-clock-skew}
 d\tau_{\rm lock} = \frac{dt}{1+z}.
\end{equation}
Therefore the loading law becomes
\begin{equation}
\label{eq:gravity-redshift-dependent-hubble}
 H_0(z)
 =
 H_{\Lambda}
 +
 \frac{1}{3(1+z)}\frac{d f_{\rm load}}{d\tau_{\rm lock}}.
\end{equation}
Equation~\eqref{eq:gravity-redshift-dependent-hubble} is the gravity-side
clock-skew derivation of a redshift-dependent Hubble intercept: once the
boundary register is finite, the same loading process is sampled at different
stages of the write cycle by low-$z$ and high-$z$ observations. A strictly
redshift-independent $H_0$ would require either
$df_{\rm load}/d\tau_{\rm lock}=0$ (no continuing boundary loading) or the
decompactification limit $N\to\infty$ in which the discrete update cycle
ceases to be physical. Neither possibility is compatible with the anomaly-free
branch, which carries a finite register and a nonzero positive vacuum loading.
The local and CMB Hubble determinations therefore probe different segments of
the same finite holographic computation rather than different fluids.

\manualappendixlabel{G.5}{app:gravity-parity-storage}
\subsection*{G.5: Gravitational Storage of B-L Parity}

The visible branch does not erase its missing antimatter complement; it stores
that complement as a topological conjugate in the same geometric sector that is
measured macroscopically by the completion residue $c_{\rm dark}$. In branch
language the visible baryon excess is admissible only because the completed
parent block remains neutral,
\begin{equation}
\label{eq:gravity-parent-neutrality}
 (B-L)_{\rm vis}+(B-L)_{\rm dark}=0,
 \qquad
 (B-L)_{\rm dark}=-(B-L)_{\rm vis}.
\end{equation}
Here $(B-L)_{\rm dark}$ is not a second thermal fluid and not a separately
propagating antimatter bath. It is the non-propagating parity sink of the
geometric completion sector: the missing antimatter of the visible projection
is stored as a topological conjugate in the $c_{\rm dark}$ branch residue.

This is why the same Holographic Stabilizer can play two roles without
contradiction. The residue $c_{\rm dark}$ sources the positive vacuum loading
$\Lambda_{\rm holo}$ in the Topological EFE, and the same residue stores the
conjugate $B-L$ bits needed to keep the projected boundary neutral. The
gravitational field and the parity sink are therefore two macroscopic images of
one and the same completion sector.

The finite de Sitter register of Eq.~\eqref{eq:gravity-horizon-bits} gives the
storage capacity for this parity bookkeeping. In that sense gravitational
storage is literal within the benchmark logic: the branch does not hide the
missing antimatter in an untracked particle sector, but in the geometric
closure data of the completed boundary block. The bulk is neutral because the
same $c_{\rm dark}$ sector that closes $\mathcal E_{\mu\nu}$ also stores the
conjugate parity information required by Eq.~\eqref{eq:gravity-parent-neutrality}.

\ifgravitystandalone
\begin{thebibliography}{99}
\bibitem{georgi1979}
H.~Georgi and C.~Jarlskog,
``A new lepton-quark mass relation in a unified theory,''
\emph{Phys. Lett. B}~86 (1979) 297--300.

\bibitem{babu1993}
K.~S.~Babu and R.~N.~Mohapatra,
``Predictive neutrino spectrum in minimal $SO(10)$ grand unification,''
\emph{Phys. Rev. Lett.}~70 (1993) 2845--2848.

\bibitem{friedan1987}
D.~Friedan and S.~Shenker,
``The Analytic Geometry of Two-Dimensional Conformal Field Theory,''
\emph{Nucl. Phys. B}~281 (1987) 509--545.

\bibitem{witten1989}
E.~Witten,
``Quantum field theory and the Jones polynomial,''
\emph{Commun. Math. Phys.}~121 (1989) 351--399.

\bibitem{gko1985}
P.~Goddard, A.~Kent, and D.~Olive,
``Virasoro algebras and coset space models,''
\emph{Phys. Lett. B}~152 (1985) 88--92.
\end{thebibliography}
\fi

\gravityenddocument
}{\newif\ifgravitystandalone
\let\gravityenddocument\relax
\ifdefined\tnmainflag
\gravitystandalonefalse
\else
\gravitystandalonetrue
\documentclass[11pt]{article}
\usepackage[margin=1in]{geometry}
\usepackage{amsmath,amssymb}
\usepackage{physics}
\usepackage{bm}

\setlength{\emergencystretch}{1.5em}
\raggedbottom

\newcommand{\pubinput}[1]{%
  \IfFileExists{#1}{\input{#1}}{\input{pub/#1}}%
}

\pubinput{physics_constants.tex}

\IfFileExists{tn.aux}{%
  \usepackage{xr}\externaldocument[tn-]{tn}%
}{%
  \IfFileExists{pub/tn.aux}{\usepackage{xr}\externaldocument[tn-]{pub/tn}}{}%
}

\title{Geometric Emergence and the Topological EFE.}
\author{Static Holographic Boundary Theory}
\date{\today}
\def\gravityenddocument{\end{document}}

\begin{document}
\maketitle
\fi

\makeatletter
\providecommand{\manualappendixlabel}[2]{%
  \begingroup
  \def\@currentlabel{#1}%
  \label{#2}%
  \endgroup
}
\newcommand{\gravitymainref}[1]{%
  \@ifundefined{r@#1}{%
    \@ifundefined{r@tn-#1}{\ref{#1}}{\ref{tn-#1}}%
  }{%
    \ref{#1}%
  }%
}
\makeatother

\ifgravitystandalone
\section{Geometric Emergence and the Topological EFE.}\label{app:gravity}\label{sec:gravity-einstein}
\else
\section*{Appendix G (continued): Geometric Emergence and the Topological EFE.}\manualappendixlabel{G}{app:gravity}\manualappendixlabel{G}{sec:gravity-einstein}
\fi

\noindent\textbf{Scope.} This appendix records the gravity-side derivations retained as backup for the anomaly-free $(26,8,312)$ benchmark branch. It is not an independent tuning sector and does not introduce new scan directions: the same modular-$T$ closure data that fix the positive completion residue $c_{\rm dark}$ also fix the positive vacuum loading $\Lambda_{\rm holo}$, the finite horizon register $N$, and the branch-local bulk closure tensor used below as the gravity-side consistency diagnostic. In this appendix the matter-asymmetry bookkeeping is classified as a \emph{Primary Boundary Residue}: the $c_{\rm dark}$ completion sector is read simultaneously as the source of the effective geometric action and as the $B-L$ parity sink that stores the missing conjugate parity bits of the visible projection. The derivation chain below is therefore conditional on the retained SHBT assumptions---static coarse-grained bulk bookkeeping, finite one-copy capacity, and the anomaly-free framing condition $\Delta_{\rm fr}=0$---rather than a claim of a complete dynamical quantum-gravity construction.

\noindent\textbf{Conceptual Scope.} The gravity-side picture used here is static in the deliberately narrow SHBT sense: it describes a ``frozen'' coarse-grained bulk fixed by the completed boundary data, not a time-dependent evaporation model. The finite bit budget $N$ inferred from $\Lambda_{\rm obs}$ is treated throughout as an \emph{Observational Boundary Condition}, not as an internal fit direction. Accordingly, any Page-curve-style identities are used only as bookkeeping constraints on finite information capacity and code-subspace reconstruction. They are not claimed to provide an independent microscopic mechanism for dynamical black-hole evaporation.

\noindent\textbf{Benchmark status.} Unless explicitly stated otherwise, the detuning arguments below are internal stress tests of the retained anomaly-free branch. They show how the displayed closure identities fail under one-step shifts of the benchmark integers or under loss of finite-capacity saturation; they are not presented as no-go theorems for every non-minimal completion outside the one-copy SHBT bookkeeping.

\manualappendixlabel{G.1}{app:curvature-sign-test}
\subsection*{G.1: Bulk Closure Tensor, Holographic Surface Tension, Finite Bit Budget, and the Topological EFE}

On the anomaly-free branch the modular complement required for fixed-parent modular-$T$ closure is the positive completion residue,
\begin{equation}
\label{eq:gravity-closure-residue}
 c_{\rm dark}=\cDarkCompletion\approx \cDarkCompletionFull>0,
\end{equation}
and the same branch data fix the positive vacuum loading
\begin{equation}
\label{eq:gravity-lambda-positive}
 \Lambda_{\rm holo}=+\frac{3\pi}{L_P^2N}=+\lambdaHoloMetersInverseSquared\,\mathrm{m}^{-2}\approx +1.08913883\times10^{-52}\,\mathrm{m}^{-2}>0.
\end{equation}
In the present appendix, the quantity $\Lambda_{\rm holo}$ is not interpreted
as a separately propagating dark-energy fluid. It is the positive
\emph{Holographic Surface Tension} of the completed $(26,8,312)$ branch: the
geometric tension term required to keep the boundary partition function
normalizable on a finite horizon register.
Write the associated de Sitter horizon scale as
\begin{equation}
\label{eq:gravity-horizon-radius}
 R_H
 \equiv
 H_{\Lambda}^{-1}
 =
 \sqrt{\frac{3}{\Lambda_{\rm holo}}},
 \qquad
 A_H=4\pi R_H^2.
\end{equation}
In the frozen-bulk bookkeeping used throughout this appendix, the completed
$SO(10)_{312}$ branch must close on a finite Euclidean $S^3$ of radius $R_H$.
Its holographic screen is the associated horizon $S^2$ of area $A_H$, so the
one-copy capacity is fixed by the saturated holographic count
\begin{equation}
\label{eq:gravity-bit-budget-derived}
 N_{\rm sat}
 \equiv
 \frac{A_H}{4L_P^2}
 =
 \frac{\pi R_H^2}{L_P^2}
 =
 \frac{3\pi}{L_P^2\Lambda_{\rm holo}}.
\end{equation}
This finite-capacity statement is the normalizability condition for the
completed branch partition function. If the horizon size were allowed to run to
infinity, then $N_{\rm sat}\to\infty$, the one-copy character sum would no
longer define a finite-capacity boundary dictionary, and the completed
$SO(10)_{312}$ partition function would cease to be physically admissible as a
normalizable branch state.
These two quantities are not independent outputs of the branch. The same
$c_{\rm dark}$ sector that sources the effective geometric action through the
positive vacuum loading also acts as the storage register for the missing
$B-L$ parity bits of the visible projection. Vacuum loading and parity-sink
bookkeeping are therefore two boundary images of the same completion residue.
This also removes the usual fine-tuning reading of the cosmological term in the
present framework. Here $\Lambda_{\rm holo}$ is not a tiny vacuum counter-term
balanced against a separate large bare contribution; it is the
\emph{Holographic Surface Tension} required to hold the topological moir\'e
pattern of the $(26,8)$ cell on a finite branch-compatible horizon and to bind
the completed $SO(10)_{312}$ partition function to the finite one-copy
register $N_{\rm sat}$. If one were to send $\Lambda_{\rm holo}\to0^+$, then
Eq.~\eqref{eq:gravity-bit-budget-derived} would force $N_{\rm sat}\to\infty$,
so the completed partition function would become divergent rather than
natural. Equivalently, the mass bridge of
Eq.~\eqref{eq:gravity-neutrino-scale-bridge} would collapse the neutrino floor
to zero as $N^{-1/4}\to0$. The positive sign in
Eq.~\eqref{eq:gravity-lambda-positive} is therefore part of the Bianchi lock
itself: finite positive surface tension, finite register, and non-vanishing
neutrino floor are one closure statement viewed in curvature, information, and
mass units.
The same closure is recorded in two branch-fixed matrices that will be used
repeatedly below. Define the \emph{Topological Identity Matrix}
\begin{equation}
\label{eq:gravity-topological-identity-matrix}
 I_{\rm topo}
 \equiv
 \mathbf I_{\rm topo}
 \equiv
 \operatorname{diag}\!\left(
 \,k_{\ell},
 \,k_q,
 \,K
 \right)
 =
 \operatorname{diag}\!\left(
 \,26,
 \,8,
 \,312
 \right).
\end{equation}
This matrix records only the discrete branch labels of the anomaly-free cell.
Separately, define the \emph{Matching Matrix}
\begin{equation}
\label{eq:gravity-topological-identity-basis}
 \mathbf M_{\rm match}
 =
 \operatorname{diag}\!\left(
 \,\kappa_{D_5},
 \,\mathcal R_{\rm GUT},
 \,r_*
 \right)
 =
 \operatorname{diag}\!\left(
 \,\simplexPrefactor,
 \,\frac{8}{28},
 \,\frac{64}{312}
 \right),
 \qquad
 r_*\equiv\frac{\langle\Sigma_{126}\rangle}{\langle\phi_{10}\rangle}=\frac{64}{312}.
\end{equation}
The matching matrix collects the branch-fixed residues that connect the rigid
topological data to the finite-capacity mass bridge, the CKM threshold
closure, and the benchmark scalar alignment. The scalar ratio rescales the
singular spectrum without perturbing the code-subspace isometries. Define the
associated matching defect by
\begin{equation}
\label{eq:gravity-topological-identity-defect}
 \Delta \mathbf M_{\rm match}
 \equiv
 \mathbf M_{\rm match}
 -
 \operatorname{diag}\!\left(
 \,\simplexPrefactor,
 \,\frac{8}{28},
 \,\frac{64}{312}
 \right).
\end{equation}
Then
\begin{equation}
\label{eq:gravity-topological-identity-saturation}
 \Delta \mathbf M_{\rm match}=0
 \qquad\Longleftrightarrow\qquad
 \mathbf M_{\rm match}
 =
 \operatorname{diag}\!\left(
 \,\simplexPrefactor,
 \,\frac{8}{28},
 \,\frac{64}{312}
 \right),
\end{equation}
while $\mathbf I_{\rm topo}$ remains fixed to the discrete branch labels of
Eq.~\eqref{eq:gravity-topological-identity-matrix}. Saturation is therefore
componentwise: $\Delta\mathbf M_{\rm match}=0$ if and only if each matching
residue attains its branch-fixed benchmark value. That is the precise sense in
which the anomaly-free branch is a topological identity equipped with a fixed
matching prescription rather than a continuously deformable family~\cite{georgi1979,babu1993,friedan1987,witten1989}.
Define the branch-closure tensor by
\begin{equation}
\label{eq:gravity-closure-tensor-local}
 \mathcal E_{\mu\nu}
 \equiv
 \frac{c_{\rm dark}M_P^2}{12}\bigl(G_{\mu\nu}+\Lambda_{\rm holo}g_{\mu\nu}\bigr)
 -T^{\rm vis}_{\mu\nu},
 \qquad
 T^{\rm vis}_{\mu\nu}
 \equiv
 T^{(SU(2)_{26})}_{\mu\nu}+T^{(SU(3)_8)}_{\mu\nu}.
\end{equation}
Equation~\eqref{eq:gravity-closure-tensor-local} defines the Bulk Closure Tensor $\mathcal E_{\mu\nu}$. The vanishing of this tensor is the necessary and sufficient condition for a framing-anomaly-free ($\Delta_{\rm fr}=0$) boundary. In the present appendix, $\mathcal E_{\mu\nu}=0$ should be read as a consistency check confirming that the selected RCFT satisfies the standard AdS/CFT isometry requirements needed for code-subspace reconstruction on the anomaly-free branch. The bulk image of modular-$T$ closure is that this tensor can only differ from zero by the same scalar framing defect that measures failure of integer-spectrum closure,
\begin{equation}
\label{eq:gravity-topological-identity-local}
 \mathcal E_{\mu\nu}=\mathcal Q_{\rm iso}\,\Delta_{\rm fr}\,g_{\mu\nu},
\end{equation}
with $\mathcal Q_{\rm iso}\neq 0$ the branch-fixed isotropic anomaly density determined by the parent, visible, and completion central charges. The branch identity therefore reads
\begin{equation}
 \mathcal E_{\mu\nu}=0
 \qquad\Longleftrightarrow\qquad
 \Delta_{\rm fr}=0.
\end{equation}
Bulk diffeomorphism invariance is therefore not an independent postulate in this benchmark appendix: the contracted Bianchi identity is the bulk Ward identity of the same boundary framing condition that sets $\Delta_{\rm fr}=0$.
Equivalently, the effective Einstein system is not an extra dynamical postulate but the covariant restatement of the branch closure,
\begin{equation}
\label{eq:gravity-effective-einstein-local}
 G_{\mu\nu}+\Lambda_{\rm holo}g_{\mu\nu}
 =
 \frac{12}{c_{\rm dark}M_P^2}T^{\rm vis}_{\mu\nu}.
\end{equation}
This is the precise sense in which a four-dimensional bulk description emerges on the anomaly-free branch: the effective Einstein system is the covariant image of the discrete closure data encoded in $\mathbf I_{\rm topo}$, while $\mathbf M_{\rm match}$ enters only through flavor-sector singular-value matching. Equivalently, the branch carries the explicit \textit{Topological EFE Identity}
\begin{equation}
\label{eq:gravity-topological-efe-identity}
 \Delta_{\rm fr}=0
 \qquad\Longleftrightarrow\qquad
 \mathcal E_{\mu\nu}=0
 \qquad\Longleftrightarrow\qquad
 G_{\mu\nu}+\Lambda_{\rm holo}g_{\mu\nu}
 =
 8\pi G_N\,T^{\rm vis}_{\mu\nu},
\end{equation}
with the Newton coupling locked by the completion residue,
\begin{equation}
\label{eq:gravity-newton-lock-local}
 8\pi G_N
 \equiv
 \frac{12}{c_{\rm dark}M_P^2}.
\end{equation}
In that sense Newton's constant is not an independent infrared dial: the gravitational coupling exists on the anomaly-free branch precisely to carry the central-charge deficit $c_{\rm dark}\approx\cDarkCompletion$ left by the visible projection. The parity-storage role of the same residue is isolated separately in Sec.~G.5.

\paragraph{Saturating Information Capacity.}
Evaluating Eq.~\eqref{eq:gravity-bit-budget-derived} on the branch-fixed
vacuum loading of Eq.~\eqref{eq:gravity-lambda-positive} gives the unique
finite register
\begin{equation}
\label{eq:gravity-saturating-information-capacity}
 N_{\rm sat}
 =
 \frac{3\pi}{L_P^2\Lambda_{\rm holo}}
 \approx
 \modularHorizonBits\ \text{bits}.
\end{equation}
This number is not merely an upper bound; it is the unique one-copy storage
capacity on which the completed boundary modular-$T$ spectrum closes. To make
that necessity explicit, write
\begin{equation}
\label{eq:gravity-capacity-detuning-parameter}
 N=N_{\rm sat}+\delta N.
\end{equation}
Holding the completed $(26,8,312)\oplus c_{\rm dark}$ block fixed, a register
detuning changes the surface-tension term by
\begin{equation}
\label{eq:gravity-capacity-lambda-detuning}
 \delta_N\Lambda_{\rm holo}
 =
 \frac{3\pi}{L_P^2}\left(\frac{1}{N_{\rm sat}+\delta N}-\frac{1}{N_{\rm sat}}\right)
 =
 -\frac{3\pi\,\delta N}{L_P^2N_{\rm sat}(N_{\rm sat}+\delta N)}.
\end{equation}
Because the completed block data are held fixed, $\delta_NT^{\rm vis}_{\mu\nu}=0$.
Equation~\eqref{eq:gravity-closure-tensor-local} therefore gives
\begin{equation}
\label{eq:gravity-capacity-closure-detuning}
 \delta_N\mathcal E_{\mu\nu}
 =
 \frac{c_{\rm dark}M_P^2}{12}\,\delta_N\Lambda_{\rm holo}\,g_{\mu\nu}.
\end{equation}
Hence $\delta_N\mathcal E_{\mu\nu}=0$ if and only if $\delta N=0$. If
$\delta N<0$, part of the completed modular-$T$ spectrum lacks storage support
on the one-copy screen; if $\delta N>0$, the excess register is not sourced by
the completed block and the partition sum leaves the one-copy normalizable
sector. Therefore
\begin{equation}
\label{eq:gravity-capacity-iff-closure}
 \mathcal E_{\mu\nu}=0
 \qquad\Longrightarrow\qquad
 N=N_{\rm sat}\approx\modularHorizonBits,
\end{equation}
so $N\approx3.42\times10^{122}$ is the \emph{Saturating Information Capacity}
required for modular-$T$ spectrum closure of the completed block~\cite{friedan1987,witten1989,gko1985}.

\paragraph{Equivalence-Principle proof by discrete detuning.}
Let
\begin{equation}
\label{eq:gravity-discrete-coordinate-vector}
 X\equiv(G_{\rm SM},k_{\ell},k_q,K),
 \qquad
 \delta X_a\equiv\delta_1 X_a\in\{\pm1\}
\end{equation}
 denote the smallest admissible detuning of a single discrete branch coordinate.
For the fixed-parent visible embedding, integer modular-$T$ closure is
equivalent to the descendant-integrality conditions
\begin{equation}
\label{eq:gravity-discrete-branch-integrality}
 \Delta_{\rm fr}=0
 \qquad\Longleftrightarrow\qquad
 \frac{K}{2k_{\ell}}\in\mathbb Z
 \quad\text{and}\quad
 \frac{K}{3k_q}\in\mathbb Z.
\end{equation}
Because these coordinates are discrete rather than continuous, proving
$\delta_1\mathcal E_{\mu\nu}\neq0$ for such one-step moves is the bulk form of
rigidity: there is no infinitesimal deformation direction that preserves the
Einstein identity.
Linearizing Eq.~\eqref{eq:gravity-closure-tensor-local} around the benchmark
cell gives
\begin{equation}
\label{eq:gravity-discrete-closure-variation}
 \delta_1\mathcal E_{\mu\nu}
 =
 \frac{c_{\rm dark}M_P^2}{12}\,\delta_1\Lambda_{\rm holo}\,g_{\mu\nu}
 -\delta_1 T^{\rm vis}_{\mu\nu}.
\end{equation}
For a one-unit leptonic-level detuning with the parent held fixed,
\begin{equation}
\label{eq:gravity-kl-unit-detuning}
 k_{\ell}:26\to26\pm1
 \qquad\Longrightarrow\qquad
 \frac{K}{2k_{\ell}}
 \in
 \left\{\frac{312}{50},\frac{312}{54}\right\}
 =
 \left\{\frac{156}{25},\frac{52}{9}\right\}
 \notin\mathbb Z,
\end{equation}
so the leptonic descendant multiplicity ceases to be integral and therefore
$\delta_1\Delta_{\rm fr}\neq0$. By Eq.~\eqref{eq:gravity-topological-identity-local}
this implies
\begin{equation}
\label{eq:gravity-kl-detuning-nonzero-closure}
 \delta_1\mathcal E_{\mu\nu}
 =
 \mathcal Q_{\rm iso}\,\delta_1\Delta_{\rm fr}\,g_{\mu\nu}
 \neq0.
\end{equation}
For a one-unit quark-level detuning with the parent held fixed,
\begin{equation}
\label{eq:gravity-kq-unit-detuning}
 k_q:8\to8\pm1
 \qquad\Longrightarrow\qquad
 \frac{K}{3k_q}
 \in
 \left\{\frac{312}{21},\frac{312}{27}\right\}
 =
 \left\{\frac{104}{7},\frac{104}{9}\right\}
 \notin\mathbb Z,
\end{equation}
so the descendant multiplicity ceases to be integral and therefore
$\Delta_{\rm fr}\neq0$. By Eq.~\eqref{eq:gravity-topological-identity-local}
this implies
\begin{equation}
\label{eq:gravity-kq-detuning-nonzero-closure}
 \delta_1\mathcal E_{\mu\nu}
 =
 \mathcal Q_{\rm iso}\,\delta_1\Delta_{\rm fr}\,g_{\mu\nu}
 \neq0.
\end{equation}
For a one-unit parent-level detuning with the visible branch held fixed,
\begin{equation}
\label{eq:gravity-parent-unit-detuning}
 K:312\to312\pm1
 \qquad\Longrightarrow\qquad
 K\pm1\not\equiv0\pmod{52}
 \quad\text{and}\quad
 K\pm1\not\equiv0\pmod{24},
\end{equation}
so neither visible descendant multiplicity remains integral and again
$\delta_1\Delta_{\rm fr}\neq0$. Therefore
\begin{equation}
\label{eq:gravity-parent-detuning-nonzero-closure}
 \delta_1\mathcal E_{\mu\nu}
 =
 \mathcal Q_{\rm iso}\,\delta_1\Delta_{\rm fr}\,g_{\mu\nu}
 \neq0.
\end{equation}
For a one-unit visible-current detuning with the branch levels held fixed,
\begin{equation}
\label{eq:gravity-gsm-unit-detuning}
 G_{\rm SM}:15\to15\pm1
 \qquad\Longrightarrow\qquad
 \delta_1\alpha^{-1}_{\rm surf}
 =
 \pm\frac{K}{k_{\ell}+k_q}
 =
 \pm\frac{312}{34}
 =
 \pm\frac{156}{17}
 \neq0,
\end{equation}
so the gauge part of the visible stress tensor shifts,
$\delta_1T^{\rm vis}_{\mu\nu}\supset\delta_1T^{\rm gauge}_{\mu\nu}\neq0$,
and Eq.~\eqref{eq:gravity-discrete-closure-variation} again gives
$\delta_1\mathcal E_{\mu\nu}\neq0$. Thus every one-step move of a
load-bearing coordinate $X_a\in\{G_{\rm SM},k_{\ell},k_q,K\}$ reopens the bulk
closure defect. Because the benchmark branch already satisfies
$\mathcal E_{\mu\nu}^{\rm bench}=0$, the detuned branch obeys
\begin{equation}
\label{eq:gravity-discrete-detuned-closure}
 \mathcal E_{\mu\nu}^{\rm det}
 =
 \delta_1\mathcal E_{\mu\nu}
 \neq 0
 \qquad\text{for every}\qquad
 \delta X_a=\pm1.
\end{equation}
This is the bulk-closure form of the reviewer trap: there is no hidden
compensating direction that can restore $\mathcal E_{\mu\nu}=0$ after a
one-step shift of any load-bearing integer coordinate. Writing the detuned
field equation as
\begin{equation}
\label{eq:gravity-detuned-einstein-identity}
 G_{\mu\nu}+\Lambda_{\rm holo}g_{\mu\nu}
 =
 8\pi G_N\,T^{\rm vis}_{\mu\nu}
 +\mathcal J_{\mu\nu}^{(a)},
 \qquad
 \mathcal J_{\mu\nu}^{(a)}
 \equiv
 \frac{12}{c_{\rm dark}M_P^2}\,\delta_1\mathcal E_{\mu\nu},
\end{equation}
the extra source satisfies $\mathcal J_{\mu\nu}^{(a)}\neq0$ for every
$\delta_1X_a=\pm1$. Because $\mathcal J_{\mu\nu}^{(a)}$ depends on which
coordinate was detuned, the response is no longer universal and cannot be
absorbed into the single branch-fixed Newton coupling of
Eq.~\eqref{eq:gravity-newton-lock-local}. The Einstein identity is therefore
violated and the bulk Equivalence Principle fails under any $\pm1$ discrete
coordinate shift~\cite{friedan1987,witten1989,gko1985}. Equivalently,
\begin{equation}
 \delta X_a=\pm1
 \qquad\Longrightarrow\qquad
 \delta\mathcal E_{\mu\nu}\neq0
 \qquad\Longrightarrow\qquad
 \mathcal J_{\mu\nu}^{(a)}\neq0,
\end{equation}
so freely falling bulk probes no longer couple to a universal geometric source.
Universality of free fall is therefore lost for every one-step detuning of a
load-bearing branch coordinate.
Within this branch bookkeeping, the Bianchi lock and the holographic bit-budget
saturation are the same statement viewed from the bulk and boundary sides. The
vanishing of the Bulk Closure Tensor requires not only $\Delta_{\rm fr}=0$ but
also that the horizon register sit at its saturated value
$N=N_{\rm sat}=\pi R_H^2/L_P^2$; otherwise $\Lambda_{\rm holo}$ is detuned,
the packing residue is no longer fully absorbed by the isotropic vacuum load,
and the bulk Ward identity reopens. In that sense
\begin{equation}
\label{eq:gravity-bianchi-bit-lock}
 \mathcal E_{\mu\nu}=0
 \qquad\Longrightarrow\qquad
 \nabla^{\mu}\mathcal E_{\mu\nu}=0
 \quad\text{with}\quad
 N=N_{\rm sat}\equiv\frac{\pi R_H^2}{L_P^2},
\end{equation}
while Sec.~G.3 spells out the same lock in FRW language.

\paragraph{The Bianchi rebuttal.}
The neighboring fixed-parent cells make the framing/diffeomorphism link explicit. Holding $(k_q,K)=(8,312)$ fixed while detuning the leptonic level away from $k_{\ell}=26$ gives the fractional vacuum modular-$T$ powers
\begin{equation}
\label{eq:gravity-detuned-vacuum-T}
 \frac{K}{2k_{\ell}}-6
 =
 \begin{cases}
 0.24, & k_{\ell}=25,\\[2pt]
 -\dfrac{2}{9}, & k_{\ell}=27,
 \end{cases}
 \qquad\Longrightarrow\qquad
 e^{2\pi i\left(\frac{K}{2k_{\ell}}-6\right)}\neq 1.
\end{equation}
Equivalently, the detuned branches carry non-vanishing framing defects,
\begin{equation}
\label{eq:gravity-detuned-framing-defect}
 \Delta_{\rm fr}(25,8)=0.24,
 \qquad
 \Delta_{\rm fr}(27,8)=\frac{2}{9},
\end{equation}
so the vacuum $T$-power is fractional rather than integer-valued. In bulk language Eq.~\eqref{eq:gravity-topological-identity-local} no longer closes to zero; instead the effective Einstein equations acquire the anomalous isotropic counter-term
\begin{equation}
\label{eq:gravity-detuned-einstein}
 G_{\mu\nu}+\Lambda_{\rm holo}g_{\mu\nu}
 =
 \frac{12}{c_{\rm dark}M_P^2}\Bigl(T^{\rm vis}_{\mu\nu}+T^{\rm anom}_{\mu\nu}\Bigr),
 \qquad
 T^{\rm anom}_{\mu\nu}
 \equiv
 \mathcal Q_{\rm iso}\,\Delta_{\rm fr}\,g_{\mu\nu}.
\end{equation}
The rebuttal point is then sharp: detuning $k_{\ell}$ to $25$ or $27$ does not merely perturb a fit metric; it injects an anomalous counter-term into the covariant bulk equations because the boundary vacuum $T$-power no longer closes integrally. The benchmark branch $(26,8,312)$ is distinguished precisely because $\Delta_{\rm fr}=0$ removes this extra source and leaves the bulk equations counter-term free.

The same branch data also carry the residual packing pressure
\begin{equation}
\label{eq:gravity-delta-topo}
 \delta_{\rm topo}
 \equiv
 1-\kappa_{D_5}
 \approx
 0.01122895,
\end{equation}
which is the non-unit $D_5$ weight-simplex remainder left after visible packing. The same positive remainder appears in the gauge-side closure language of the main text through the benchmark threshold row
\begin{equation}
\label{eq:gravity-alpha-closure}
 \alpha^{-1}(M_{\rm match})
 =
 137.035999\ldots+\delta_{\rm topo}.
\end{equation}
The gravity-side statement is stronger: the same residue is absorbed into the positive dS$_4$ vacuum loading rather than left as a free stress defect.

\paragraph{Gauge-Emergence Cutoff.}
The same branch closure also fixes the ultraviolet gauge-density proxy carried by the holographic screen. With the visible current normalization $G_{\rm SM}=15$, the surface-tension proxy is
\begin{equation}
\label{eq:gravity-surface-gauge-proxy}
 \alpha^{-1}_{\rm surf}
 =
 G_{\rm SM}\frac{K}{k_{\ell}+k_q}
 =
 \frac{15K}{k_{\ell}+k_q}.
\end{equation}
On the retained visible branch $(k_{\ell},k_q)=\benchmarkVisiblePair$ the Diophantine descendant rule forces the parent tower
\begin{equation}
\label{eq:gravity-parent-tower}
 K\in \operatorname{lcm}(2k_{\ell},3k_q)\,\mathbb Z_{>0}
 =\benchmarkParentLevel\,\mathbb Z_{>0},
 \qquad
 K=\benchmarkParentLevel n,
 \quad n\in\mathbb Z_{>0},
\end{equation}
so the surface datum becomes
\begin{equation}
\label{eq:gravity-surface-gauge-tower}
 \alpha^{-1}_{\rm surf}
 =
 15\frac{\benchmarkParentLevel n}{\benchmarkLeptonLevel+\benchmarkQuarkLevel}
=
 \alphaSurfBenchmarkExact\,n.
\end{equation}
The first admissible parent therefore gives
\begin{equation}
\label{eq:gravity-surface-gauge-benchmark}
 K=\benchmarkParentLevel
 \qquad\Longrightarrow\qquad
 \alpha^{-1}_{\rm surf}
 =
 \alphaSurfBenchmarkExact
 \approx
 \alphaSurfBenchmarkDecimal,
\end{equation}
which is the physical benchmark value of the ultraviolet surface gauge datum. By contrast every higher admissible parent lies in the sparse-boundary regime,
\begin{equation}
\label{eq:gravity-sparse-boundary}
 K>312
 \quad\Longrightarrow\quad
 n\ge 2
 \quad\Longrightarrow\quad
 \alpha^{-1}_{\rm surf}
 \ge
 2\alphaSurfBenchmarkExact
 \approx
 275.294118,
\end{equation}
or equivalently
\begin{equation}
\label{eq:gravity-sparse-boundary-direct}
 \alpha_{\rm surf}
 \le
 \frac{17}{4680}
 \approx
 3.63\times10^{-3}
 =
 \frac12\,\alpha_{\rm surf}^{\rm bench}.
\end{equation}
Thus increasing $K$ above the first admissible parent does not create new visible gauge channels; it dilutes the same fifteen visible current slots across a larger parent cycle and makes the boundary gauge network progressively sparser. In the SHBT dictionary bulk gauge fields are the continuum image of these boundary current-current correlators, so the sparse-boundary tower drives the emergent electromagnetic sector toward decoupling. In the formal limit $K\to\infty$ one has $\alpha_{\rm surf}\to0$, and the bulk gauge field vanishes altogether. Already the first higher multiple $K=624$ enters the regime $\alpha^{-1}_{\rm surf}\ge275$, where the boundary gauge density is too weak to support an observationally viable four-dimensional bulk with detectable electromagnetic interactions.

The exclusion of $K>\benchmarkParentLevel$ is therefore geometric rather than merely numerological. A branch with $K>\benchmarkParentLevel$ would produce a bulk in which electromagnetic forces are effectively absent at the emergence scale, so charged matter would fail to organize into the observed four-dimensional sector. The retained parent level is consequently locked not only by framing closure but also by gauge detectability:
\begin{equation}
\label{eq:gravity-gauge-detectability-lock}
 \text{detectable bulk gauge fields}
 \qquad\Longleftrightarrow\qquad
 K=\benchmarkParentLevel
\end{equation}
within the canonical one-copy $SO(10)\supset SU(3)\times SU(2)$ branch used here. In that precise sense the stability of $K$ is anchored by the very existence of measurable gauge fields in the emergent bulk.

Equivalently, define the gauge-support ratio on the admissible parent tower by
\begin{equation}
\label{eq:gravity-gauge-support-ratio}
 \rho_{\rm gauge}(K)
 \equiv
 \frac{\alpha_{\rm surf}(K)}{\alpha_{\rm surf}(\benchmarkParentLevel)}
 =
 \frac{\benchmarkParentLevel}{K}
 =
 \frac{1}{n},
 \qquad
 K=\benchmarkParentLevel n.
\end{equation}
Avoiding the Sparse Boundary onset is the strict support condition
\begin{equation}
\label{eq:gravity-gauge-support-condition}
 \alpha^{-1}_{\rm surf}(K)<2\alphaSurfBenchmarkExact
 \qquad\Longleftrightarrow\qquad
 \rho_{\rm gauge}(K)>\frac12.
\end{equation}
Because the admissible tower has $n\in\mathbb Z_{>0}$, Eq.~\eqref{eq:gravity-gauge-support-condition} implies $n<2$, hence necessarily $n=1$ and therefore
\begin{equation}
\label{eq:gravity-unique-gauge-supporting-parent}
 K=\benchmarkParentLevel=312.
\end{equation}
Thus $K=312$ is the \emph{Unique Gauge-Supporting Parent}: every higher multiple $K\ge624$ lies in the Sparse Boundary regime, where the visible current network is too dilute to sustain an emergent four-dimensional electromagnetic sector~\cite{friedan1987,witten1989,gko1985}.

\manualappendixlabel{G.2}{app:proof-of-mass}
\manualappendixlabel{G.2}{app:holographic-surface-tension}
\manualappendixlabel{G.2}{app:first-law-efe}
\subsection*{G.2: Holographic Surface Tension, Newton Lock, Hubble Friction, and Information Return}

The positive vacuum loading in Eq.~\eqref{eq:gravity-lambda-positive} is the
branch-fixed holographic surface tension of the completed $(26,8,312)$ cell.
Equations~\eqref{eq:gravity-horizon-radius} and
\eqref{eq:gravity-bit-budget-derived} then fix the branch-local de Sitter
register through
\begin{equation}
\label{eq:gravity-horizon-bits}
 N_{\rm holo}\equiv N
 =
 \frac{3\pi}{L_P^2\Lambda_{\rm holo}}
 \approx \maxComplexityCapacity\ \text{bits}.
\end{equation}
The same finite register is the coarse-grained storage capacity for the
$B-L$ parity sink carried by $c_{\rm dark}$. In that sense the de Sitter
horizon count is simultaneously the vacuum-loading register of the geometric
action and the bookkeeping register that stores the missing conjugate parity
bits required for parent-block neutrality.
Combining this with the branch-fixed mass relation
\begin{equation}
\label{eq:gravity-neutrino-scale-bridge}
 m_{\nu}=\kappa_{D_5}M_PN^{-1/4}
\end{equation}
gives
\begin{equation}
\label{eq:gravity-neutrino-scale-bitcount}
 N
 =
 \frac{\kappa_{D_5}^4M_P^4}{m_{\nu}^4}.
\end{equation}
Substituting Eq.~\eqref{eq:gravity-neutrino-scale-bitcount} into
Eq.~\eqref{eq:gravity-horizon-bits}, and using $G_N=L_P^2=M_P^{-2}$ in the
present normalization, yields the \emph{Unity of Scale Identity}
\begin{equation}
\label{eq:gravity-unity-of-scale-identity}
 \Lambda_{\rm holo}
 =
 \frac{3\pi}{\kappa_{D_5}^4}
 G_N m_{\nu}^4.
\end{equation}
In the present appendix Eq.~\eqref{eq:gravity-unity-of-scale-identity} is not
just a derived curiosity. It is the benchmark's \emph{primary internal
falsification metric}: once the anomaly-free branch fixes $\kappa_{D_5}$ and
the observed cosmological loading fixes $\Lambda_{\rm holo}$ through
Eq.~\eqref{eq:gravity-horizon-bits}, the lightest neutrino scale is no longer
free to drift independently. Any empirical failure of
Eq.~\eqref{eq:gravity-unity-of-scale-identity} would be an internal failure of
the branch closure itself rather than a poor phenomenological fit that could be
rescued by retuning hidden parameters. Operationally, once
$\Lambda_{\rm holo}$ and $G_N$ are fixed, Eq.~\eqref{eq:gravity-unity-of-scale-identity}
is the first gravity-side number to check: the lightest neutrino mass must
land in the branch-compatible few-meV regime or the anomaly-free dictionary
fails internally.

\paragraph{Formal proof of the Unity of Scale Identity.}
Starting from Eq.~\eqref{eq:gravity-neutrino-scale-bridge},
\begin{equation}
 m_{\nu}^4
 =
 \kappa_{D_5}^4 M_P^4 N^{-1},
 \qquad\Longrightarrow\qquad
 N
 =
 \frac{\kappa_{D_5}^4 M_P^4}{m_{\nu}^4}.
\end{equation}
Substituting this expression for $N$ into Eq.~\eqref{eq:gravity-horizon-bits}
gives
\begin{equation}
 \Lambda_{\rm holo}
 =
 \frac{3\pi}{L_P^2}
 \frac{m_{\nu}^4}{\kappa_{D_5}^4 M_P^4}.
\end{equation}
Using the branch normalization $G_N=L_P^2=M_P^{-2}$, one has
\begin{equation}
 L_P^2 M_P^4 = G_N^{-1},
\end{equation}
so the previous line becomes
\begin{equation}
 \Lambda_{\rm holo}
 =
 \frac{3\pi}{\kappa_{D_5}^4}
 G_N m_{\nu}^4,
\end{equation}
which is exactly Eq.~\eqref{eq:gravity-unity-of-scale-identity}. Because the
anomaly-free branch has $\kappa_{D_5}>0$, $G_N>0$, and $m_{\nu}>0$, the right-
hand side is strictly positive; the Unity of Scale Identity therefore proves
that the branch-compatible $\Lambda_{\rm holo}$ is a positive surface-tension
term rather than a freely choosable fluid contribution.
The cosmological surface tension, Newton coupling, and neutrino mass floor are
therefore not independent phenomenological dials in this benchmark; they are
the same saturated branch datum expressed in curvature, gravitational, and
defect-mass units.

\paragraph{The scale trap.}
This identity is also the theory's primary internal falsification metric and
simplest internal trap door. If the lightest
neutrino mass were measured significantly above the branch-fixed meV regime---
for example near $10\,\mathrm{meV}$---then
Eq.~\eqref{eq:gravity-neutrino-scale-bitcount} would force a horizon register
$N\propto m_{\nu}^{-4}$ smaller by a factor of roughly
$\sim1.6\times10^2$ than the saturation value fixed by the observed
$\Lambda_{\rm holo}$. Feeding that detuned register back into
Eq.~\eqref{eq:gravity-horizon-bits} would in turn require a surface tension
$\Lambda_{\rm holo}$ larger by the same factor, incompatible with the observed
cosmological constant. By Eq.~\eqref{eq:gravity-capacity-closure-detuning}
this implies $\delta_N\mathcal E_{\mu\nu}\neq0$; by
Eq.~\eqref{eq:gravity-topological-identity-local} the same mismatch is exactly
the reopening of framing closure, i.e. $\Delta_{\rm fr}\neq0$. In that sense a
ten-meV lightest neutrino would not merely shift the benchmark; it would
falsify the anomaly-free branch by forcing a bit budget incompatible with the
observed vacuum loading.

The corresponding geometric surface-tension action is
\begin{equation}
\label{eq:gravity-surface-action}
 S_{\rm surf}[g]
 =
 -\sigma_{\rm holo}\int d^4x\,\sqrt{-g},
 \qquad
 \sigma_{\rm holo}\equiv\frac{\Lambda_{\rm holo}}{8\pi G_N}.
\end{equation}
Varying with respect to the metric gives the effective geometric stress-energy tensor
\begin{equation}
\label{eq:gravity-geom-stress}
 T_{\mu\nu}^{(\mathrm{geom})}
 \equiv
 -\frac{2}{\sqrt{-g}}\frac{\delta S_{\rm surf}}{\delta g^{\mu\nu}}
 =
 -\sigma_{\rm holo}g_{\mu\nu}
 =
 -\frac{\Lambda_{\rm holo}}{8\pi G_N}\,g_{\mu\nu}.
\end{equation}
Equations~\eqref{eq:gravity-effective-einstein-local} and
\eqref{eq:gravity-newton-lock-local} therefore identify the geometric surface
stress with the same completion residue that closes the boundary partition
function. The Newton coefficient is branch-fixed rather than inserted by hand:
\begin{equation}
\label{eq:gravity-newton-lock-restated}
 8\pi G_N
 =
 \frac{12}{c_{\rm dark}M_P^2}
 \approx
 \frac{12}{\cDarkCompletion M_P^2}.
\end{equation}
Gravity is therefore present on the selected branch specifically because the
positive central-charge deficit $c_{\rm dark}$ must be carried by a covariant
bulk channel; without that completion residue there is no Topological EFE to
support the visible stress tensor.

\paragraph{Hubble-Friction Lock.}
The same branch-fixed surface tension also determines the unique isotropic
friction channel that can absorb the packing deficit
$\delta_{\rm topo}=1-\kappa_{D_5}$. On the anomaly-free background,
\begin{equation}
\label{eq:gravity-hubble-lock-preview}
 H_{\Lambda}
 \equiv
 \sqrt{\frac{\Lambda_{\rm holo}}{3}},
 \qquad
 3H_{\Lambda}^2=\Lambda_{\rm holo},
\end{equation}
and the Bianchi lock from Eq.~\eqref{eq:gravity-bianchi-lock-local} with
$\partial_{\nu}\Lambda_{\rm holo}=0$ requires
\begin{equation}
\label{eq:gravity-frw-friction-preview}
 \dot\rho_{\rm vis}+3H_{\Lambda}(\rho_{\rm vis}+p_{\rm vis})=0.
\end{equation}
If one nevertheless postulates a different isotropic damping rate $\Gamma$,
\begin{equation}
\label{eq:gravity-general-friction-preview}
 \dot\rho_{\rm vis}+3\Gamma(\rho_{\rm vis}+p_{\rm vis})=0,
\end{equation}
then subtracting Eq.~\eqref{eq:gravity-frw-friction-preview} gives
\begin{equation}
\label{eq:gravity-friction-uniqueness-preview}
 3(\Gamma-H_{\Lambda})(\rho_{\rm vis}+p_{\rm vis})=0.
\end{equation}
On the matter-carrying branch $\rho_{\rm vis}+p_{\rm vis}\neq0$, and once
$\Delta_{\rm fr}=0$ no second isotropic counter-term survives. Hence the only
solution is
\begin{equation}
\label{eq:gravity-friction-solution-preview}
 \Gamma=H_{\Lambda}.
\end{equation}
Therefore the isotropic Hubble friction $3H_{\Lambda}$ is the unique mechanism
that absorbs the packing deficit $\delta_{\rm topo}$ into the positive vacuum
loading rather than leaving a residual source in $\nabla^{\mu}T^{\rm vis}_{\mu\nu}$.
Section~G.3 gives the expanded contradiction proof in the moir\'e-consistency
language, but the lock itself is already forced by the Topological EFE identity.

Let $B$ be a ball-shaped boundary region, let $H_{\rm mod}(B)$ be its modular Hamiltonian, and let $\xi_B$ be the associated modular-flow vector field that acts as the local boost/isometry at the entangling cut. The first law of entanglement requires
\begin{equation}
\label{eq:gravity-first-law}
 \delta S_{EE}(B)=\delta\langle H_{\rm mod}(B)\rangle.
\end{equation}
For perturbations around the anomaly-free background, the gravitational side of the same identity is the canonical constraint on the bulk Cauchy slice $\Sigma_B$ bounded by the extremal surface of $B$,
\begin{equation}
\label{eq:gravity-first-law-constraint}
 \delta S_{EE}(B)-\delta\langle H_{\rm mod}(B)\rangle
 =
 \frac{6}{c_{\rm dark}M_P^2}
 \int_{\Sigma_B}
 \xi_B^{\nu}\,\delta\mathcal E_{\mu\nu}
 \, d\Sigma^{\mu}.
\end{equation}
Preserving the branch identity means that Eq.~\eqref{eq:gravity-first-law} must hold for every ball $B$ and for every allowed linearized perturbation. The only way for the right-hand side of Eq.~\eqref{eq:gravity-first-law-constraint} to vanish for arbitrary modular-flow generators $\xi_B$ is
\begin{equation}
\label{eq:gravity-linearized-closure}
 \delta\mathcal E_{\mu\nu}=0.
\end{equation}
Substituting the definition of $\mathcal E_{\mu\nu}$ then gives the linearized Einstein field equations on the branch,
\begin{equation}
\label{eq:gravity-linearized-efe}
 \delta G_{\mu\nu}+\Lambda_{\rm holo}\,\delta g_{\mu\nu}
 =
 \frac{12}{c_{\rm dark}M_P^2}\,\delta T^{\rm vis}_{\mu\nu}.
\end{equation}
Thus the entanglement first law is the linearized expression of the same branch closure already encoded in Eq.~\eqref{eq:gravity-effective-einstein-local}. In this appendix, the vanishing of $\delta\mathcal E_{\mu\nu}$ is used as a consistency check: it confirms that the selected RCFT obeys the standard AdS/CFT modular-flow/isometry matching required for the holographic bridge on the anomaly-free branch. Any failure of the linearized Einstein equations would reopen a nonzero $\delta\mathcal E_{\mu\nu}$ and therefore signal a breakdown of that consistency condition.
Equivalently, linearized information return is only compatible with the
branch-fixed background when the finite register remains saturated. A putative
decompactification of the finite $S^3$ background would change
$\Lambda_{\rm holo}$ through Eq.~\eqref{eq:gravity-bit-budget-derived}, violate
the Unity of Scale Identity in Eq.~\eqref{eq:gravity-unity-of-scale-identity},
and reintroduce the same bulk closure defect that the Bianchi lock excludes.

The same information-return statement can be written directly in terms of the bare topological flavor overlap block. On the selected branch the restricted three-channel overlap matrix built from the branch-fixed $SU(2)_{26}$ modular data is
\begin{equation}
\label{eq:gravity-flavor-overlap}
 S_{\rm flav}
 \equiv
 \bigl[S^{(26)}_{q_{\alpha}g_i}\bigr]_{\alpha,i=1}^{3},
 \qquad
 \det S_{\rm flav}\neq0,
 \qquad
 \operatorname{rank}S_{\rm flav}=3.
\end{equation}
Hence the Gram matrix $S_{\rm flav}^{\dagger}S_{\rm flav}$ is positive definite and its polar unitary exists,
\begin{equation}
\label{eq:gravity-flavor-polar}
 U_{\rm flav}
 \equiv
 S_{\rm flav}\bigl(S_{\rm flav}^{\dagger}S_{\rm flav}\bigr)^{-1/2},
 \qquad
 U_{\rm flav}^{\dagger}U_{\rm flav}=\mathbf 1_3.
\end{equation}
If $S_{\rm flav}=W_{\rm flav}\Sigma_{\rm topo}V_{\rm flav}^{\dagger}$ is the
singular-value decomposition of the bare overlap block, then scalar matching is
implemented in the singular basis as
\begin{equation}
\label{eq:gravity-flavor-dressed-overlap}
 S_{\rm dress}
 \equiv
 W_{\rm flav}\bigl(\Sigma_{\rm topo}\mathbf M_{\rm match}\bigr)V_{\rm flav}^{\dagger}.
\end{equation}
Because $\mathbf M_{\rm match}$ is positive diagonal, it rescales only the
singular values and leaves the polar unitary $U_{\rm flav}=W_{\rm flav}V_{\rm flav}^{\dagger}$ unchanged.
If $\det S_{\rm flav}=0$, some nonzero flavor vector would lie in the kernel of $S_{\rm flav}$ and would be projected out before bulk reconstruction. Information loss would then occur already in the flavor sector. Because the benchmark block is rank-$3$, no such null direction exists.

\paragraph{Theorem IV: Holographic Unitary Lock.}
Let
\begin{equation}
\label{eq:gravity-information-return-map}
 \mathcal U_{\rm holo}:\mathcal H_{\rm bulk}^{\rm code}\to\mathcal H_{\partial}^{\rm comp}
\end{equation}
be the branch-fixed bulk-to-boundary reconstruction map from the anomaly-free
code subspace to the completed boundary Hilbert space. On the selected anomaly-free branch, the restriction of $\mathcal U_{\rm holo}$ to the visible three-channel flavor code is the polar unitary $U_{\rm flav}$ extracted from the bare overlap block $S_{\rm flav}$. The Holographic Unitary Lock is therefore secured by the affine level $k_{\ell}=26$ encoded in $S^{(26)}$ and is independent of the Higgs VEV scale, because scalar matching acts only through $\mathbf M_{\rm match}$ on the singular spectrum. Information Return is therefore not an auxiliary postulate but the direct consequence of the branch-fixed unitary lock enforced by the non-singular flavor overlap matrix,
\begin{equation}
\label{eq:gravity-information-return-isometry}
 \langle \mathcal U_{\rm holo}\psi,\mathcal U_{\rm holo}\phi\rangle_{\partial}
 =
 \langle \psi,\phi\rangle_{\rm bulk}
 \qquad\Longleftrightarrow\qquad
 \det S_{\rm flav}\neq0
 \quad
 \text{with}\quad
 \Delta_{\rm fr}=0.
\end{equation}
Equivalently, this theorem should be read as the precise structural requirement for Information
Return in the present framework rather than as a separate dynamical
information-paradox claim: on the anomaly-free branch,
$\mathcal E_{\mu\nu}=0$ together with $\delta\mathcal E_{\mu\nu}=0$ confirms
that the chosen RCFT satisfies the standard AdS/CFT isometry conditions needed
for code-subspace reconstruction.

\paragraph{Proof.}
Because $\det S_{\rm flav}\neq0$, the Hermitian matrix
$S_{\rm flav}^{\dagger}S_{\rm flav}$ is positive definite and admits the
inverse square root appearing in Eq.~\eqref{eq:gravity-flavor-polar}. Its
polar part $U_{\rm flav}$ is therefore unitary, so on the visible code subspace
\begin{equation}
 U_{\rm flav}^{\dagger}U_{\rm flav}=\mathbf 1_3
 \qquad\Longrightarrow\qquad
 \langle U_{\rm flav}\psi,U_{\rm flav}\phi\rangle
 =
 \langle \psi,\phi\rangle.
\end{equation}
Thus no flavor amplitude is projected out and the reconstruction map preserves
inner products. Equation~\eqref{eq:gravity-flavor-dressed-overlap} shows that
any Higgs-sector matching rescales only the positive singular spectrum and does
not alter the polar unitary extracted from the bare $k_{\ell}=26$ overlap
block. This is the direct linear-algebraic content of Theorem~IV: Holographic
Unitary Lock.

If instead $\det S_{\rm flav}=0$, then there exists a nonzero vector
$v\in\mathbb C^3$ with $S_{\rm flav}v=0$. The corresponding reconstructed state
has vanishing boundary norm while the bulk state still has nonzero norm, so the
map loses information already in the flavor sector and cannot be isometric.
The companion background condition is $\Delta_{\rm fr}=0$: by
Eq.~\eqref{eq:gravity-topological-identity-local} this removes the residual
scalar framing source, while Eqs.~\eqref{eq:gravity-first-law-constraint} and
\eqref{eq:gravity-linearized-closure} ensure that no linearized defect reopens
under modular flow. Once the anomaly filter is imposed, the rank-$3$
non-singularity of $S_{\rm flav}$ is precisely what guarantees Information
Return, because the bulk reconstruction map is then isometric on the completed
code subspace.

\manualappendixlabel{G.3}{app:gwb-tilt-corollary}
\manualappendixlabel{G.3}{app:gwb-tilt-requirement}
\manualappendixlabel{G.3}{app:topological-moire-consistency}
\subsection*{G.3: Topological Moir\'e Consistency, Vacuum Loading, and the Bianchi Lock}

The scalar departure from exact scale invariance is the long-wavelength
\emph{Topological Moir\'e Pattern} generated when the discrete $K=312$
boundary lattice is mapped into a continuous four-dimensional horizon
geometry. Writing the scalar mismatch as
\begin{equation}
\label{eq:gravity-moire-scalar}
 \Delta_{\rm Moire}
 \equiv
 1-n_s
 \approx
 0.0352,
 \qquad
 n_s-1\approx -0.0352,
\end{equation}
the observed red tilt is the continuum beat envelope of the branch-fixed
lattice spacing. The same branch also carries the residual packing deficit of
Eq.~\eqref{eq:gravity-delta-topo}, which supplies the tensor-side
\emph{Geometric Friction} of the $D_5$ weight-simplex packing. Defining
\begin{equation}
\label{eq:gravity-geometric-friction}
 \mathcal F_{\rm geom}
 \equiv
 1-\kappa_{D_5}
 =
 \delta_{\rm topo}
 \approx
 0.01122895,
 \qquad
 \kappa_{D_5}=\simplexPrefactor\approx 0.98877,
\end{equation}
the branch-fixed tensor tilt is
\begin{equation}
\label{eq:gravity-tensor-friction}
 n_t\approx -\mathcal F_{\rm geom}=-(1-\kappa_{D_5})\approx -0.0112.
\end{equation}
The positive vacuum loading supplies the exact isotropic damping scale needed
to absorb this packing remainder into the de Sitter background. Writing
\begin{equation}
\label{eq:gravity-hubble-scale}
 H_{\Lambda}\equiv \sqrt{\frac{\Lambda_{\rm holo}}{3}},
 \qquad
 3H_{\Lambda}^2=\Lambda_{\rm holo},
\end{equation}
the Bianchi identity for the closure tensor is
\begin{equation}
\label{eq:gravity-bianchi-lock-local}
 \nabla^{\mu}\mathcal E_{\mu\nu}
 =
 \frac{c_{\rm dark}M_P^2}{12}\,\partial_{\nu}\Lambda_{\rm holo}
 -\nabla^{\mu}T^{\rm vis}_{\mu\nu}
 =0.
\end{equation}
Because $\partial_{\nu}\Lambda_{\rm holo}=0$ on the completed $(26,8,312)$
branch, Eq.~\eqref{eq:gravity-bianchi-lock-local} reduces to the visible
stress-tensor conservation law
\begin{equation}
\label{eq:gravity-visible-stress-conservation}
 \nabla^{\mu}T^{\rm vis}_{\mu\nu}=0.
\end{equation}
Equivalently, $\nabla^{\mu}T_{\mu\nu}=0$ is a structural requirement of the
$(26,8,312)$ branch, with $T_{\mu\nu}$ understood as the branch-completed
stress tensor whose visible projection is $T^{\rm vis}_{\mu\nu}$.
On the isotropic anomaly-free branch $\partial_{\nu}\Lambda_{\rm holo}=0$, so
the only allowed residue is the scalar friction $\mathcal F_{\rm geom}$ of
Eq.~\eqref{eq:gravity-geometric-friction}; no open vector stress survives.
Equivalently, in the effective FRW rewriting the continuity equation takes the
form
\begin{equation}
\label{eq:gravity-frw-friction}
 \dot\rho_{\rm vis}+3H_{\Lambda}(\rho_{\rm vis}+p_{\rm vis})=0,
\end{equation}
with $H_{\Lambda}$ fixed entirely by the positive vacuum loading
$\Lambda_{\rm holo}$. Because the anomaly-free branch supplies no second
isotropic damping scale, Eq.~\eqref{eq:gravity-frw-friction} can be read as an
exact cancellation statement for the lifted packing remainder. Writing a
putative mismatch as
\begin{equation}
\label{eq:gravity-frw-lock-mismatch}
 \dot\rho_{\rm vis}+3H_{\Lambda}(\rho_{\rm vis}+p_{\rm vis})
 =
 \Delta_{\rm lock}(\rho_{\rm vis}+p_{\rm vis}),
\end{equation}
the anomaly-free branch requires $\Delta_{\rm lock}=0$. Here
$\Delta_{\rm lock}$ is the uncancelled FRW lift of the branch-fixed packing
deficit $\delta_{\rm topo}=1-\kappa_{D_5}$. If the Hubble friction $3H_{\Lambda}$
did not exactly match the geometric cancellation required by
$\delta_{\rm topo}$, then $\Delta_{\rm lock}\neq0$ and therefore
\paragraph{Geometric Friction Identity.}
To make the uniqueness claim explicit, replace the de Sitter damping coefficient
by a general isotropic rate $\Gamma$ and write
\begin{equation}
\label{eq:gravity-general-friction}
 \dot\rho_{\rm vis}+3\Gamma(\rho_{\rm vis}+p_{\rm vis})=0.
\end{equation}
Subtracting Eq.~\eqref{eq:gravity-frw-friction} gives the identity
\begin{equation}
\label{eq:gravity-friction-identity}
 3(\Gamma-H_{\Lambda})(\rho_{\rm vis}+p_{\rm vis})
 =
 \Delta_{\rm lock}(\rho_{\rm vis}+p_{\rm vis}).
\end{equation}
On the exact level-matched branch,
\begin{equation}
\label{eq:gravity-level-matching-local}
 K=2I_{\ell}^{\star}k_{\ell}=3I_q^{\star}k_q
 \qquad\Longrightarrow\qquad
 (k_{\ell},k_q,K)=(26,8,312),
\end{equation}
so $\Delta_{\rm fr}=0$ and the only surviving scalar residue is the packing
deficit $\delta_{\rm topo}=\mathcal F_{\rm geom}$. No second isotropic
counter-term is allowed once the framing anomaly is removed. Therefore the
Bianchi lock requires $\Delta_{\rm lock}=0$, and on the matter-carrying visible
branch $\rho_{\rm vis}+p_{\rm vis}$ is not identically zero. Equation
\eqref{eq:gravity-friction-identity} then has the unique solution
\begin{equation}
\label{eq:gravity-unique-friction}
 \Gamma=H_{\Lambda}.
\end{equation}
Hence the Hubble friction $3H_{\Lambda}$ is the unique isotropic sink that
absorbs $\delta_{\rm topo}$ into the positive vacuum loading. Detuning
$k_{\ell}$ away from $26$ reopens $\Delta_{\rm fr}\neq0$, activates the
anomalous counter-term in Eq.~\eqref{eq:gravity-detuned-einstein}, and
therefore destroys this uniqueness.
Thus,
\begin{equation}
\label{eq:gravity-visible-stress-mismatch}
 \nabla^{\mu}T^{\rm vis}_{\mu 0}
 =
 \dot\rho_{\rm vis}+3H_{\Lambda}(\rho_{\rm vis}+p_{\rm vis})
 =
 \Delta_{\rm lock}(\rho_{\rm vis}+p_{\rm vis})
 \neq 0,
\end{equation}
so the visible stress tensor would fail to be conserved. The term $3H_{\Lambda}$
is therefore the exact Hubble friction that absorbs the packing deficit
$\delta_{\rm topo}\approx0.0112$ into the positive vacuum loading rather than
leaving a Bianchi-violating source term. The same residual therefore reappears as the finite tensor drag $n_t\approx-\delta_{\rm topo}$,
while the scalar sector records the longer-wavelength moir\'e beat
$\Delta_{\rm Moire}=1-n_s$. Together they express the same branch-fixed bulk
consistency: vacuum loading, geometric friction, and the Bianchi lock are one
and the same closure statement viewed at curvature, transport, and tilt level.
At the level of the disclosed branch residues this statement is an exact
Topological Identity: $\mathcal F_{\rm geom}\approx0.01122895$ is the scalar remainder
that the continuity equation must absorb, and the tensor tilt
$n_t\approx-\mathcal F_{\rm geom}\approx-0.0112$ is the tensor-side imprint of
that same conservation law. If either relation failed, the packing deficit
would reappear as $\nabla^{\mu}T_{\mu\nu}\neq0$ and the Bulk Closure Tensor
$\mathcal E_{\mu\nu}$ would no longer vanish on the selected branch.

\paragraph{Gauge-Emergence Exclusion Principle.}
The Bianchi lock also carries a gauge-side exclusion statement. The
Gauge-Emergence Cutoff
\begin{equation}
\label{eq:gravity-gauge-emergence-cutoff-local}
 K_{\rm GE}\equiv\benchmarkParentLevel=312
\end{equation}
is a physical exclusion principle rather than a bookkeeping convention. Because
the admissible parent tower is $K=\benchmarkParentLevel n$, any
\begin{equation}
\label{eq:gravity-sparse-boundary-cutoff-local}
 K\ge624
 \qquad\Longrightarrow\qquad
 n\ge2
 \qquad\Longrightarrow\qquad
 \alpha^{-1}_{\rm surf}\ge 2\alphaSurfBenchmarkExact,
 \quad
 \alpha_{\rm surf}\le\frac{17}{4680},
\end{equation}
so the branch enters the \emph{Sparse Boundary} regime already identified in
Sec.~G.1. In that regime the visible current-current network is too dilute to
sustain an emergent four-dimensional electromagnetic sector: the gauge density
falls below the branch-local support threshold, and the bulk no longer carries
observationally viable electromagnetism. Thus $K=312$ is not merely the first
allowed parent level; it is the unique parent level that is simultaneously
Bianchi-locked and gauge-supporting, while every $K\ge624$ is physically
excluded by sparse-boundary gauge dilution~\cite{friedan1987,witten1989,gko1985}.

\paragraph{Forbidden-Void Proposition.}
A vanishing cosmological constant ($\Lambda = 0$) is topologically forbidden
on the $(26,8,312)$ branch. It necessitates a divergent bit-budget
($N\to\infty$), rendering the boundary partition function physically
non-normalizable. Through Eq.~\eqref{eq:gravity-neutrino-scale-bridge} the same
limit also collapses the neutrino floor~\cite{friedan1987,witten1989,gko1985}. Eq.~\eqref{eq:gravity-bit-budget-derived} makes this implication
explicit on the one-copy lattice. In implication form,
\begin{equation}
\label{eq:gravity-forbidden-void-chain}
 \Lambda_{\rm holo}=0
 \Longrightarrow
 H_{\Lambda}=0
 \Longrightarrow
 N_{\rm sat}=\infty
 \Longrightarrow
 Z_{\partial}^{(26,8,312)}\ \text{non-normalizable}
 \Longrightarrow
 m_{\nu}=\kappa_{D_5}M_PN^{-1/4}=0.
\end{equation}
This proposition is the boundary-normalizability form of the Bianchi lock:
$\Lambda_{\rm holo}=0$ is not a soft limit inside the admissible branch family,
but an exit from the finite-capacity Hilbert sector of the completed
$(26,8,312)$ block.
The proof-by-contradiction is then immediate. Assume that the completed
boundary modular-$T$ matrix still closed on the anomaly-free branch while
$\Lambda_{\rm holo}=0$. Then Eq.~\eqref{eq:gravity-hubble-scale} gives
$H_{\Lambda}=0$, and Eq.~\eqref{eq:gravity-bit-budget-derived} gives
$N_{\rm sat}\to\infty$, so the finite one-copy register decompactifies and the
partition function leaves the normalizable finite-capacity sector. The mass
bridge in Eq.~\eqref{eq:gravity-neutrino-scale-bridge} then forces
$m_{\nu}=\kappa_{D_5}M_PN^{-1/4}\to0$, i.e. the neutrino floor collapses. The
FRW friction channel disappears, but the branch-fixed packing residue
$\delta_{\rm topo}=1-\kappa_{D_5}>0$ remains. Equation
\eqref{eq:gravity-frw-lock-mismatch} would then force
$\Delta_{\rm lock}\neq0$, hence
\begin{equation}
\label{eq:gravity-bianchi-lock-contradiction}
 \nabla^{\mu}T^{\rm vis}_{\mu\nu}\neq0
 \qquad\Longrightarrow\qquad
 \nabla^{\mu}\mathcal E_{\mu\nu}\neq0,
\end{equation}
in contradiction with the Bianchi lock in
Eq.~\eqref{eq:gravity-bianchi-lock-local}. Therefore $\Lambda_{\rm holo}=0$
cannot satisfy both conditions simultaneously: it would require an infinite
bit-budget and thereby violate the finite-capacity premise of the
$(26,8,312)$ lattice. The only branch-compatible solution for which the
completed boundary modular-$T$ matrix closes and the bulk Bianchi identity is
satisfied is the positive vacuum loading already fixed in
Eq.~\eqref{eq:gravity-lambda-positive},
\begin{equation}
 \Lambda_{\rm holo}>0.
\end{equation}

\paragraph{Trap note.}
This is the Bianchi-lock form of the same scale trap derived in
Sec.~G.2. If the lightest neutrino mass were measured significantly
outside the predicted few-meV regime, the resulting bit-budget $N$
would be incompatible with the observed expansion rate, forcing a
non-vanishing framing anomaly $\Delta_{\rm fr}\neq0$ and collapsing the
holographic dictionary~\cite{friedan1987,witten1989,gko1985}.

\manualappendixlabel{G.4}{app:holographic-computational-gradients}
\manualappendixlabel{G.4}{app:clock-skew-derivation}
\subsection*{G.4: Holographic Computational Gradients}

The anomaly-free branch carries a finite horizon register,
$N=3\pi/(L_P^2\Lambda_{\rm holo})$, and that finite register is not only a
storage bound but also a finite update budget. Let $N_{\rm load}(t)$ denote the
cumulative number of boundary bits that have been loaded onto the completion
sector by cosmic time $t$, and define the normalized completion fraction
\begin{equation}
\label{eq:gravity-loading-fraction}
 f_{\rm load}(t)
 \equiv
 \frac{N_{\rm load}(t)}{N},
 \qquad
 0\leq f_{\rm load}(t)\leq 1.
\end{equation}
The physical interpretation of the cosmological expansion is then explicit:
\emph{The expansion of the universe is the loading speed of new boundary bits
onto the $c_{\rm dark}$ stabilizer.}
In the FRW rewriting this means that the late-time expansion rate is the
coarse-grained write speed of the completion fraction,
\begin{equation}
\label{eq:gravity-loading-hubble}
 H_{\rm eff}(t)
 =
 H_{\Lambda}
 +
 \frac{1}{3}\frac{d f_{\rm load}}{dt}.
\end{equation}

The clock-skew statement follows because the loading process is read on the
past light cone, where shells at different redshift do not share the same
proper write clock. If $\tau_{\rm lock}$ is the proper time of a local
boundary-update cycle, then the observer-frame and source-frame ticks are
related by the usual redshift factor,
\begin{equation}
\label{eq:gravity-clock-skew}
 d\tau_{\rm lock} = \frac{dt}{1+z}.
\end{equation}
Therefore the loading law becomes
\begin{equation}
\label{eq:gravity-redshift-dependent-hubble}
 H_0(z)
 =
 H_{\Lambda}
 +
 \frac{1}{3(1+z)}\frac{d f_{\rm load}}{d\tau_{\rm lock}}.
\end{equation}
Equation~\eqref{eq:gravity-redshift-dependent-hubble} is the gravity-side
clock-skew derivation of a redshift-dependent Hubble intercept: once the
boundary register is finite, the same loading process is sampled at different
stages of the write cycle by low-$z$ and high-$z$ observations. A strictly
redshift-independent $H_0$ would require either
$df_{\rm load}/d\tau_{\rm lock}=0$ (no continuing boundary loading) or the
decompactification limit $N\to\infty$ in which the discrete update cycle
ceases to be physical. Neither possibility is compatible with the anomaly-free
branch, which carries a finite register and a nonzero positive vacuum loading.
The local and CMB Hubble determinations therefore probe different segments of
the same finite holographic computation rather than different fluids.

\manualappendixlabel{G.5}{app:gravity-parity-storage}
\subsection*{G.5: Gravitational Storage of B-L Parity}

The visible branch does not erase its missing antimatter complement; it stores
that complement as a topological conjugate in the same geometric sector that is
measured macroscopically by the completion residue $c_{\rm dark}$. In branch
language the visible baryon excess is admissible only because the completed
parent block remains neutral,
\begin{equation}
\label{eq:gravity-parent-neutrality}
 (B-L)_{\rm vis}+(B-L)_{\rm dark}=0,
 \qquad
 (B-L)_{\rm dark}=-(B-L)_{\rm vis}.
\end{equation}
Here $(B-L)_{\rm dark}$ is not a second thermal fluid and not a separately
propagating antimatter bath. It is the non-propagating parity sink of the
geometric completion sector: the missing antimatter of the visible projection
is stored as a topological conjugate in the $c_{\rm dark}$ branch residue.

This is why the same Holographic Stabilizer can play two roles without
contradiction. The residue $c_{\rm dark}$ sources the positive vacuum loading
$\Lambda_{\rm holo}$ in the Topological EFE, and the same residue stores the
conjugate $B-L$ bits needed to keep the projected boundary neutral. The
gravitational field and the parity sink are therefore two macroscopic images of
one and the same completion sector.

The finite de Sitter register of Eq.~\eqref{eq:gravity-horizon-bits} gives the
storage capacity for this parity bookkeeping. In that sense gravitational
storage is literal within the benchmark logic: the branch does not hide the
missing antimatter in an untracked particle sector, but in the geometric
closure data of the completed boundary block. The bulk is neutral because the
same $c_{\rm dark}$ sector that closes $\mathcal E_{\mu\nu}$ also stores the
conjugate parity information required by Eq.~\eqref{eq:gravity-parent-neutrality}.

\ifgravitystandalone
\begin{thebibliography}{99}
\bibitem{georgi1979}
H.~Georgi and C.~Jarlskog,
``A new lepton-quark mass relation in a unified theory,''
\emph{Phys. Lett. B}~86 (1979) 297--300.

\bibitem{babu1993}
K.~S.~Babu and R.~N.~Mohapatra,
``Predictive neutrino spectrum in minimal $SO(10)$ grand unification,''
\emph{Phys. Rev. Lett.}~70 (1993) 2845--2848.

\bibitem{friedan1987}
D.~Friedan and S.~Shenker,
``The Analytic Geometry of Two-Dimensional Conformal Field Theory,''
\emph{Nucl. Phys. B}~281 (1987) 509--545.

\bibitem{witten1989}
E.~Witten,
``Quantum field theory and the Jones polynomial,''
\emph{Commun. Math. Phys.}~121 (1989) 351--399.

\bibitem{gko1985}
P.~Goddard, A.~Kent, and D.~Olive,
``Virasoro algebras and coset space models,''
\emph{Phys. Lett. B}~152 (1985) 88--92.
\end{thebibliography}
\fi

\gravityenddocument
}

\section*{Appendix K: Terminology Convention}\manualappendixlabel{K}{app:interpretive-mappings}

This appendix fixes the terminology used for the completion residue. The
language below does not alter branch selection, PMNS/CKM transport, the RG
Consistency Audit, the reported $\chi^2$ values, or the fundamental
consistency of the $SO(10)_{312}$ benchmark.

\paragraph{Terminology convention.}
The main text uses standard phrases such as \emph{boundary reconstruction map},
\emph{support map}, \emph{modular overlap}, and \emph{topological map}. When a
compact label is useful, the completion residue may also be named explicitly.

\paragraph{Topological Gravitational Ghost (TGG).}
The label \emph{Topological Gravitational Ghost of the framing defect} denotes
the non-propagating completion residue $c_{\rm dark}\equiv c_{\rm comp}$ and
its coarse-grained gravitational footprint. It does not denote a WIMP
candidate, a particle fluid, or an additional freely scanning matter sector.

\paragraph{Benchmark status.}
The benchmark is evaluated on one anomaly-free branch with fixed support data
and without any observer-dependent tuning rule. The terminology introduced here
is purely nominative: it does not enlarge the benchmark tuple or alter the
disclosed frequentist consistency statistics.

\begin{thebibliography}{99}

\bibitem{cohen1999}
A.~G.~Cohen, D.~B.~Kaplan, and A.~E.~Nelson,
``Effective field theory, black holes, and the cosmological constant,''
\emph{Phys. Rev. Lett.}~82 (1999) 4971--4974.

\bibitem{feruglio2017}
F.~Feruglio,
``Are neutrino masses modular forms?,''
arXiv:1706.08749.

\bibitem{kobayashi2018}
T.~Kobayashi, K.~Tanaka, and T.~H.~Tatsuishi,
``Neutrino mixing from finite modular groups,''
\emph{Phys. Rev. D}~98 (2018) 016004.

\bibitem{criado2018}
J.~C.~Criado and F.~Feruglio,
``Modular invariance faces precision neutrino data,''
\emph{SciPost Phys.}~5 (2018) 042.

\bibitem{kobayashi2017}
T.~Kobayashi, S.~Nagamoto, and S.~Uemura,
``Modular symmetry in magnetized/intersecting D-brane models,''
\emph{PTEP}~2017 (2017) 023B02.

\bibitem{heckman2010}
J.~J.~Heckman and C.~Vafa,
``Flavor hierarchy from F-theory,''
\emph{Nucl. Phys. B}~837 (2010) 137--151.

\bibitem{bouchard2010}
V.~Bouchard, J.~Heckman, J.~Seo, and C.~Vafa,
``F-theory and neutrinos: Kaluza--Klein dilution of flavor hierarchy,''
\emph{JHEP}~01 (2010) 061.

\bibitem{georgi1979}
H.~Georgi and C.~Jarlskog,
``A new lepton-quark mass relation in a unified theory,''
\emph{Phys. Lett. B}~86 (1979) 297--300.

\bibitem{babu1993}
K.~S.~Babu and R.~N.~Mohapatra,
``Predictive neutrino spectrum in minimal $SO(10)$ grand unification,''
\emph{Phys. Rev. Lett.}~70 (1993) 2845--2848.

\bibitem{ooguri2007}
H.~Ooguri and C.~Vafa,
``On the Geometry of the String Landscape and the Swampland,''
\emph{Nucl. Phys. B}~766 (2007) 21--33.

\bibitem{almheiri2015}
A.~Almheiri, X.~Dong, and D.~Harlow,
``Bulk Locality and Quantum Error Correction in AdS/CFT,''
\emph{JHEP}~04 (2015) 163.

\bibitem{pastawski2015}
F.~Pastawski, B.~Yoshida, D.~Harlow, and J.~Preskill,
``Holographic quantum error-correcting codes: Toy models for the bulk/boundary correspondence,''
\emph{JHEP}~06 (2015) 149.

\bibitem{friedan1987}
D.~Friedan and S.~Shenker,
``The Analytic Geometry of Two-Dimensional Conformal Field Theory,''
\emph{Nucl. Phys. B}~281 (1987) 509--545.

\bibitem{witten1989}
E.~Witten,
``Quantum field theory and the Jones polynomial,''
\emph{Commun. Math. Phys.}~121 (1989) 351--399.

\bibitem{antusch2005}
S.~Antusch, J.~Kersten, M.~Lindner, M.~Ratz, and M.~A.~Schmidt,
``Running neutrino mass parameters in see-saw scenarios,''
\emph{JHEP}~03 (2005) 024.

\bibitem{bpz1984}
A.~A.~Belavin, A.~M.~Polyakov, and A.~B.~Zamolodchikov,
``Infinite conformal symmetry in two-dimensional quantum field theory,''
\emph{Nucl. Phys. B}~241 (1984) 333--380.

\bibitem{gko1985}
P.~Goddard, A.~Kent, and D.~Olive,
``Virasoro algebras and coset space models,''
\emph{Phys. Lett. B}~152 (1985) 88--92.

\end{thebibliography}

\end{document}
